% GENERATED FILE -- DO NOT EDIT.
%
% Produced by harness/assemble_contract.py from module.yaml and claims/*.tex.
% Edit the claim files and regenerate; `--check` fails if this file is stale.
\documentclass[11pt,a4paper]{article}
\pdftrailerid{e15bbad2dd79964c2d7f665747bac719}
%% --- 1. Mathematics (order matters: physics after braket wins \ket, \tr) ---
\usepackage{amsmath}
\usepackage{amssymb}
\usepackage{amsfonts}
\usepackage{amsthm}
\usepackage{mathtools}
\usepackage{dsfont}
\usepackage{braket}
\usepackage{physics}
%% --- 2. Page furniture ------------------------------------------------------
\usepackage{xcolor}
\usepackage{array}    % >{...} column prefixes in the wide tables of some claims
\usepackage{enumitem}
\usepackage{graphicx}
\usepackage{booktabs}
\usepackage{longtable}
\usepackage{xltabular}              % front- and back-matter tables across pages
\usepackage{needspace}
\usepackage{xurl}
\usepackage[noautomatic]{imakeidx}  % defined-terms index; latexmk runs makeindex
\makeindex[title=Index of defined terms,intoc]
\usepackage{tcolorbox}
\tcbuselibrary{skins,breakable}
%% --- 3. Links and navigation (hyperref late; bookmark, refcount after it) ---
\usepackage{hyperref}
\usepackage{bookmark}
\usepackage{refcount}
%% --- 4. Colours: every distinction also has a shape (rule, caps, size, weight)
\definecolor{nnGrey}{gray}{0.40}      % non-normative text
\definecolor{metaGrey}{gray}{0.30}    % the generated meta line
\definecolor{techGrey}{gray}{0.45}    % heading of a technical claim
\definecolor{mainRule}{RGB}{0,110,140}
\definecolor{reqAccent}{RGB}{165,70,0}
\definecolor{linkBlue}{RGB}{0,62,150}
\definecolor{extLink}{RGB}{120,40,140}  % links into another PDF
\hypersetup{colorlinks=true, linkcolor=linkBlue, citecolor=linkBlue,
  urlcolor=linkBlue, filecolor=extLink, bookmarksnumbered=true,
  bookmarksopen=true, bookmarksopenlevel=1, pdfdisplaydoctitle=true}
% A link into another PDF (\href to a .pdf, \hyperref{file}..., xr-hyper) is a
% relative URI, <file>#nameddest=<destination>, in the file-link colour:
% browser viewers (Edge, Chrome) ignore hyperref's default GoToR action.
\makeatletter
\renewcommand*{\hyper@linkfile}[3]{% anchor text, file name, destination
  \begingroup
    \ifx\\#3\\\def\CQ@target{#2}\else\def\CQ@target{#2\hyper@hash nameddest=#3}\fi
    \let\@urlcolor\@filecolor
    \hyper@linkurl{#1}{\CQ@target}%
  \endgroup}
\makeatother
%% --- 5. The mapped Unicode set (harness/contract_navigation.py, UNICODE_TEX)
\DeclareUnicodeCharacter{03B1}{\ensuremath{\alpha}}
\DeclareUnicodeCharacter{03B2}{\ensuremath{\beta}}
\DeclareUnicodeCharacter{03B3}{\ensuremath{\gamma}}
\DeclareUnicodeCharacter{03B4}{\ensuremath{\delta}}
\DeclareUnicodeCharacter{03B5}{\ensuremath{\varepsilon}}
\DeclareUnicodeCharacter{03F5}{\ensuremath{\epsilon}}
\DeclareUnicodeCharacter{03B6}{\ensuremath{\zeta}}
\DeclareUnicodeCharacter{03B7}{\ensuremath{\eta}}
\DeclareUnicodeCharacter{03B8}{\ensuremath{\theta}}
\DeclareUnicodeCharacter{03D1}{\ensuremath{\vartheta}}
\DeclareUnicodeCharacter{03B9}{\ensuremath{\iota}}
\DeclareUnicodeCharacter{03BA}{\ensuremath{\kappa}}
\DeclareUnicodeCharacter{03BB}{\ensuremath{\lambda}}
\DeclareUnicodeCharacter{03BC}{\ensuremath{\mu}}
\DeclareUnicodeCharacter{03BD}{\ensuremath{\nu}}
\DeclareUnicodeCharacter{03BE}{\ensuremath{\xi}}
\DeclareUnicodeCharacter{03C0}{\ensuremath{\pi}}
\DeclareUnicodeCharacter{03C1}{\ensuremath{\rho}}
\DeclareUnicodeCharacter{03F1}{\ensuremath{\varrho}}
\DeclareUnicodeCharacter{03C3}{\ensuremath{\sigma}}
\DeclareUnicodeCharacter{03C2}{\ensuremath{\varsigma}}
\DeclareUnicodeCharacter{03C4}{\ensuremath{\tau}}
\DeclareUnicodeCharacter{03C5}{\ensuremath{\upsilon}}
\DeclareUnicodeCharacter{03C6}{\ensuremath{\varphi}}
\DeclareUnicodeCharacter{03D5}{\ensuremath{\phi}}
\DeclareUnicodeCharacter{03C7}{\ensuremath{\chi}}
\DeclareUnicodeCharacter{03C8}{\ensuremath{\psi}}
\DeclareUnicodeCharacter{03C9}{\ensuremath{\omega}}
\DeclareUnicodeCharacter{0393}{\ensuremath{\Gamma}}
\DeclareUnicodeCharacter{0394}{\ensuremath{\Delta}}
\DeclareUnicodeCharacter{0398}{\ensuremath{\Theta}}
\DeclareUnicodeCharacter{039B}{\ensuremath{\Lambda}}
\DeclareUnicodeCharacter{039E}{\ensuremath{\Xi}}
\DeclareUnicodeCharacter{03A0}{\ensuremath{\Pi}}
\DeclareUnicodeCharacter{03A3}{\ensuremath{\Sigma}}
\DeclareUnicodeCharacter{03A5}{\ensuremath{\Upsilon}}
\DeclareUnicodeCharacter{03A6}{\ensuremath{\Phi}}
\DeclareUnicodeCharacter{03A8}{\ensuremath{\Psi}}
\DeclareUnicodeCharacter{03A9}{\ensuremath{\Omega}}
\DeclareUnicodeCharacter{00E1}{\'a}
\DeclareUnicodeCharacter{00E0}{\`a}
\DeclareUnicodeCharacter{00E2}{\^a}
\DeclareUnicodeCharacter{00E4}{\"a}
\DeclareUnicodeCharacter{00E5}{\aa{}}
\DeclareUnicodeCharacter{00E9}{\'e}
\DeclareUnicodeCharacter{00E8}{\`e}
\DeclareUnicodeCharacter{00EA}{\^e}
\DeclareUnicodeCharacter{00EB}{\"e}
\DeclareUnicodeCharacter{00ED}{\'{\i}}
\DeclareUnicodeCharacter{00EF}{\"{\i}}
\DeclareUnicodeCharacter{00F3}{\'o}
\DeclareUnicodeCharacter{00F2}{\`o}
\DeclareUnicodeCharacter{00F4}{\^o}
\DeclareUnicodeCharacter{00F6}{\"o}
\DeclareUnicodeCharacter{00F8}{\o{}}
\DeclareUnicodeCharacter{00FA}{\'u}
\DeclareUnicodeCharacter{00F9}{\`u}
\DeclareUnicodeCharacter{00FB}{\^u}
\DeclareUnicodeCharacter{00FC}{\"u}
\DeclareUnicodeCharacter{00F1}{\~n}
\DeclareUnicodeCharacter{00E7}{\c{c}}
\DeclareUnicodeCharacter{00DF}{\ss{}}
\DeclareUnicodeCharacter{00C4}{\"A}
\DeclareUnicodeCharacter{00C5}{\AA{}}
\DeclareUnicodeCharacter{00C9}{\'E}
\DeclareUnicodeCharacter{00D6}{\"O}
\DeclareUnicodeCharacter{00DC}{\"U}
\DeclareUnicodeCharacter{2013}{\textendash{}}
\DeclareUnicodeCharacter{2014}{\textemdash{}}
\DeclareUnicodeCharacter{2018}{`}
\DeclareUnicodeCharacter{2019}{'}
\DeclareUnicodeCharacter{201C}{``}
\DeclareUnicodeCharacter{201D}{''}
\DeclareUnicodeCharacter{2026}{\ldots{}}
\DeclareUnicodeCharacter{00A0}{~}
\DeclareUnicodeCharacter{00A7}{\S{}}
\DeclareUnicodeCharacter{2264}{\ensuremath{\le}}
\DeclareUnicodeCharacter{2265}{\ensuremath{\ge}}
\DeclareUnicodeCharacter{2260}{\ensuremath{\neq}}
\DeclareUnicodeCharacter{2248}{\ensuremath{\approx}}
\DeclareUnicodeCharacter{2261}{\ensuremath{\equiv}}
\DeclareUnicodeCharacter{223C}{\ensuremath{\sim}}
\DeclareUnicodeCharacter{2243}{\ensuremath{\simeq}}
\DeclareUnicodeCharacter{2245}{\ensuremath{\cong}}
\DeclareUnicodeCharacter{221D}{\ensuremath{\propto}}
\DeclareUnicodeCharacter{226A}{\ensuremath{\ll}}
\DeclareUnicodeCharacter{226B}{\ensuremath{\gg}}
\DeclareUnicodeCharacter{2192}{\ensuremath{\to}}
\DeclareUnicodeCharacter{2190}{\ensuremath{\leftarrow}}
\DeclareUnicodeCharacter{2194}{\ensuremath{\leftrightarrow}}
\DeclareUnicodeCharacter{21D2}{\ensuremath{\Rightarrow}}
\DeclareUnicodeCharacter{21D0}{\ensuremath{\Leftarrow}}
\DeclareUnicodeCharacter{21D4}{\ensuremath{\Leftrightarrow}}
\DeclareUnicodeCharacter{21A6}{\ensuremath{\mapsto}}
\DeclareUnicodeCharacter{2208}{\ensuremath{\in}}
\DeclareUnicodeCharacter{2209}{\ensuremath{\notin}}
\DeclareUnicodeCharacter{2286}{\ensuremath{\subseteq}}
\DeclareUnicodeCharacter{2282}{\ensuremath{\subset}}
\DeclareUnicodeCharacter{2287}{\ensuremath{\supseteq}}
\DeclareUnicodeCharacter{2283}{\ensuremath{\supset}}
\DeclareUnicodeCharacter{2229}{\ensuremath{\cap}}
\DeclareUnicodeCharacter{222A}{\ensuremath{\cup}}
\DeclareUnicodeCharacter{2216}{\ensuremath{\setminus}}
\DeclareUnicodeCharacter{2205}{\ensuremath{\emptyset}}
\DeclareUnicodeCharacter{2200}{\ensuremath{\forall}}
\DeclareUnicodeCharacter{2203}{\ensuremath{\exists}}
\DeclareUnicodeCharacter{2115}{\ensuremath{\mathbb{N}}}
\DeclareUnicodeCharacter{2124}{\ensuremath{\mathbb{Z}}}
\DeclareUnicodeCharacter{211A}{\ensuremath{\mathbb{Q}}}
\DeclareUnicodeCharacter{211D}{\ensuremath{\mathbb{R}}}
\DeclareUnicodeCharacter{2102}{\ensuremath{\mathbb{C}}}
\DeclareUnicodeCharacter{00B7}{\ensuremath{\cdot}}
\DeclareUnicodeCharacter{00D7}{\ensuremath{\times}}
\DeclareUnicodeCharacter{2212}{\ensuremath{-}}
\DeclareUnicodeCharacter{00B1}{\ensuremath{\pm}}
\DeclareUnicodeCharacter{2213}{\ensuremath{\mp}}
\DeclareUnicodeCharacter{221A}{\ensuremath{\surd}}
\DeclareUnicodeCharacter{2211}{\ensuremath{\sum}}
\DeclareUnicodeCharacter{220F}{\ensuremath{\prod}}
\DeclareUnicodeCharacter{222B}{\ensuremath{\int}}
\DeclareUnicodeCharacter{221E}{\ensuremath{\infty}}
\DeclareUnicodeCharacter{2202}{\ensuremath{\partial}}
\DeclareUnicodeCharacter{2207}{\ensuremath{\nabla}}
\DeclareUnicodeCharacter{2297}{\ensuremath{\otimes}}
\DeclareUnicodeCharacter{2295}{\ensuremath{\oplus}}
\DeclareUnicodeCharacter{2218}{\ensuremath{\circ}}
\DeclareUnicodeCharacter{2308}{\ensuremath{\lceil}}
\DeclareUnicodeCharacter{2309}{\ensuremath{\rceil}}
\DeclareUnicodeCharacter{230A}{\ensuremath{\lfloor}}
\DeclareUnicodeCharacter{230B}{\ensuremath{\rfloor}}
\DeclareUnicodeCharacter{27E8}{\ensuremath{\langle}}
\DeclareUnicodeCharacter{27E9}{\ensuremath{\rangle}}
\DeclareUnicodeCharacter{2016}{\ensuremath{\|}}
\DeclareUnicodeCharacter{2113}{\ensuremath{\ell}}
\DeclareUnicodeCharacter{00B0}{\ensuremath{^{\circ}}}
\DeclareUnicodeCharacter{00B2}{\ensuremath{^{2}}}
\DeclareUnicodeCharacter{00B3}{\ensuremath{^{3}}}
\DeclareUnicodeCharacter{00B9}{\ensuremath{^{1}}}
\DeclareUnicodeCharacter{2080}{\ensuremath{_{0}}}
\DeclareUnicodeCharacter{2081}{\ensuremath{_{1}}}
\DeclareUnicodeCharacter{2082}{\ensuremath{_{2}}}
\DeclareUnicodeCharacter{207F}{\ensuremath{^{n}}}
\DeclareUnicodeCharacter{2096}{\ensuremath{_{k}}}
%% --- 6. Numbering: one counter per section -----------------------
\numberwithin{equation}{section}
\setcounter{tocdepth}{1}
%% --- 7. The project macro set, then the claim environments -----------------
% spec/macros.tex, read relative to this module's directory (latexmk -cd).
% =============================================================================
%  spec/macros.tex -- the one macro set every generated contract inputs.
%
%  Two parts: part A, the macros spec/notation.md catalogues, and part B,
%  the macros local to one module, grouped by module.
%  harness/assemble_contract.py emits % =============================================================================
%  spec/macros.tex -- the one macro set every generated contract inputs.
%
%  Two parts: part A, the macros spec/notation.md catalogues, and part B,
%  the macros local to one module, grouped by module.
%  harness/assemble_contract.py emits % =============================================================================
%  spec/macros.tex -- the one macro set every generated contract inputs.
%
%  Two parts: part A, the macros spec/notation.md catalogues, and part B,
%  the macros local to one module, grouped by module.
%  harness/assemble_contract.py emits \input{../../macros} for it, so the
%  file is read from spec/modules/<id>/.
%
%  RULES
%    * A macro that spec/notation.md catalogues is a verbatim copy of it
%      (part A). Changing one is a notation change: bump spec/notation.md.
%    * Part B holds the module-local macros. A macro defined by two modules
%      with the same expansion appears once, under the first module (in the
%      order below) that defines it. Two macros with the same name and
%      DIFFERENT expansions are both kept under a module prefix, and the
%      claims write the prefixed names:
%        \Chain  -> \redChain (reductions, \mathfrak{C}: the reduction chain)
%                   \pcChain  (protocol-compiler, \mathfrak{R}: the two-segment
%                             chain; \Chsrc and \Chrec are built on it)
%        \Nmsm   -> \redNmsm  (reductions, \mathcal{N}: a POVM measurement
%                             channel)
%                   \pcNmsm   (\mathcal{N}_{\mathsf A}: Alice's measurement
%                             channel, written by
%                             transmitter/thm:certificate-sufficiency,
%                             transmitter/rem:xreg and
%                             reductions/lem:source-transfer)
%    * Deliberately omitted: a \provenance{tag}{details} macro. Provenance
%      is a header field of every claim file (spec/claim-format.md section
%      3); a claim body that calls such a macro fails to compile, which is
%      the intended signal.
%
%  REQUIRED PACKAGES, loaded by the assembler's preamble before this file:
%    amsmath, amssymb, amsfonts, mathtools, dsfont, braket, physics, xcolor
%    (braket before physics, so that physics wins for \ket, \bra, \tr, ...).
%
%  This file contains only definitions: no \usepackage, no document body.
% =============================================================================


% =============================================================================
%  PART A -- notation macros, the ones spec/notation.md catalogues, each a
%  verbatim copy of its catalogue row; grouped under numbered headings.
% =============================================================================

% -----------------------------------------------------------------------------
% 1. State spaces and operator sets (a "system label" -> spaces of operators)
% -----------------------------------------------------------------------------
\newcommand{\Hs}[1]{\mathcal{H}(#1)}        % Hilbert space of a system label
\newcommand{\Lin}[1]{\mathcal{L}(#1)}       % bounded linear operators on H(#1)
\newcommand{\Linn}[2]{\mathcal{L}(#1,#2)}   % linear maps H(#1) -> H(#2)
\newcommand{\Pos}[1]{\mathrm{Pos}(#1)}      % positive semidefinite operators on H(#1)
\newcommand{\Density}{\mathrm{D}}           % density operators: \Density(A) = D(A)
\newcommand{\Seq}[1]{S_{=}(#1)}             % normalized states (trace = 1)
\newcommand{\Sleq}[1]{S_{\leq}(#1)}         % subnormalized states (0 <= trace <= 1)
\newcommand{\id}{\mathds{1}}                % identity operator
\newcommand{\idon}[1]{\mathds{1}_{#1}}      % identity on a named system, e.g. \idon{C}

% DEPRECATED: applies \mathcal to both arguments itself, so it takes bare
% letters, not system labels. Prefer \Linn; do not use in new contracts.
\providecommand{\Lop}[2]{\mathcal{L}(\mathcal{#1}, \mathcal{#2})}


% -----------------------------------------------------------------------------
% 2. Composite systems
% -----------------------------------------------------------------------------
\newcommand{\dsum}{\oplus}                  % direct sum of system labels
% Tensor product: use \otimes directly.
% n-fold direct sum of an operator space: \Lin{A}^{\oplus n}


% -----------------------------------------------------------------------------
% 3. Events, conditioning and classical registers (for STATES)
%    Convention:  \wedge  -> subnormalized "and-event" component, e.g. \rho_{A\wedge x}
%                 |        -> normalized conditional state,        e.g. \rho_{A|x}
% (Write these directly in subscripts; no extra macros needed.)
% -----------------------------------------------------------------------------


% -----------------------------------------------------------------------------
% 4. Photon-number subspaces / direct-sum blocks
%    Convention:  operators carry the restriction "on top", in brackets;
%                 states carry it as a \wedge subscript (see Sec. 3).
% -----------------------------------------------------------------------------
\newcommand{\Pnum}[1]{P_{#1}}               % photon-number projector (system from context)
\newcommand{\Pnumon}[2]{P_{#1}^{#2}}        % with explicit system superscript, e.g. P_n^C
\newcommand{\blk}[2]{#1^{(#2)}}             % operator block, e.g. \blk{M}{\leqN} = M^{(<=N)}
\newcommand{\leqN}{\leq N}                  % use as \blk{M}{\leqN}
\newcommand{\gtN}{> N}                      % use as \blk{M}{\gtN}
\newcommand{\eqn}{= n}                      % use as \blk{M}{\eqn}


% -----------------------------------------------------------------------------
% 5. System labels used in the protocol
% -----------------------------------------------------------------------------
\newcommand{\Chan}{C}                       % quantum-channel (Eve) input to Bob
\newcommand{\Prb}{P}                        % probe input
\newcommand{\Flag}{{\mathrm{Flg}}}          % flag system (EXISTS ONLY AFTER squashing)
\newcommand{\Loss}{{\mathrm{Lss}}}          % loss system
\newcommand{\Vc}{{\mathrm{Vcm}}}            % vacuum system
\newcommand{\Nout}{N_{\mathrm{out}}}
\newcommand{\Nin}{N_{\mathrm{in}}}


% -----------------------------------------------------------------------------
% 6. Alphabets
% -----------------------------------------------------------------------------
\newcommand{\Yset}{\mathcal{Y}}             % Bob's (coarse-grained) detection alphabet
\newcommand{\Yraw}{\mathcal{Y}_{\mathrm{raw}}}% raw click patterns {0,1}^4
\newcommand{\Yscset}{\mathcal{Y}_{\singlec}}% single-click outcomes
\newcommand{\Ymcset}{\mathcal{Y}_{\mc}}     % multi-click outcomes
\newcommand{\Tset}{\mathcal{T}_{\mathrm{set}}}% private setting set (the probe settings in the optical instance)
\newcommand{\Xset}{\mathcal{X}}             % Alice's alphabet
\newcommand{\Dset}{\{H,V,+,-\}}             % polarization / detector labels
\newcommand{\pass}{\top}                    % pass sifting


% -----------------------------------------------------------------------------
% 7. Mode / outcome abbreviations
% -----------------------------------------------------------------------------
\newcommand{\on}{{\mathrm{on}}}
\newcommand{\off}{{\mathrm{off}}}
\newcommand{\dc}{{\mathrm{dc}}}             % dark count
\newcommand{\sq}{\mathrm{sq}}               % squashed
\newcommand{\mc}{\mathrm{mc}}               % multi-click
\newcommand{\nc}{\mathrm{nc}}               % no-click
\newcommand{\singlec}{\mathrm{sc}}          % single-click
\newcommand{\vac}{\mathrm{vac}}             % vacuum
\newcommand{\kout}{{k_\mathrm{out}}}
\newcommand{\kin}{{k_\mathrm{in}}}


% -----------------------------------------------------------------------------
% 8. Maps, channels and miscellaneous operators
% -----------------------------------------------------------------------------
\newcommand{\Prob}{\mathrm{Prob}}           % probability (use \Prob[\cdot] or \Prob(\cdot))
\newcommand{\lmin}{\lambda_{\min}}          % minimum eigenvalue
\newcommand{\Wflag}{W_{\Flag}}              % flag-space weight
\newcommand{\transpose}{\mathsf{T}}
\newcommand{\innerp}[2]{\left\langle #1 \vert #2 \right\rangle}
\renewcommand{\vec}[1]{\boldsymbol{#1}}     % NB: overrides the LaTeX arrow accent


% -----------------------------------------------------------------------------
% 9. Key-rate numerics
% -----------------------------------------------------------------------------
\newcommand{\cG}{\mathcal{G}}               % protocol map
\newcommand{\cZ}{\mathcal{Z}}               % pinching map
\newcommand{\Amap}{\mathcal{A}_{\mathrm{ann}}} % announcement + measurement map
\newcommand{\Amapnb}{\mathcal{A}_{\mathrm{ann}|\mathrm{nb}}}
\newcommand{\simpl}{}                       % empty suffix for simplified block maps
\newcommand{\acSet}{\mathcal{S}_{\mathrm{acc}}} % acceptance set
\newcommand{\testE}{\mathcal{E}_{\mathrm{test}}} % test-statistics map
\newcommand{\ecLeak}{\mathrm{leak}^{\mathrm{EC}}_{\mathrm{obs}}}
\newcommand{\pgen}{p_{\mathrm{gen}}}
\newcommand{\askeyr}{R_{\infty}}             % asymptotic key rate
\newcommand{\ppass}{p_{\mathrm{pass}}}       % sifting pass probability


% -----------------------------------------------------------------------------
% 10. Finite-size analysis / Renyi entropies
% -----------------------------------------------------------------------------

% --- Round types: test vs. key generation (register T, alphabet \Fset) --------
\newcommand{\tst}{\mathrm{test}}            % test round
\newcommand{\gen}{\mathrm{gen}}             % key-generation round
\newcommand{\Fset}{\mathcal{T}}             % round-type alphabet {test, gen} (not \Tset)
\newcommand{\Chat}{\widehat{C}}             % acceptance-test register
\newcommand{\Chatset}{\widehat{\mathcal{C}}}% its alphabet

% --- Basis / key extraction functions ----------------------------------------
\newcommand{\bA}{b_{A}}                     % Alice's basis-extraction function
\newcommand{\bB}{b_{B}}                     % Bob's basis-extraction function
\newcommand{\kB}{k_{B}}                     % Bob's key-extraction function

% --- Maps --------------------------------------------------------------------
\newcommand{\Mprot}{\mathcal{M}}            % single-round protocol map
\newcommand{\Phimap}{\boldsymbol{\Phi}}     % test-statistics map
\newcommand{\Pdeph}{\mathcal{P}}            % dephasing (pinching) map on \Blreg

% --- Isometries --------------------------------------------------------------
\newcommand{\Vmsm}{V_{\mathrm{msm}}}        % measurement isometry
\newcommand{\Vann}{V_{\mathrm{ann}}}        % announcement isometry
\newcommand{\VI}{V_{I}}                     % Eve's copy of the announcements
\newcommand{\VChat}{V_{\Chat}}              % acceptance-data isometry
\newcommand{\VR}{V_{R}}                     % sifting + key map -> raw key R
\newcommand{\VZ}{V_{Z}}                     % raw key -> key register Z
\newcommand{\Wprot}{W}                      % full protocol isometry

% --- Renyi entropies ---------------------------------------------------------
% SECURITY-CRITICAL: see spec/notation.md §4. Never change an arrow to
% make an inequality work out; never convert silently between the sandwiched
% (tilde) and Petz (overline) families.
\newcommand{\Hsandup}{\widetilde{H}^{\uparrow}_{\alpha}}   % sandwiched, up
\newcommand{\Hsanddown}{\widetilde{H}^{\downarrow}_{\alpha}}% sandwiched, down
\newcommand{\Hpetzup}{\overline{H}^{\uparrow}_{\alpha}}    % Petz, up
\newcommand{\Hpetzdown}{\overline{H}^{\downarrow}_{\alpha}}% Petz, down
\newcommand{\hup}{h^{\uparrow}_{\alpha}}                   % single-round rate
\newcommand{\hqkd}{h^{\uparrow,\mathrm{QKD}}_{\alpha}}     % single-round QKD rate
\newcommand{\hlow}{\underline{h}_{\alpha}}                 % lower bound on \hqkd
% --- spec/notation.md 3.5 and 4.5 -------------------------------------------
% \hlb  is the first-order linearization lower bound on \hlow: the headline
%       key rate (the `first-order` tier of foundations/def:tiers). NOT \hlow
%       itself.
% \tnum is the instance's one declared numerical tolerance. Its
%       near-collision with the declared marginal \tau_A is disclosed in
%       spec/notation.md 3.5; bare \tau is reserved and denotes neither.
\newcommand{\hlb}{\underline{h}^{\mathrm{FO}}_{\alpha}}    % first-order lower bound on \hlow
\newcommand{\tnum}{\tau_{\mathrm{num}}}                    % declared numerical tolerance


% =============================================================================
%  PART B -- module-local macros, grouped by module. Deprecated aliases,
%  defined so that a text written with them compiles, used in no new contract: \Lop,
%  \testE, \Pdeph, \pass, \ecLeak (spec/notation.md section 3.9).
% =============================================================================

% --- meat-core -----------------------------------------------------------------
\newcommand{\SecSet}{\mathbb{S}}            % sector index set of Q
\newcommand{\Qsys}{Q}                       % declared sector-carrying measured system
\newcommand{\Vset}{\mathcal{V}}             % declared joint outcome alphabet of the Q-POVM
\newcommand{\CAset}{\mathcal{C}_{A}}        % Alice announcement alphabet
\newcommand{\CBset}{\mathcal{C}_{B}}        % Bob announcement alphabet
\newcommand{\Party}{\mathfrak{P}}           % declared party partition
\newcommand{\Oset}[1]{\mathcal{O}_{#1}}     % readout alphabet of party #1
\newcommand{\rd}[1]{\pi_{#1}}               % readout map of party #1
\newcommand{\Sadm}{\mathcal{S}_{A}}         % admissible Alice subspace
\newcommand{\nomarg}{\mathsf{none}}         % "no trusted marginal declared"
\newcommand{\Zset}{\mathcal{Z}_{\perp}}     % raw-key alphabet incl. \perp
\newcommand{\Zval}{\mathcal{Z}}             % key alphabet without \perp
\newcommand{\Etil}{\widetilde{E}}           % single-round purifying register
\newcommand{\Ehat}{\widehat{E}}             % Eve's n-round register
\newcommand{\OmAT}{\Omega_{\mathrm{AT}}}    % acceptance event
\newcommand{\epa}{\varepsilon_{\mathrm{PA}}}
\newcommand{\eev}{\varepsilon_{\mathrm{EV}}}
\newcommand{\lEC}{\lambda_{\mathrm{EC}}}
\newcommand{\Decl}{\mathfrak{D}}            % protocol declaration
\newcommand{\CList}{\mathfrak{L}}           % declared constraint list
\newcommand{\PinS}{\mathcal{P}_{\SecSet}}   % sector pinching on Q
\newcommand{\Feas}{\mathcal{W}}             % feasible set of block states
\newcommand{\Chkey}{\Chatset_{\mathrm{key}}}% key-live acceptance values
\newcommand{\Kann}{\mathcal{K}}             % announcement-class alphabet
\newcommand{\hold}{\mathsf{h}}              % declared key holder (A or B)
\newcommand{\sfD}{D}                        % classical KL divergence (base 2)
\newcommand{\Dtwo}{D}                       % quantum relative entropy (base 2)
\newcommand{\Dsand}{\widetilde{D}_{\alpha}} % sandwiched Renyi divergence
\newcommand{\idmap}{\mathrm{id}}            % identity channel
\newcommand{\Aev}{\mathcal{A}}              % Eve's per-round attack channel
\newcommand{\AttSet}{\mathfrak{A}}          % licensed attack-channel set
\newcommand{\Chacc}{\Chatset_{\mathrm{acc}}}% acceptance support of \acSet
\newcommand{\epsstar}{\varepsilon^{\star}}      % distance to an exact POVM
\newcommand{\dfH}{\Delta_{\mathrm{herm}}}       % Hermiticity defect
\newcommand{\dfP}{\Delta_{\mathrm{pos}}}        % negative part
\newcommand{\dfC}{\Delta_{\mathrm{cpl}}}        % completeness defect
\newcommand{\lmoq}{g}                       % acceptance-side gradient block
\newcommand{\lmoJ}{G}                       % feasible-side gradient block
\newcommand{\bacc}{\beta_{\mathrm{acc}}}    % certified acceptance-side minimum
\newcommand{\Itset}{\mathfrak{I}}           % finite index set of iterates (C22)

% --- reductions ----------------------------------------------------------------
\newcommand{\Bsys}{B}                        % physical (pre-squashing) system
\newcommand{\Tcl}[1]{\mathcal{T}(#1)}        % trace-class operators on H(#1)
\newcommand{\Ins}{\mathcal{I}}               % generalized-instrument branch map
\newcommand{\redNmsm}{\mathcal{N}}              % POVM measurement channel
\newcommand{\Eprot}{\mathcal{E}}             % n-round protocol channel
\newcommand{\Grnd}{\mathsf{G}}               % one round map of a protocol
\newcommand{\Kfin}{\mathsf{K}}               % end-of-protocol processing
\newcommand{\Crit}{\mathsf{sec}}             % value of a security criterion
\newcommand{\Pflag}{\Pi_{\Flag}}             % flag-sector projector on Q
\newcommand{\Pleq}{\Pi_{\leqN}}              % preserved-block projector
\newcommand{\Pgt}{\Pi_{\gtN}}                % above-cutoff projector
\newcommand{\Sph}{S^{\mathrm{ph}}}           % family of physical inputs
\newcommand{\Srch}{S^{\mathrm{rch}}}         % reachable squashed states
\newcommand{\Ssq}{S^{\sq}}                   % squashed admissible set
\newcommand{\Ffss}{F^{\mathrm{FSS}}}         % flag-state-squasher target
\newcommand{\Fwp}{F^{\mathrm{WP}}}           % multi-bound WPFSS target
\newcommand{\Mwp}{M^{\mathrm{WP}}}           % renormalized above-cutoff POVM
\newcommand{\Mdec}{{M'}^{\mathrm{ph}}}       % declared physical family on B
\newcommand{\Ppvm}{P^{\Flag}}                % flag PVM
\newcommand{\Lmp}{{\Lambda_{\mathrm{mp}}}}   % measure-and-prepare channel
\newcommand{\Lwp}{{\Lambda^{\mathrm{WP}}}}   % multi-bound WPFSS channel
\newcommand{\Lfss}{{\Lambda^{\mathrm{FSS}}}} % flag-state-squasher channel
\newcommand{\lmax}{\lambda_{\max}}           % maximum eigenvalue
\newcommand{\Lnoise}{\widetilde{\Lambda}}    % flag-noise channel
\newcommand{\cg}{\varpi}                     % declared coarse-graining map
\newcommand{\MCK}{\mathsf{MC}}               % multiclick raw-pattern set
\newcommand{\redChain}{\mathfrak{C}}            % reduction chain
\newcommand{\Lnk}{\ell}                      % a single typed link
\newcommand{\dsq}{\preceq}                      % the squashing relation
\newcommand{\opn}[1]{\left\lVert #1\right\rVert}           % operator (spectral) norm
\newcommand{\mxn}[1]{\left\lVert #1\right\rVert_{\max}}    % entrywise max modulus
\newcommand{\Fld}{\mathbb{F}}                   % exact ordered substrate field
\newcommand{\Fldi}{\mathbb{F}_{i}}              % its Gaussian extension F + iF
\newcommand{\lat}[1]{\big\lfloor #1\big\rfloor_{\kappa}} % round DOWN onto the lattice

% --- protocol-compiler ---------------------------------------------------------
\newcommand{\Spec}{\mathfrak{S}}             % protocol specification (compiler input)
\newcommand{\pcChain}{\mathfrak{R}}            % reduction chain (two typed segments)
\newcommand{\Chsrc}{\pcChain_{\mathrm{src}}}   % its Alice-side source-map segment
\newcommand{\Chrec}{\pcChain_{\mathrm{rec}}}   % its receiver-side squashing segment
\newcommand{\rvirt}{\rho^{\mathrm{v}}}       % virtual (PRE-source-map) signal state
\newcommand{\Fph}{F^{\mathrm{ph}}}           % physical (pre-reduction) POVM element
\newcommand{\Vph}{\mathcal{V}^{\mathrm{ph}}} % physical outcome alphabet
\newcommand{\Tpub}{\mathcal{T}_{\mathrm{pub}}}% public-choice alphabet
\newcommand{\Uset}{\mathcal{U}}              % residual outcome alphabet
\newcommand{\idxm}{\mathsf{idx}}             % label-to-index map
\newcommand{\Asig}{A'}                       % transmitted (signal) system -- the REAL one
\newcommand{\Avirt}{A''}                     % virtual transmitted system (source descriptor of Xi_1)
\newcommand{\Abar}{\bar{A}}                  % Alice's measured register
\newcommand{\Ahat}{\hat{A}}                  % Alice's shield register
\newcommand{\pcNmsm}{\mathcal{N}_{\mathsf A}}  % Alice's measurement channel (see the RULES block)
\newcommand{\Phip}{\Phi^{+}}                 % maximally entangled vector label
\newcommand{\Smap}{\mathfrak{s}}             % the emitted sector map (fraktur: a compile-time object)

% --- receiver ------------------------------------------------------------------
\newcommand{\Modes}{\mathbb{M}}              % declared ordered mode set
\newcommand{\Inp}{\mathbb{I}}                % declared input modes
\newcommand{\InpC}{\mathbb{I}_{\Chan}}       % channel input modes
\newcommand{\InpS}{\mathbb{I}_{\Prb}}        % probe (source) input modes
\newcommand{\Env}{\mathbb{E}}                % dilation (environment) modes
\newcommand{\Port}[1]{\mathbb{P}_{#1}}       % mode set of detector port #1
\newcommand{\Emd}[1]{E_{#1}}                 % coordinate projector on modes #1
\newcommand{\Gmat}{G}                        % instrument (transfer) matrix
\newcommand{\Fock}{\mathfrak{F}}             % Fock / second-quantization functor
\newcommand{\Symn}[1]{\mathrm{Sym}^{#1}}     % symmetric tensor power
\newcommand{\normord}[1]{{:}\mspace{2mu}#1\mspace{2mu}{:}}   % normal ordering
\newcommand{\rk}{\operatorname{rank}}        % rank (physics.sty owns \rank)
\newcommand{\SCK}{\mathsf{SC}}               % single-click pattern set
\newcommand{\Recv}{\mathfrak{R}}             % receiver instrument tuple
\newcommand{\Cat}{\mathfrak{C}}              % component catalogue
\newcommand{\Grph}{\mathfrak{E}}             % declared element graph
\newcommand{\defic}{\mathbb{L}}              % loss deficiency operator
\newcommand{\Pvac}[1]{\Pi_{\vac}(#1)}        % joint no-click element on ports #1
\newcommand{\pmc}{p_{\mc}}                   % uniform high-photon mc weight
\newcommand{\Bnd}{\beta}                     % the scalar bound function beta(n)
\newcommand{\Rot}{R}                         % planar rotation matrix
\newcommand{\Wplate}{W}                      % retarder (wave-plate) matrix
\newcommand{\Bsp}{B}                         % beamsplitter core matrix

% --- statistics ----------------------------------------------------------------
\newcommand{\honest}{p^{\mathrm{hon}}}          % honest per-round distribution
\newcommand{\ecomp}{\varepsilon_{\mathrm{comp}}}% completeness budget
\newcommand{\simplex}[1]{\Delta(#1)}            % probability simplex
\newcommand{\Chatcon}{\Chatset_{\mathrm{con}}}  % constrained outcomes
\newcommand{\Ireg}[3]{\mathcal{I}_{#1}(#2,#3)}  % regularized incomplete beta
\newcommand{\Ncount}[1]{\mathsf{N}_{#1}}        % per-outcome count
\newcommand{\qlo}{\mathbf{q}^{-}}               % box lower bounds
\newcommand{\qhi}{\mathbf{q}^{+}}               % box upper bounds
 for it, so the
%  file is read from spec/modules/<id>/.
%
%  RULES
%    * A macro that spec/notation.md catalogues is a verbatim copy of it
%      (part A). Changing one is a notation change: bump spec/notation.md.
%    * Part B holds the module-local macros. A macro defined by two modules
%      with the same expansion appears once, under the first module (in the
%      order below) that defines it. Two macros with the same name and
%      DIFFERENT expansions are both kept under a module prefix, and the
%      claims write the prefixed names:
%        \Chain  -> \redChain (reductions, \mathfrak{C}: the reduction chain)
%                   \pcChain  (protocol-compiler, \mathfrak{R}: the two-segment
%                             chain; \Chsrc and \Chrec are built on it)
%        \Nmsm   -> \redNmsm  (reductions, \mathcal{N}: a POVM measurement
%                             channel)
%                   \pcNmsm   (\mathcal{N}_{\mathsf A}: Alice's measurement
%                             channel, written by
%                             transmitter/thm:certificate-sufficiency,
%                             transmitter/rem:xreg and
%                             reductions/lem:source-transfer)
%    * Deliberately omitted: a \provenance{tag}{details} macro. Provenance
%      is a header field of every claim file (spec/claim-format.md section
%      3); a claim body that calls such a macro fails to compile, which is
%      the intended signal.
%
%  REQUIRED PACKAGES, loaded by the assembler's preamble before this file:
%    amsmath, amssymb, amsfonts, mathtools, dsfont, braket, physics, xcolor
%    (braket before physics, so that physics wins for \ket, \bra, \tr, ...).
%
%  This file contains only definitions: no \usepackage, no document body.
% =============================================================================


% =============================================================================
%  PART A -- notation macros, the ones spec/notation.md catalogues, each a
%  verbatim copy of its catalogue row; grouped under numbered headings.
% =============================================================================

% -----------------------------------------------------------------------------
% 1. State spaces and operator sets (a "system label" -> spaces of operators)
% -----------------------------------------------------------------------------
\newcommand{\Hs}[1]{\mathcal{H}(#1)}        % Hilbert space of a system label
\newcommand{\Lin}[1]{\mathcal{L}(#1)}       % bounded linear operators on H(#1)
\newcommand{\Linn}[2]{\mathcal{L}(#1,#2)}   % linear maps H(#1) -> H(#2)
\newcommand{\Pos}[1]{\mathrm{Pos}(#1)}      % positive semidefinite operators on H(#1)
\newcommand{\Density}{\mathrm{D}}           % density operators: \Density(A) = D(A)
\newcommand{\Seq}[1]{S_{=}(#1)}             % normalized states (trace = 1)
\newcommand{\Sleq}[1]{S_{\leq}(#1)}         % subnormalized states (0 <= trace <= 1)
\newcommand{\id}{\mathds{1}}                % identity operator
\newcommand{\idon}[1]{\mathds{1}_{#1}}      % identity on a named system, e.g. \idon{C}

% DEPRECATED: applies \mathcal to both arguments itself, so it takes bare
% letters, not system labels. Prefer \Linn; do not use in new contracts.
\providecommand{\Lop}[2]{\mathcal{L}(\mathcal{#1}, \mathcal{#2})}


% -----------------------------------------------------------------------------
% 2. Composite systems
% -----------------------------------------------------------------------------
\newcommand{\dsum}{\oplus}                  % direct sum of system labels
% Tensor product: use \otimes directly.
% n-fold direct sum of an operator space: \Lin{A}^{\oplus n}


% -----------------------------------------------------------------------------
% 3. Events, conditioning and classical registers (for STATES)
%    Convention:  \wedge  -> subnormalized "and-event" component, e.g. \rho_{A\wedge x}
%                 |        -> normalized conditional state,        e.g. \rho_{A|x}
% (Write these directly in subscripts; no extra macros needed.)
% -----------------------------------------------------------------------------


% -----------------------------------------------------------------------------
% 4. Photon-number subspaces / direct-sum blocks
%    Convention:  operators carry the restriction "on top", in brackets;
%                 states carry it as a \wedge subscript (see Sec. 3).
% -----------------------------------------------------------------------------
\newcommand{\Pnum}[1]{P_{#1}}               % photon-number projector (system from context)
\newcommand{\Pnumon}[2]{P_{#1}^{#2}}        % with explicit system superscript, e.g. P_n^C
\newcommand{\blk}[2]{#1^{(#2)}}             % operator block, e.g. \blk{M}{\leqN} = M^{(<=N)}
\newcommand{\leqN}{\leq N}                  % use as \blk{M}{\leqN}
\newcommand{\gtN}{> N}                      % use as \blk{M}{\gtN}
\newcommand{\eqn}{= n}                      % use as \blk{M}{\eqn}


% -----------------------------------------------------------------------------
% 5. System labels used in the protocol
% -----------------------------------------------------------------------------
\newcommand{\Chan}{C}                       % quantum-channel (Eve) input to Bob
\newcommand{\Prb}{P}                        % probe input
\newcommand{\Flag}{{\mathrm{Flg}}}          % flag system (EXISTS ONLY AFTER squashing)
\newcommand{\Loss}{{\mathrm{Lss}}}          % loss system
\newcommand{\Vc}{{\mathrm{Vcm}}}            % vacuum system
\newcommand{\Nout}{N_{\mathrm{out}}}
\newcommand{\Nin}{N_{\mathrm{in}}}


% -----------------------------------------------------------------------------
% 6. Alphabets
% -----------------------------------------------------------------------------
\newcommand{\Yset}{\mathcal{Y}}             % Bob's (coarse-grained) detection alphabet
\newcommand{\Yraw}{\mathcal{Y}_{\mathrm{raw}}}% raw click patterns {0,1}^4
\newcommand{\Yscset}{\mathcal{Y}_{\singlec}}% single-click outcomes
\newcommand{\Ymcset}{\mathcal{Y}_{\mc}}     % multi-click outcomes
\newcommand{\Tset}{\mathcal{T}_{\mathrm{set}}}% private setting set (the probe settings in the optical instance)
\newcommand{\Xset}{\mathcal{X}}             % Alice's alphabet
\newcommand{\Dset}{\{H,V,+,-\}}             % polarization / detector labels
\newcommand{\pass}{\top}                    % pass sifting


% -----------------------------------------------------------------------------
% 7. Mode / outcome abbreviations
% -----------------------------------------------------------------------------
\newcommand{\on}{{\mathrm{on}}}
\newcommand{\off}{{\mathrm{off}}}
\newcommand{\dc}{{\mathrm{dc}}}             % dark count
\newcommand{\sq}{\mathrm{sq}}               % squashed
\newcommand{\mc}{\mathrm{mc}}               % multi-click
\newcommand{\nc}{\mathrm{nc}}               % no-click
\newcommand{\singlec}{\mathrm{sc}}          % single-click
\newcommand{\vac}{\mathrm{vac}}             % vacuum
\newcommand{\kout}{{k_\mathrm{out}}}
\newcommand{\kin}{{k_\mathrm{in}}}


% -----------------------------------------------------------------------------
% 8. Maps, channels and miscellaneous operators
% -----------------------------------------------------------------------------
\newcommand{\Prob}{\mathrm{Prob}}           % probability (use \Prob[\cdot] or \Prob(\cdot))
\newcommand{\lmin}{\lambda_{\min}}          % minimum eigenvalue
\newcommand{\Wflag}{W_{\Flag}}              % flag-space weight
\newcommand{\transpose}{\mathsf{T}}
\newcommand{\innerp}[2]{\left\langle #1 \vert #2 \right\rangle}
\renewcommand{\vec}[1]{\boldsymbol{#1}}     % NB: overrides the LaTeX arrow accent


% -----------------------------------------------------------------------------
% 9. Key-rate numerics
% -----------------------------------------------------------------------------
\newcommand{\cG}{\mathcal{G}}               % protocol map
\newcommand{\cZ}{\mathcal{Z}}               % pinching map
\newcommand{\Amap}{\mathcal{A}_{\mathrm{ann}}} % announcement + measurement map
\newcommand{\Amapnb}{\mathcal{A}_{\mathrm{ann}|\mathrm{nb}}}
\newcommand{\simpl}{}                       % empty suffix for simplified block maps
\newcommand{\acSet}{\mathcal{S}_{\mathrm{acc}}} % acceptance set
\newcommand{\testE}{\mathcal{E}_{\mathrm{test}}} % test-statistics map
\newcommand{\ecLeak}{\mathrm{leak}^{\mathrm{EC}}_{\mathrm{obs}}}
\newcommand{\pgen}{p_{\mathrm{gen}}}
\newcommand{\askeyr}{R_{\infty}}             % asymptotic key rate
\newcommand{\ppass}{p_{\mathrm{pass}}}       % sifting pass probability


% -----------------------------------------------------------------------------
% 10. Finite-size analysis / Renyi entropies
% -----------------------------------------------------------------------------

% --- Round types: test vs. key generation (register T, alphabet \Fset) --------
\newcommand{\tst}{\mathrm{test}}            % test round
\newcommand{\gen}{\mathrm{gen}}             % key-generation round
\newcommand{\Fset}{\mathcal{T}}             % round-type alphabet {test, gen} (not \Tset)
\newcommand{\Chat}{\widehat{C}}             % acceptance-test register
\newcommand{\Chatset}{\widehat{\mathcal{C}}}% its alphabet

% --- Basis / key extraction functions ----------------------------------------
\newcommand{\bA}{b_{A}}                     % Alice's basis-extraction function
\newcommand{\bB}{b_{B}}                     % Bob's basis-extraction function
\newcommand{\kB}{k_{B}}                     % Bob's key-extraction function

% --- Maps --------------------------------------------------------------------
\newcommand{\Mprot}{\mathcal{M}}            % single-round protocol map
\newcommand{\Phimap}{\boldsymbol{\Phi}}     % test-statistics map
\newcommand{\Pdeph}{\mathcal{P}}            % dephasing (pinching) map on \Blreg

% --- Isometries --------------------------------------------------------------
\newcommand{\Vmsm}{V_{\mathrm{msm}}}        % measurement isometry
\newcommand{\Vann}{V_{\mathrm{ann}}}        % announcement isometry
\newcommand{\VI}{V_{I}}                     % Eve's copy of the announcements
\newcommand{\VChat}{V_{\Chat}}              % acceptance-data isometry
\newcommand{\VR}{V_{R}}                     % sifting + key map -> raw key R
\newcommand{\VZ}{V_{Z}}                     % raw key -> key register Z
\newcommand{\Wprot}{W}                      % full protocol isometry

% --- Renyi entropies ---------------------------------------------------------
% SECURITY-CRITICAL: see spec/notation.md §4. Never change an arrow to
% make an inequality work out; never convert silently between the sandwiched
% (tilde) and Petz (overline) families.
\newcommand{\Hsandup}{\widetilde{H}^{\uparrow}_{\alpha}}   % sandwiched, up
\newcommand{\Hsanddown}{\widetilde{H}^{\downarrow}_{\alpha}}% sandwiched, down
\newcommand{\Hpetzup}{\overline{H}^{\uparrow}_{\alpha}}    % Petz, up
\newcommand{\Hpetzdown}{\overline{H}^{\downarrow}_{\alpha}}% Petz, down
\newcommand{\hup}{h^{\uparrow}_{\alpha}}                   % single-round rate
\newcommand{\hqkd}{h^{\uparrow,\mathrm{QKD}}_{\alpha}}     % single-round QKD rate
\newcommand{\hlow}{\underline{h}_{\alpha}}                 % lower bound on \hqkd
% --- spec/notation.md 3.5 and 4.5 -------------------------------------------
% \hlb  is the first-order linearization lower bound on \hlow: the headline
%       key rate (the `first-order` tier of foundations/def:tiers). NOT \hlow
%       itself.
% \tnum is the instance's one declared numerical tolerance. Its
%       near-collision with the declared marginal \tau_A is disclosed in
%       spec/notation.md 3.5; bare \tau is reserved and denotes neither.
\newcommand{\hlb}{\underline{h}^{\mathrm{FO}}_{\alpha}}    % first-order lower bound on \hlow
\newcommand{\tnum}{\tau_{\mathrm{num}}}                    % declared numerical tolerance


% =============================================================================
%  PART B -- module-local macros, grouped by module. Deprecated aliases,
%  defined so that a text written with them compiles, used in no new contract: \Lop,
%  \testE, \Pdeph, \pass, \ecLeak (spec/notation.md section 3.9).
% =============================================================================

% --- meat-core -----------------------------------------------------------------
\newcommand{\SecSet}{\mathbb{S}}            % sector index set of Q
\newcommand{\Qsys}{Q}                       % declared sector-carrying measured system
\newcommand{\Vset}{\mathcal{V}}             % declared joint outcome alphabet of the Q-POVM
\newcommand{\CAset}{\mathcal{C}_{A}}        % Alice announcement alphabet
\newcommand{\CBset}{\mathcal{C}_{B}}        % Bob announcement alphabet
\newcommand{\Party}{\mathfrak{P}}           % declared party partition
\newcommand{\Oset}[1]{\mathcal{O}_{#1}}     % readout alphabet of party #1
\newcommand{\rd}[1]{\pi_{#1}}               % readout map of party #1
\newcommand{\Sadm}{\mathcal{S}_{A}}         % admissible Alice subspace
\newcommand{\nomarg}{\mathsf{none}}         % "no trusted marginal declared"
\newcommand{\Zset}{\mathcal{Z}_{\perp}}     % raw-key alphabet incl. \perp
\newcommand{\Zval}{\mathcal{Z}}             % key alphabet without \perp
\newcommand{\Etil}{\widetilde{E}}           % single-round purifying register
\newcommand{\Ehat}{\widehat{E}}             % Eve's n-round register
\newcommand{\OmAT}{\Omega_{\mathrm{AT}}}    % acceptance event
\newcommand{\epa}{\varepsilon_{\mathrm{PA}}}
\newcommand{\eev}{\varepsilon_{\mathrm{EV}}}
\newcommand{\lEC}{\lambda_{\mathrm{EC}}}
\newcommand{\Decl}{\mathfrak{D}}            % protocol declaration
\newcommand{\CList}{\mathfrak{L}}           % declared constraint list
\newcommand{\PinS}{\mathcal{P}_{\SecSet}}   % sector pinching on Q
\newcommand{\Feas}{\mathcal{W}}             % feasible set of block states
\newcommand{\Chkey}{\Chatset_{\mathrm{key}}}% key-live acceptance values
\newcommand{\Kann}{\mathcal{K}}             % announcement-class alphabet
\newcommand{\hold}{\mathsf{h}}              % declared key holder (A or B)
\newcommand{\sfD}{D}                        % classical KL divergence (base 2)
\newcommand{\Dtwo}{D}                       % quantum relative entropy (base 2)
\newcommand{\Dsand}{\widetilde{D}_{\alpha}} % sandwiched Renyi divergence
\newcommand{\idmap}{\mathrm{id}}            % identity channel
\newcommand{\Aev}{\mathcal{A}}              % Eve's per-round attack channel
\newcommand{\AttSet}{\mathfrak{A}}          % licensed attack-channel set
\newcommand{\Chacc}{\Chatset_{\mathrm{acc}}}% acceptance support of \acSet
\newcommand{\epsstar}{\varepsilon^{\star}}      % distance to an exact POVM
\newcommand{\dfH}{\Delta_{\mathrm{herm}}}       % Hermiticity defect
\newcommand{\dfP}{\Delta_{\mathrm{pos}}}        % negative part
\newcommand{\dfC}{\Delta_{\mathrm{cpl}}}        % completeness defect
\newcommand{\lmoq}{g}                       % acceptance-side gradient block
\newcommand{\lmoJ}{G}                       % feasible-side gradient block
\newcommand{\bacc}{\beta_{\mathrm{acc}}}    % certified acceptance-side minimum
\newcommand{\Itset}{\mathfrak{I}}           % finite index set of iterates (C22)

% --- reductions ----------------------------------------------------------------
\newcommand{\Bsys}{B}                        % physical (pre-squashing) system
\newcommand{\Tcl}[1]{\mathcal{T}(#1)}        % trace-class operators on H(#1)
\newcommand{\Ins}{\mathcal{I}}               % generalized-instrument branch map
\newcommand{\redNmsm}{\mathcal{N}}              % POVM measurement channel
\newcommand{\Eprot}{\mathcal{E}}             % n-round protocol channel
\newcommand{\Grnd}{\mathsf{G}}               % one round map of a protocol
\newcommand{\Kfin}{\mathsf{K}}               % end-of-protocol processing
\newcommand{\Crit}{\mathsf{sec}}             % value of a security criterion
\newcommand{\Pflag}{\Pi_{\Flag}}             % flag-sector projector on Q
\newcommand{\Pleq}{\Pi_{\leqN}}              % preserved-block projector
\newcommand{\Pgt}{\Pi_{\gtN}}                % above-cutoff projector
\newcommand{\Sph}{S^{\mathrm{ph}}}           % family of physical inputs
\newcommand{\Srch}{S^{\mathrm{rch}}}         % reachable squashed states
\newcommand{\Ssq}{S^{\sq}}                   % squashed admissible set
\newcommand{\Ffss}{F^{\mathrm{FSS}}}         % flag-state-squasher target
\newcommand{\Fwp}{F^{\mathrm{WP}}}           % multi-bound WPFSS target
\newcommand{\Mwp}{M^{\mathrm{WP}}}           % renormalized above-cutoff POVM
\newcommand{\Mdec}{{M'}^{\mathrm{ph}}}       % declared physical family on B
\newcommand{\Ppvm}{P^{\Flag}}                % flag PVM
\newcommand{\Lmp}{{\Lambda_{\mathrm{mp}}}}   % measure-and-prepare channel
\newcommand{\Lwp}{{\Lambda^{\mathrm{WP}}}}   % multi-bound WPFSS channel
\newcommand{\Lfss}{{\Lambda^{\mathrm{FSS}}}} % flag-state-squasher channel
\newcommand{\lmax}{\lambda_{\max}}           % maximum eigenvalue
\newcommand{\Lnoise}{\widetilde{\Lambda}}    % flag-noise channel
\newcommand{\cg}{\varpi}                     % declared coarse-graining map
\newcommand{\MCK}{\mathsf{MC}}               % multiclick raw-pattern set
\newcommand{\redChain}{\mathfrak{C}}            % reduction chain
\newcommand{\Lnk}{\ell}                      % a single typed link
\newcommand{\dsq}{\preceq}                      % the squashing relation
\newcommand{\opn}[1]{\left\lVert #1\right\rVert}           % operator (spectral) norm
\newcommand{\mxn}[1]{\left\lVert #1\right\rVert_{\max}}    % entrywise max modulus
\newcommand{\Fld}{\mathbb{F}}                   % exact ordered substrate field
\newcommand{\Fldi}{\mathbb{F}_{i}}              % its Gaussian extension F + iF
\newcommand{\lat}[1]{\big\lfloor #1\big\rfloor_{\kappa}} % round DOWN onto the lattice

% --- protocol-compiler ---------------------------------------------------------
\newcommand{\Spec}{\mathfrak{S}}             % protocol specification (compiler input)
\newcommand{\pcChain}{\mathfrak{R}}            % reduction chain (two typed segments)
\newcommand{\Chsrc}{\pcChain_{\mathrm{src}}}   % its Alice-side source-map segment
\newcommand{\Chrec}{\pcChain_{\mathrm{rec}}}   % its receiver-side squashing segment
\newcommand{\rvirt}{\rho^{\mathrm{v}}}       % virtual (PRE-source-map) signal state
\newcommand{\Fph}{F^{\mathrm{ph}}}           % physical (pre-reduction) POVM element
\newcommand{\Vph}{\mathcal{V}^{\mathrm{ph}}} % physical outcome alphabet
\newcommand{\Tpub}{\mathcal{T}_{\mathrm{pub}}}% public-choice alphabet
\newcommand{\Uset}{\mathcal{U}}              % residual outcome alphabet
\newcommand{\idxm}{\mathsf{idx}}             % label-to-index map
\newcommand{\Asig}{A'}                       % transmitted (signal) system -- the REAL one
\newcommand{\Avirt}{A''}                     % virtual transmitted system (source descriptor of Xi_1)
\newcommand{\Abar}{\bar{A}}                  % Alice's measured register
\newcommand{\Ahat}{\hat{A}}                  % Alice's shield register
\newcommand{\pcNmsm}{\mathcal{N}_{\mathsf A}}  % Alice's measurement channel (see the RULES block)
\newcommand{\Phip}{\Phi^{+}}                 % maximally entangled vector label
\newcommand{\Smap}{\mathfrak{s}}             % the emitted sector map (fraktur: a compile-time object)

% --- receiver ------------------------------------------------------------------
\newcommand{\Modes}{\mathbb{M}}              % declared ordered mode set
\newcommand{\Inp}{\mathbb{I}}                % declared input modes
\newcommand{\InpC}{\mathbb{I}_{\Chan}}       % channel input modes
\newcommand{\InpS}{\mathbb{I}_{\Prb}}        % probe (source) input modes
\newcommand{\Env}{\mathbb{E}}                % dilation (environment) modes
\newcommand{\Port}[1]{\mathbb{P}_{#1}}       % mode set of detector port #1
\newcommand{\Emd}[1]{E_{#1}}                 % coordinate projector on modes #1
\newcommand{\Gmat}{G}                        % instrument (transfer) matrix
\newcommand{\Fock}{\mathfrak{F}}             % Fock / second-quantization functor
\newcommand{\Symn}[1]{\mathrm{Sym}^{#1}}     % symmetric tensor power
\newcommand{\normord}[1]{{:}\mspace{2mu}#1\mspace{2mu}{:}}   % normal ordering
\newcommand{\rk}{\operatorname{rank}}        % rank (physics.sty owns \rank)
\newcommand{\SCK}{\mathsf{SC}}               % single-click pattern set
\newcommand{\Recv}{\mathfrak{R}}             % receiver instrument tuple
\newcommand{\Cat}{\mathfrak{C}}              % component catalogue
\newcommand{\Grph}{\mathfrak{E}}             % declared element graph
\newcommand{\defic}{\mathbb{L}}              % loss deficiency operator
\newcommand{\Pvac}[1]{\Pi_{\vac}(#1)}        % joint no-click element on ports #1
\newcommand{\pmc}{p_{\mc}}                   % uniform high-photon mc weight
\newcommand{\Bnd}{\beta}                     % the scalar bound function beta(n)
\newcommand{\Rot}{R}                         % planar rotation matrix
\newcommand{\Wplate}{W}                      % retarder (wave-plate) matrix
\newcommand{\Bsp}{B}                         % beamsplitter core matrix

% --- statistics ----------------------------------------------------------------
\newcommand{\honest}{p^{\mathrm{hon}}}          % honest per-round distribution
\newcommand{\ecomp}{\varepsilon_{\mathrm{comp}}}% completeness budget
\newcommand{\simplex}[1]{\Delta(#1)}            % probability simplex
\newcommand{\Chatcon}{\Chatset_{\mathrm{con}}}  % constrained outcomes
\newcommand{\Ireg}[3]{\mathcal{I}_{#1}(#2,#3)}  % regularized incomplete beta
\newcommand{\Ncount}[1]{\mathsf{N}_{#1}}        % per-outcome count
\newcommand{\qlo}{\mathbf{q}^{-}}               % box lower bounds
\newcommand{\qhi}{\mathbf{q}^{+}}               % box upper bounds
 for it, so the
%  file is read from spec/modules/<id>/.
%
%  RULES
%    * A macro that spec/notation.md catalogues is a verbatim copy of it
%      (part A). Changing one is a notation change: bump spec/notation.md.
%    * Part B holds the module-local macros. A macro defined by two modules
%      with the same expansion appears once, under the first module (in the
%      order below) that defines it. Two macros with the same name and
%      DIFFERENT expansions are both kept under a module prefix, and the
%      claims write the prefixed names:
%        \Chain  -> \redChain (reductions, \mathfrak{C}: the reduction chain)
%                   \pcChain  (protocol-compiler, \mathfrak{R}: the two-segment
%                             chain; \Chsrc and \Chrec are built on it)
%        \Nmsm   -> \redNmsm  (reductions, \mathcal{N}: a POVM measurement
%                             channel)
%                   \pcNmsm   (\mathcal{N}_{\mathsf A}: Alice's measurement
%                             channel, written by
%                             transmitter/thm:certificate-sufficiency,
%                             transmitter/rem:xreg and
%                             reductions/lem:source-transfer)
%    * Deliberately omitted: a \provenance{tag}{details} macro. Provenance
%      is a header field of every claim file (spec/claim-format.md section
%      3); a claim body that calls such a macro fails to compile, which is
%      the intended signal.
%
%  REQUIRED PACKAGES, loaded by the assembler's preamble before this file:
%    amsmath, amssymb, amsfonts, mathtools, dsfont, braket, physics, xcolor
%    (braket before physics, so that physics wins for \ket, \bra, \tr, ...).
%
%  This file contains only definitions: no \usepackage, no document body.
% =============================================================================


% =============================================================================
%  PART A -- notation macros, the ones spec/notation.md catalogues, each a
%  verbatim copy of its catalogue row; grouped under numbered headings.
% =============================================================================

% -----------------------------------------------------------------------------
% 1. State spaces and operator sets (a "system label" -> spaces of operators)
% -----------------------------------------------------------------------------
\newcommand{\Hs}[1]{\mathcal{H}(#1)}        % Hilbert space of a system label
\newcommand{\Lin}[1]{\mathcal{L}(#1)}       % bounded linear operators on H(#1)
\newcommand{\Linn}[2]{\mathcal{L}(#1,#2)}   % linear maps H(#1) -> H(#2)
\newcommand{\Pos}[1]{\mathrm{Pos}(#1)}      % positive semidefinite operators on H(#1)
\newcommand{\Density}{\mathrm{D}}           % density operators: \Density(A) = D(A)
\newcommand{\Seq}[1]{S_{=}(#1)}             % normalized states (trace = 1)
\newcommand{\Sleq}[1]{S_{\leq}(#1)}         % subnormalized states (0 <= trace <= 1)
\newcommand{\id}{\mathds{1}}                % identity operator
\newcommand{\idon}[1]{\mathds{1}_{#1}}      % identity on a named system, e.g. \idon{C}

% DEPRECATED: applies \mathcal to both arguments itself, so it takes bare
% letters, not system labels. Prefer \Linn; do not use in new contracts.
\providecommand{\Lop}[2]{\mathcal{L}(\mathcal{#1}, \mathcal{#2})}


% -----------------------------------------------------------------------------
% 2. Composite systems
% -----------------------------------------------------------------------------
\newcommand{\dsum}{\oplus}                  % direct sum of system labels
% Tensor product: use \otimes directly.
% n-fold direct sum of an operator space: \Lin{A}^{\oplus n}


% -----------------------------------------------------------------------------
% 3. Events, conditioning and classical registers (for STATES)
%    Convention:  \wedge  -> subnormalized "and-event" component, e.g. \rho_{A\wedge x}
%                 |        -> normalized conditional state,        e.g. \rho_{A|x}
% (Write these directly in subscripts; no extra macros needed.)
% -----------------------------------------------------------------------------


% -----------------------------------------------------------------------------
% 4. Photon-number subspaces / direct-sum blocks
%    Convention:  operators carry the restriction "on top", in brackets;
%                 states carry it as a \wedge subscript (see Sec. 3).
% -----------------------------------------------------------------------------
\newcommand{\Pnum}[1]{P_{#1}}               % photon-number projector (system from context)
\newcommand{\Pnumon}[2]{P_{#1}^{#2}}        % with explicit system superscript, e.g. P_n^C
\newcommand{\blk}[2]{#1^{(#2)}}             % operator block, e.g. \blk{M}{\leqN} = M^{(<=N)}
\newcommand{\leqN}{\leq N}                  % use as \blk{M}{\leqN}
\newcommand{\gtN}{> N}                      % use as \blk{M}{\gtN}
\newcommand{\eqn}{= n}                      % use as \blk{M}{\eqn}


% -----------------------------------------------------------------------------
% 5. System labels used in the protocol
% -----------------------------------------------------------------------------
\newcommand{\Chan}{C}                       % quantum-channel (Eve) input to Bob
\newcommand{\Prb}{P}                        % probe input
\newcommand{\Flag}{{\mathrm{Flg}}}          % flag system (EXISTS ONLY AFTER squashing)
\newcommand{\Loss}{{\mathrm{Lss}}}          % loss system
\newcommand{\Vc}{{\mathrm{Vcm}}}            % vacuum system
\newcommand{\Nout}{N_{\mathrm{out}}}
\newcommand{\Nin}{N_{\mathrm{in}}}


% -----------------------------------------------------------------------------
% 6. Alphabets
% -----------------------------------------------------------------------------
\newcommand{\Yset}{\mathcal{Y}}             % Bob's (coarse-grained) detection alphabet
\newcommand{\Yraw}{\mathcal{Y}_{\mathrm{raw}}}% raw click patterns {0,1}^4
\newcommand{\Yscset}{\mathcal{Y}_{\singlec}}% single-click outcomes
\newcommand{\Ymcset}{\mathcal{Y}_{\mc}}     % multi-click outcomes
\newcommand{\Tset}{\mathcal{T}_{\mathrm{set}}}% private setting set (the probe settings in the optical instance)
\newcommand{\Xset}{\mathcal{X}}             % Alice's alphabet
\newcommand{\Dset}{\{H,V,+,-\}}             % polarization / detector labels
\newcommand{\pass}{\top}                    % pass sifting


% -----------------------------------------------------------------------------
% 7. Mode / outcome abbreviations
% -----------------------------------------------------------------------------
\newcommand{\on}{{\mathrm{on}}}
\newcommand{\off}{{\mathrm{off}}}
\newcommand{\dc}{{\mathrm{dc}}}             % dark count
\newcommand{\sq}{\mathrm{sq}}               % squashed
\newcommand{\mc}{\mathrm{mc}}               % multi-click
\newcommand{\nc}{\mathrm{nc}}               % no-click
\newcommand{\singlec}{\mathrm{sc}}          % single-click
\newcommand{\vac}{\mathrm{vac}}             % vacuum
\newcommand{\kout}{{k_\mathrm{out}}}
\newcommand{\kin}{{k_\mathrm{in}}}


% -----------------------------------------------------------------------------
% 8. Maps, channels and miscellaneous operators
% -----------------------------------------------------------------------------
\newcommand{\Prob}{\mathrm{Prob}}           % probability (use \Prob[\cdot] or \Prob(\cdot))
\newcommand{\lmin}{\lambda_{\min}}          % minimum eigenvalue
\newcommand{\Wflag}{W_{\Flag}}              % flag-space weight
\newcommand{\transpose}{\mathsf{T}}
\newcommand{\innerp}[2]{\left\langle #1 \vert #2 \right\rangle}
\renewcommand{\vec}[1]{\boldsymbol{#1}}     % NB: overrides the LaTeX arrow accent


% -----------------------------------------------------------------------------
% 9. Key-rate numerics
% -----------------------------------------------------------------------------
\newcommand{\cG}{\mathcal{G}}               % protocol map
\newcommand{\cZ}{\mathcal{Z}}               % pinching map
\newcommand{\Amap}{\mathcal{A}_{\mathrm{ann}}} % announcement + measurement map
\newcommand{\Amapnb}{\mathcal{A}_{\mathrm{ann}|\mathrm{nb}}}
\newcommand{\simpl}{}                       % empty suffix for simplified block maps
\newcommand{\acSet}{\mathcal{S}_{\mathrm{acc}}} % acceptance set
\newcommand{\testE}{\mathcal{E}_{\mathrm{test}}} % test-statistics map
\newcommand{\ecLeak}{\mathrm{leak}^{\mathrm{EC}}_{\mathrm{obs}}}
\newcommand{\pgen}{p_{\mathrm{gen}}}
\newcommand{\askeyr}{R_{\infty}}             % asymptotic key rate
\newcommand{\ppass}{p_{\mathrm{pass}}}       % sifting pass probability


% -----------------------------------------------------------------------------
% 10. Finite-size analysis / Renyi entropies
% -----------------------------------------------------------------------------

% --- Round types: test vs. key generation (register T, alphabet \Fset) --------
\newcommand{\tst}{\mathrm{test}}            % test round
\newcommand{\gen}{\mathrm{gen}}             % key-generation round
\newcommand{\Fset}{\mathcal{T}}             % round-type alphabet {test, gen} (not \Tset)
\newcommand{\Chat}{\widehat{C}}             % acceptance-test register
\newcommand{\Chatset}{\widehat{\mathcal{C}}}% its alphabet

% --- Basis / key extraction functions ----------------------------------------
\newcommand{\bA}{b_{A}}                     % Alice's basis-extraction function
\newcommand{\bB}{b_{B}}                     % Bob's basis-extraction function
\newcommand{\kB}{k_{B}}                     % Bob's key-extraction function

% --- Maps --------------------------------------------------------------------
\newcommand{\Mprot}{\mathcal{M}}            % single-round protocol map
\newcommand{\Phimap}{\boldsymbol{\Phi}}     % test-statistics map
\newcommand{\Pdeph}{\mathcal{P}}            % dephasing (pinching) map on \Blreg

% --- Isometries --------------------------------------------------------------
\newcommand{\Vmsm}{V_{\mathrm{msm}}}        % measurement isometry
\newcommand{\Vann}{V_{\mathrm{ann}}}        % announcement isometry
\newcommand{\VI}{V_{I}}                     % Eve's copy of the announcements
\newcommand{\VChat}{V_{\Chat}}              % acceptance-data isometry
\newcommand{\VR}{V_{R}}                     % sifting + key map -> raw key R
\newcommand{\VZ}{V_{Z}}                     % raw key -> key register Z
\newcommand{\Wprot}{W}                      % full protocol isometry

% --- Renyi entropies ---------------------------------------------------------
% SECURITY-CRITICAL: see spec/notation.md §4. Never change an arrow to
% make an inequality work out; never convert silently between the sandwiched
% (tilde) and Petz (overline) families.
\newcommand{\Hsandup}{\widetilde{H}^{\uparrow}_{\alpha}}   % sandwiched, up
\newcommand{\Hsanddown}{\widetilde{H}^{\downarrow}_{\alpha}}% sandwiched, down
\newcommand{\Hpetzup}{\overline{H}^{\uparrow}_{\alpha}}    % Petz, up
\newcommand{\Hpetzdown}{\overline{H}^{\downarrow}_{\alpha}}% Petz, down
\newcommand{\hup}{h^{\uparrow}_{\alpha}}                   % single-round rate
\newcommand{\hqkd}{h^{\uparrow,\mathrm{QKD}}_{\alpha}}     % single-round QKD rate
\newcommand{\hlow}{\underline{h}_{\alpha}}                 % lower bound on \hqkd
% --- spec/notation.md 3.5 and 4.5 -------------------------------------------
% \hlb  is the first-order linearization lower bound on \hlow: the headline
%       key rate (the `first-order` tier of foundations/def:tiers). NOT \hlow
%       itself.
% \tnum is the instance's one declared numerical tolerance. Its
%       near-collision with the declared marginal \tau_A is disclosed in
%       spec/notation.md 3.5; bare \tau is reserved and denotes neither.
\newcommand{\hlb}{\underline{h}^{\mathrm{FO}}_{\alpha}}    % first-order lower bound on \hlow
\newcommand{\tnum}{\tau_{\mathrm{num}}}                    % declared numerical tolerance


% =============================================================================
%  PART B -- module-local macros, grouped by module. Deprecated aliases,
%  defined so that a text written with them compiles, used in no new contract: \Lop,
%  \testE, \Pdeph, \pass, \ecLeak (spec/notation.md section 3.9).
% =============================================================================

% --- meat-core -----------------------------------------------------------------
\newcommand{\SecSet}{\mathbb{S}}            % sector index set of Q
\newcommand{\Qsys}{Q}                       % declared sector-carrying measured system
\newcommand{\Vset}{\mathcal{V}}             % declared joint outcome alphabet of the Q-POVM
\newcommand{\CAset}{\mathcal{C}_{A}}        % Alice announcement alphabet
\newcommand{\CBset}{\mathcal{C}_{B}}        % Bob announcement alphabet
\newcommand{\Party}{\mathfrak{P}}           % declared party partition
\newcommand{\Oset}[1]{\mathcal{O}_{#1}}     % readout alphabet of party #1
\newcommand{\rd}[1]{\pi_{#1}}               % readout map of party #1
\newcommand{\Sadm}{\mathcal{S}_{A}}         % admissible Alice subspace
\newcommand{\nomarg}{\mathsf{none}}         % "no trusted marginal declared"
\newcommand{\Zset}{\mathcal{Z}_{\perp}}     % raw-key alphabet incl. \perp
\newcommand{\Zval}{\mathcal{Z}}             % key alphabet without \perp
\newcommand{\Etil}{\widetilde{E}}           % single-round purifying register
\newcommand{\Ehat}{\widehat{E}}             % Eve's n-round register
\newcommand{\OmAT}{\Omega_{\mathrm{AT}}}    % acceptance event
\newcommand{\epa}{\varepsilon_{\mathrm{PA}}}
\newcommand{\eev}{\varepsilon_{\mathrm{EV}}}
\newcommand{\lEC}{\lambda_{\mathrm{EC}}}
\newcommand{\Decl}{\mathfrak{D}}            % protocol declaration
\newcommand{\CList}{\mathfrak{L}}           % declared constraint list
\newcommand{\PinS}{\mathcal{P}_{\SecSet}}   % sector pinching on Q
\newcommand{\Feas}{\mathcal{W}}             % feasible set of block states
\newcommand{\Chkey}{\Chatset_{\mathrm{key}}}% key-live acceptance values
\newcommand{\Kann}{\mathcal{K}}             % announcement-class alphabet
\newcommand{\hold}{\mathsf{h}}              % declared key holder (A or B)
\newcommand{\sfD}{D}                        % classical KL divergence (base 2)
\newcommand{\Dtwo}{D}                       % quantum relative entropy (base 2)
\newcommand{\Dsand}{\widetilde{D}_{\alpha}} % sandwiched Renyi divergence
\newcommand{\idmap}{\mathrm{id}}            % identity channel
\newcommand{\Aev}{\mathcal{A}}              % Eve's per-round attack channel
\newcommand{\AttSet}{\mathfrak{A}}          % licensed attack-channel set
\newcommand{\Chacc}{\Chatset_{\mathrm{acc}}}% acceptance support of \acSet
\newcommand{\epsstar}{\varepsilon^{\star}}      % distance to an exact POVM
\newcommand{\dfH}{\Delta_{\mathrm{herm}}}       % Hermiticity defect
\newcommand{\dfP}{\Delta_{\mathrm{pos}}}        % negative part
\newcommand{\dfC}{\Delta_{\mathrm{cpl}}}        % completeness defect
\newcommand{\lmoq}{g}                       % acceptance-side gradient block
\newcommand{\lmoJ}{G}                       % feasible-side gradient block
\newcommand{\bacc}{\beta_{\mathrm{acc}}}    % certified acceptance-side minimum
\newcommand{\Itset}{\mathfrak{I}}           % finite index set of iterates (C22)

% --- reductions ----------------------------------------------------------------
\newcommand{\Bsys}{B}                        % physical (pre-squashing) system
\newcommand{\Tcl}[1]{\mathcal{T}(#1)}        % trace-class operators on H(#1)
\newcommand{\Ins}{\mathcal{I}}               % generalized-instrument branch map
\newcommand{\redNmsm}{\mathcal{N}}              % POVM measurement channel
\newcommand{\Eprot}{\mathcal{E}}             % n-round protocol channel
\newcommand{\Grnd}{\mathsf{G}}               % one round map of a protocol
\newcommand{\Kfin}{\mathsf{K}}               % end-of-protocol processing
\newcommand{\Crit}{\mathsf{sec}}             % value of a security criterion
\newcommand{\Pflag}{\Pi_{\Flag}}             % flag-sector projector on Q
\newcommand{\Pleq}{\Pi_{\leqN}}              % preserved-block projector
\newcommand{\Pgt}{\Pi_{\gtN}}                % above-cutoff projector
\newcommand{\Sph}{S^{\mathrm{ph}}}           % family of physical inputs
\newcommand{\Srch}{S^{\mathrm{rch}}}         % reachable squashed states
\newcommand{\Ssq}{S^{\sq}}                   % squashed admissible set
\newcommand{\Ffss}{F^{\mathrm{FSS}}}         % flag-state-squasher target
\newcommand{\Fwp}{F^{\mathrm{WP}}}           % multi-bound WPFSS target
\newcommand{\Mwp}{M^{\mathrm{WP}}}           % renormalized above-cutoff POVM
\newcommand{\Mdec}{{M'}^{\mathrm{ph}}}       % declared physical family on B
\newcommand{\Ppvm}{P^{\Flag}}                % flag PVM
\newcommand{\Lmp}{{\Lambda_{\mathrm{mp}}}}   % measure-and-prepare channel
\newcommand{\Lwp}{{\Lambda^{\mathrm{WP}}}}   % multi-bound WPFSS channel
\newcommand{\Lfss}{{\Lambda^{\mathrm{FSS}}}} % flag-state-squasher channel
\newcommand{\lmax}{\lambda_{\max}}           % maximum eigenvalue
\newcommand{\Lnoise}{\widetilde{\Lambda}}    % flag-noise channel
\newcommand{\cg}{\varpi}                     % declared coarse-graining map
\newcommand{\MCK}{\mathsf{MC}}               % multiclick raw-pattern set
\newcommand{\redChain}{\mathfrak{C}}            % reduction chain
\newcommand{\Lnk}{\ell}                      % a single typed link
\newcommand{\dsq}{\preceq}                      % the squashing relation
\newcommand{\opn}[1]{\left\lVert #1\right\rVert}           % operator (spectral) norm
\newcommand{\mxn}[1]{\left\lVert #1\right\rVert_{\max}}    % entrywise max modulus
\newcommand{\Fld}{\mathbb{F}}                   % exact ordered substrate field
\newcommand{\Fldi}{\mathbb{F}_{i}}              % its Gaussian extension F + iF
\newcommand{\lat}[1]{\big\lfloor #1\big\rfloor_{\kappa}} % round DOWN onto the lattice

% --- protocol-compiler ---------------------------------------------------------
\newcommand{\Spec}{\mathfrak{S}}             % protocol specification (compiler input)
\newcommand{\pcChain}{\mathfrak{R}}            % reduction chain (two typed segments)
\newcommand{\Chsrc}{\pcChain_{\mathrm{src}}}   % its Alice-side source-map segment
\newcommand{\Chrec}{\pcChain_{\mathrm{rec}}}   % its receiver-side squashing segment
\newcommand{\rvirt}{\rho^{\mathrm{v}}}       % virtual (PRE-source-map) signal state
\newcommand{\Fph}{F^{\mathrm{ph}}}           % physical (pre-reduction) POVM element
\newcommand{\Vph}{\mathcal{V}^{\mathrm{ph}}} % physical outcome alphabet
\newcommand{\Tpub}{\mathcal{T}_{\mathrm{pub}}}% public-choice alphabet
\newcommand{\Uset}{\mathcal{U}}              % residual outcome alphabet
\newcommand{\idxm}{\mathsf{idx}}             % label-to-index map
\newcommand{\Asig}{A'}                       % transmitted (signal) system -- the REAL one
\newcommand{\Avirt}{A''}                     % virtual transmitted system (source descriptor of Xi_1)
\newcommand{\Abar}{\bar{A}}                  % Alice's measured register
\newcommand{\Ahat}{\hat{A}}                  % Alice's shield register
\newcommand{\pcNmsm}{\mathcal{N}_{\mathsf A}}  % Alice's measurement channel (see the RULES block)
\newcommand{\Phip}{\Phi^{+}}                 % maximally entangled vector label
\newcommand{\Smap}{\mathfrak{s}}             % the emitted sector map (fraktur: a compile-time object)

% --- receiver ------------------------------------------------------------------
\newcommand{\Modes}{\mathbb{M}}              % declared ordered mode set
\newcommand{\Inp}{\mathbb{I}}                % declared input modes
\newcommand{\InpC}{\mathbb{I}_{\Chan}}       % channel input modes
\newcommand{\InpS}{\mathbb{I}_{\Prb}}        % probe (source) input modes
\newcommand{\Env}{\mathbb{E}}                % dilation (environment) modes
\newcommand{\Port}[1]{\mathbb{P}_{#1}}       % mode set of detector port #1
\newcommand{\Emd}[1]{E_{#1}}                 % coordinate projector on modes #1
\newcommand{\Gmat}{G}                        % instrument (transfer) matrix
\newcommand{\Fock}{\mathfrak{F}}             % Fock / second-quantization functor
\newcommand{\Symn}[1]{\mathrm{Sym}^{#1}}     % symmetric tensor power
\newcommand{\normord}[1]{{:}\mspace{2mu}#1\mspace{2mu}{:}}   % normal ordering
\newcommand{\rk}{\operatorname{rank}}        % rank (physics.sty owns \rank)
\newcommand{\SCK}{\mathsf{SC}}               % single-click pattern set
\newcommand{\Recv}{\mathfrak{R}}             % receiver instrument tuple
\newcommand{\Cat}{\mathfrak{C}}              % component catalogue
\newcommand{\Grph}{\mathfrak{E}}             % declared element graph
\newcommand{\defic}{\mathbb{L}}              % loss deficiency operator
\newcommand{\Pvac}[1]{\Pi_{\vac}(#1)}        % joint no-click element on ports #1
\newcommand{\pmc}{p_{\mc}}                   % uniform high-photon mc weight
\newcommand{\Bnd}{\beta}                     % the scalar bound function beta(n)
\newcommand{\Rot}{R}                         % planar rotation matrix
\newcommand{\Wplate}{W}                      % retarder (wave-plate) matrix
\newcommand{\Bsp}{B}                         % beamsplitter core matrix

% --- statistics ----------------------------------------------------------------
\newcommand{\honest}{p^{\mathrm{hon}}}          % honest per-round distribution
\newcommand{\ecomp}{\varepsilon_{\mathrm{comp}}}% completeness budget
\newcommand{\simplex}[1]{\Delta(#1)}            % probability simplex
\newcommand{\Chatcon}{\Chatset_{\mathrm{con}}}  % constrained outcomes
\newcommand{\Ireg}[3]{\mathcal{I}_{#1}(#2,#3)}  % regularized incomplete beta
\newcommand{\Ncount}[1]{\mathsf{N}_{#1}}        % per-outcome count
\newcommand{\qlo}{\mathbf{q}^{-}}               % box lower bounds
\newcommand{\qhi}{\mathbf{q}^{+}}               % box upper bounds

\makeatletter
\newcommand*{\ClaimHeadFont}{\bfseries}   % switched by prominence
\newcommand*{\ClaimNoteFont}{\mdseries}
\newcommand*{\ClaimShortId}{}
\newcommand*{\SetShortId}[1]{\renewcommand*{\ClaimShortId}{#1}}
\newcommand*{\RestatedName}{}
% Upright bodies throughout: italic is reserved for non-normative text.
% \@endparpenalty\@M keeps the meta line on the statement's page.
\newtheoremstyle{claim}{\topsep}{\topsep}{\normalfont\@endparpenalty\@M}{}%
  {\ClaimHeadFont}{.}{.5em}{\thmname{#1}\thmnumber{ #2}\thmnote{ {\ClaimNoteFont(#3)}}}
\newtheoremstyle{claimass}{\topsep}{\topsep}{\normalfont\@endparpenalty\@M}{}%
  {\ClaimHeadFont}{.}{.5em}{\thmname{#1}\thmnumber{ #2}%
  \ifx\ClaimShortId\@empty\else\ [\textsf{\ClaimShortId}]\fi%
  \thmnote{ {\ClaimNoteFont(#3)}}}
\newtheoremstyle{claimremark}{.5\topsep}{.5\topsep}{\small\@endparpenalty\@M}{}%
  {\itshape}{.}{.5em}{\thmname{#1}\thmnumber{ #2}\thmnote{ {\upshape(#3)}}}
\newtheoremstyle{claimreq}{\topsep}{\topsep}{\normalfont\@endparpenalty\@M}{}%
  {\color{reqAccent}\scshape}{.}{.5em}%
  {\thmname{#1}\thmnumber{ #2}\thmnote{ {\upshape(#3)}}}
\makeatother
\theoremstyle{claim}
\newtheorem{theorem}{Theorem}[section]
\newtheorem{lemma}[theorem]{Lemma}
\newtheorem{proposition}[theorem]{Proposition}
\newtheorem{corollary}[theorem]{Corollary}
\newtheorem{fact}[theorem]{Fact}
\newtheorem{definition}[theorem]{Definition}
\newtheorem*{restatedclaim}{\RestatedName}
\theoremstyle{claimass}
\newtheorem{assumption}[theorem]{Assumption}
\newtheorem*{restatedassumption}{\RestatedName}
\theoremstyle{claimremark}
\newtheorem{remark}[theorem]{Remark}
\newtheorem{example}[theorem]{Example}
\theoremstyle{claimreq}
\newtheorem{requirement}[theorem]{Requirement}
\newtheorem*{restatedrequirement}{\RestatedName}
%% --- 8. Prominence: main results a margin rule and a bold heading ----------
\newtcolorbox{mainclaim}{enhanced jigsaw, breakable, blanker,
  borderline west={1.5pt}{-8pt}{mainRule}, left=0pt, right=0pt, top=0pt, bottom=0pt,
  boxsep=0pt, before skip=\topsep, after skip=\topsep, parbox=false,
  before upper={\renewcommand*{\ClaimNoteFont}{\bfseries}}}
\newenvironment{technicalclaim}{\renewcommand*{\ClaimHeadFont}{\mdseries\color{techGrey}}%
  \renewcommand*{\ClaimNoteFont}{}}{}
%% --- 9. Non-normative text, the meta line, navigation ------------------------
\makeatletter
% Tie what follows to this paragraph: no page break, no list penalty before
% the next theorem head (the \@nobreak route a section heading uses).
\newcommand{\keepwithnext}{\par\nobreak\vspace{3pt}\@nobreaktrue
  \everypar{\if@nobreak\@nobreakfalse\clubpenalty\@M\else\clubpenalty\@clubpenalty\everypar{}\fi}}
\newcommand{\claimsummary}[2]{% #1 claim id (named destination), #2 summary
  \par\addvspace{\bigskipamount}\Needspace{6\baselineskip}%
  {\small\itshape\color{nnGrey}\noindent\hypertarget{claim-#1}{}#2\par}\keepwithnext}
\newcommand{\claimanchor}[1]{% a claim without a summary: the destination alone
  \par\addvspace{\bigskipamount}\Needspace{6\baselineskip}%
  \phantomsection\hypertarget{claim-#1}{}\keepwithnext}
\newcommand{\claimleadin}[1]{{\color{nnGrey}\noindent#1\par}\keepwithnext}
\newcommand{\claimmeta}[1]{\par\nobreak\vspace{2pt}%
  {\footnotesize\color{metaGrey}\raggedright\interlinepenalty\@M\noindent#1\par}%
  \addvspace{\medskipamount}}
% After an assumption: the sentence generated from its verification field; after
% any claim with one: its plain-text reading (user_statement). Both
% non-normative, both grey.
\newcommand{\claimverification}[1]{\par\nobreak\vspace{2pt}%
  {\footnotesize\color{metaGrey}\noindent#1\par}\nobreak}
\newcommand{\claimreading}[1]{\par\nobreak\vspace{2pt}%
  {\footnotesize\color{nnGrey}\noindent\emph{Reading:}\enspace#1\par}\nobreak}
\makeatother
\newcommand{\claimbookmark}[4]{% #1 id, #2 kind word, #3 label, #4 plain title
  \bookmark[dest=claim-#1,level=2]{#2 \getrefnumber{#3} #4}}
\newcommand{\metasep}{\unskip\enspace\textperiodcentered\enspace\ignorespaces}
\newcommand{\metalabel}[1]{\emph{#1}\enspace}
\newcommand{\sectionintro}[1]{{\color{nnGrey}\noindent#1\par}\smallskip}
\newcommand{\sectionimports}[1]{{\footnotesize\color{metaGrey}\raggedright\noindent\metalabel{Imports}#1\par}\smallskip}
\newenvironment{overview}{\par\color{nnGrey}\itshape\noindent\ignorespaces}{\par}
% References: the whole "Def. 2.1" is the link; ids are the fallback text of
% every link into another PDF (a relative URI to a named destination).
\newcommand{\mref}[2]{\hyperref[#2]{#1~\ref*{#2}}}
\newcommand{\aref}[2]{\hyperref[#2]{\textsf{#1}}}
\DeclareUrlCommand{\claimid}{\urlstyle{sf}\def\UrlBreaks{\do\/\do\:\do\-\do\.\do\_}%
  \def\UrlBigBreaks{}\def\UrlNoBreaks{}}
\newcommand{\extref}[2]{\hyperref{#1}{}{claim-#2}{\claimid{#2}}}
\newcommand{\extaref}[3]{\hyperref{#1}{}{claim-#2}{\textsf{#3}}}
\newcommand{\coderef}[2]{\href{#1}{\claimid{#2}}}
% Claim-body commands (spec/claim-format.md section 4): \xref{<full id>} prints
% the kind, title and id of a claim of another contract, linked to its stable
% destination in that contract's PDF (one \DeclareXref per cited claim below);
% \term[gloss]{phrase} marks a defined phrase for the index (the gloss feeds
% the Terms table); \tref[phrase]{text} is a use of a defined term. The
% assembler rewrites both in the body into \termdef (the definition's named
% destination term-<slug>) and \termref / \exttermref (a link to it, in this
% or another contract's PDF); the fallbacks below serve text it copies
% elsewhere, such as the digest. \iface is an interface key.
\makeatletter
\newcommand{\DeclareXref}[4]{% id, kind word, title, contract PDF
  \@namedef{xref@#1}{\hyperref{#4}{}{claim-#1}{#2 ``#3'' ({\small\claimid{#1}})}}}
\newcommand{\xref}[1]{\@ifundefined{xref@#1}%
  {\textbf{[unresolved \nolinkurl{#1}]}\@latex@warning{Unresolved xref #1}}%
  {\@nameuse{xref@#1}}}
\makeatother
\DeclareUrlCommand{\iface}{\urlstyle{tt}}
\newcommand{\term}[2][]{\emph{#2}\index{#2}}
\newcommand{\tref}[2][]{#2}
\newcommand{\termdef}[2]{\hypertarget{#1}{}\emph{#2}\index{#2}}
\newcommand{\termref}[2]{\hyperlink{#1}{#2}}
\newcommand{\exttermref}[3]{\hyperref{#1}{}{#2}{#3}}
\hypersetup{pdftitle={Acceptance statistics (statistics) -- contract 0.10.0}}

% Claims of other contracts cited here (\xref, [[...]]).
\DeclareXref{foundations/def:exact-island}{Definition}{The exact island}{../foundations/contract.pdf}
\DeclareXref{foundations/req:consumer-decides-exactly}{Requirement}{Exact decision by every consumer}{../foundations/contract.pdf}
\DeclareXref{foundations/req:rounding-direction}{Requirement}{Rounding direction of emitted bounds}{../foundations/contract.pdf}
\DeclareXref{meat-core/def:admissible-acceptance-set}{Definition}{Admissible acceptance set}{../meat-core/contract.pdf}
\DeclareXref{meat-core/def:feasible}{Definition}{Feasible set and computable objective}{../meat-core/contract.pdf}
\DeclareXref{meat-core/def:postprocessing-model}{Definition}{Fixed-length post-processing model}{../meat-core/contract.pdf}
\DeclareXref{meat-core/frank-wolfe/def:first-order-setting}{Definition}{Ambient program, pairings, acceptance support and linear minimisation}{../meat-core/contract.pdf}
\DeclareXref{meat-core/frank-wolfe/rem:acc-side}{Remark}{Exact minimisation on the acceptance side}{../meat-core/contract.pdf}
\DeclareXref{meat-core/prop:minimal-divergence}{Proposition}{The minimal divergence as the value with every pair dropped}{../meat-core/contract.pdf}
\DeclareXref{meat-core/rem:hlow-range}{Remark}{Reading an infinite value}{../meat-core/contract.pdf}
\DeclareXref{meat-core/req:infinite-value}{Requirement}{Propagation of an infinite value}{../meat-core/contract.pdf}
\DeclareXref{meat-core/req:unreachable-acceptance-set}{Requirement}{The check of the acceptance set against the feasible statistics}{../meat-core/contract.pdf}
\DeclareXref{protocol-compiler/ass:coin}{Assumption}{Public choices as private coins of declared bias}{../protocol-compiler/contract.pdf}
\DeclareXref{protocol-compiler/def:declaration}{Definition}{Protocol declaration}{../protocol-compiler/contract.pdf}
\DeclareXref{protocol-compiler/def:feasible}{Definition}{Feasible set}{../protocol-compiler/contract.pdf}
\DeclareXref{protocol-compiler/def:spec}{Definition}{Protocol specification}{../protocol-compiler/contract.pdf}
\DeclareXref{reductions/ass:tsm}{Assumption}{Output-state test statistics}{../reductions/contract.pdf}

\begin{document}
\begin{center}
{\LARGE\bfseries Acceptance statistics (statistics)}\\[4pt]
{\large Contract, version 0.10.0}
\end{center}
\begin{overview}
This module builds, before the run and from a declared honest model, the acceptance set against which a protocol's observed outcome frequencies are tested, and states how often an honest run passes that test; the key-rate engine of the meat-core module reads the set as the acceptance set of the protocol declaration it solves. Its core fixes the \termref{term-frequency-vector}{frequency vector} of a run and the two admissible geometries, a single frequency vector or a rational polytope in the probability simplex, the polytope possibly in \termref{term-canonical-box-form}{canonical box form} with recorded bounds per outcome. The core's main results are that every geometry passing its \termref{term-validation-predicate}{validation predicate} is a non-empty, compact and convex subset of the simplex, as the engine's acceptance assumption requires (\hyperlink{claim-statistics/thm:geometry-validation}{Theorem~\ref*{thm:geometry-validation}}), and that over a non-empty canonical box, filling the coordinates of smallest coefficient first minimises a rational linear functional exactly, with a dual witness of the same value in the format the engine's certificate expects, while the outcomes with a positive upper bound include every outcome on which a point of the box puts weight (\hyperlink{claim-statistics/thm:box-oracle}{Theorem~\ref*{thm:box-oracle}}). The core also supplies the binomial tails in closed form and the \termref{term-honest-envelope}{honest envelope}, a per-outcome range of honest probabilities of which one honest law is the degenerate case. The binomial-box submodule constructs the box and bounds the probability that an honest run fails the test; the singleton submodule evaluates a one-point test exactly. The declared premises are the honest-round model, an acceptance value recorded in every round, an honest law that some feasible state of the engine's declaration produces, and the output-state model of the test statistics (\hyperref{../reductions/contract.pdf}{}{claim-reductions/ass:tsm}{A-TSM}), under which the acceptance outcome carries every per-round public choice of known bias made at a measuring station, a letter-side choice being a component of the letter drawn as declared (\hyperref{../protocol-compiler/contract.pdf}{}{claim-protocol-compiler/ass:coin}{A-COIN}); that the alphabet and round count are those of that declaration, and that the data are exact and the tails conservatively evaluated, are conditions the engine decides. The honest-round model and the premise on that feasible state bear on completeness, rate and usability, never on the soundness of a key.
\end{overview}
\medskip
\noindent{\small\textcolor{nnGrey}{\itshape Grey text}, here and throughout, is non-normative: it explains, and nothing in it is assumed or proved.}

\smallskip
\noindent{\small The \href{https://github.com/fraunhofersg/certiqs-key/blob/v0.1.0/spec/modules/statistics/design.md}{design notes} (non-normative) explain why the module is built as it is and draw its processing flowchart. The dependencies between the sections are drawn in \href{graphs/overview.svg}{the section dependency matrix}, and \href{https://github.com/fraunhofersg/certiqs-key/blob/v0.1.0/spec/modules/statistics/graphs/README.md}{the graph index} opens every section's dependency graph. The \href{digest.pdf}{statements digest} holds every definition, assumption, result and requirement without the proofs. \href{https://github.com/fraunhofersg/certiqs-key/blob/v0.1.0/harness/contract-reading.md}{How to read a contract} explains this layout and its links.}

{\small\tableofcontents}

\section*{Main results}\addcontentsline{toc}{section}{Main results}
\sectionintro{The module's most important results, those its overview names as what it establishes, with every assumption each rests on, directly or through the claims it uses. Each links to its statement and proof; the statements digest holds them in full.}
\par\noindent\hyperref[thm:geometry-validation]{\textbf{Theorem~\ref*{thm:geometry-validation}}}\enspace Validity of the admissible geometries\par\nobreak
{\small\itshape\color{nnGrey}\noindent An acceptance set whose data pass validation (a probability vector, or a feasible linear system) is a non-empty, compact, admissible acceptance set (\xref{meat-core/def:admissible-acceptance-set}). For a box of per-outcome bounds, validation means the lower sum to at most one and the upper to at least one.\par}\nobreak
\claimmeta{\metalabel{Rests on}\aref{A-ALPHABET}{ass:alphabet}}
\par\noindent\hyperref[thm:box-oracle]{\textbf{Theorem~\ref*{thm:box-oracle}}}\enspace Acceptance support and the exact box oracle\par\nobreak
{\small\itshape\color{nnGrey}\noindent Over a non-empty canonical box, filling the coordinates of smallest coefficient first minimises a rational linear functional exactly, with a dual witness of the same value. The outcomes with a positive upper bound include every outcome on which a point of the box puts weight, and can include more.\par}\nobreak
\claimmeta{\metalabel{Rests on}\aref{A-ALPHABET}{ass:alphabet}}
\par\noindent\hyperref[thm:box-construction]{\textbf{Theorem~\ref*{thm:box-construction}}}\enspace Properties of the constructed box\par\nobreak
{\small\itshape\color{nnGrey}\noindent For any admissible endpoint rule (count bounds meeting binomial tail budgets), the acceptance box has exact rational bounds and contains every law within the honest envelope's bounds. It is non-empty and compact and, by construction, an admissible acceptance set (\xref{meat-core/def:admissible-acceptance-set}).\par}\nobreak
\claimmeta{\metalabel{Rests on}\aref{A-ALPHABET}{ass:alphabet}}
\par\noindent\hyperref[thm:completeness-bound]{\textbf{Theorem~\ref*{thm:completeness-bound}}}\enspace The completeness bound\par\nobreak
{\small\itshape\color{nnGrey}\noindent If each honest round's conditional outcome law stays inside the envelope, an honest run fails the acceptance test with probability at most \ensuremath{\varepsilon}*, the sum of the binomial tails at the envelope's ends beyond the box's count bounds, and \ensuremath{\varepsilon}* is at most the completeness budget. No independence is needed.\par}\nobreak
\claimmeta{\metalabel{Rests on}\aref{A-ALPHABET}{ass:alphabet}, \aref{A-HON-ENV}{ass:hon-env}}
\par\noindent\hyperref[thm:singleton-completeness]{\textbf{Theorem~\ref*{thm:singleton-completeness}}}\enspace Completeness of a singleton acceptance set\par\nobreak
{\small\itshape\color{nnGrey}\noindent Under i.i.d. honest rounds a singleton acceptance set accepts with its point's multinomial probability. That is 0, with no completeness budget below one, if n times a coordinate is no integer or the point leaves the law's support; else the least budget is one minus it, zero only for one point mass.\par}\nobreak
\claimmeta{\metalabel{Rests on}\aref{A-ALPHABET}{ass:alphabet}, \aref{A-HON}{ass:hon}}

\section*{Assumptions at a glance}\addcontentsline{toc}{section}{Assumptions at a glance}
\sectionintro{Every assumption this contract states or its claims assume, with its category and the claims of this contract that assume it directly. An assumption of another contract links into that contract's PDF.}
{\footnotesize\begin{xltabular}{\linewidth}{@{}>{\raggedright\arraybackslash}p{1.6cm}>{\raggedright\arraybackslash}p{1.9cm}>{\raggedright\arraybackslash}X>{\raggedright\arraybackslash}p{2.3cm}>{\raggedright\arraybackslash}p{2.8cm}@{}}
\toprule
Assumption & Category & Title & Stated as & Assumed directly by\\
\midrule
\endhead
\aref{A-ALPHABET}{ass:alphabet} & protocol structure & An acceptance value recorded in every round & \mref{Assumption}{ass:alphabet} & \mref{Thm.}{thm:geometry-validation}, \mref{Thm.}{thm:box-oracle}, \mref{Rem.}{rem:guard-activation}, \mref{Thm.}{thm:box-construction}, \mref{Thm.}{thm:completeness-bound}, \mref{Thm.}{thm:honest-containment}, \mref{Thm.}{thm:monotonicity-limits}, \mref{Thm.}{thm:singleton-completeness}\\
\addlinespace
\extaref{../reductions/contract.pdf}{reductions/ass:tsm}{A-TSM} & protocol structure & Output-state test statistics & \extref{../reductions/contract.pdf}{reductions/ass:tsm} & \mref{Rem.}{rem:guard-activation}, \mref{Rem.}{rem:sound-indep}\\
\addlinespace
\aref{A-HON}{ass:hon} & honest model & Honest i.i.d. rounds & \mref{Assumption}{ass:hon} & \mref{Rem.}{rem:sound-indep}, \mref{Thm.}{thm:singleton-completeness}\\
\addlinespace
\aref{A-HON-ENV}{ass:hon-env} & honest model & Honest rounds inside a declared envelope & \mref{Assumption}{ass:hon-env} & \mref{Thm.}{thm:completeness-bound}, \mref{Rem.}{rem:sound-indep}\\
\addlinespace
\aref{A-REACH}{ass:reach} & honest model & An honest law produced by a feasible state & \mref{Assumption}{ass:reach} & \mref{Thm.}{thm:honest-containment}, \mref{Rem.}{rem:hlow-consequence}\\
\bottomrule
\end{xltabular}}

\section*{Notation}\addcontentsline{toc}{section}{Notation}
\sectionintro{The symbols of the repository's notation catalogue (\href{../../notation.md}{spec/notation.md}) that this contract uses, in the catalogue's order, with their meaning.}
{\small\begin{xltabular}{\linewidth}{@{}>{\raggedright\arraybackslash}p{3.6cm}>{\raggedright\arraybackslash}X@{}}
\toprule
Symbol & Meaning\\
\midrule
\endhead
$\Decl$ & Protocol declaration $\big(A,\tau_A;\,Q,\mathbb{S};\,\{M_x\};\,\{F_v\};\,\mathfrak{P};\,a_A,a_B;\,\phi;\,r,\mathsf{h};\,\mathfrak{L}\big)$ \textemdash{} the tuple of \extref{../meat-core/contract.pdf}{meat-core/def:declaration}: the party partition $\mathfrak{P}$ is a slot, and the marginal slot $\tau_A$ takes a state or the no-marginal value (trusted- vs untrusted-source mode)\\
$\Chat$, $\Chatset$ & The acceptance-statistic register $\widehat{C}$ and its alphabet $\widehat{\mathcal{C}}$\\
$\acSet$ & Acceptance set\\
$\Phimap$ & The test-statistics map $\mathcal{L}(AQ)\to\mathbb{R}^{\widehat{\mathcal{C}}}$\\
$\Feas$ & Feasible set of sector-block-diagonal marginals: $\operatorname{tr}\omega=1$ always; $\omega_A=\tau_A$ in trusted-source mode only; $\mathfrak{L}$ if declared (\extref{../meat-core/contract.pdf}{meat-core/def:feasible})\\
$\hlow$ & Computable relaxation: lower bound on $h^{\uparrow,\mathrm{QKD}}_{\alpha}$; written \texttt{h\_low} in plain text\\
$\OmAT$ & Acceptance event $\operatorname{freq}_{\hat c_1^n}\in\mathcal{S}_{\mathrm{acc}}$\\
$\sfD$ & Classical KL divergence, base 2; also the quantum relative entropy on operators\\
$\transpose$ & Transpose, $X^{\mathsf T}$\\
\bottomrule
\end{xltabular}}

\section*{Terms}\addcontentsline{toc}{section}{Terms}
\sectionintro{The phrases this contract defines, and the terms of other contracts it uses, each with a one-line gloss and the definition that fixes it. In the text, a linked term opens its definition.}
{\small\begin{xltabular}{\linewidth}{@{}>{\raggedright\arraybackslash}p{3.6cm}>{\raggedright\arraybackslash}X>{\raggedright\arraybackslash}p{2.6cm}@{}}
\toprule
Term & Meaning & Defined in\\
\midrule
\endhead
\termref{term-admissible-endpoint-rule}{admissible endpoint rule} & A map from a round count, a success probability and two tail budgets to two count endpoints whose binomial tails meet those budgets, with fixed values at a zero budget and at a degenerate mass. & \mref{Definition}{def:endpoint-rule}\\
\addlinespace[2pt]
\termref{term-admissible-geometry}{admissible geometry} & A tag, singleton or polytope, together with the data of that tag; an acceptance set is declared as one. & \mref{Definition}{def:taxonomy}\\
\addlinespace[2pt]
\termref{term-admissible-input}{admissible input} & The data the construction accepts: alphabet, round count, honest model, budget, constrained outcomes, allocation, and choice of rule and evaluation, the numbers exact and the alphabet and round count those of the consuming declaration. & \mref{Definition}{def:endpoint-rule}\\
\addlinespace[2pt]
\termref{term-canonical-box-form}{canonical box form} & A polytope geometry that records a lower and an upper bound per outcome and whose rows are exactly the two bounds of each outcome. & \mref{Definition}{def:taxonomy}\\
\addlinespace[2pt]
\termref{term-conservative-evaluation}{conservative evaluation} & An evaluation of the acceptance-set construction on exact rational data that decides every tail comparison exactly or against a bound rounding the tail upward. & \mref{Definition}{def:conservative-evaluation}\\
\addlinespace[2pt]
\termref{term-construction-map}{construction map} & The map that turns an admissible input and an admissible endpoint rule into the acceptance box: the guarded count endpoints over n on the constrained outcomes, 0 and 1 elsewhere, intersected with the simplex. & \mref{Definition}{def:construction}\\
\addlinespace[2pt]
\exttermref{../foundations/contract.pdf}{term-consumer}{consumer} & a module other than a supplier that reads declared data & \extref{../foundations/contract.pdf}{foundations/def:declared-data}\\
\addlinespace[2pt]
\termref{term-count}{count} & The number of rounds of an outcome string in which a given outcome occurs. & \mref{Definition}{def:freq}\\
\addlinespace[2pt]
\termref{term-degenerate-envelope}{degenerate envelope} & An honest envelope whose lower and upper vectors are equal, hence one honest law. & \mref{Definition}{def:envelope}\\
\addlinespace[2pt]
\termref{term-envelope-set}{envelope set} & The laws on the simplex that lie between the lower and the upper vector of an honest envelope. & \mref{Definition}{def:envelope}\\
\addlinespace[2pt]
\termref{term-exact-binomial-quantile-rule}{exact binomial quantile rule} & The endpoint rule that returns the extremal endpoints, the extremal binomial quantiles of the two tail budgets. & \mref{Definition}{def:interval}\\
\addlinespace[2pt]
\exttermref{../foundations/contract.pdf}{term-exact-island}{exact island} & the computations of the engine, its instances' declarations included, by which a consumer of data or the declarer about to declare them decides whether they satisfy the conditions their declaration states, a declarer computes and decides the reference of a measurement family, a consumer's replacement of an emitted scalar bound is decided to lie on the bound's weakening side, a producer constructs the operator data it emits, the acceptance side of the optimisation is constructed and certified, or a lower bound on a linear minimum over the feasible side of the optimisation is certified & \extref{../foundations/contract.pdf}{foundations/def:exact-island}\\
\addlinespace[2pt]
\termref{term-extremal-endpoints}{extremal endpoints} & The largest lower count and the smallest upper count in 0 to n whose binomial tails meet their allocated budgets. & \mref{Definition}{def:interval}\\
\addlinespace[2pt]
\termref{term-frequency-vector}{frequency vector} & The counts of an outcome string divided by its length n, one coordinate per outcome. & \mref{Definition}{def:freq}\\
\addlinespace[2pt]
\termref{term-guarded-endpoints}{guarded endpoints} & A rule's two count endpoints after the guard: the lower one capped at the floor of n times the envelope's lower value, the upper one raised to at least the ceiling of n times its upper value. & \mref{Definition}{def:endpoint-rule}\\
\addlinespace[2pt]
\termref{term-honest-envelope}{honest envelope} & A lower and an upper exact rational probability per outcome, the declared range of honest probabilities. & \mref{Definition}{def:envelope}\\
\addlinespace[2pt]
\termref{term-honest-containment-guard}{honest-containment guard} & The widening of a rule's count interval that makes it contain n times the whole envelope interval of the outcome. & \mref{Definition}{def:endpoint-rule}\\
\addlinespace[2pt]
\exttermref{../foundations/contract.pdf}{term-instance}{instance} & the party that fixes one protocol's device data, protocol data and run configuration & \extref{../foundations/contract.pdf}{foundations/def:declared-data}\\
\addlinespace[2pt]
\termref{term-matched-alphabet-and-round-count}{matched alphabet and round count} & An acceptance alphabet, with a fixed enumeration, and a round count that are those of the declaration the consumer solves. & \mref{Definition}{def:matched-alphabet}\\
\addlinespace[2pt]
\termref{term-monotone-allocation-rule}{monotone allocation rule} & A map assigning to each completeness budget an allocation within it, every allocated tail budget non-decreasing in the budget. & \mref{Definition}{def:monotone-rule}\\
\addlinespace[2pt]
\termref{term-monotone-in-the-allocation}{monotone in the allocation} & Said of an endpoint rule whose lower endpoint never falls and whose upper endpoint never rises as its own tail budget grows. & \mref{Definition}{def:endpoint-rule}\\
\addlinespace[2pt]
\exttermref{../reductions/contract.pdf}{term-outcome-alphabet}{outcome alphabet} & The finite ordered set of device outcomes that a private setting is joined to. & \extref{../reductions/contract.pdf}{reductions/def:joint-family}\\
\addlinespace[2pt]
\termref{term-point-form}{point form} & The honest model given as one law rather than as an envelope. & \mref{Definition}{def:envelope}\\
\addlinespace[2pt]
\exttermref{../foundations/contract.pdf}{term-round-homogeneous}{round-homogeneous} & one fixed list of objects describes every round, and every round behaves as they describe whatever the other rounds do & \extref{../foundations/contract.pdf}{foundations/def:round-homogeneity}\\
\addlinespace[2pt]
\termref{term-side-wise-reading}{side-wise reading} & The convention that a single-value claim of this submodule is read at the envelope's lower value for the lower endpoint and at its upper value for the upper endpoint. & \mref{Definition}{def:endpoint-rule}\\
\addlinespace[2pt]
\termref{term-uniform-rule}{uniform rule} & The allocation rule that splits the completeness budget equally over the two tails of every constrained outcome. & \mref{Definition}{def:monotone-rule}\\
\addlinespace[2pt]
\termref{term-validation-predicate}{validation predicate} & The condition on a geometry's data that validation checks: a probability vector for a singleton, a feasible linear system for a polytope. & \mref{Definition}{def:taxonomy}\\
\bottomrule
\end{xltabular}}

\clearpage
\part{Mathematics}

\section{Frequencies and acceptance geometries}\label{group-statistics/def:freq}
\sectionintro{A run's outcome counts divided by the round count give a \termref{term-frequency-vector}{frequency vector} in the probability simplex, and the acceptance test asks whether that vector lies in a declared set. The alphabet and the round count are those of the declaration the consuming engine solves. The set is an \termref{term-admissible-geometry}{admissible geometry}, a single frequency vector or a rational polytope in the simplex, the polytope possibly in \termref{term-canonical-box-form}{canonical box form} with recorded lower and upper bounds per outcome. The main result, \hyperlink{claim-statistics/thm:geometry-validation}{Theorem~\ref*{thm:geometry-validation}}, is that a geometry passing its \termref{term-validation-predicate}{validation predicate} is a non-empty, compact and convex subset of the simplex; for the canonical box form the predicate is a closed-form condition on the sums of the bounds.}

% claim statistics/def:freq
\claimsummary{statistics/def:freq}{An outcome's count is the number of rounds in which it occurs, and the frequency vector of a run of n rounds is its counts divided by n. The acceptance test asks whether this vector lies in the acceptance set.}
\claimbookmark{statistics/def:freq}{Definition}{def:freq}{Frequency vector and counts}
\claimleadin{%
The event $\OmAT=\{\operatorname{freq}_{\hat c_1^n}\in\acSet\}$ that a run
of $n$ rounds passes the acceptance test, the acceptance event of the
key-length post-processing (\xref{meat-core/def:postprocessing-model}; the symbols are listed in
\href{../../../spec/notation.md\#35-optimization-and-key-length-objects}{\nolinkurl{spec/notation.md}\texttt{\#35-optimization-and-key-length-objects}}), is defined
on the frequency vector fixed here.
}
\begin{definition}[Frequency vector and counts]\label{def:freq}
For an outcome string $\hat c_1^n=(\hat c_1,\dots,\hat c_n)\in\Chatset^{\,n}$
and $\hat c\in\Chatset$ put
\begin{equation}\label{eq:freq}
  \Ncount{\hat c}(\hat c_1^n)\coloneqq
  \big|\{\,j\in\{1,\dots,n\}:\hat c_j=\hat c\,\}\big|,
  \qquad
  \operatorname{freq}_{\hat c_1^n}(\hat c)\coloneqq
  \tfrac{1}{n}\,\Ncount{\hat c}(\hat c_1^n):
\end{equation}
the \termdef{term-count}{count} of the outcome $\hat c$ and the
\termdef{term-frequency-vector}{frequency vector} of the string. Then
$\operatorname{freq}_{\hat c_1^n}\in\simplex{\Chatset}$ for every string, and
its coordinates lie in $\tfrac1n\{0,1,\dots,n\}$.
\end{definition}
\claimmeta{\metalabel{Used by}\mref{Lem.}{lem:freq-in-simplex}, \mref{Def.}{def:endpoint-rule}, \mref{Thm.}{thm:completeness-bound}, \mref{Lem.}{lem:widening}, \mref{Lem.}{lem:bernstein-tail}, \mref{Thm.}{thm:singleton-completeness}
\metasep newly derived
\metasep accepted}

% claim statistics/lem:freq-in-simplex
\claimsummary{statistics/lem:freq-in-simplex}{Every frequency vector has non-negative coordinates summing to one, since each round contributes exactly one outcome. Intersecting a constraint set with the probability simplex therefore never changes the probability that a run passes the test, so simplex constraints cost no completeness.}
\claimbookmark{statistics/lem:freq-in-simplex}{Lemma}{lem:freq-in-simplex}{Frequency vectors in the probability simplex}
\begin{lemma}[Frequency vectors in the probability simplex]\label{lem:freq-in-simplex}
Let $n\ge1$, and write $\simplex{\Chatset}$ for the set of
$\mathbf{q}\in\mathbb{R}^{\Chatset}$ with non-negative coordinates summing
to one.
\begin{enumerate}[label=(\roman*),leftmargin=2.2em]
\item For every outcome string $\hat c_1^n\in\Chatset^{\,n}$, every count
$\Ncount{\hat c}(\hat c_1^n)$ of Definition~\ref{def:freq} is a non-negative
integer, at most $n$ of them are non-zero, and
$\sum_{\hat c\in\Chatset}\Ncount{\hat c}(\hat c_1^n)=n$. Hence
$\operatorname{freq}_{\hat c_1^n}\in\simplex{\Chatset}$.
\item For every set $S\subseteq\mathbb{R}^{\Chatset}$ and every random string
$\Chat_1^n$ with values in $\Chatset^{\,n}$, the sets
$\{\operatorname{freq}_{\Chat_1^n}\in S\cap\simplex{\Chatset}\}$ and
$\{\operatorname{freq}_{\Chat_1^n}\in S\}$ of sample points are equal. In
particular, whenever either is an event,
\begin{equation}\label{eq:simplex-free}
  \Pr\big[\operatorname{freq}_{\Chat_1^n}\in S\cap\simplex{\Chatset}\big]
  =\Pr\big[\operatorname{freq}_{\Chat_1^n}\in S\big].
\end{equation}
\end{enumerate}
\end{lemma}
\claimmeta{\metalabel{Uses}\mref{Def.}{def:freq}
\metasep \metalabel{Used by}\mref{Thm.}{thm:completeness-bound}
\metasep newly derived
\metasep accepted}

\begin{proof}
\emph{(i).} Fix $\hat c_1^n$. By \eqref{eq:freq},
$\Ncount{\hat c}(\hat c_1^n)$ is the cardinality of the set
$J_{\hat c}\coloneqq\{j\in\{1,\dots,n\}:\hat c_j=\hat c\}$, a non-negative
integer. Each round $j\in\{1,\dots,n\}$ lies in $J_{\hat c}$ for exactly one
$\hat c$, namely $\hat c=\hat c_j$, so the sets $J_{\hat c}$ for
$\hat c\in\Chatset$ are pairwise disjoint with union $\{1,\dots,n\}$. Hence
at most $n$ of them are non-empty, and summing their cardinalities gives
$\sum_{\hat c\in\Chatset}\Ncount{\hat c}(\hat c_1^n)=n$. Dividing by $n\ge1$,
the coordinates $\operatorname{freq}_{\hat c_1^n}(\hat c)
=\tfrac1n\Ncount{\hat c}(\hat c_1^n)$ are non-negative and sum to one, which
is $\operatorname{freq}_{\hat c_1^n}\in\simplex{\Chatset}$.

\emph{(ii).} The inclusion
$\{\operatorname{freq}_{\Chat_1^n}\in S\cap\simplex{\Chatset}\}
\subseteq\{\operatorname{freq}_{\Chat_1^n}\in S\}$ holds because
$S\cap\simplex{\Chatset}\subseteq S$. Conversely, at every sample point the
value of $\Chat_1^n$ is a string in $\Chatset^{\,n}$, so by (i) the frequency
vector there lies in $\simplex{\Chatset}$; if it also lies in $S$, it lies in
$S\cap\simplex{\Chatset}$. The two sets are therefore equal, and equal events
have equal probability, which is \eqref{eq:simplex-free}.
\end{proof}

% claim statistics/def:matched-alphabet
\claimsummary{statistics/def:matched-alphabet}{An alphabet and a round count are matched when the alphabet, with a fixed enumeration, is the acceptance alphabet of the declaration the consumer solves and the round count is the same. The acceptance set then lives on the consumer's coordinates and is normalised by its round count.}
\claimbookmark{statistics/def:matched-alphabet}{Definition}{def:matched-alphabet}{Alphabet and round count matched to the consuming declaration}
\claimleadin{%
The definition relates this module's inputs to the declaration the consumer
solves, so no input of this module decides it; the agreement of
\exttermref{../reductions/contract.pdf}{term-outcome-alphabet}{outcome alphabets} across artefacts is checked where
the consumer validates its inputs
(\href{../../../spec/architecture.md\#36-bundle-validation-is-input-validation-not-a-module}{\nolinkurl{spec/architecture.md}\texttt{\#36-bundle-validation-is-input-validation-not-a-module}};
Requirement~\ref{req:matched-alphabet}), and the round count is one run
parameter that this module and the consumer both read.
}
\begin{definition}[Alphabet and round count matched to the consuming declaration]\label{def:matched-alphabet}
An acceptance alphabet $\Chatset$ and a round count $n$ are a
\termdef{term-matched-alphabet-and-round-count}{matched alphabet and
round count} when $\Chatset$ is the finite, non-empty acceptance alphabet of
the same declaration $\Decl$ (\xref{protocol-compiler/def:declaration}) the
\exttermref{../foundations/contract.pdf}{term-consumer}{consumer} solves, the codomain of $\phi$, with a fixed enumeration, and
$n$ is the same total round count.
\end{definition}
\claimmeta{\metalabel{Used by}\mref{Thm.}{thm:geometry-validation}, \mref{Thm.}{thm:box-oracle}, \mref{Req.}{req:matched-alphabet}, \mref{Def.}{def:endpoint-rule}, \mref{Rem.}{rem:guard-activation}, \mref{Thm.}{thm:honest-containment}, \mref{Thm.}{thm:singleton-completeness}
\metasep newly derived
\metasep accepted}

% claim statistics/def:taxonomy
\claimsummary{statistics/def:taxonomy}{An admissible geometry is a singleton, one rational vector, or a polytope, rational linear inequalities intersected with the probability simplex, each with its validation predicate. A polytope in canonical box form is given by recorded lower and upper bounds per outcome.}
\claimbookmark{statistics/def:taxonomy}{Definition}{def:taxonomy}{Admissible geometries and the canonical box form}
\claimleadin{%
The acceptance test asks whether the frequency vector of a run lies in a
declared set. The set is declared in one of two forms, and a polytope may in
addition record per-outcome bounds from which its rows are built.
}
\begin{definition}[Admissible geometries and the canonical box form]\label{def:taxonomy}
An \termdef{term-admissible-geometry}{admissible geometry} is a pair consisting of a \emph{tag}
and its \emph{data}, drawn from exactly the following two cases, each with
its data, the set $\Gamma$ the data denote and its
\termdef{term-validation-predicate}{validation predicate}. There is no third tag.
\begin{description}[leftmargin=1.9em]
\item[$\mathrm{singleton}$.] Data: a vector
$\mathbf{q}^{\ast}\in\mathbb{Q}^{\Chatset}$. Denotation
$\Gamma\coloneqq\{\mathbf{q}^{\ast}\}$. \emph{Validation predicate:}
$q^{\ast}(\hat c)\ge0$ for every $\hat c$ and $\sum_{\hat c}q^{\ast}(\hat
c)=1$.
\item[$\mathrm{polytope}$.] Data: an integer row count $m_{\mathrm{row}}\ge0$,
a matrix $B\in\mathbb{Q}^{\,m_{\mathrm{row}}\times\Chatset}$ and a vector
$\mathbf{b}\in\mathbb{Q}^{\,m_{\mathrm{row}}}$. Denotation
\begin{equation}\label{eq:polytope}
  \Gamma\coloneqq\{\mathbf{q}\in\simplex{\Chatset}:\ B\mathbf{q}\le\mathbf{b}\},
\end{equation}
the inequality componentwise. \emph{Validation predicate:} feasibility of the
linear system $\{\mathbf{q}\ge0,\ \sum_{\hat c}q(\hat c)=1,\
B\mathbf{q}\le\mathbf{b}\}$ over the rationals.
\end{description}
A $\mathrm{polytope}$ geometry is in \termdef{term-canonical-box-form}{canonical box form} when it carries, in addition, a
pair $\qlo,\qhi\in\mathbb{Q}^{\Chatset}$ with $0\le\qlo\le\qhi$
componentwise, and its $(B,\mathbf{b})$ are exactly
\begin{equation}\label{eq:boxrows}
  m_{\mathrm{row}}=2m,\qquad
  \begin{aligned}
    &\text{row }i:\ \ \delta_{\hat c^{(i)}}^{\transpose},
      &&\text{right-hand side } q^{+}(\hat c^{(i)}),\\
    &\text{row }m+i:\ \ -\delta_{\hat c^{(i)}}^{\transpose},
      &&\text{right-hand side } -q^{-}(\hat c^{(i)}),
  \end{aligned}
  \qquad i=1,\dots,m,
\end{equation}
where $m\coloneqq|\Chatset|$, $\hat c^{(1)},\dots,\hat c^{(m)}$ is the
declaration's enumeration of $\Chatset$ and $\delta_{\hat c}$ is the standard
basis vector of $\mathbb{R}^{\Chatset}$ at $\hat c$. Row indices are integers,
never alphabet elements; equivalently
\begin{equation}\label{eq:boxset}
  \Gamma=\{\mathbf{q}\in\simplex{\Chatset}:\
  q^{-}(\hat c)\le q(\hat c)\le q^{+}(\hat c)\ \ \forall\hat c\in\Chatset\}.
\end{equation}
\end{definition}
\claimmeta{\metalabel{Used by}\mref{Thm.}{thm:geometry-validation}, \mref{Thm.}{thm:box-oracle}, \mref{Def.}{def:envelope}, \mref{Req.}{req:geometry-tag}, \mref{Req.}{req:box-bounds-key}, \mref{Def.}{def:construction}, \mref{Thm.}{thm:box-construction}, \mref{Lem.}{lem:widening}
\metasep newly derived
\metasep implemented
\metasep \metalabel{Code}\coderef{https://github.com/fraunhofersg/certiqs-key/blob/v0.1.0/src/certiqs_key/statistics/taxonomy.py}{certiqs_key.statistics.taxonomy:GeometryTag}}

% claim statistics/thm:geometry-validation
\claimsummary{statistics/thm:geometry-validation}{An acceptance set whose data pass validation (a probability vector, or a feasible linear system) is a non-empty, compact, admissible acceptance set (\xref{meat-core/def:admissible-acceptance-set}). For a box of per-outcome bounds, validation means the lower sum to at most one and the upper to at least one.}
\claimbookmark{statistics/thm:geometry-validation}{Theorem}{thm:geometry-validation}{Validity of the admissible geometries}
\claimleadin{%
The consuming engine assumes that its acceptance set is a closed convex subset
of the probability simplex. The alphabet is the finite, enumerated one of
Definition~\ref{def:matched-alphabet}, so the simplex is compact; a polytope cuts it
with finitely many closed half-spaces, and along the segment from the lower to
the upper bound vector of a canonical box the coordinate sum moves
continuously from the sum of the lower bounds to that of the upper bounds,
which gives the predicate in closed form.
}
\begin{mainclaim}
\begin{theorem}[Validity of the admissible geometries]\label{thm:geometry-validation}
Let $\Chatset$ and $n$ be a \termref{term-matched-alphabet-and-round-count}{matched alphabet and round count}
(Definition~\ref{def:matched-alphabet}), and let $(\text{tag},\text{data})$
be an admissible geometry with denotation $\Gamma$.
\begin{enumerate}[label=(\roman*),leftmargin=2.2em]
\item If the validation predicate holds, $\Gamma$ is a non-empty, closed,
convex and compact subset of $\simplex{\Chatset}$. In particular
\xref{meat-core/def:admissible-acceptance-set}, that $\acSet$ is a closed convex subset of the
probability simplex on $\Chatset$, holds by construction for every admissible
geometry whose validation predicate holds.
\item For the tag $\mathrm{singleton}$, the validation predicate is necessary
and sufficient for $\Gamma$ to be a non-empty subset of $\simplex{\Chatset}$.
\item For the tag $\mathrm{polytope}$ in canonical box form with
$0\le\qlo\le\qhi$, the validation predicate, the condition
$\Gamma\neq\varnothing$ and the closed-form condition
\begin{equation}\label{eq:boxfeas}
  \sum_{\hat c\in\Chatset}q^{-}(\hat c)\ \le\ 1\ \le\
  \sum_{\hat c\in\Chatset}q^{+}(\hat c)
\end{equation}
are equivalent.
\end{enumerate}
\end{theorem}
\end{mainclaim}
\claimmeta{\metalabel{Uses}\mref{Def.}{def:matched-alphabet}, \mref{Def.}{def:taxonomy}
\metasep \metalabel{Assumes}\aref{A-ALPHABET}{ass:alphabet}
\metasep \metalabel{Used by}\mref{Thm.}{thm:box-oracle}, \mref{Thm.}{thm:box-construction}, \mref{Thm.}{thm:honest-containment}
\metasep newly derived
\metasep implemented (effective: accepted)
\metasep \metalabel{Code}\coderef{https://github.com/fraunhofersg/certiqs-key/blob/v0.1.0/src/certiqs_key/statistics/taxonomy.py}{certiqs_key.statistics.taxonomy:validate_geometry}
\metasep \metalabel{Oracles}\coderef{https://github.com/fraunhofersg/certiqs-key/blob/v0.1.0/tests/oracles/registry_statistics_acceptance.py}{registry_statistics_acceptance.py}}

\begin{proof}
\emph{(i).} $\simplex{\Chatset}$ is the intersection of the closed half-spaces
$\{q(\hat c)\ge0\}$ with the hyperplane $\{\sum_{\hat c}q(\hat c)=1\}$, hence
closed and convex; it is bounded, since $0\le q(\hat c)\le1$ for every
coordinate, hence compact ($\Chatset$ is finite, so
$\mathbb{R}^{\Chatset}$ is finite-dimensional and Heine--Borel applies). For
$\mathrm{polytope}$, $\{B\mathbf{q}\le\mathbf{b}\}$ is a finite intersection
of closed half-spaces, so closed and convex, and $\Gamma$ is the intersection
of two closed convex sets, hence closed and convex; being a closed subset of
the compact $\simplex{\Chatset}$, it is compact. For $\mathrm{singleton}$, a
one-point set is closed, convex and compact, and lies in $\simplex{\Chatset}$
exactly when the validation predicate holds. Non-emptiness follows from the
validation predicate in every case: for $\mathrm{singleton}$ by (ii); for a
general $\mathrm{polytope}$ because the predicate is rational feasibility of
$\{\mathbf{q}\ge0,\ \sum_{\hat c}q(\hat c)=1,\ B\mathbf{q}\le\mathbf{b}\}$ and
a feasible point is by \eqref{eq:polytope} an element of $\Gamma$; and in
canonical box form equivalently by the closed form (iii).

\emph{(ii).} $\{\mathbf{q}^{\ast}\}\subseteq\simplex{\Chatset}$ and
$\{\mathbf{q}^{\ast}\}\neq\varnothing$ hold if and only if
$\mathbf{q}^{\ast}\in\simplex{\Chatset}$, which is the predicate verbatim.

\emph{(iii).} We show that the predicate gives $\Gamma\neq\varnothing$, that
$\Gamma\neq\varnothing$ gives \eqref{eq:boxfeas}, and that
\eqref{eq:boxfeas} gives the predicate. By \eqref{eq:boxrows} the rows
$B\mathbf{q}\le\mathbf{b}$ read $\qlo\le\mathbf{q}\le\qhi$ componentwise.

\emph{Predicate $\Rightarrow$ non-emptiness.} A rational solution of the
system of the predicate is by \eqref{eq:polytope} an element of $\Gamma$.

\emph{Non-emptiness $\Rightarrow$ \eqref{eq:boxfeas}.} If
$\mathbf{q}\in\Gamma$ then $\qlo\le\mathbf{q}\le\qhi$ componentwise and
$\sum_{\hat c}q(\hat c)=1$; summing the two componentwise inequalities gives
\eqref{eq:boxfeas}.

\emph{\eqref{eq:boxfeas} $\Rightarrow$ predicate.} Put, for
$\theta\in[0,1]$, $\mathbf{q}(\theta)\coloneqq\qlo+\theta(\qhi-\qlo)$. Since
$\qlo\le\qhi$, every $\mathbf{q}(\theta)$ satisfies
$\qlo\le\mathbf{q}(\theta)\le\qhi$ and, using $\qlo\ge0$, also
$\mathbf{q}(\theta)\ge0$. Write $S^{-}\coloneqq\sum_{\hat c}q^{-}(\hat c)$ and
$S^{+}\coloneqq\sum_{\hat c}q^{+}(\hat c)$, so that $S^{-}\le1\le S^{+}$ by
\eqref{eq:boxfeas} and $\sum_{\hat c}q(\theta)(\hat c)=S^{-}+\theta(S^{+}-S^{-})$.
Put $\theta^{\ast}\coloneqq(1-S^{-})/(S^{+}-S^{-})$ if $S^{+}>S^{-}$, and
$\theta^{\ast}\coloneqq0$ if $S^{+}=S^{-}$, in which case $S^{-}=S^{+}=1$. In
both cases $\theta^{\ast}\in[0,1]$ and
$\sum_{\hat c}q(\theta^{\ast})(\hat c)=1$, and $\theta^{\ast}$ is rational
because $\qlo$ and $\qhi$ are (Definition~\ref{def:taxonomy}). So
$\mathbf{q}(\theta^{\ast})$ is a rational vector with
$\mathbf{q}(\theta^{\ast})\ge0$, $\sum_{\hat c}q(\theta^{\ast})(\hat c)=1$ and
$\qlo\le\mathbf{q}(\theta^{\ast})\le\qhi$, a rational solution of the system
of the predicate.
\end{proof}

% claim statistics/rem:two-tags
\claimsummary{statistics/rem:two-tags}{The two tags are exactly the two geometries over which the consuming engine certifies the minimum of a linear functional, and the canonical box form is a sub-form of the polytope, not a third tag. A consumer that ignores the recorded bounds still reads a well-formed polytope.}
\claimbookmark{statistics/rem:two-tags}{Remark}{rem:two-tags}{The two tags and the consumer's certified geometries}
\begin{remark}[The two tags and the consumer's certified geometries]\label{rem:two-tags}
The two tags of Definition~\ref{def:taxonomy} correspond one to one to the two
instantiations of \xref{meat-core/frank-wolfe/rem:acc-side}, the geometries
for which that remark certifies the acceptance-side minimum, the minimum of a
linear functional over the acceptance set: (a) singleton, (b)
polytope $\{\mathbf{q}\ge0,\ \sum_{\hat c}q(\hat c)=1,\
B\mathbf{q}\le\mathbf{b}\}$ for exact $B$ and $\mathbf{b}$. That remark
certifies no set of any other form. The canonical box form is a
\emph{sub-form} of $\mathrm{polytope}$, not a third tag: it adds a recorded
pair $(\qlo,\qhi)$ (Requirement~\ref{req:box-bounds-key}) that licenses the
closed-form oracle of Theorem~\ref{thm:box-oracle}, and a \exttermref{../foundations/contract.pdf}{term-consumer}{consumer} that
ignores the recorded pair still reads a well-formed $\mathrm{polytope}$.
\end{remark}
\claimmeta{newly derived
\metasep accepted}

\section{Binomial tails}\label{group-statistics/thm:binomial-tails}
\sectionintro{The count of one outcome over independent rounds with one law is binomial, and the box construction reads its two tails. They are regularised incomplete beta integrals, with every boundary case of the count stated and the upper tail non-decreasing in the success probability (\hyperlink{claim-statistics/thm:binomial-tails}{Theorem~\ref*{thm:binomial-tails}}). The two strict tails the construction compares with a budget follow by a shift of the count (\hyperlink{claim-statistics/cor:tail-forms}{Corollary~\ref*{cor:tail-forms}}), and at a rational success probability both can be decided exactly (\hyperlink{claim-statistics/lem:exact-strict-tails}{Lemma~\ref*{lem:exact-strict-tails}}).}

% claim statistics/thm:binomial-tails
\claimsummary{statistics/thm:binomial-tails}{The upper and lower tails of a binomial count are regularised incomplete beta integrals, with every boundary case of the count stated. The upper tail is non-decreasing in the success probability.}
\claimbookmark{statistics/thm:binomial-tails}{Theorem}{thm:binomial-tails}{Binomial tails in closed form}
\claimleadin{%
The count of one outcome over independent rounds with one law is binomial,
and the box construction compares its two tails with their budgets. The
theorem writes both tails as regularised incomplete beta integrals, deriving
the normalisation of the beta integral inside the proof.
}
\begin{theorem}[Binomial tails in closed form]\label{thm:binomial-tails}
Let $n\ge1$, $p\in[0,1]$, and let $\Ncount{}$ be a binomial random variable
with parameters $(n,p)$, i.e.\
$\Pr[\Ncount{}=i]=\binom{n}{i}p^{i}(1-p)^{n-i}$ for $0\le i\le n$. Write
\begin{equation}\label{eq:ireg}
  \Ireg{x}{a}{b}\coloneqq
  \frac{\int_0^{x}t^{a-1}(1-t)^{b-1}\,\mathrm{d}t}
       {\int_0^{1}t^{a-1}(1-t)^{b-1}\,\mathrm{d}t},
  \qquad x\in[0,1],\ a,b\ge1 .
\end{equation}
Then
\begin{align}
  \Pr[\Ncount{}\ge k]&=\frac{n!}{(k-1)!\,(n-k)!}
    \int_0^{p}t^{k-1}(1-t)^{n-k}\,\mathrm{d}t
    =\Ireg{p}{k}{n-k+1}, &&1\le k\le n,
    \label{eq:uppertail}\\
  \Pr[\Ncount{}\le k]&=1-\Pr[\Ncount{}\ge k+1]
    =\Ireg{1-p}{n-k}{k+1}, &&0\le k\le n-1,
    \label{eq:lowertail}
\end{align}
together with the boundary values $\Pr[\Ncount{}\ge k]=1$ for $k\le0$,
$\Pr[\Ncount{}\ge k]=0$ for $k>n$, $\Pr[\Ncount{}\le k]=0$ for $k<0$ and
$\Pr[\Ncount{}\le k]=1$ for $k\ge n$. Moreover $p\mapsto\Pr[\Ncount{}\ge k]$
is non-decreasing on $[0,1]$ for every $k\ge1$.
\end{theorem}
\claimmeta{\metalabel{Used by}\mref{Cor.}{cor:tail-forms}, \mref{Lem.}{lem:exact-strict-tails}, \mref{Rem.}{rem:guard-activation}, \mref{Lem.}{lem:pmf-enclosure}, \mref{Lem.}{lem:tail-majorant}
\metasep newly derived
\metasep implemented
\metasep \metalabel{Code}\coderef{https://github.com/fraunhofersg/certiqs-key/blob/v0.1.0/src/certiqs_key/statistics/tails.py}{certiqs_key.statistics.tails:binomial_cdf}
\metasep \metalabel{Oracles}\coderef{https://github.com/fraunhofersg/certiqs-key/blob/v0.1.0/tests/oracles/registry_statistics_acceptance.py}{registry_statistics_acceptance.py}}

\begin{proof}
\emph{Idea.} The upper tail is a polynomial in $p$ vanishing at $p=0$ whose
derivative telescopes to a single term; integrating that term gives the first
form of \eqref{eq:uppertail}, and its value at $p=1$ gives the normalisation.
The lower tail is the upper tail of $n-\Ncount{}$.

\emph{Step 1 (the derivative of the upper tail).} Fix $k$ with $1\le k\le n$
and define the polynomial
$G(p)\coloneqq\sum_{i=k}^{n}\binom{n}{i}p^{i}(1-p)^{n-i}$, so that
$G(p)=\Pr[\Ncount{}\ge k]$ and $G(0)=0$ (every summand carries $p^{i}$ with
$i\ge k\ge1$). Set
\begin{equation*}
  a_i\coloneqq n\binom{n-1}{i-1}p^{i-1}(1-p)^{n-i},
  \qquad 1\le i\le n ,
\end{equation*}
a polynomial in $p$ for every such $i$: the exponent $n-i$ is non-negative
throughout, so no negative power of $1-p$ occurs anywhere in this proof and no
division is performed. Differentiate the summands of $G$ separately in the two
cases.

For $k\le i\le n-1$, the product rule and the absorption identities
$i\binom{n}{i}=n\binom{n-1}{i-1}$ and $(n-i)\binom{n}{i}=n\binom{n-1}{i}$ give
\begin{equation*}
  \frac{\mathrm{d}}{\mathrm{d}p}\Big[\binom{n}{i}p^{i}(1-p)^{n-i}\Big]
  =i\binom{n}{i}p^{i-1}(1-p)^{n-i}-(n-i)\binom{n}{i}p^{i}(1-p)^{n-i-1}
  =a_i-a_{i+1},
\end{equation*}
where $i\ge k\ge1$ makes the first power non-negative, $i\le n-1$ makes the
second one non-negative, and $a_{i+1}=n\binom{n-1}{i}p^{i}(1-p)^{n-i-1}$ by
substituting $i+1$ for $i$ in the definition of $a_i$, which is legitimate
because $i+1\le n$. For $i=n$ the summand is $p^{n}$, whose derivative is
$np^{n-1}=a_n$; this is the case the general formula would have written with a
factor $(1-p)^{-1}$, and it is treated separately for exactly that reason.
Summing the two cases, the sum telescopes,
\begin{equation*}
  G'(p)=\sum_{i=k}^{n-1}\big(a_i-a_{i+1}\big)+a_n
       =\big(a_k-a_n\big)+a_n
       =a_k ,
\end{equation*}
which is also correct at $k=n$, where the telescoping sum is empty and
$G'(p)=a_n=a_k$.

\emph{Step 2 (the integral form and its normalisation).} Spelling $a_k$ out,
$G'(p)=n\binom{n-1}{k-1}p^{k-1}(1-p)^{n-k}$, and by the fundamental
theorem of calculus with $G(0)=0$,
\begin{equation}\label{eq:Gint}
  G(p)=n\binom{n-1}{k-1}\int_0^{p}t^{k-1}(1-t)^{n-k}\,\mathrm{d}t,
  \qquad
  n\binom{n-1}{k-1}=\frac{n!}{(k-1)!\,(n-k)!}.
\end{equation}
This is the first equality of \eqref{eq:uppertail}. For the second, evaluate
\eqref{eq:Gint} at $p=1$: there $\Ncount{}=n$ with probability one and $k\le
n$, so $G(1)=1$, whence
$\int_0^{1}t^{k-1}(1-t)^{n-k}\,\mathrm{d}t=\frac{(k-1)!\,(n-k)!}{n!}$;
dividing \eqref{eq:Gint} by this value gives exactly
$\Ireg{p}{k}{n-k+1}$ as defined in \eqref{eq:ireg}, with $a=k\ge1$ and
$b=n-k+1\ge1$.

\emph{Step 3 (the lower tail).} For \eqref{eq:lowertail}, note first that
$\{\Ncount{}\le k\}$ and $\{\Ncount{}\ge k+1\}$ partition the sample space.
Next, $n-\Ncount{}$ is binomial with parameters $(n,1-p)$, because
$\Pr[n-\Ncount{}=i]=\Pr[\Ncount{}=n-i]
=\binom{n}{n-i}p^{n-i}(1-p)^{i}=\binom{n}{i}(1-p)^{i}p^{n-i}$. Hence, for
$0\le k\le n-1$ (so that $n-k\ge1$),
$\Pr[\Ncount{}\le k]=\Pr[n-\Ncount{}\ge n-k]=\Ireg{1-p}{n-k}{k+1}$ by
\eqref{eq:uppertail} applied to $n-\Ncount{}$.

\emph{Step 4 (boundary values and monotonicity).} The boundary values are
immediate from $0\le\Ncount{}\le n$. Monotonicity in $p$ follows from
$G'(p)=n\binom{n-1}{k-1}p^{k-1}(1-p)^{n-k}\ge0$ on $[0,1]$ for $1\le k\le n$;
for $k>n$ the tail is the constant $0$ by the boundary values.
\end{proof}

% claim statistics/cor:tail-forms
\claimsummary{statistics/cor:tail-forms}{For every count k from 0 to n the strict tails of a binomial count N, the probabilities that N lies below k and above k, are regularised incomplete beta integrals, zero at k = 0 and at k = n respectively. These are the two tails the endpoint rules compare with their budgets.}
\claimbookmark{statistics/cor:tail-forms}{Corollary}{cor:tail-forms}{The two strict binomial tails}
\claimleadin{%
An endpoint rule compares the probability that the count falls strictly below
a lower endpoint, or strictly above an upper one, with a budget. Both strict
tails follow from Theorem~\ref{thm:binomial-tails} by a shift of the count.
}
\begin{corollary}[The two strict binomial tails]\label{cor:tail-forms}
Let $n\ge1$, $p\in[0,1]$ and let $\Ncount{}$ be binomial with parameters
$(n,p)$, as in Theorem~\ref{thm:binomial-tails}. Then, for $0\le k\le n$,
\begin{equation}\label{eq:tails-used}
\begin{aligned}
  \Pr[\Ncount{}<k]&=\Pr[\Ncount{}\le k-1]
  =\begin{cases}0,&k=0,\\[2pt] \Ireg{1-p}{n-k+1}{k},&1\le k\le n,\end{cases}\\[4pt]
  \Pr[\Ncount{}>k]&=\Pr[\Ncount{}\ge k+1]
  =\begin{cases}\Ireg{p}{k+1}{n-k},&0\le k\le n-1,\\[2pt] 0,&k=n.\end{cases}
\end{aligned}
\end{equation}
\end{corollary}
\claimmeta{\metalabel{Uses}\mref{Thm.}{thm:binomial-tails}
\metasep \metalabel{Used by}\mref{Lem.}{lem:exact-strict-tails}, \mref{Prop.}{prop:rule-exact}, \mref{Rem.}{rem:guard-activation}, \mref{Rem.}{rem:epsstar}
\metasep newly derived
\metasep implemented
\metasep \metalabel{Code}\coderef{https://github.com/fraunhofersg/certiqs-key/blob/v0.1.0/src/certiqs_key/statistics/tails.py}{certiqs_key.statistics.tails:upper_tail}}

\begin{proof}
Both are obtained from \eqref{eq:uppertail}--\eqref{eq:lowertail} by
substituting $k-1$ and $k+1$ respectively, $\Ncount{}$ being integer-valued.
For the first form: at $k=0$ the boundary value $\Pr[\Ncount{}\le k-1]=0$ for
$k-1<0$ of Theorem~\ref{thm:binomial-tails} applies; for $1\le k\le n$ one has
$0\le k-1\le n-1$, and \eqref{eq:lowertail} at $k-1$ gives
$\Ireg{1-p}{n-(k-1)}{(k-1)+1}=\Ireg{1-p}{n-k+1}{k}$. For the second form: for
$0\le k\le n-1$ one has $1\le k+1\le n$, and \eqref{eq:uppertail} at $k+1$
gives $\Ireg{p}{k+1}{n-(k+1)+1}=\Ireg{p}{k+1}{n-k}$; at $k=n$ the boundary
value $\Pr[\Ncount{}\ge k+1]=0$ for $k+1>n$ of Theorem~\ref{thm:binomial-tails} applies.
\end{proof}

% claim statistics/lem:exact-strict-tails
\claimsummary{statistics/lem:exact-strict-tails}{At a rational success probability both strict tails of a binomial count are finite sums of binomial probabilities, hence rational numbers, and each comparison of a tail with a rational budget is decided exactly. The endpoint rules can therefore be evaluated without rounding.}
\claimbookmark{statistics/lem:exact-strict-tails}{Lemma}{lem:exact-strict-tails}{Exact strict tails at a rational success probability}
\claimleadin{%
Corollary~\ref{cor:tail-forms} writes the two strict tails as regularised
incomplete beta integrals, the form a numerical evaluation of the incomplete
beta function uses. At a rational success probability the finite sum of the
binomial probabilities gives each tail exactly instead; which of the two
routes an evaluation takes is constrained only by the requirement that it be
a \termref{term-conservative-evaluation}{conservative evaluation}
(Definition~\ref{def:conservative-evaluation}).
}
\begin{lemma}[Exact strict tails at a rational success probability]\label{lem:exact-strict-tails}
Let $n\ge1$, let $p\in[0,1]\cap\mathbb{Q}$ and let $\Ncount{}$ be binomial
with parameters $(n,p)$, as in Theorem~\ref{thm:binomial-tails}. For
$0\le k\le n$,
\begin{equation}\label{eq:tails-sums}
  \Pr[\Ncount{}<k]=\sum_{i=0}^{k-1}\binom{n}{i}p^{i}(1-p)^{n-i},
  \qquad
  \Pr[\Ncount{}>k]=\sum_{i=k+1}^{n}\binom{n}{i}p^{i}(1-p)^{n-i},
\end{equation}
an empty sum being $0$, and:
\begin{enumerate}[label=(\roman*),leftmargin=2.2em]
\item both strict tails are rational numbers, obtained from $n$, $k$ and $p$
by finitely many additions and multiplications of rationals, and they are
equal to the closed forms \eqref{eq:tails-used} of
Corollary~\ref{cor:tail-forms};
\item for every rational $\varepsilon$, whether
$\Pr[\Ncount{}<k]\le\varepsilon$ holds, and whether
$\Pr[\Ncount{}>k]\le\varepsilon$ holds, is decided exactly by evaluating the
rational number \eqref{eq:tails-sums} and comparing two rationals, which is a
comparison of two integers.
\end{enumerate}
\end{lemma}
\claimmeta{\metalabel{Uses}\mref{Thm.}{thm:binomial-tails}, \mref{Cor.}{cor:tail-forms}
\metasep \metalabel{Used by}\mref{Prop.}{prop:rule-exact}, \mref{Rem.}{rem:guard-activation}, \mref{Rem.}{rem:epsstar}
\metasep newly derived
\metasep accepted}

\begin{proof}
\emph{The sums.} $\Ncount{}$ takes values in $\{0,\dots,n\}$ with
$\Pr[\Ncount{}=i]=\binom{n}{i}p^{i}(1-p)^{n-i}$ (Theorem~\ref{thm:binomial-tails},
with $0^{0}=1$, the convention under which this law is a probability law at
$p\in\{0,1\}$). For $0\le k\le n$ the event $\{\Ncount{}<k\}$ is the disjoint
union of the events $\{\Ncount{}=i\}$ for the integers $0\le i\le k-1$, and
$\{\Ncount{}>k\}$ is the disjoint union of the events $\{\Ncount{}=i\}$ for
$k+1\le i\le n$; at $k=0$, respectively $k=n$, the union is empty and the
event has probability $0$. Finite additivity of the probability gives
\eqref{eq:tails-sums}.

\emph{(i).} Since $p$ is rational, so is $1-p$; every $\binom{n}{i}$ is an
integer; and $p^{i}(1-p)^{n-i}$ is a product of $n$ rational factors (the
empty product being $1$). Each summand of \eqref{eq:tails-sums} is therefore
a product of finitely many rationals, and each tail a sum of at most $n+1$
such summands. The rationals are closed under addition and multiplication, so
both tails are rational. They equal the forms \eqref{eq:tails-used} because
Corollary~\ref{cor:tail-forms} states that the same two probabilities
$\Pr[\Ncount{}<k]$ and $\Pr[\Ncount{}>k]$ equal those forms, for the same
$n$, $p$ and $k$.

\emph{(ii).} By (i) each tail is a rational number $a/b$ with integers $a$
and $b>0$, computed in finitely many exact operations. Write the rational
$\varepsilon$ as $c/d$ with integers $c$ and $d>0$. Then
$a/b\le c/d$ holds if and only if $ad\le cb$, multiplying by $bd>0$, which is
a comparison of two integers and is decided exactly.
\end{proof}

\section{The exact box oracle}\label{group-statistics/thm:box-oracle}
\sectionintro{The consuming engine minimises linear functionals over the acceptance set and needs to know on which outcomes an accepted frequency can put weight. The main result, \hyperlink{claim-statistics/thm:box-oracle}{Theorem~\ref*{thm:box-oracle}}, solves the minimisation over a non-empty canonical box exactly by filling the coordinates of smallest coefficient first, from their lower towards their upper bounds, supplies a dual witness of the same value in the format the engine's certificate expects, and shows that the outcomes with a positive upper bound include every outcome on which a point of the box puts weight, and can include more.}

% claim statistics/thm:box-oracle
\claimsummary{statistics/thm:box-oracle}{Over a non-empty canonical box, filling the coordinates of smallest coefficient first minimises a rational linear functional exactly, with a dual witness of the same value. The outcomes with a positive upper bound include every outcome on which a point of the box puts weight, and can include more.}
\claimbookmark{statistics/thm:box-oracle}{Theorem}{thm:box-oracle}{Acceptance support and the exact box oracle}
\claimleadin{%
The consuming engine minimises linear functionals over the acceptance set and
needs to know on which outcomes an accepted frequency can put weight. Over a
canonical box the minimum has a closed form, and the outcomes with a positive
upper bound form a closed-form superset of the outcomes that can carry weight.
The dual witness of (iii) has the format of instantiation~(b) of
\xref{meat-core/frank-wolfe/rem:acc-side}: \eqref{eq:dualfeas} implies, at
every outcome, the dual-feasibility inequality that instantiation requires,
and the witness
attains the minimum, \eqref{eq:dualvalue}, which is why the bound it gives is
tight. Rounded down, that bound remains a lower bound, the direction of
\xref{foundations/req:rounding-direction}, so the remark's certified
acceptance-side bound $\bacc$ may be taken as any downward-rounded evaluation
of $\langle g,\mathbf{q}^{\mathrm{opt}}\rangle$.
}
\begin{mainclaim}
\begin{theorem}[Acceptance support and the exact box oracle]\label{thm:box-oracle}
Let $\Chatset$ and $n$ be a \termref{term-matched-alphabet-and-round-count}{matched alphabet and round count}
(Definition~\ref{def:matched-alphabet}), let $\acSet$ be a non-empty
\termref{term-canonical-box-form}{canonical box}, \eqref{eq:boxset}, with
$0\le\qlo\le\qhi$, and let $g\in\mathbb{Q}^{\Chatset}$.
\begin{enumerate}[label=(\roman*),leftmargin=2.2em]
\item \emph{(acceptance support)} The set
$\Chatset^{+}\coloneqq\{\hat c\in\Chatset: q^{+}(\hat c)>0\}$ contains
$\operatorname{supp}\mathbf{q}$ for every $\mathbf{q}\in\acSet$; hence it
contains the acceptance support
$\Chacc=\bigcup_{\mathbf{q}\in\acSet}\operatorname{supp}\mathbf{q}$ of
\xref{meat-core/frank-wolfe/def:first-order-setting}, which it need not equal: the box
constraints together with the simplex constraint can force $q(\hat c)=0$ at an
outcome with $q^{+}(\hat c)>0$. For a $\mathrm{singleton}$ geometry, which
declares no box, the union of the supports is
$\operatorname{supp}\mathbf{q}^{\ast}$, the acceptance support itself.
\item \emph{(exact minimum in closed form)} Enumerate
$\Chatset=\{c_1,\dots,c_m\}$ so that $g(c_1)\le\dots\le g(c_m)$. Put
$r\coloneqq1-\sum_{\hat c}q^{-}(\hat c)$ and
$D_i\coloneqq q^{+}(c_i)-q^{-}(c_i)\ge0$, let $j^{\ast}$ be the least index
with $\sum_{i\le j^{\ast}}D_i\ge r$, and define
\begin{equation}\label{eq:greedy}
  q^{\mathrm{opt}}(c_i)\coloneqq
  \begin{cases}
    q^{+}(c_i), & i<j^{\ast},\\[2pt]
    q^{-}(c_i)+\big(r-\sum_{i'<j^{\ast}}D_{i'}\big), & i=j^{\ast},\\[2pt]
    q^{-}(c_i), & i>j^{\ast}.
  \end{cases}
\end{equation}
Then $j^{\ast}$ exists, $\mathbf{q}^{\mathrm{opt}}\in\acSet$, and
\begin{equation}\label{eq:exactmin}
  \langle g,\mathbf{q}^{\mathrm{opt}}\rangle
  \;=\;\min_{\mathbf{q}\in\acSet}\ \langle g,\mathbf{q}\rangle ,
\end{equation}
the minimum being attained. All quantities are exact rationals when
$g,\qlo,\qhi$ are.
\item \emph{(dual witness)} With $B,\mathbf{b}$ as in
\eqref{eq:boxrows}, $\lambda\coloneqq g(c_{j^{\ast}})$, $v\coloneqq-\lambda$ and
$\mathbf{u}\coloneqq(\mathbf{u}^{+},\mathbf{u}^{-})$ given by
$u^{+}(\hat c)\coloneqq\max\{\lambda-g(\hat c),0\}$ and
$u^{-}(\hat c)\coloneqq\max\{g(\hat c)-\lambda,0\}$, one has
$\mathbf{u}\ge0$,
\begin{equation}\label{eq:dualfeas}
  g(\hat c)+(B^{\transpose}\mathbf{u})(\hat c)+v=0
  \qquad\text{for every }\hat c\in\Chatset,
\end{equation}
in particular the left-hand side is $\ge0$, and
\begin{equation}\label{eq:dualvalue}
  -\langle\mathbf{b},\mathbf{u}\rangle-v
  \;=\;\langle g,\mathbf{q}^{\mathrm{opt}}\rangle
  \;=\;\min_{\mathbf{q}\in\acSet}\langle g,\mathbf{q}\rangle .
\end{equation}
The lower bound $-\langle\mathbf{b},\mathbf{u}\rangle-v$ on the minimum that
this witness gives is therefore \emph{tight} for the canonical box, and every
downward-rounded evaluation of $\langle g,\mathbf{q}^{\mathrm{opt}}\rangle$ is
a lower bound on $\min_{\mathbf{q}\in\acSet}\langle g,\mathbf{q}\rangle$.
\end{enumerate}
\end{theorem}
\end{mainclaim}
\claimmeta{\metalabel{Uses}\mref{Def.}{def:matched-alphabet}, \mref{Def.}{def:taxonomy}, \mref{Thm.}{thm:geometry-validation}
\metasep \metalabel{Assumes}\aref{A-ALPHABET}{ass:alphabet}
\metasep \metalabel{Used by}\mref{Rem.}{rem:box-oracle-consumer}, \mref{Req.}{req:acceptance-support}
\metasep newly derived
\metasep implemented (effective: accepted)
\metasep \metalabel{Code}\coderef{https://github.com/fraunhofersg/certiqs-key/blob/v0.1.0/src/certiqs_key/statistics/minimization.py}{certiqs_key.statistics.minimization:box_minimum}
\metasep \metalabel{Oracles}\coderef{https://github.com/fraunhofersg/certiqs-key/blob/v0.1.0/tests/oracles/registry_statistics_acceptance.py}{registry_statistics_acceptance.py}}

\begin{proof}
\emph{(i).} If $\mathbf{q}\in\acSet$ and $q^{+}(\hat c)=0$ then
$0\le q(\hat c)\le q^{+}(\hat c)=0$, so $\hat c\notin\operatorname{supp}
\mathbf{q}$; contrapositively $\operatorname{supp}\mathbf{q}\subseteq\Chatset^{+}$,
and taking the union over $\mathbf{q}\in\acSet$ gives
$\Chacc\subseteq\Chatset^{+}$. That the inclusion can be strict: on
$\Chatset=\{c_1,c_2,c_3\}$ with $\qlo=(\tfrac12,\tfrac12,0)$ and
$\qhi=(\tfrac12,\tfrac12,\tfrac1{10})$ the box is non-empty by
\eqref{eq:boxfeas}, the simplex constraint forces $q(c_3)=0$ for every
$\mathbf{q}\in\acSet$, so $\Chacc=\{c_1,c_2\}$ while
$\Chatset^{+}=\{c_1,c_2,c_3\}$. For a singleton the union
$\bigcup_{\mathbf{q}\in\acSet}\operatorname{supp}\mathbf{q}$ is
$\operatorname{supp}\mathbf{q}^{\ast}$ by inspection, which is therefore the
acceptance support $\Chacc$.

\emph{(ii).} \emph{Existence of $j^{\ast}$.} $\acSet\neq\varnothing$ gives
$\sum_{\hat c}q^{+}(\hat c)\ge1$ by Theorem~\ref{thm:geometry-validation}\,(iii), hence
$\sum_{i\le m}D_i=\sum_{\hat c}q^{+}(\hat c)-\sum_{\hat c}q^{-}(\hat c)\ge
1-\sum_{\hat c}q^{-}(\hat c)=r$, so the index set is non-empty and has a least
element. Also $r\ge0$, again by Theorem~\ref{thm:geometry-validation}\,(iii).

\emph{Feasibility of $\mathbf{q}^{\mathrm{opt}}$.} For $i\neq j^{\ast}$ the value is
one of the two box endpoints. At $i=j^{\ast}$ the increment
$w\coloneqq r-\sum_{i'<j^{\ast}}D_{i'}$ satisfies $w>0$ when
$j^{\ast}>1$, by minimality of $j^{\ast}$, and $w=r\ge0$ when
$j^{\ast}=1$; and $w\le D_{j^{\ast}}$ by the defining inequality
$\sum_{i\le j^{\ast}}D_i\ge r$. Hence
$q^{-}(c_{j^{\ast}})\le q^{\mathrm{opt}}(c_{j^{\ast}})\le q^{+}(c_{j^{\ast}})$, and
$\mathbf{q}^{\mathrm{opt}}\ge\qlo\ge0$. Summing,
$\sum_i q^{\mathrm{opt}}(c_i)=\sum_i q^{-}(c_i)+\sum_{i<j^{\ast}}D_i+w
=\sum_i q^{-}(c_i)+r=1$. So $\mathbf{q}^{\mathrm{opt}}\in\acSet$.

\emph{Optimality.} Let $\mathbf{q}\in\acSet$ and $\lambda=g(c_{j^{\ast}})$.
Since $\sum_i\big(q(c_i)-q^{\mathrm{opt}}(c_i)\big)=1-1=0$,
\begin{equation*}
  \langle g,\mathbf{q}\rangle-\langle g,\mathbf{q}^{\mathrm{opt}}\rangle
  =\sum_{i=1}^{m}\big(g(c_i)-\lambda\big)
   \big(q(c_i)-q^{\mathrm{opt}}(c_i)\big).
\end{equation*}
Each summand is non-negative: for $i<j^{\ast}$ the ordering gives
$g(c_i)-\lambda\le0$ while $q(c_i)-q^{\mathrm{opt}}(c_i)=q(c_i)-q^{+}(c_i)\le0$; for
$i>j^{\ast}$, $g(c_i)-\lambda\ge0$ while
$q(c_i)-q^{\mathrm{opt}}(c_i)=q(c_i)-q^{-}(c_i)\ge0$; and at $i=j^{\ast}$ the first
factor is $0$. Hence
$\langle g,\mathbf{q}\rangle\ge\langle g,\mathbf{q}^{\mathrm{opt}}\rangle$ for every
feasible $\mathbf{q}$, which is \eqref{eq:exactmin} with attainment at
$\mathbf{q}^{\mathrm{opt}}$.

\emph{(iii).} Non-negativity of $\mathbf{u}$ is clear. With the row structure
\eqref{eq:boxrows}, $B^{\transpose}\mathbf{u}=\mathbf{u}^{+}-\mathbf{u}^{-}$,
and for every real $x$ one has $\max\{x,0\}-\max\{-x,0\}=x$; applying this at
$x=\lambda-g(\hat c)$ gives
$u^{+}(\hat c)-u^{-}(\hat c)=\lambda-g(\hat c)$, so
$g(\hat c)+(B^{\transpose}\mathbf{u})(\hat c)+v
=g(\hat c)+\lambda-g(\hat c)-\lambda=0$, which is \eqref{eq:dualfeas}.

For \eqref{eq:dualvalue}, note first that the sorted order and
$\lambda=g(c_{j^{\ast}})$ give $u^{-}(c_i)=0$ for $i\le j^{\ast}$ and
$u^{+}(c_i)=0$ for $i\ge j^{\ast}$. Then
\begin{align*}
  -\langle\mathbf{b},\mathbf{u}\rangle-v
  &=-\langle\qhi,\mathbf{u}^{+}\rangle+\langle\qlo,\mathbf{u}^{-}\rangle
    +\lambda\\
  &=-\sum_{i<j^{\ast}}q^{+}(c_i)\big(\lambda-g(c_i)\big)
    +\sum_{i>j^{\ast}}q^{-}(c_i)\big(g(c_i)-\lambda\big)+\lambda\\
  &=\lambda\Big(1-\sum_{i<j^{\ast}}q^{+}(c_i)-\sum_{i>j^{\ast}}q^{-}(c_i)\Big)
    +\sum_{i<j^{\ast}}g(c_i)q^{+}(c_i)+\sum_{i>j^{\ast}}g(c_i)q^{-}(c_i)\\
  &=\lambda\,q^{\mathrm{opt}}(c_{j^{\ast}})+\sum_{i\neq j^{\ast}}g(c_i)q^{\mathrm{opt}}(c_i)
   \;=\;\langle g,\mathbf{q}^{\mathrm{opt}}\rangle ,
\end{align*}
where the fourth line uses $q^{\mathrm{opt}}(c_i)=q^{+}(c_i)$ for $i<j^{\ast}$,
$q^{\mathrm{opt}}(c_i)=q^{-}(c_i)$ for $i>j^{\ast}$, $\sum_i q^{\mathrm{opt}}(c_i)=1$ and
$\lambda=g(c_{j^{\ast}})$. Combining with \eqref{eq:exactmin} gives
\eqref{eq:dualvalue}. In particular $-\langle\mathbf{b},\mathbf{u}\rangle-v$
is a lower bound on the minimum that it equals, and a downward-rounded
evaluation of $\langle g,\mathbf{q}^{\mathrm{opt}}\rangle$ is at most
$\langle g,\mathbf{q}^{\mathrm{opt}}\rangle$, which is the minimum by
\eqref{eq:exactmin}.
\end{proof}

% claim statistics/rem:box-oracle-consumer
\claimsummary{statistics/rem:box-oracle-consumer}{Over the canonical box the minimum of a linear functional that the consuming engine certifies is computable exactly in closed form, with no linear-programming solve and a dual witness for re-verification. A linear minimisation oracle over the box can be implemented as a sort plus a prefix sum.}
\claimbookmark{statistics/rem:box-oracle-consumer}{Remark}{rem:box-oracle-consumer}{The closed-form box minimum for the consuming engine}
\begin{remark}[The closed-form box minimum for the consuming engine]\label{rem:box-oracle-consumer}
\xref{meat-core/frank-wolfe/rem:acc-side} lets the certificate consume, in
place of the pairing with a point attaining the minimum, a number $\bacc$
with
\begin{equation}\label{eq:bacc-shape}
  \bacc\ \le\ \min_{\mathbf{q}\in\acSet}\ \sum_{\hat c\in\Chacc}g(\hat c)\,
  q(\hat c)
\end{equation}
for a given exact $g$, and for a polytope it computes the minimum by a
linear program solved exactly over the rationals. For the representation
this module emits the minimum is computable in closed form, exactly, by
\eqref{eq:greedy}: no dual search and no linear-programming solve is needed,
and the dual witness of Theorem~\ref{thm:box-oracle}\,(iii) is available for
independent re-verification in exactly the shape that remark's
instantiation~(b) expects. Two consequences follow. First, the acceptance-side
minimisation is not a source of numerical error for this geometry: besides the
consumer's outward rounding of the box bounds when it reads the set
(\xref{foundations/req:consumer-decides-exactly}), which leaves a non-empty
canonical box containing the emitted one, to which
Theorem~\ref{thm:box-oracle} applies again, the only rounding is the final
downward evaluation of a rational number, which keeps a lower bound a lower
bound (\xref{foundations/req:rounding-direction}). Second,
a consumer's linear minimisation oracle over $\acSet$ may be implemented as a
sort plus a prefix sum.
\end{remark}
\claimmeta{\metalabel{Uses}\mref{Thm.}{thm:box-oracle}
\metasep newly derived
\metasep accepted}

\section{The honest envelope}\label{group-statistics/def:envelope}
\sectionintro{The honest model is an \termref{term-honest-envelope}{honest envelope}: a lower and an upper honest probability per outcome, whose \termref{term-envelope-set}{envelope set} holds every law between them. One honest law is the \termref{term-degenerate-envelope}{degenerate envelope}, and assuming that every round's conditional outcome law lies in a degenerate envelope is the same as assuming independent rounds with that law (\hyperlink{claim-statistics/lem:hon-degenerate}{Lemma~\ref*{lem:hon-degenerate}}). A \termref{term-conservative-evaluation}{conservative evaluation} takes the envelope, the budget and its allocation as exact rationals and decides every tail comparison exactly or against an upward-rounded bound (\hyperlink{claim-statistics/def:conservative-evaluation}{Definition~\ref*{def:conservative-evaluation}}).}

% claim statistics/def:envelope
\claimsummary{statistics/def:envelope}{An honest envelope is a lower and an upper exact rational probability per outcome, and its envelope set holds every law in the simplex between them. A single honest law is the degenerate envelope, whose envelope set is that law alone.}
\claimbookmark{statistics/def:envelope}{Definition}{def:envelope}{Honest envelope and its envelope set}
\claimleadin{%
The honest model is a range of honest probabilities per outcome rather than a
single law, and a single law is the special case of ranges of width zero.
}
\begin{definition}[Honest envelope and its envelope set]\label{def:envelope}
An \termdef{term-honest-envelope}{honest envelope} on $\Chatset$ is a pair
${\honest}^{-},{\honest}^{+}\in\mathbb{Q}^{\Chatset}$ of exact rational
vectors with
\begin{equation}\label{eq:envelope}
  0\ \le\ {\honest}^{-}(\hat c)\ \le\ {\honest}^{+}(\hat c)\ \le\ 1
  \quad\forall\hat c\in\Chatset,
  \qquad
  \sum_{\hat c\in\Chatset}{\honest}^{-}(\hat c)\ \le\ 1\ \le\
  \sum_{\hat c\in\Chatset}{\honest}^{+}(\hat c).
\end{equation}
Its \termdef{term-envelope-set}{envelope set} is
\begin{equation}\label{eq:envelope-set}
  \mathsf{H}\ \coloneqq\ \big\{\mathbf{p}\in\simplex{\Chatset}:\
  {\honest}^{-}(\hat c)\le p(\hat c)\le{\honest}^{+}(\hat c)\ \
  \forall\hat c\in\Chatset\big\},
\end{equation}
the set of honest laws the envelope admits. An honest envelope with
${\honest}^{-}={\honest}^{+}$ is a \termdef{term-degenerate-envelope}{degenerate envelope}; then
\eqref{eq:envelope} forces the common value to be a law
$\honest\in\simplex{\Chatset}\cap\mathbb{Q}^{\Chatset}$ and
$\mathsf{H}=\{\honest\}$. A single honest law $\honest$ (the
\termdef{term-point-form}{point form}) is read as the degenerate envelope
$(\honest,\honest)$. The pair $({\honest}^{-},{\honest}^{+})$ is a
\termref{term-canonical-box-form}{canonical box form} in the sense of Definition~\ref{def:taxonomy}
(rational, $0\le{\honest}^{-}\le{\honest}^{+}$), $\mathsf{H}$ is its
denotation \eqref{eq:boxset}, and the second condition of
\eqref{eq:envelope} is its closed-form feasibility condition
\eqref{eq:boxfeas}.
\end{definition}
\claimmeta{\metalabel{Uses}\mref{Def.}{def:taxonomy}
\metasep \metalabel{Used by}\mref{Lem.}{lem:hon-degenerate}, \mref{Def.}{def:conservative-evaluation}, \mref{Ass.}{ass:hon-env}, \mref{Ass.}{ass:reach}, \mref{Def.}{def:endpoint-rule}, \mref{Lem.}{lem:guard}, \mref{Def.}{def:interval}, \mref{Def.}{def:construction}, \mref{Thm.}{thm:box-construction}, \mref{Thm.}{thm:completeness-bound}, \mref{Rem.}{rem:envelope}, \mref{Thm.}{thm:honest-containment}, \mref{Lem.}{lem:widening}, \mref{Thm.}{thm:monotonicity-limits}, \mref{Rem.}{rem:singleton-honest}
\metasep newly derived
\metasep implemented
\metasep \metalabel{Code}\coderef{https://github.com/fraunhofersg/certiqs-key/blob/v0.1.0/src/certiqs_key/statistics/construction.py}{certiqs_key.statistics.construction:HonestEnvelope}}

% claim statistics/lem:hon-degenerate
\claimsummary{statistics/lem:hon-degenerate}{Every round's conditional honest law lying in a degenerate envelope is the same as independent, identically distributed honest rounds with that one law. The envelope assumption therefore generalises the single-law assumption.}
\claimbookmark{statistics/lem:hon-degenerate}{Lemma}{lem:hon-degenerate}{The envelope assumption at a degenerate envelope}
\claimleadin{%
The completeness statements read one of two honest-round premises: independent
rounds with one law (Assumption~\ref{ass:hon}), or rounds whose conditional
laws stay inside a declared envelope (Assumption~\ref{ass:hon-env}). At a
degenerate envelope the two premises coincide.
}
\begin{lemma}[The envelope assumption at a degenerate envelope]\label{lem:hon-degenerate}
Let $\honest\in\simplex{\Chatset}\cap\mathbb{Q}^{\Chatset}$ and let
$\Chat_1,\dots,\Chat_n$ be $\Chatset$-valued random variables on one
probability space. The following are equivalent:
\begin{enumerate}[label=(\alph*),leftmargin=2.2em]
\item the statement of Assumption~\ref{ass:hon-env} holds for the degenerate
envelope $(\honest,\honest)$ of Definition~\ref{def:envelope}: there is a
filtration $\mathcal{F}_0\subseteq\dots\subseteq\mathcal{F}_n$ with $\Chat_j$
$\mathcal{F}_j$-measurable and
$\Pr[\Chat_j=\hat c\,|\,\mathcal{F}_{j-1}]=\honest(\hat c)$ almost surely,
for every $j$ and $\hat c$;
\item the statement of Assumption~\ref{ass:hon} holds: $\Chat_1,\dots,\Chat_n$
are independent and identically distributed with law $\honest$.
\end{enumerate}
\end{lemma}
\claimmeta{\metalabel{Uses}\mref{Def.}{def:envelope}
\metasep newly derived
\metasep accepted}

\begin{proof}
With a degenerate envelope, \eqref{eq:hon-env} reads
$\honest(\hat c)\le\Pr[\Chat_j=\hat c\,|\,\mathcal{F}_{j-1}]\le\honest(\hat
c)$, i.e.\ the equality displayed in (a).

\emph{(a)$\Rightarrow$(b).} Fix $j$. For every $A\in\mathcal{F}_{j-1}$ and
every $\hat c$, the defining property of conditional probability gives
\begin{equation*}
  \Pr\big[\{\Chat_j=\hat c\}\cap A\big]
  =\mathbb{E}\big[\mathbf{1}_A\,\Pr[\Chat_j=\hat c\,|\,\mathcal{F}_{j-1}]\big]
  =\honest(\hat c)\,\Pr[A].
\end{equation*}
Taking $A$ the whole space shows $\Pr[\Chat_j=\hat c]=\honest(\hat c)$, so
the display says $\Pr[\{\Chat_j=\hat c\}\cap A]=\Pr[\Chat_j=\hat c]\Pr[A]$:
$\Chat_j$ is independent of $\mathcal{F}_{j-1}$ and has law $\honest$. Since
$\Chat_i$ is $\mathcal{F}_i$-measurable and $\mathcal{F}_i\subseteq
\mathcal{F}_{j-1}$ for $i<j$, the event
$A=\{\Chat_1=\hat c_1,\dots,\Chat_{j-1}=\hat c_{j-1}\}$ lies in
$\mathcal{F}_{j-1}$, whence
$\Pr[\Chat_1=\hat c_1,\dots,\Chat_j=\hat c_j]
=\honest(\hat c_j)\,\Pr[\Chat_1=\hat c_1,\dots,\Chat_{j-1}=\hat c_{j-1}]$
for all $\hat c_1,\dots,\hat c_j\in\Chatset$. Induction on $j$ from $1$ to $n$
gives $\Pr[\Chat_1=\hat c_1,\dots,\Chat_n=\hat c_n]=\prod_{j=1}^{n}\honest(\hat
c_j)$, the product law: the $\Chat_j$ are independent, each with law
$\honest$.

\emph{(b)$\Rightarrow$(a).} Take the natural filtration: $\mathcal{F}_0$ the
trivial $\sigma$-algebra and $\mathcal{F}_j\coloneqq\sigma(\Chat_1,\dots,
\Chat_j)$. Each $\Chat_j$ is $\mathcal{F}_j$-measurable. By independence,
$\Chat_j$ is independent of $\mathcal{F}_{j-1}$, so for every
$A\in\mathcal{F}_{j-1}$,
$\mathbb{E}[\mathbf{1}_A\mathbf{1}\{\Chat_j=\hat c\}]=\honest(\hat c)\Pr[A]
=\mathbb{E}[\mathbf{1}_A\,\honest(\hat c)]$; the constant $\honest(\hat c)$
is $\mathcal{F}_{j-1}$-measurable, hence a version of
$\Pr[\Chat_j=\hat c\,|\,\mathcal{F}_{j-1}]$, and (a) holds.
\end{proof}

% claim statistics/def:conservative-evaluation
\claimsummary{statistics/def:conservative-evaluation}{An evaluation of the construction is conservative when the honest envelope, the completeness budget and its allocation are exact rationals and every tail comparison is decided exactly or against an upward-rounded tail. It is the arithmetic under which rounding can only widen the acceptance set.}
\claimbookmark{statistics/def:conservative-evaluation}{Definition}{def:conservative-evaluation}{Exact data and conservative evaluation}
\claimleadin{%
With exact data the construction of the acceptance set lies on the
\exttermref{../foundations/contract.pdf}{term-exact-island}{exact island}, the computations the engine must carry out exactly
(\xref{foundations/def:exact-island}\,(ii)). Rounding in
the safe direction enlarges $\acSet$, which lowers the consumer's $\hlow$, an
infimum over a superset being at most the infimum over a subset
(Lemma~\ref{lem:widening}); it costs rate, never soundness. The consuming
engine's outward rounding of the emitted bounds is such a rounding: it is the
one alteration of emitted data that
\xref{foundations/req:consumer-decides-exactly} allows a consumer, which
still decides every predicate on a replaced bound exactly. The engine refuses
inexact data (Requirement~\ref{req:input-refusal}) and decides its tail
comparisons in this way on each evaluation
(Requirements~\ref{req:rule-in-force} and~\ref{req:enclosure-substrate}).
}
\begin{definition}[Exact data and conservative evaluation]\label{def:conservative-evaluation}
An evaluation of the acceptance-set construction is a
\termdef{term-conservative-evaluation}{conservative evaluation} when the honest envelope's vectors
${\honest}^{-},{\honest}^{+}$ (Definition~\ref{def:envelope}; for a single
declared law, $\honest$ itself), $\ecomp$ and the allocation are exact
rationals, and every tail comparison an endpoint rule makes
(Definition~\ref{def:endpoint-rule}; for the exact rule, those of
Definition~\ref{def:interval}) is decided either exactly or against a bound
that rounds the tail \emph{upward}.
\end{definition}
\claimmeta{\metalabel{Uses}\mref{Def.}{def:envelope}
\metasep \metalabel{Used by}\mref{Prop.}{prop:rule-exact}, \mref{Thm.}{thm:box-construction}, \mref{Thm.}{thm:completeness-bound}, \mref{Rem.}{rem:epsstar}, \mref{Thm.}{thm:honest-containment}, \mref{Thm.}{thm:monotonicity-limits}, \mref{Prop.}{prop:rule-bernstein}, \mref{Rem.}{rem:completeness-on-rules}, \mref{Prop.}{prop:bracketed-endpoints}, \mref{Rem.}{rem:local-evaluation}
\metasep newly derived
\metasep implemented
\metasep \metalabel{Code}\coderef{https://github.com/fraunhofersg/certiqs-key/blob/v0.1.0/src/certiqs_key/statistics/exact.py}{certiqs_key.statistics.exact:as_exact_rational}}

\section{Assumptions}\label{group-statistics/ass:hon}
\sectionintro{These are the premises this module declares, and none is evidence about an adversary: honest rounds, independent with one law or conditionally inside a declared envelope; an acceptance value recorded in every round; and a law of the envelope that some feasible state of the consuming declaration produces. A failure of the honest-round premises costs completeness, rate or usability, never the soundness of a key. A failure of that last premise means that the declared devices do not describe what produced the law, and a key length computed for them can grow with the mismatch; the consuming engine checks a run whose honest law is supplied or whose declaration carries a constraint the user declared, unless the run opts out, states the cap a failure it proves puts on the probability that devices as declared pass the acceptance test, and refuses the run where that cap is at most a threshold, one half by default. A round missing from the record normalises the counts against the wrong round count.}

% claim statistics/ass:hon
\claimsummary{statistics/ass:hon}{On an honest run the acceptance statistics are independent and identically distributed with the declared law. The premise bears on the advertised honest abort probability and restricts no adversary.}
\claimbookmark{statistics/ass:hon}{Assumption}{ass:hon}{Honest i.i.d. rounds}
\claimleadin{%
The premise fixes one honest law for every round and narrows
Assumption~\ref{ass:hon-env} to the \termref{term-degenerate-envelope}{degenerate envelope} of that law. It
is asserted about the honest channel only, never against an adversary.
Physically it says that the honest channel acts identically and independently
in every round and that there is no memory across rounds: the sameness and
the no-memory clause of a \exttermref{../foundations/contract.pdf}{term-round-homogeneous}{round-homogeneous} family, asserted of the
honest channel.
}
\SetShortId{A-HON}
\begin{assumption}[Honest i.i.d. rounds]\label{ass:hon}
(i) On an honest run the acceptance statistics $\Chat_1,\dots,\Chat_n$ are
independent and identically distributed with law $\honest$ on $\Chatset$.
\end{assumption}
\claimverification{The engine does not verify this assumption; a run that rests on it records it as unverified.}
\claimreading{Without an eavesdropper, every round's acceptance outcome is drawn independently from the same declared distribution. If you expect drift between or within runs, declare an honest envelope instead. This premise only bears on the advertised probability that an honest run aborts: if it fails, an honest run may abort more or less often than advertised. Security is unaffected.}
\claimmeta{\metalabel{Category}honest model
\metasep \metalabel{Affects}completeness
\metasep \metalabel{Refines}\aref{A-HON-ENV}{ass:hon-env}
\metasep \metalabel{Assumed by}\mref{Rem.}{rem:sound-indep}, \mref{Thm.}{thm:singleton-completeness}
\metasep \metalabel{Machine key}\iface{honest_rounds_iid}
\metasep newly derived
\metasep implemented
\metasep \metalabel{Code}\coderef{https://github.com/fraunhofersg/certiqs-key/blob/v0.1.0/src/certiqs_key/statistics/construction.py}{certiqs_key.statistics.construction:HONEST_MODEL_TAG}}

% claim statistics/ass:hon-env
\claimsummary{statistics/ass:hon-env}{On an honest run each round's probability of each acceptance outcome, given everything before it, stays inside the declared honest envelope. The premise bears on the advertised honest abort probability and restricts no adversary.}
\claimbookmark{statistics/ass:hon-env}{Assumption}{ass:hon-env}{Honest rounds inside a declared envelope}
\claimleadin{%
The premise is asserted about the honest channel only, never against an
adversary. Under it the honest channel may change from round to round within
a run, from run to run, and in a way that depends on past outcomes or on
further honest variables recorded in $\mathcal{F}_{j-1}$, as long as each
round's conditional law stays inside the envelope.
}
\SetShortId{A-HON-ENV}
\begin{assumption}[Honest rounds inside a declared envelope]\label{ass:hon-env}
Let $({\honest}^{-},{\honest}^{+})$ be the declared \termref{term-honest-envelope}{honest envelope}
(Definition~\ref{def:envelope}). On an honest run there is a filtration
$\mathcal{F}_0\subseteq\mathcal{F}_1\subseteq\dots\subseteq\mathcal{F}_n$ on
the underlying probability space such that
(i) $\Chat_j$ is $\mathcal{F}_j$-measurable for every $j\in\{1,\dots,n\}$;
(ii) for every such $j$ and every $\hat c\in\Chatset$, almost surely,
\begin{equation}\label{eq:hon-env}
  {\honest}^{-}(\hat c)\ \le\
  \Pr\big[\Chat_j=\hat c\ \big|\ \mathcal{F}_{j-1}\big]\ \le\
  {\honest}^{+}(\hat c).
\end{equation}
\end{assumption}
\claimverification{The engine does not verify this assumption; a run that rests on it records it as unverified.}
\claimreading{When no attacker is present, each round's probability of each acceptance outcome, given everything that happened before, stays inside the honest model: between the lower and upper bounds of the honest envelope where you declared one, and equal to the honest distribution where you gave a single one instead (an envelope of zero width). With an envelope the channel may drift within and between runs as long as it stays inside it; with a single distribution it may not drift at all. This premise only bears on the advertised probability that an honest run aborts; it never affects security. If drift can leave the honest model, declare a wider envelope.}
\claimmeta{\metalabel{Uses}\mref{Def.}{def:envelope}
\metasep \metalabel{Category}honest model
\metasep \metalabel{Affects}completeness
\metasep \metalabel{Assumed by}\mref{Thm.}{thm:completeness-bound}, \mref{Rem.}{rem:sound-indep}
\metasep \metalabel{Machine key}\iface{honest_rounds_within_envelope}
\metasep newly derived
\metasep implemented
\metasep \metalabel{Code}\coderef{https://github.com/fraunhofersg/certiqs-key/blob/v0.1.0/src/certiqs_key/statistics/construction.py}{certiqs_key.statistics.construction:ENVELOPE_MODEL_TAG}}

% claim statistics/ass:alphabet
\claimsummary{statistics/ass:alphabet}{In every one of the n rounds of the run an acceptance value of the alphabet is recorded. The counts and the frequency vector then describe the whole run, normalised by its round count.}
\claimbookmark{statistics/ass:alphabet}{Assumption}{ass:alphabet}{An acceptance value recorded in every round}
\SetShortId{A-ALPHABET}
\begin{assumption}[An acceptance value recorded in every round]\label{ass:alphabet}
(i) In every one of the $n$ rounds of the run, an acceptance value in
$\Chatset$ is recorded.
\end{assumption}
\claimverification{The engine does not verify this assumption; a run that rests on it records it as unverified.}
\claimreading{Every one of the n rounds of your run records an acceptance value from the declared acceptance alphabet: no round counted in n is missing from the record. If rounds go missing (for example through lost time tags), the counts are normalised against the wrong number of rounds and the acceptance test does not describe the run. A round missing from the record also puts security at stake: the acceptance test must be evaluated on the acceptance values of all n rounds (\hyperref{../meat-core/contract.pdf}{}{claim-meat-core/ass:conform}{A-CONFORM}), and a measured round the record drops is a sifting step the declaration does not describe (\hyperref{../protocol-compiler/contract.pdf}{}{claim-protocol-compiler/ass:spec}{A-SPEC}). Where the engine builds the acceptance set, it checks that the alphabet, its order and n are those of your protocol declaration; a set carried in the declaration is read in the declaration's order, so writing it in that order is yours to check.}
\claimmeta{\metalabel{Category}protocol structure
\metasep \metalabel{Affects}secrecy, completeness
\metasep \metalabel{Assumed by}\mref{Thm.}{thm:geometry-validation}, \mref{Thm.}{thm:box-oracle}, \mref{Rem.}{rem:guard-activation}, \mref{Thm.}{thm:box-construction}, \mref{Thm.}{thm:completeness-bound}, \mref{Thm.}{thm:honest-containment}, \mref{Thm.}{thm:monotonicity-limits}, \mref{Thm.}{thm:singleton-completeness}
\metasep \metalabel{Machine key}\iface{acceptance_recorded_every_round}
\metasep newly derived
\metasep accepted}

% claim statistics/ass:reach
\claimsummary{statistics/ass:reach}{Some feasible state of the consuming declaration produces a law of the declared honest envelope, for a single declared law that law itself. It rules out an incompatibility of the honest model with the declaration.}
\claimbookmark{statistics/ass:reach}{Assumption}{ass:reach}{An honest law produced by a feasible state}
\claimleadin{%
The premise is an obligation of whoever declares the honest model, together
with the declaration: it is a joint property of the declared
\termref{term-honest-envelope}{honest envelope} and of the protocol declaration the \exttermref{../foundations/contract.pdf}{term-consumer}{consumer}
solves.
}
\SetShortId{A-REACH}
\begin{assumption}[An honest law produced by a feasible state]\label{ass:reach}
(i) There exists $\omega^{\mathrm{hon}}\in\Feas$, the feasible set of the
consuming declaration (\xref{protocol-compiler/def:feasible}), with
$\Phimap[\omega^{\mathrm{hon}}]\in\mathsf{H}$, the \termref{term-envelope-set}{envelope set} of the
declared honest envelope (Definition~\ref{def:envelope}); for a single
declared law $\honest$ (the degenerate envelope, $\mathsf{H}=\{\honest\}$)
this reads $\Phimap[\omega^{\mathrm{hon}}]=\honest$.
\end{assumption}
\claimverification{The engine does not verify this assumption; a run that rests on it records it as unverified.}
\claimreading{Your honest model fits your protocol declaration: a state the declaration admits as feasible produces acceptance statistics in your declared honest envelope, or equal to your single declared honest law. If it fails, the declared devices do not describe the statistics and a key length for them is not warranted. The engine's own honest models are built to match the declared Alice marginal; nothing proves this premise for them. Unless switched off, the engine checks every run whose honest statistics you supply or whose declaration carries a constraint you declared: from a proven mismatch it states a cap on the probability that devices behaving as the declaration says pass your acceptance test, and refuses the run where that cap is at most one half, a threshold you can set. On an unchecked run a mismatch goes undetected and this premise is yours to vouch for.}
\claimmeta{\metalabel{Uses}\mref{Def.}{def:envelope}
\metasep \metalabel{Category}honest model
\metasep \metalabel{Affects}secrecy, completeness
\metasep \metalabel{Assumed by}\mref{Thm.}{thm:honest-containment}, \mref{Rem.}{rem:hlow-consequence}
\metasep \metalabel{Machine key}\iface{honest_reachable}
\metasep newly derived
\metasep accepted}

\clearpage
\part{Engine requirements}
\sectionintro{What an implementation of this module must do, record or refuse; each requirement names the result that motivates it.}

\section{Requirements: the emitted geometry}\label{group-statistics/req:geometry-tag}
\sectionintro{The engine emits every acceptance set under one of the two geometry tags with that tag's data, and a consumer refuses a set under any other tag. It records a canonical box's lower and upper bounds beside the rows they determine. As the acceptance support it emits the outcomes with a positive upper bound, a superset never substituted for the support itself, and for a singleton the support of its vector. A set is docked only on the alphabet and round count of the declaration its consumer solves: a built set is compared with it, a set the declaration carries is read in its enumeration, and a mismatch produces no result.}

% claim statistics/req:geometry-tag
\claimsummary{statistics/req:geometry-tag}{Every acceptance set is emitted under one of the two geometry tags with that tag's data. A consumer presented with any other tag refuses the set rather than guessing what it means.}
\claimbookmark{statistics/req:geometry-tag}{Requirement}{req:geometry-tag}{Acceptance sets under exactly the two geometry tags}
\claimleadin{%
Definition~\ref{def:taxonomy} admits two tags and no third,
Theorem~\ref{thm:geometry-validation} shows what a validated set under either tag is, and
\xref{meat-core/frank-wolfe/rem:acc-side} certifies the minimum of a linear
functional over the acceptance set for exactly these two forms.
}
\begin{requirement}[Acceptance sets under exactly the two geometry tags]\label{req:geometry-tag}
The engine emits every acceptance set under one of the two tags of
Definition~\ref{def:taxonomy} (\iface{geometry_tag}) together with the data of
that tag, and a \exttermref{../foundations/contract.pdf}{term-consumer}{consumer} presented with an acceptance set under any
other tag refuses it rather than guessing.
\end{requirement}
\claimmeta{\metalabel{Motivated by}\mref{Def.}{def:taxonomy}, \mref{Thm.}{thm:geometry-validation}, \extref{../meat-core/contract.pdf}{meat-core/frank-wolfe/rem:acc-side}
\metasep implemented
\metasep \metalabel{Code}\coderef{https://github.com/fraunhofersg/certiqs-key/blob/v0.1.0/src/certiqs_key/statistics/taxonomy.py}{certiqs_key.statistics.taxonomy:GeometryTag}}

% claim statistics/req:box-bounds-key
\claimsummary{statistics/req:box-bounds-key}{A geometry in canonical box form is emitted with its lower and upper bound vectors recorded beside the rows they determine. The recorded pair is what licenses the closed-form box oracle.}
\claimbookmark{statistics/req:box-bounds-key}{Requirement}{req:box-bounds-key}{Recorded bounds of a canonical box}
\claimleadin{%
A consumer that reads the recorded pair can minimise over the box in closed
form (Theorem~\ref{thm:box-oracle}); one that ignores it still reads a well-formed
polytope (Remark~\ref{rem:two-tags}).
}
\begin{requirement}[Recorded bounds of a canonical box]\label{req:box-bounds-key}
The engine emits a geometry in \termref{term-canonical-box-form}{canonical box form} with its pair
$(\qlo,\qhi)$ recorded under \iface{box_bounds}, beside the rows
\eqref{eq:boxrows} that the pair determines.
\end{requirement}
\claimmeta{\metalabel{Motivated by}\mref{Def.}{def:taxonomy}, \mref{Rem.}{rem:two-tags}, \mref{Thm.}{thm:box-oracle}
\metasep implemented
\metasep \metalabel{Code}\coderef{https://github.com/fraunhofersg/certiqs-key/blob/v0.1.0/src/certiqs_key/statistics/taxonomy.py}{certiqs_key.statistics.taxonomy:canonical_box}}

% claim statistics/req:acceptance-support
\claimsummary{statistics/req:acceptance-support}{For a canonical box the engine emits the outcomes with a positive upper bound, a superset of the acceptance support that is never substituted for it; for a singleton it emits the support of the vector, which is the acceptance support itself.}
\claimbookmark{statistics/req:acceptance-support}{Requirement}{req:acceptance-support}{The emitted acceptance support}
\claimleadin{%
Theorem~\ref{thm:box-oracle}\,(i) shows that the outcomes with a positive upper bound
contain the support of every point of the box, and that they can contain more
than the union of those supports.
}
\begin{requirement}[The emitted acceptance support]\label{req:acceptance-support}
The engine emits under \iface{acceptance_support} the set
$\Chatset^{+}=\{\hat c\in\Chatset:q^{+}(\hat c)>0\}$ of
Theorem~\ref{thm:box-oracle}\,(i) for a geometry in \termref{term-canonical-box-form}{canonical box form}, and
the set $\operatorname{supp}\mathbf{q}^{\ast}$ for a $\mathrm{singleton}$
geometry. For a box the emitted set is a superset of the acceptance support
$\Chacc$ of \xref{meat-core/frank-wolfe/def:first-order-setting} and is never to be
substituted for it; for a singleton the emitted set and $\Chacc$ coincide.
\end{requirement}
\claimmeta{\metalabel{Motivated by}\mref{Thm.}{thm:box-oracle}, \extref{../meat-core/contract.pdf}{meat-core/frank-wolfe/def:first-order-setting}
\metasep implemented (effective: accepted)
\metasep \metalabel{Code}\coderef{https://github.com/fraunhofersg/certiqs-key/blob/v0.1.0/src/certiqs_key/statistics/taxonomy.py}{certiqs_key.statistics.taxonomy:acceptance_support}}

% claim statistics/req:matched-alphabet
\claimsummary{statistics/req:matched-alphabet}{An acceptance set is docked only on the alphabet, in its declared enumeration, and the round count of the declaration its consumer solves: a built set is compared with the declaration first, a set the declaration carries is read in that enumeration. A disagreement produces no result.}
\claimbookmark{statistics/req:matched-alphabet}{Requirement}{req:matched-alphabet}{Alphabet and round count taken from the consuming declaration}
\claimleadin{%
Every result of this module that reads an alphabet and a round count reads
them as a \termref{term-matched-alphabet-and-round-count}{matched alphabet and round count}
(Definition~\ref{def:matched-alphabet}); the match is a relation between
artefacts and not a property of an input, so it is decided where both
artefacts are available.
}
\begin{requirement}[Alphabet and round count taken from the consuming declaration]\label{req:matched-alphabet}
The engine docks an acceptance set into the declaration a \exttermref{../foundations/contract.pdf}{term-consumer}{consumer}
solves only on a \termref{term-matched-alphabet-and-round-count}{matched alphabet and round count}
(Definition~\ref{def:matched-alphabet}).
\begin{enumerate}[label=(\roman*),leftmargin=2.2em]
\item A set it builds, it builds on the acceptance alphabet, in its declared
  enumeration, of the declaration the consumer solves and on the round count
  of the one run parameter the consumer reads, and before the set is docked
  it compares the set's alphabet and enumeration with the declaration's. A
  disagreement is a hard error of the consumer's input validation: it is a
  refusal code of neither module, nothing is repaired, and no result is
  produced.
\item A set that the declaration itself carries, with no alphabet of its
  own, it reads positionally: its $k$-th coordinate is the $k$-th acceptance
  value in the declaration's declared enumeration, and its round count is
  the one run parameter the consumer reads, as in (i). A set whose number of coordinates differs
  from the size of the acceptance alphabet is refused as ill formed, and no
  result is produced.
\end{enumerate}
\end{requirement}
\claimmeta{\metalabel{Motivated by}\mref{Def.}{def:matched-alphabet}
\metasep accepted
\metasep \metalabel{Code}\coderef{https://github.com/fraunhofersg/certiqs-key/blob/v0.1.0/src/certiqs_key/runs/pipeline.py}{certiqs_key.runs.pipeline:_require_matched_alphabet}}

\clearpage
\part{Binomial box construction}
\sectionintro{This submodule turns an honest envelope, a completeness budget and its allocation over selected outcomes into a box of accepted frequencies. An admissible endpoint rule returns, per outcome, a lower count at the envelope's lower value and an upper count at its upper value whose tails meet their budgets, and a guard widens them to contain the whole envelope interval. Its main results are that for every such rule the construction map emits a canonical box of exact rationals that contains every law of the envelope set and meets the consuming engine's acceptance assumption (\hyperlink{claim-statistics/binomial-box/thm:box-construction}{Theorem~\ref*{thm:box-construction}}), and that an honest run whose per-round conditional outcome law stays inside the envelope fails the test with probability at most the sum of the binomial tails at the envelope's ends beyond the guarded endpoints, which the allocation keeps within the budget (\hyperlink{claim-statistics/binomial-box/thm:completeness-bound}{Theorem~\ref*{thm:completeness-bound}}); neither independent rounds nor independent outcome counts are assumed. The construction never fails on an admissible input, so every refusal comes from the input. Under a monotone allocation rule and an endpoint rule monotone in the allocation, which the enclosure evaluation of the exact rule is not claimed to be, a larger budget gives a smaller or equal box; for the exact rule the box is not always nested as the round count grows. The exact binomial quantile rule is the default; the Bernstein rule gives closed-form endpoints whose unrounded half-width decreases strictly with the round count, and a local evaluation of the exact rule decides its tail comparisons on enclosures in log space. Completeness is a usability guarantee, and soundness never rests on the honest-round assumptions.}

\section{Admissible inputs, endpoint rules and the guard}\label{group-statistics/binomial-box/def:endpoint-rule}
\sectionintro{An \termref{term-admissible-input}{admissible input} fixes the alphabet, the round count, the \termref{term-honest-envelope}{honest envelope}, the completeness budget, its allocation over the constrained outcomes, and the choice of rule and evaluation, the numbers exact. An endpoint rule returns, for one binomial count, a lower and an upper endpoint whose tails meet their budgets, the lower taken at the envelope's lower value and the upper at its upper value, and the \termref{term-honest-containment-guard}{honest-containment guard} widens them to contain the whole envelope interval. \hyperlink{claim-statistics/binomial-box/lem:guard}{Lemma~\ref*{lem:guard}} proves for every rule that the \termref{term-guarded-endpoints}{guarded endpoints} are ordered around that interval, still meet their budgets, widen the rule's interval by an exactly known amount, and lie within 1/n of the envelope value wherever the guard acts.}

% claim statistics/binomial-box/def:endpoint-rule
\claimsummary{statistics/binomial-box/def:endpoint-rule}{An admissible endpoint rule returns count endpoints whose binomial tails meet their budgets. On an admissible input the lower endpoint is taken at the envelope's lower value, the upper at its upper value, and a guard widens both to contain the envelope interval.}
\claimbookmark{statistics/binomial-box/def:endpoint-rule}{Definition}{def:endpoint-rule}{Admissible inputs, endpoint rules and guarded endpoints}
\claimleadin{%
The construction turns the two tail budgets of each constrained outcome into
an interval of counts. An endpoint rule does this for one binomial count; the
definition fixes what every rule must guarantee, the inputs on which the
construction runs, and the guard that widens a rule's interval to contain the
\termref{term-honest-envelope}{honest envelope}. The lower endpoint is taken at the envelope's lower
value and the upper endpoint at its upper value because, under
Assumption~\ref{ass:hon-env}, the binomial laws at those two values bound the
count of Definition~\ref{def:freq} from below and from above
(Lemma~\ref{lem:envelope-dominance}). The construction and the theorems about
it take the admissible endpoint rule $R$ as a parameter, independent of the
choice of rule and evaluation an input carries in clause (viii), and state as
a hypothesis any further property of $R$ they need; that choice selects the
rule in force (Requirement~\ref{req:rule-in-force}), and a choice outside
clause (viii) is violation (11) of
Theorem~\ref{thm:honest-containment}\,(iii).
}
\begin{definition}[Admissible inputs, endpoint rules and guarded endpoints]\label{def:endpoint-rule}
Fix $n\ge1$ and let $\Ncount{}$ be binomial with parameters $(n,p)$, where
$p\in[0,1]$ is rational. An \termdef{term-admissible-endpoint-rule}{admissible endpoint rule} $R$ assigns to every
$(n,p,\varepsilon^{-},\varepsilon^{+})$ with rational
$\varepsilon^{-},\varepsilon^{+}\in[0,1)$ two integers
$R^{-}=R^{-}(n,p,\varepsilon^{-})$ and $R^{+}=R^{+}(n,p,\varepsilon^{+})$
such that
\begin{enumerate}[label=(\alph*),leftmargin=2.2em]
\item \emph{(admissibility)} $R^{-},R^{+}\in\{0,\dots,n\}$ and
\begin{equation}\label{eq:rule-admissible}
  \Pr\big[\Ncount{}<R^{-}\big]\le\varepsilon^{-},
  \qquad
  \Pr\big[\Ncount{}>R^{+}\big]\le\varepsilon^{+};
\end{equation}
\item \emph{(zero budget)} if $0<p<1$, then $R^{-}=0$ at $\varepsilon^{-}=0$
and $R^{+}=n$ at $\varepsilon^{+}=0$;
\item \emph{(degenerate mass)} $R^{-}=R^{+}=0$ at $p=0$ and $R^{-}=R^{+}=n$
at $p=1$, for every $\varepsilon^{\pm}$.
\end{enumerate}
The rule is \termdef{term-monotone-in-the-allocation}{monotone in the allocation} when, the other arguments fixed,
$\varepsilon^{-}\mapsto R^{-}$ is non-decreasing and
$\varepsilon^{+}\mapsto R^{+}$ is non-increasing.

An \termdef{term-admissible-input}{admissible input} consists of
\begin{enumerate}[label=(\roman*),leftmargin=2.2em]
\item a finite, non-empty acceptance alphabet $\Chatset$ with a fixed
enumeration;
\item a round count, an integer $n\ge1$;
\item an honest input on $\Chatset$: exactly one of a law
$\honest\in\simplex{\Chatset}$, the \termref{term-point-form}{point form}, and an honest envelope
$({\honest}^{-},{\honest}^{+})$ (Definition~\ref{def:envelope}); a law in the
point form is read as the \termref{term-degenerate-envelope}{degenerate envelope} $(\honest,\honest)$;
\item a completeness budget $\ecomp\in[0,1)$;
\item a set $\Chatcon\subseteq\Chatset$ of constrained outcomes;
\item an allocation
$\varepsilon^{-},\varepsilon^{+}:\Chatcon\to\mathbb{Q}_{\ge0}$ with
$\sum_{\hat c\in\Chatcon}(\varepsilon^{-}(\hat c)+\varepsilon^{+}(\hat
c))\le\ecomp$;
\item exactness: every coordinate of the honest input, $\ecomp$ and every
value of the allocation is an exact rational;
\item a choice of rule, the \termref{term-exact-binomial-quantile-rule}{exact binomial quantile rule} of
Definition~\ref{def:interval} or the Bernstein rule of
Proposition~\ref{prop:rule-bernstein}, and a choice of one of two
evaluations, by the closed forms or by the enclosures of
Proposition~\ref{prop:bracketed-endpoints}, which on the exact rule decides
its tail comparisons: by the enclosures, it gives the rule
$R_{\mathrm{encl}}$ of that proposition in place of the exact rule;
\item consistency with the \exttermref{../foundations/contract.pdf}{term-consumer}{consumer}: $\Chatset$, its enumeration and
$n$ are those of the declaration the consumer solves, a
\termref{term-matched-alphabet-and-round-count}{matched alphabet and round count} (Definition~\ref{def:matched-alphabet}).
\end{enumerate}
The clauses fix the form of the data --- an enumerated set $\Chatset$, an
integer $n$, one honest input that is a vector (the point form) or a pair of
vectors (the envelope form), a number $\ecomp$, a set $\Chatcon$ with a pair
of vectors indexed by it, and a named rule and evaluation --- and conditions
on data of that form.

Fix an admissible input, with its honest envelope
$({\honest}^{-},{\honest}^{+})$. For $\hat c\in\Chatcon$ write
$p^{-}_{\hat c}\coloneqq{\honest}^{-}(\hat c)$ and
$p^{+}_{\hat c}\coloneqq{\honest}^{+}(\hat c)$, let $\Ncount{\hat c}^{-}$ and
$\Ncount{\hat c}^{+}$ be binomial with parameters $(n,p^{-}_{\hat c})$ and
$(n,p^{+}_{\hat c})$ respectively, and abbreviate
$R^{-}\coloneqq R^{-}(n,p^{-}_{\hat c},\varepsilon^{-}(\hat c))$ and
$R^{+}\coloneqq R^{+}(n,p^{+}_{\hat c},\varepsilon^{+}(\hat c))$: the lower
endpoint is the rule's at the envelope's lower value, the upper endpoint the
rule's at its upper value. The \termdef{term-guarded-endpoints}{guarded endpoints} of $R$ at $\hat c$, after the
\termdef{term-honest-containment-guard}{honest-containment guard}, are
\begin{equation}\label{eq:guard}
  k^{-}(\hat c)\coloneqq\min\big\{R^{-},\ \lfloor n\,p^{-}_{\hat c}\rfloor\big\},
  \qquad
  k^{+}(\hat c)\coloneqq\max\big\{R^{+},\ \lceil n\,p^{+}_{\hat c}\rceil\big\}.
\end{equation}
With the degenerate envelope, $p^{-}_{\hat c}=p^{+}_{\hat c}=\honest(\hat c)
\eqqcolon p_{\hat c}$, $\Ncount{\hat c}^{-}$ and $\Ncount{\hat c}^{+}$ have
the one law of $\Ncount{\hat c}$, binomial $(n,p_{\hat c})$, and
\eqref{eq:guard} reads $k^{-}(\hat c)=\min\{R^{-},\lfloor np_{\hat
c}\rfloor\}$, $k^{+}(\hat c)=\max\{R^{+},\lceil np_{\hat c}\rceil\}$ with both
rule endpoints at $p_{\hat c}$.

The \termdef{term-side-wise-reading}{side-wise reading} is the following. By
\eqref{eq:guard}, $k^{-}(\hat c)$ depends on the envelope only through
$p^{-}_{\hat c}$ and $k^{+}(\hat c)$ only through $p^{+}_{\hat c}$. A claim of
this submodule that fixes one value $p_{\hat c}$ and one count
$\Ncount{\hat c}$ binomial $(n,p_{\hat c})$ and speaks of the guarded
endpoints --- Definition~\ref{def:interval} and
Propositions~\ref{prop:rule-exact}, \ref{prop:rule-bernstein} and
\ref{prop:bracketed-endpoints} do --- states each assertion about the lower
endpoint (one naming $k^{-}$, $R^{-}$, $\varepsilon^{-}$ and the lower tail)
with $p_{\hat c}\coloneqq p^{-}_{\hat c}$ and
$\Ncount{\hat c}\coloneqq\Ncount{\hat c}^{-}$, and each assertion about the
upper endpoint with $p_{\hat c}\coloneqq p^{+}_{\hat c}$ and
$\Ncount{\hat c}\coloneqq\Ncount{\hat c}^{+}$; with the degenerate envelope
the two readings coincide and are the claim as written.
\end{definition}
\claimmeta{\metalabel{Uses}\mref{Def.}{def:freq}, \mref{Def.}{def:matched-alphabet}, \mref{Def.}{def:envelope}
\metasep \metalabel{Used by}\mref{Lem.}{lem:guard}, \mref{Def.}{def:interval}, \mref{Prop.}{prop:rule-exact}, \mref{Rem.}{rem:guard}, \mref{Rem.}{rem:guard-activation}, \mref{Def.}{def:construction}, \mref{Thm.}{thm:box-construction}, \mref{Thm.}{thm:completeness-bound}, \mref{Rem.}{rem:epsstar}, \mref{Thm.}{thm:honest-containment}, \mref{Lem.}{lem:widening}, \mref{Thm.}{thm:monotonicity-limits}, \mref{Prop.}{prop:rule-bernstein}, \mref{Rem.}{rem:completeness-on-rules}, \mref{Prop.}{prop:bracketed-endpoints}, \mref{Rem.}{rem:local-evaluation}, \mref{Req.}{req:rule-in-force}, \mref{Req.}{req:input-refusal}
\metasep newly derived
\metasep implemented (effective: accepted)
\metasep \metalabel{Code}\coderef{https://github.com/fraunhofersg/certiqs-key/blob/v0.1.0/src/certiqs_key/statistics/rules.py}{certiqs_key.statistics.rules:EndpointRule}}

% claim statistics/binomial-box/lem:guard
\claimsummary{statistics/binomial-box/lem:guard}{For every admissible endpoint rule the guarded endpoints enclose the envelope interval, still meet both tail budgets and widen the rule's interval by an exactly known number of counts. Where the guard acts, the emitted frequency endpoint lies within 1/n of that side's envelope value.}
\claimbookmark{statistics/binomial-box/lem:guard}{Lemma}{lem:guard}{The honest-containment guard for every rule}
\claimleadin{%
The guard acts on whatever endpoints a rule returns, so its properties hold
for every rule at once; what is particular to the exact rule is
Proposition~\ref{prop:rule-exact}.
}
\begin{lemma}[The honest-containment guard for every rule]\label{lem:guard}
Let $R$ be an admissible endpoint rule (Definition~\ref{def:endpoint-rule})
and $\hat c\in\Chatcon$; write $p^{\mp}\coloneqq p^{\mp}_{\hat c}$ for the
envelope's values at $\hat c$,
$R^{-}\coloneqq R^{-}(n,p^{-},\varepsilon^{-}(\hat c))$,
$R^{+}\coloneqq R^{+}(n,p^{+},\varepsilon^{+}(\hat c))$, and let
$k^{\mp}(\hat c)$ be the guarded endpoints \eqref{eq:guard}. Then:
\begin{enumerate}[label=(\roman*),leftmargin=2.2em]
\item \emph{(ordering)}
\begin{equation}\label{eq:ordering}
  0\ \le\ k^{-}(\hat c)\ \le\ \lfloor n p^{-}\rfloor
  \ \le\ \lceil n p^{+}\rceil\ \le\ k^{+}(\hat c)\ \le\ n ;
\end{equation}
\item \emph{(the budgets are met)} the maps $k\mapsto\Pr[\Ncount{\hat c}^{-}<k]$
and $k\mapsto\Pr[\Ncount{\hat c}^{+}>k]$ are non-decreasing and non-increasing
respectively on the integers; consequently
\begin{equation}\label{eq:budget-met}
  \Pr\big[\Ncount{\hat c}^{-}<k^{-}(\hat c)\big]\le\varepsilon^{-}(\hat c),
  \qquad
  \Pr\big[\Ncount{\hat c}^{+}>k^{+}(\hat c)\big]\le\varepsilon^{+}(\hat c),
\end{equation}
and more generally \emph{every} pair $(\tilde k^{-},\tilde k^{+})$ with
$0\le\tilde k^{-}\le R^{-}$ and $R^{+}\le\tilde k^{+}\le n$ satisfies
\eqref{eq:budget-met} in place of $k^{\mp}(\hat c)$: enlarging the rule's
interval is always admissible;
\item \emph{(exact cost of the guard)} relative to the rule's own interval
$[R^{-},R^{+}]$, the guard \eqref{eq:guard} enlarges the interval by exactly
\begin{equation}\label{eq:guard-cost}
  R^{-}-k^{-}(\hat c)
    =\max\big\{0,\ R^{-}-\lfloor np^{-}\rfloor\big\},
  \qquad
  k^{+}(\hat c)-R^{+}
    =\max\big\{0,\ \lceil np^{+}\rceil-R^{+}\big\}
\end{equation}
counts on the lower and the upper side respectively;
\item \emph{(distance from the honest value)} measured from the envelope's
value on each side, the guard costs at most $1/n$ \emph{unconditionally}:
\begin{equation}\label{eq:guard-from-p}
  0\le p^{-}-\frac{k^{-}(\hat c)}{n}
    \le\max\Big\{p^{-}-\frac{R^{-}}{n},\ \frac1n\Big\},
  \qquad
  0\le \frac{k^{+}(\hat c)}{n}-p^{+}
    \le\max\Big\{\frac{R^{+}}{n}-p^{+},\ \frac1n\Big\},
\end{equation}
so on whichever side the guard is active, the emitted endpoint lies within
$1/n$ of that side's envelope value.
\end{enumerate}
With the degenerate envelope, $p^{-}=p^{+}=p_{\hat c}$, both
$\Ncount{\hat c}^{\mp}$ are binomial $(n,p_{\hat c})$, and (i)--(iv) are
statements about that single law.
\end{lemma}
\claimmeta{\metalabel{Uses}\mref{Def.}{def:envelope}, \mref{Def.}{def:endpoint-rule}
\metasep \metalabel{Used by}\mref{Prop.}{prop:rule-exact}, \mref{Rem.}{rem:guard}, \mref{Rem.}{rem:guard-activation}, \mref{Thm.}{thm:box-construction}, \mref{Thm.}{thm:completeness-bound}, \mref{Rem.}{rem:envelope}, \mref{Lem.}{lem:widening}, \mref{Thm.}{thm:monotonicity-limits}, \mref{Prop.}{prop:rule-bernstein}, \mref{Rem.}{rem:completeness-on-rules}, \mref{Prop.}{prop:bracketed-endpoints}
\metasep newly derived
\metasep implemented (effective: accepted)
\metasep \metalabel{Code}\coderef{https://github.com/fraunhofersg/certiqs-key/blob/v0.1.0/src/certiqs_key/statistics/intervals.py}{certiqs_key.statistics.intervals:guarded_endpoints}
\metasep \metalabel{Oracles}\coderef{https://github.com/fraunhofersg/certiqs-key/blob/v0.1.0/tests/oracles/registry_statistics_acceptance.py}{registry_statistics_acceptance.py}}

\begin{proof}
\emph{(i).} The chain \eqref{eq:ordering} is immediate from \eqref{eq:guard}:
$k^{-}(\hat c)$ is a minimum one of whose arguments is $\lfloor np^{-}\rfloor$,
$k^{+}(\hat c)$ a maximum one of whose arguments is $\lceil np^{+}\rceil$, and
$\lfloor np^{-}\rfloor\le np^{-}\le np^{+}\le\lceil np^{+}\rceil$ because
$p^{-}\le p^{+}$ (Definition~\ref{def:envelope}); the outer bounds
hold because $R^{-}\ge0$, $\lfloor np^{-}\rfloor\ge0$, $R^{+}\le n$ and
$\lceil np^{+}\rceil\le n$ (Definition~\ref{def:endpoint-rule}\,(a) and
$0\le p^{-}\le p^{+}\le1$).

\emph{(ii).} The function $k\mapsto\Pr[\Ncount{\hat c}^{-}<k]$ is
non-decreasing, because the events $\{\Ncount{\hat c}^{-}<k\}$ increase with
$k$; similarly $k\mapsto\Pr[\Ncount{\hat c}^{+}>k]$ is non-increasing. Hence
for $0\le\tilde k^{-}\le R^{-}$,
$\Pr[\Ncount{\hat c}^{-}<\tilde k^{-}]\le\Pr[\Ncount{\hat c}^{-}<R^{-}]
\le\varepsilon^{-}(\hat c)$ by \eqref{eq:rule-admissible} at
$p=p^{-}$, and for $R^{+}\le\tilde k^{+}\le n$,
$\Pr[\Ncount{\hat c}^{+}>\tilde k^{+}]\le\Pr[\Ncount{\hat c}^{+}>R^{+}]
\le\varepsilon^{+}(\hat c)$ by \eqref{eq:rule-admissible} at $p=p^{+}$.
Since \eqref{eq:guard} gives $0\le k^{-}(\hat c)\le R^{-}$ and
$R^{+}\le k^{+}(\hat c)\le n$ (the latter by (i)), \eqref{eq:budget-met} is
the case $(\tilde k^{-},\tilde k^{+})=(k^{-}(\hat c),k^{+}(\hat c))$.

\emph{(iii).} The two identities \eqref{eq:guard-cost} are \eqref{eq:guard}
together with $\kappa-\min\{\kappa,x\}=\max\{0,\kappa-x\}$ and
$\max\{\kappa,x\}-\kappa=\max\{0,x-\kappa\}$, valid for all reals
$\kappa,x$.

\emph{(iv).} By \eqref{eq:guard},
\begin{equation*}
  p^{-}-\frac{k^{-}(\hat c)}{n}
  =p^{-}-\frac{\min\{R^{-},\lfloor np^{-}\rfloor\}}{n}
  =\max\Big\{p^{-}-\frac{R^{-}}{n},\
              \frac{np^{-}-\lfloor np^{-}\rfloor}{n}\Big\},
\end{equation*}
and $np^{-}-\lfloor np^{-}\rfloor\in[0,1)$, so the second entry lies in
$[0,1/n)$ and the maximum is bounded by $\max\{p^{-}-R^{-}/n,\,1/n\}$;
non-negativity holds because the second entry is non-negative. The upper
display of \eqref{eq:guard-from-p} is the mirror statement at $p^{+}$, using
$\lceil np^{+}\rceil-np^{+}\in[0,1)$.
\end{proof}

\section{The exact rule}\label{group-statistics/binomial-box/def:interval}
\sectionintro{The default rule takes the extremal binomial quantiles: the largest lower count and the smallest upper count whose tails meet their budgets. \hyperlink{claim-statistics/binomial-box/prop:rule-exact}{Proposition~\ref*{prop:rule-exact}} proves that this is an \termref{term-admissible-endpoint-rule}{admissible endpoint rule}, \termref{term-monotone-in-the-allocation}{monotone in the allocation}, and that the guard may widen its raw interval by far more than one count while the emitted endpoint stays within 1/n of the honest value. Two remarks weigh what the guard buys and costs, and show that it acts at ordinary budgets.}

% claim statistics/binomial-box/def:interval
\claimsummary{statistics/binomial-box/def:interval}{The exact rule takes the largest lower count and the smallest upper count whose binomial tails meet their budgets, the lower at the envelope's lower value and the upper at its upper value. The guard then widens this raw interval to contain the envelope interval.}
\claimbookmark{statistics/binomial-box/def:interval}{Definition}{def:interval}{The exact binomial quantile rule}
\claimleadin{%
The exact binomial quantile rule takes the extremal binomial quantiles. It is an
\termref{term-admissible-endpoint-rule}{admissible endpoint rule}, \termref{term-monotone-in-the-allocation}{monotone in the allocation}, by
Proposition~\ref{prop:rule-exact}\,(ii), and it is the default rule
(Requirement~\ref{req:rule-in-force}).
}
\begin{definition}[The exact binomial quantile rule]\label{def:interval}
Fix an \termref{term-admissible-input}{admissible input} (Definition~\ref{def:endpoint-rule}), with its
round count $n\ge1$, its \termref{term-honest-envelope}{honest envelope} $({\honest}^{-},{\honest}^{+})$
(Definition~\ref{def:envelope}; a single law $\honest$ is the
\termref{term-degenerate-envelope}{degenerate envelope}), its constrained outcomes
$\Chatcon\subseteq\Chatset$, its completeness budget $\ecomp\in[0,1)$ and its
budget allocation $\varepsilon^{-},\varepsilon^{+}:\Chatcon\to\mathbb{Q}_{\ge0}$
with $\sum_{\hat c\in\Chatcon}(\varepsilon^{-}(\hat c)+\varepsilon^{+}(\hat
c))\le\ecomp$; in particular $0\le\varepsilon^{\pm}(\hat c)\le\ecomp<1$ for
every $\hat c\in\Chatcon$. For $\hat c\in\Chatcon$ let $p_{\hat c}$ be the
rule's argument --- by the \termref{term-side-wise-reading}{side-wise reading} of
Definition~\ref{def:endpoint-rule}, $p_{\hat c}=p^{-}_{\hat c}$ in
\eqref{eq:kappaminus} and $p_{\hat c}=p^{+}_{\hat c}$ in \eqref{eq:kappaplus},
both equal to $\honest(\hat c)$ for the degenerate envelope --- and let
$\Ncount{\hat c}$ be binomial with parameters $(n,p_{\hat c})$. Define the
\termdef{term-extremal-endpoints}{extremal endpoints}
\begin{align}
  \kappa^{-}(\hat c)&\coloneqq\max\big\{k\in\{0,\dots,n\}:\
    \Pr[\Ncount{\hat c}<k]\le\varepsilon^{-}(\hat c)\big\},
    \label{eq:kappaminus}\\
  \kappa^{+}(\hat c)&\coloneqq\min\big\{k\in\{0,\dots,n\}:\
    \Pr[\Ncount{\hat c}>k]\le\varepsilon^{+}(\hat c)\big\}.
    \label{eq:kappaplus}
\end{align}
The \termdef{term-exact-binomial-quantile-rule}{exact binomial quantile
rule} is the endpoint rule
$R^{\mp}_{\mathrm{ex}}(n,p_{\hat c},\varepsilon^{\mp}(\hat c))
\coloneqq\kappa^{\mp}(\hat c)$. Its \termref{term-guarded-endpoints}{guarded endpoints} are those of
\eqref{eq:guard} with $R^{\mp}=\kappa^{\mp}(\hat c)$, $\kappa^{-}(\hat c)$
taken at $p^{-}_{\hat c}$ and $\kappa^{+}(\hat c)$ at $p^{+}_{\hat c}$, i.e.\
\begin{equation*}
  k^{-}(\hat c)=\min\big\{\kappa^{-}(\hat c),\ \lfloor n\,p^{-}_{\hat c}\rfloor\big\},
  \qquad
  k^{+}(\hat c)=\max\big\{\kappa^{+}(\hat c),\ \lceil n\,p^{+}_{\hat c}\rceil\big\}.
\end{equation*}
\end{definition}
\claimmeta{\metalabel{Uses}\mref{Def.}{def:envelope}, \mref{Def.}{def:endpoint-rule}
\metasep \metalabel{Used by}\mref{Prop.}{prop:rule-exact}, \mref{Rem.}{rem:guard}, \mref{Rem.}{rem:guard-activation}, \mref{Def.}{def:construction}, \mref{Lem.}{lem:widening}, \mref{Thm.}{thm:monotonicity-limits}, \mref{Prop.}{prop:rule-bernstein}, \mref{Rem.}{rem:completeness-on-rules}, \mref{Prop.}{prop:bracketed-endpoints}, \mref{Rem.}{rem:local-evaluation}, \mref{Req.}{req:rule-in-force}
\metasep newly derived
\metasep implemented (effective: accepted)
\metasep \metalabel{Code}\coderef{https://github.com/fraunhofersg/certiqs-key/blob/v0.1.0/src/certiqs_key/statistics/intervals.py}{certiqs_key.statistics.intervals:extremal_endpoints}}

% claim statistics/binomial-box/prop:rule-exact
\claimsummary{statistics/binomial-box/prop:rule-exact}{The exact quantiles are the extremal counts meeting their budgets and form an admissible endpoint rule, monotone in the allocation, as do counts decided against upper bounds on the tails. The guard may widen the interval by many counts, yet a moved endpoint stays within 1/n of the honest value.}
\claimbookmark{statistics/binomial-box/prop:rule-exact}{Proposition}{prop:rule-exact}{Properties of the exact binomial quantile rule}
\begin{proposition}[Properties of the exact binomial quantile rule]\label{prop:rule-exact}
Assume a \termref{term-conservative-evaluation}{conservative evaluation}
(Definition~\ref{def:conservative-evaluation}). With the notation of
Definition~\ref{def:interval}, for every $\hat c\in\Chatcon$:
\begin{enumerate}[label=(\roman*),leftmargin=2.2em]
\item the two sets in \eqref{eq:kappaminus}--\eqref{eq:kappaplus} are non-empty
subsets of the finite set $\{0,\dots,n\}$, so $\kappa^{\pm}(\hat c)$ exist and
are unique;
\item \emph{(the exact rule is an admissible endpoint rule)} the pair
$R^{\mp}_{\mathrm{ex}}=\kappa^{\mp}(\hat c)$ satisfies
Definition~\ref{def:endpoint-rule}\,(a)--(c): $\kappa^{\pm}(\hat c)\in
\{0,\dots,n\}$ and
\begin{equation}\label{eq:exact-admissible}
  \Pr\big[\Ncount{\hat c}<\kappa^{-}(\hat c)\big]\le\varepsilon^{-}(\hat c),
  \qquad
  \Pr\big[\Ncount{\hat c}>\kappa^{+}(\hat c)\big]\le\varepsilon^{+}(\hat c);
\end{equation}
if $0<p_{\hat c}<1$ then $\kappa^{-}(\hat c)=0$ at $\varepsilon^{-}(\hat c)=0$
and $\kappa^{+}(\hat c)=n$ at $\varepsilon^{+}(\hat c)=0$; and
$\kappa^{-}(\hat c)=\kappa^{+}(\hat c)=0$ at $p_{\hat c}=0$,
$\kappa^{-}(\hat c)=\kappa^{+}(\hat c)=n$ at $p_{\hat c}=1$. Moreover the
rule is monotone in the allocation: $\kappa^{-}(\hat c)$ is non-decreasing in
$\varepsilon^{-}(\hat c)$ and $\kappa^{+}(\hat c)$ is non-increasing in
$\varepsilon^{+}(\hat c)$. Consequently Lemma~\ref{lem:guard} applies to the
exact rule, and its guarded endpoints $k^{\mp}(\hat c)$ satisfy
\eqref{eq:ordering} and \eqref{eq:budget-met};
\item \emph{(extremality and admissible relaxation)} $\kappa^{-}(\hat c)$ is
the largest and $\kappa^{+}(\hat c)$ the smallest integer in $\{0,\dots,n\}$
satisfying the respective inequality of \eqref{eq:exact-admissible}, and
\emph{every} pair $(\tilde k^{-},\tilde k^{+})$ with
$0\le\tilde k^{-}\le\kappa^{-}(\hat c)$ and
$\kappa^{+}(\hat c)\le\tilde k^{+}\le n$ satisfies \eqref{eq:budget-met} in
place of $k^{\mp}(\hat c)$. Enlarging the interval is therefore always
admissible; shrinking it beyond $[\kappa^{-},\kappa^{+}]$ is not;
\item the defining inequalities are the closed forms
\eqref{eq:tails-used} of Corollary~\ref{cor:tail-forms}, so $\kappa^{\pm}(\hat c)$
are computable from $n$, $p_{\hat c}$ and the allocation alone
(Lemma~\ref{lem:exact-strict-tails});
\item \emph{(degenerate honest mass)} if $p_{\hat c}=0$ then
$k^{-}(\hat c)=k^{+}(\hat c)=0$, and if $p_{\hat c}=1$ then
$k^{-}(\hat c)=k^{+}(\hat c)=n$; in both cases $\varepsilon^{-}(\hat c)$ and
$\varepsilon^{+}(\hat c)$ are unspent, both tails in \eqref{eq:budget-met}
being exactly zero;
\item \emph{(exact cost of the guard for this rule)} write $p=p_{\hat c}$.
Relative to the raw quantile interval
$[\kappa^{-}(\hat c),\kappa^{+}(\hat c)]$, the guard \eqref{eq:guard}
enlarges the interval by exactly \eqref{eq:guard-cost} with
$R^{\mp}=\kappa^{\mp}(\hat c)$, and neither enlargement is bounded by one
count: for every $n\ge1$ and every $p\in(0,1)$, an allocation with
$\varepsilon^{-}(\hat c)\ge1-p^{\,n}$ --- admissible whenever $\ecomp$
permits, since $1-p^{\,n}<1$ --- gives $\kappa^{-}(\hat c)=n$ and a lower
enlargement of $n-\lfloor np\rfloor$ counts, that is, at least $1-p$ in
frequency \emph{uniformly in $n$}; symmetrically
$\varepsilon^{+}(\hat c)\ge1-(1-p)^{n}$ gives $\kappa^{+}(\hat c)=0$ and an
upper enlargement of $\lceil np\rceil$ counts. The exact conditions are
\begin{align}
  \kappa^{-}(\hat c)-k^{-}(\hat c)=0
    &\iff \varepsilon^{-}(\hat c)
      <\Pr\big[\Ncount{\hat c}\le\lfloor np\rfloor\big],
    \label{eq:guard-lo0}\\
  \kappa^{-}(\hat c)-k^{-}(\hat c)\le1
    &\iff \varepsilon^{-}(\hat c)
      <\Pr\big[\Ncount{\hat c}\le\lfloor np\rfloor+1\big],
    \label{eq:guard-lo1}\\
  k^{+}(\hat c)-\kappa^{+}(\hat c)=0
    &\iff \varepsilon^{+}(\hat c)
      <\Pr\big[\Ncount{\hat c}\ge\lceil np\rceil\big],
    \label{eq:guard-up0}\\
  k^{+}(\hat c)-\kappa^{+}(\hat c)\le1
    &\iff \varepsilon^{+}(\hat c)
      <\Pr\big[\Ncount{\hat c}\ge\lceil np\rceil-1\big].
    \label{eq:guard-up1}
\end{align}
Measured instead from the honest value, the guard costs at most $1/n$
unconditionally, \eqref{eq:guard-from-p} of Lemma~\ref{lem:guard}\,(iv);
\item \emph{(evaluation against upper bounds)} let $U^{-}_k$ and $U^{+}_k$,
for $k\in\{0,\dots,n\}$, be rationals determined by $n$, $p_{\hat c}$ and
$k$ alone, with
$U^{-}_k\ge\Pr[\Ncount{\hat c}<k]$, $U^{+}_k\ge\Pr[\Ncount{\hat c}>k]$,
$U^{-}_0=0$ and $U^{+}_n=0$, and with every $U^{\mp}_k$ equal to its tail
when $p_{\hat c}\in\{0,1\}$. Put
\begin{equation*}
  \tilde\kappa^{-}(\hat c)\coloneqq\max\big\{k\in\{0,\dots,n\}:\
    U^{-}_{j}\le\varepsilon^{-}(\hat c)\text{ for }1\le j\le k\big\},
  \qquad
  \tilde\kappa^{+}(\hat c)\coloneqq\min\big\{k\in\{0,\dots,n\}:\
    U^{+}_{k}\le\varepsilon^{+}(\hat c)\big\}.
\end{equation*}
Then $\tilde R^{\mp}=\tilde\kappa^{\mp}(\hat c)$ is an admissible endpoint
rule (Definition~\ref{def:endpoint-rule}\,(a)--(c)), monotone in the
allocation, with $\tilde\kappa^{-}(\hat c)\le\kappa^{-}(\hat c)$ and
$\tilde\kappa^{+}(\hat c)\ge\kappa^{+}(\hat c)$, and it is the exact rule
when every $U^{\mp}_k$ equals its tail. Consequently Lemma~\ref{lem:guard}
applies to it.
\end{enumerate}
\end{proposition}
\claimmeta{\metalabel{Uses}\mref{Cor.}{cor:tail-forms}, \mref{Lem.}{lem:exact-strict-tails}, \mref{Def.}{def:conservative-evaluation}, \mref{Def.}{def:endpoint-rule}, \mref{Lem.}{lem:guard}, \mref{Def.}{def:interval}
\metasep \metalabel{Used by}\mref{Rem.}{rem:guard}, \mref{Rem.}{rem:guard-activation}, \mref{Rem.}{rem:epsstar}, \mref{Prop.}{prop:bracketed-endpoints}
\metasep newly derived
\metasep implemented (effective: accepted)
\metasep \metalabel{Code}\coderef{https://github.com/fraunhofersg/certiqs-key/blob/v0.1.0/src/certiqs_key/statistics/rules.py}{certiqs_key.statistics.rules:SweepExactRule}
\metasep \metalabel{Oracles}\coderef{https://github.com/fraunhofersg/certiqs-key/blob/v0.1.0/tests/oracles/registry_statistics_acceptance.py}{registry_statistics_acceptance.py}}

\begin{proof}
\emph{(i).} $k=0$ belongs to the set in \eqref{eq:kappaminus}, since
$\Pr[\Ncount{\hat c}<0]=0\le\varepsilon^{-}(\hat c)$; and $k=n$ belongs to the
set in \eqref{eq:kappaplus}, since
$\Pr[\Ncount{\hat c}>n]=0\le\varepsilon^{+}(\hat c)$. A non-empty finite set of
integers has a unique maximum and a unique minimum.

\emph{(ii).} \emph{Admissibility.} $\kappa^{-}(\hat c)$ is an element of the
set in \eqref{eq:kappaminus} and $\kappa^{+}(\hat c)$ one of the set in
\eqref{eq:kappaplus}, so both lie in $\{0,\dots,n\}$ and satisfy the
inequalities that define those sets, which are \eqref{eq:exact-admissible}.
This is Definition~\ref{def:endpoint-rule}\,(a).

\emph{Monotonicity in the allocation.} Fix every argument but
$\varepsilon^{-}(\hat c)$. Raising $\varepsilon^{-}(\hat c)$ enlarges the set
in \eqref{eq:kappaminus}, so its maximum does not fall; raising
$\varepsilon^{+}(\hat c)$ enlarges the set in \eqref{eq:kappaplus}, so its
minimum does not rise.

\emph{Zero budget, Definition~\ref{def:endpoint-rule}\,(b).} Let
$0<p_{\hat c}<1$ and write $p=p_{\hat c}$. At $\varepsilon^{-}(\hat c)=0$:
for $k\ge1$, $\Pr[\Ncount{\hat c}<k]\ge\Pr[\Ncount{\hat c}=0]=(1-p)^{n}>0$, so
only $k=0$ lies in the set of \eqref{eq:kappaminus} and
$\kappa^{-}(\hat c)=0$. At $\varepsilon^{+}(\hat c)=0$: for $k\le n-1$,
$\Pr[\Ncount{\hat c}>k]\ge\Pr[\Ncount{\hat c}=n]=p^{n}>0$, so only $k=n$ lies
in the set of \eqref{eq:kappaplus} and $\kappa^{+}(\hat c)=n$.

\emph{Degenerate mass, Definition~\ref{def:endpoint-rule}\,(c).} If
$p_{\hat c}=0$ then $\Ncount{\hat c}=0$ with probability one, so
$\Pr[\Ncount{\hat c}<k]$ is $0$ at $k=0$ and $1$ for $k\ge1$; since
$\varepsilon^{-}(\hat c)<1$ the set in \eqref{eq:kappaminus} is $\{0\}$ and
$\kappa^{-}(\hat c)=0$; and $\Pr[\Ncount{\hat c}>k]=0$ for every $k\ge0$, so
the set in \eqref{eq:kappaplus} is all of $\{0,\dots,n\}$ and
$\kappa^{+}(\hat c)=0$. If $p_{\hat c}=1$ then $\Ncount{\hat c}=n$ with
probability one, so $\Pr[\Ncount{\hat c}<k]=0$ for every $k\le n$ and
$\kappa^{-}(\hat c)=n$; and $\Pr[\Ncount{\hat c}>k]$ is $1$ for $k\le n-1$ and
$0$ at $k=n$, so with $\varepsilon^{+}(\hat c)<1$ the set in
\eqref{eq:kappaplus} is $\{n\}$ and $\kappa^{+}(\hat c)=n$.

The consequence is Lemma~\ref{lem:guard} applied to $R=R_{\mathrm{ex}}$.

\emph{(iii).} Extremality is the definition of the maximum and the minimum
together with the monotonicity of the two tails in $k$
(Lemma~\ref{lem:guard}\,(ii)); the relaxation statement is
Lemma~\ref{lem:guard}\,(ii) with $R^{\mp}=\kappa^{\mp}(\hat c)$.

\emph{(iv).} By Corollary~\ref{cor:tail-forms}, with $p=p_{\hat c}$, the
tails in the inequalities that define the sets of \eqref{eq:kappaminus} and
\eqref{eq:kappaplus} are the closed forms \eqref{eq:tails-used}. Under a
conservative evaluation $p_{\hat c}$, a coordinate of ${\honest}^{-}$ or of
${\honest}^{+}$ (Definition~\ref{def:interval}), is rational, the envelope's
vectors being exact rationals (Definition~\ref{def:conservative-evaluation});
and $\varepsilon^{\mp}(\hat c)$ are rational by
Definition~\ref{def:interval}, which also gives $n\ge1$ and
$p_{\hat c}\in[0,1]$. By
Lemma~\ref{lem:exact-strict-tails}\,(ii) each inequality, one for each
$k\in\{0,\dots,n\}$ in each set, is then decided exactly from $n$,
$p_{\hat c}$, $k$ and $\varepsilon^{\mp}(\hat c)$. So both sets, and with
them their maximum $\kappa^{-}(\hat c)$ and minimum $\kappa^{+}(\hat c)$, are
computed exactly from $n$, $p_{\hat c}$ and the allocation alone.

\emph{(v).} By (ii), $\kappa^{-}(\hat c)=\kappa^{+}(\hat c)=0$ at
$p_{\hat c}=0$, and $\lfloor n\cdot0\rfloor=\lceil n\cdot0\rceil=0$, so
\eqref{eq:guard} gives $k^{-}(\hat c)=k^{+}(\hat c)=0$; at $p_{\hat c}=1$,
$\kappa^{\pm}(\hat c)=n$ and $\lfloor n\rfloor=\lceil n\rceil=n$, so
$k^{-}(\hat c)=k^{+}(\hat c)=n$. In the first case
$\Pr[\Ncount{\hat c}<0]=0$ and $\Pr[\Ncount{\hat c}>0]=0$ because
$\Ncount{\hat c}=0$ almost surely; in the second, $\Pr[\Ncount{\hat c}<n]=0$
and $\Pr[\Ncount{\hat c}>n]=0$ because $\Ncount{\hat c}=n$ almost surely. So
both tails in \eqref{eq:budget-met} are exactly $0$.

\emph{(vi).} The identities are Lemma~\ref{lem:guard}\,(iii) with
$R^{\mp}=\kappa^{\mp}(\hat c)$.

\emph{Unboundedness.} If $\varepsilon^{-}(\hat c)\ge1-p^{\,n}$ then
$\Pr[\Ncount{\hat c}<n]=\Pr[\Ncount{\hat c}\le n-1]=1-p^{\,n}
\le\varepsilon^{-}(\hat c)$, so $k=n$ is admissible in \eqref{eq:kappaminus}
and $\kappa^{-}(\hat c)=n$; the lower enlargement is then
$n-\lfloor np\rfloor\ge n-np=n(1-p)$ counts, i.e.\ at least $1-p$ in
frequency. Such an $\varepsilon^{-}(\hat c)$ is admissible whenever $\ecomp$
permits, because $1-p^{\,n}<1$ for $p\in(0,1)$. Symmetrically, if
$\varepsilon^{+}(\hat c)\ge1-(1-p)^{n}$ then
$\Pr[\Ncount{\hat c}>0]=1-(1-p)^{n}\le\varepsilon^{+}(\hat c)$, so $k=0$ is
admissible in \eqref{eq:kappaplus} and $\kappa^{+}(\hat c)=0$, giving an upper
enlargement of $\lceil np\rceil$ counts.

\emph{The four exact conditions.} We use twice that
$\varepsilon^{\pm}(\hat c)\le\ecomp<1$, which
Definition~\ref{def:interval} records for every admissible input.

First, for every integer $K\ge0$,
\begin{equation*}
  \kappa^{-}(\hat c)\le K
  \quad\Longleftrightarrow\quad
  \Pr\big[\Ncount{\hat c}\le K\big]>\varepsilon^{-}(\hat c).
\end{equation*}
If $K+1\le n$: ($\Rightarrow$) $K+1$ is then not admissible in
\eqref{eq:kappaminus}, by maximality of $\kappa^{-}(\hat c)$, which is exactly
$\Pr[\Ncount{\hat c}<K+1]=\Pr[\Ncount{\hat c}\le K]>\varepsilon^{-}(\hat c)$;
($\Leftarrow$) for every integer $j$ with $K+1\le j\le n$, monotonicity of
$k\mapsto\Pr[\Ncount{\hat c}<k]$ (Lemma~\ref{lem:guard}\,(ii)) gives
$\Pr[\Ncount{\hat c}<j]\ge\Pr[\Ncount{\hat c}\le K]>\varepsilon^{-}(\hat c)$,
so no such $j$ is admissible and the maximum is at most $K$. If instead
$K+1>n$, both sides hold unconditionally: $\kappa^{-}(\hat c)\le n\le K$ by
(i), and $\Pr[\Ncount{\hat c}\le K]=1>\varepsilon^{-}(\hat c)$.
By \eqref{eq:guard-cost} the lower enlargement is $0$ iff
$\kappa^{-}(\hat c)\le\lfloor np\rfloor$ and is at most $1$ iff
$\kappa^{-}(\hat c)\le\lfloor np\rfloor+1$, so taking $K=\lfloor np\rfloor$
and $K=\lfloor np\rfloor+1$ gives \eqref{eq:guard-lo0} and
\eqref{eq:guard-lo1}.

Second, and symmetrically, for every integer $K\le n$,
\begin{equation*}
  \kappa^{+}(\hat c)\ge K
  \quad\Longleftrightarrow\quad
  \Pr\big[\Ncount{\hat c}\ge K\big]>\varepsilon^{+}(\hat c),
\end{equation*}
by the same argument applied to the non-increasing map
$k\mapsto\Pr[\Ncount{\hat c}>k]$ and the minimality of $\kappa^{+}(\hat c)$:
for $K-1\ge0$ the index $K-1$ is inadmissible iff
$\Pr[\Ncount{\hat c}>K-1]=\Pr[\Ncount{\hat c}\ge K]>\varepsilon^{+}(\hat c)$,
and for $K\le0$ both sides hold unconditionally. Taking $K=\lceil np\rceil$
and $K=\lceil np\rceil-1$, both of which are at most $n$, gives
\eqref{eq:guard-up0} and \eqref{eq:guard-up1}.

The distance statement is Lemma~\ref{lem:guard}\,(iv).

\emph{(vii).} Both sets are non-empty: the condition defining the first is
vacuous at $k=0$, and $k=n$ belongs to the second because $U^{+}_n=0$. Write
$\tilde k^{\mp}\coloneqq\tilde\kappa^{\mp}(\hat c)$.
\emph{Admissibility.} If $\tilde k^{-}=0$ then
$\Pr[\Ncount{\hat c}<\tilde k^{-}]=0$; otherwise
$\Pr[\Ncount{\hat c}<\tilde k^{-}]\le U^{-}_{\tilde k^{-}}
\le\varepsilon^{-}(\hat c)$, the last step being the defining condition at
$j=\tilde k^{-}$. And
$\Pr[\Ncount{\hat c}>\tilde k^{+}]\le U^{+}_{\tilde k^{+}}
\le\varepsilon^{+}(\hat c)$. Both endpoints lie in $\{0,\dots,n\}$, so this
is Definition~\ref{def:endpoint-rule}\,(a); it also places $\tilde k^{-}$
in the set of \eqref{eq:kappaminus} and $\tilde k^{+}$ in that of
\eqref{eq:kappaplus}, so $\tilde k^{-}\le\kappa^{-}(\hat c)$ and
$\tilde k^{+}\ge\kappa^{+}(\hat c)$ by the extremality of (iii).
\emph{Zero budget.} Let $0<p_{\hat c}<1$. At $\varepsilon^{-}(\hat c)=0$,
(ii) gives $\kappa^{-}(\hat c)=0$, so $0\le\tilde k^{-}\le0$; at
$\varepsilon^{+}(\hat c)=0$ it gives $\kappa^{+}(\hat c)=n$, so
$n\ge\tilde k^{+}\ge n$. This is Definition~\ref{def:endpoint-rule}\,(b).
\emph{Equality with the exact rule.} If every $U^{\mp}_k$ equals its tail,
the second set is that of \eqref{eq:kappaplus}; and since
$k\mapsto\Pr[\Ncount{\hat c}<k]$ is non-decreasing
(Lemma~\ref{lem:guard}\,(ii)), the condition holds at every $j\le k$
exactly when it holds at $j=k$, so the first set is that of
\eqref{eq:kappaminus}. Hence $\tilde k^{\mp}=\kappa^{\mp}(\hat c)$.
\emph{Degenerate mass.} At $p_{\hat c}\in\{0,1\}$ every bound equals its
tail, so the rule is the exact one there, and (ii) gives
Definition~\ref{def:endpoint-rule}\,(c).
\emph{Monotonicity in the allocation.} The bounds do not depend on the
allocation. Raising $\varepsilon^{-}(\hat c)$ enlarges the first set, so its
maximum does not fall; raising $\varepsilon^{+}(\hat c)$ enlarges the second,
so its minimum does not rise. The consequence is Lemma~\ref{lem:guard}
applied to $R=\tilde R$.
\end{proof}

% claim statistics/binomial-box/rem:guard
\claimsummary{statistics/binomial-box/rem:guard}{The guard buys honest containment for every budget and every allocation, and never raises the failure probability. Its cost is not bounded by one count against the raw quantile interval, yet an endpoint it moves stays within 1/n of the honest value.}
\claimbookmark{statistics/binomial-box/rem:guard}{Remark}{rem:guard}{Benefit and cost of the honest-containment guard}
\begin{remark}[Benefit and cost of the honest-containment guard]\label{rem:guard}
This remark reads the \termref{term-honest-containment-guard}{honest-containment guard} on the exact rule of
Definition~\ref{def:interval}, whose rule endpoints are the quantiles
$\kappa^{\pm}$; Lemma~\ref{lem:guard} states what holds for every rule. It is
written for a single honest value $p_{\hat c}=\honest(\hat c)$, the
degenerate envelope; for a general envelope every statement below holds side
by side, by the side-wise reading of Definition~\ref{def:endpoint-rule}: what
it says of the lower endpoint with $p_{\hat c}=p^{-}_{\hat c}$, and what it
says of the upper endpoint with $p_{\hat c}=p^{+}_{\hat c}$.

\emph{What the guard buys} is unconditional: the whole envelope interval
$[p^{-}_{\hat c},p^{+}_{\hat c}]$ --- for a single law, $\honest(\hat c)$ ---
lies in $[k^{-}(\hat c)/n,\,k^{+}(\hat c)/n]$ for \emph{every} admissible
budget and \emph{every} admissible allocation, including allocations so
lopsided that the raw quantile interval
$[\kappa^{-}(\hat c),\kappa^{+}(\hat c)]$ excludes the honest value.
Theorem~\ref{thm:box-construction}\,(iii) and Theorem~\ref{thm:honest-containment} rest on exactly that,
and without the guard both would carry a side condition on how large
$\ecomp$ may be. By Proposition~\ref{prop:rule-exact}\,(iii) the widening is
always admissible: it can only lower the failure probability.

\emph{What the guard costs} has two different and equally correct answers,
both stated in Lemma~\ref{lem:guard}\,(iii)--(iv) and
Proposition~\ref{prop:rule-exact}\,(vi); the second alone would overstate how
little the guard costs.

\emph{Against the raw quantile interval} the cost is \eqref{eq:guard-cost},
and it is \emph{not} bounded by one count. A concrete admissible input:
$\Chatcon=\{\hat c\}$, $n=10$, $p_{\hat c}=\tfrac12$,
$\varepsilon^{-}(\hat c)=\tfrac9{10}$ and $\varepsilon^{+}(\hat c)=\tfrac1{20}$,
the allocation summing to $\tfrac{19}{20}$ and so admissible for any
$\ecomp\in[\tfrac{19}{20},1)$. There
$\Pr[\Ncount{\hat c}<7]=\tfrac{53}{64}\le\tfrac9{10}$ while
$\Pr[\Ncount{\hat c}<8]=\tfrac{121}{128}>\tfrac9{10}$, so $\kappa^{-}=7$
against $\lfloor np_{\hat c}\rfloor=5$: the guard widens the interval by two
counts. Proposition~\ref{prop:rule-exact}\,(vi) shows the widening is
unbounded in the same direction --- it is at least $1-p_{\hat c}$ in
frequency, uniformly in $n$, once a single tail's allocation reaches
$1-p_{\hat c}^{\,n}$.

\emph{Against the honest value} the quantity bounded is the emitted
endpoint's \emph{distance} from $\honest(\hat c)$, not the enlargement, and
that distance is at most $1/n$, unconditionally: by \eqref{eq:guard-from-p},
on the side where the guard acts the emitted endpoint sits within $1/n$ of
$\honest(\hat c)$. This is the sense in which the guard is cheap, and it is
the sense that survives the comparison with Theorem~\ref{thm:monotonicity-limits}\,(iv): the
$O(n^{-1/2})$ half-width of \eqref{eq:width} is proved for the
\emph{guarded} endpoints, with no side condition, so the guard never breaks
that bound.
\end{remark}
\claimmeta{\metalabel{Uses}\mref{Def.}{def:endpoint-rule}, \mref{Lem.}{lem:guard}, \mref{Def.}{def:interval}, \mref{Prop.}{prop:rule-exact}
\metasep \metalabel{Used by}\mref{Rem.}{rem:epsstar}
\metasep newly derived
\metasep accepted}

% claim statistics/binomial-box/rem:guard-activation
\claimsummary{statistics/binomial-box/rem:guard-activation}{Four exact comparisons of the allocation with the binomial distribution function decide whether the guard acts and whether it costs more than one count. It acts even at ordinary budgets, on any outcome whose positive expected count is at most its tail budget.}
\claimbookmark{statistics/binomial-box/rem:guard-activation}{Remark}{rem:guard-activation}{Regimes in which the guard acts}
\begin{remark}[Regimes in which the guard acts]\label{rem:guard-activation}
This remark reads the guard on the exact rule of
Definition~\ref{def:interval}, written for a single honest value
$p_{\hat c}=\honest(\hat c)$, the degenerate envelope; for a general envelope
every statement holds side by side, by the side-wise reading of
Definition~\ref{def:endpoint-rule}.

The four exact conditions \eqref{eq:guard-lo0}--\eqref{eq:guard-up1} of
Proposition~\ref{prop:rule-exact}\,(vi) say where the guard acts and what it
costs there. By \eqref{eq:guard-lo0} and \eqref{eq:guard-up0} the guard is
entirely inert as soon as each allocated tail budget is strictly below the
mass the count law places at or below $\lfloor np_{\hat c}\rfloor$,
respectively at or above $\lceil np_{\hat c}\rceil$; by \eqref{eq:guard-lo1}
and \eqref{eq:guard-up1} it costs at most one count under the correspondingly
weaker conditions. All four are comparisons of an allocated budget against a
value of the binomial distribution function, hence decidable exactly in the
closed forms \eqref{eq:tails-used} of Corollary~\ref{cor:tail-forms}: each
is the negation of a comparison of Lemma~\ref{lem:exact-strict-tails}\,(ii)
at the count next to the one it names, or a boundary value of
Theorem~\ref{thm:binomial-tails} where that count leaves $\{0,\dots,n\}$,
with no machinery beyond what
Definition~\ref{def:interval} already requires. The two
pairs fail in different regimes, however, and are not to be read together.

\emph{A cost of more than one count} --- failure of \eqref{eq:guard-lo1} or of
\eqref{eq:guard-up1} --- requires a single outcome's single tail to carry a
budget at least the mass the count law places up to one step beyond
$np_{\hat c}$, that is
$\varepsilon^{-}(\hat c)\ge\Pr[\Ncount{\hat c}\le\lfloor np_{\hat c}\rfloor+1]$
respectively
$\varepsilon^{+}(\hat c)\ge\Pr[\Ncount{\hat c}\ge\lceil np_{\hat c}\rceil-1]$.
At the ends of the range that cannot happen at all: if $np_{\hat c}<1$ then
$\lceil np_{\hat c}\rceil-1\le0$ and
$\Pr[\Ncount{\hat c}\ge\lceil np_{\hat c}\rceil-1]=1>\varepsilon^{+}(\hat c)$,
and if $\lfloor np_{\hat c}\rfloor\ge n-1$ then
$\lfloor np_{\hat c}\rfloor+1\ge n$ and
$\Pr[\Ncount{\hat c}\le\lfloor np_{\hat c}\rfloor+1]=1>\varepsilon^{-}(\hat
c)$, both using $\varepsilon^{\pm}(\hat c)\le\ecomp<1$. In the example of
Remark~\ref{rem:guard} the budget is indeed near-unit: $\varepsilon^{-}(\hat
c)=\tfrac9{10}$ against $\Pr[\Ncount{\hat c}\le6]=\tfrac{53}{64}$. That is
also the regime of the $\ecomp=\tfrac45$-on-one-outcome counterexample of
Theorem~\ref{thm:monotonicity-limits}\,(v), and under an allocation that concentrates a
near-unit budget on one or two constrained outcomes the enlargement is to be
read off \eqref{eq:guard-cost}, not estimated by a rule of thumb.

\emph{Activation of the guard at all} --- failure of \eqref{eq:guard-lo0} or
of \eqref{eq:guard-up0} --- is a different matter: it is not confined to
near-unit budgets. It occurs at ordinary budgets and under the engine's
default \termref{term-uniform-rule}{uniform allocation rule}, for any outcome whose expected count over
the whole run is positive and at or below the tail budget allocated to it.
Precisely, if $0<np_{\hat c}\le\varepsilon^{+}(\hat c)$ --- which forces
$np_{\hat c}<1$ and hence $\lceil np_{\hat c}\rceil=1$, since
$\varepsilon^{+}(\hat c)<1$ --- then Bernoulli's inequality
$(1-p_{\hat c})^{n}\ge1-np_{\hat c}$ gives
\begin{equation*}
  \Pr\big[\Ncount{\hat c}\ge\lceil np_{\hat c}\rceil\big]
  =1-(1-p_{\hat c})^{n}\ \le\ np_{\hat c}\ \le\ \varepsilon^{+}(\hat c),
\end{equation*}
so \eqref{eq:guard-up0} fails, $\kappa^{+}(\hat c)=0$, and the guard raises
the emitted endpoint to $k^{+}(\hat c)=1$. Symmetrically, if
$0<n(1-p_{\hat c})\le\varepsilon^{-}(\hat c)$ then $\lfloor
np_{\hat c}\rfloor=n-1$ and $\Pr[\Ncount{\hat c}\le n-1]=1-p_{\hat c}^{\,n}\le
n(1-p_{\hat c})\le\varepsilon^{-}(\hat c)$ --- the same inequality applied at
$1-p_{\hat c}$ --- so \eqref{eq:guard-lo0} fails, $\kappa^{-}(\hat c)=n$, and
the guard lowers the emitted endpoint to $k^{-}(\hat c)=n-1$. A concrete
instance is an \termref{term-admissible-input}{admissible input} with the engine's default rule and
default allocation (Requirements~\ref{req:rule-in-force}
and~\ref{req:request-defaults}): $\Chatset=\{\hat c_0,\hat c_1\}$, $n=10$,
$\honest=(\tfrac{999}{1000},\tfrac1{1000})$, $\Chatcon=\Chatset$,
$\ecomp=\tfrac1{10}$ and the uniform allocation
$\varepsilon^{\pm}(\hat c)=\ecomp/(2|\Chatcon|)=\tfrac1{40}$; there
$np_{\hat c_1}=n(1-p_{\hat c_0})=\tfrac1{100}<\tfrac1{40}$, so the guard is
active on the upper side of $\hat c_1$ and on the lower side of $\hat c_0$
--- at a budget nowhere near $1$, and with the budget spread uniformly rather
than concentrated. Nor is this a small-$n$ artefact: at $n=10^{6}$ with
$\ecomp=\tfrac1{10}$ and $|\Chatcon|=5$ the uniform rule allocates
$\varepsilon^{\pm}(\hat c)=10^{-2}$, and a constrained outcome of honest
probability $10^{-9}$ has $np_{\hat c}=10^{-3}\le\varepsilon^{+}(\hat c)$, so
the guard forces $q^{+}(\hat c)=1/n$ where the raw quantile interval would
have given $0$. Under A-TSM (\xref{reductions/ass:tsm}), this
module's alphabet being that of the consuming declaration
(Definition~\ref{def:matched-alphabet}), the acceptance alphabet is the
joint outcome, carrying every per-round public choice: one made at a
measuring station as a component of the outcome by A-TSM, and one on the
letter side as a component of the letter
(\xref{protocol-compiler/def:spec}\,(a)), drawn as declared by A-COIN
(\xref{protocol-compiler/ass:coin}); so outcomes of
expected count at or below one are a routine feature of a fine-grained
$\Chatset$ rather than an exotic one. What the guard costs on them is
nevertheless bounded, by both measures: exactly one count --- the relevant
condition of the previous paragraph holds in each case,
$\lceil np_{\hat c}\rceil-1\le0$ giving \eqref{eq:guard-up1} on the upper
side and $\lfloor np_{\hat c}\rfloor+1\ge n$ giving \eqref{eq:guard-lo1} on
the lower side --- and at most $1/n$ from the honest value by
\eqref{eq:guard-from-p}.
\end{remark}
\claimmeta{\metalabel{Uses}\mref{Def.}{def:matched-alphabet}, \mref{Thm.}{thm:binomial-tails}, \mref{Cor.}{cor:tail-forms}, \mref{Lem.}{lem:exact-strict-tails}, \mref{Def.}{def:endpoint-rule}, \mref{Lem.}{lem:guard}, \mref{Def.}{def:interval}, \mref{Prop.}{prop:rule-exact}, \mref{Def.}{def:construction}
\metasep \metalabel{Assumes}\extaref{../reductions/contract.pdf}{reductions/ass:tsm}{A-TSM}, \aref{A-ALPHABET}{ass:alphabet}
\metasep newly derived
\metasep accepted}

\section{The construction map}\label{group-statistics/binomial-box/def:construction}
\sectionintro{The guarded counts divided by the round count bound the constrained outcomes, the unconstrained ones keep the bounds 0 and 1, and the box is intersected with the simplex. The main result, \hyperlink{claim-statistics/binomial-box/thm:box-construction}{Theorem~\ref*{thm:box-construction}}, is that for every \termref{term-admissible-endpoint-rule}{admissible endpoint rule} the output is a canonical box of exact rationals containing every law of the \termref{term-envelope-set}{envelope set}, a non-empty, compact and convex subset of the simplex by construction.}

% claim statistics/binomial-box/def:construction
\claimsummary{statistics/binomial-box/def:construction}{The construction map bounds each constrained outcome by its guarded count endpoints divided by the round count, gives every other outcome the bounds 0 and 1, and intersects the box with the probability simplex. Its output is a polytope in canonical box form built from the honest envelope.}
\claimbookmark{statistics/binomial-box/def:construction}{Definition}{def:construction}{The construction map}
\claimleadin{%
The map is defined for every \termref{term-admissible-endpoint-rule}{admissible endpoint rule}; which rule is
used, by default the \termref{term-exact-binomial-quantile-rule}{exact binomial quantile rule} of
Definition~\ref{def:interval}, is Requirement~\ref{req:rule-in-force}.
}
\begin{definition}[The construction map]\label{def:construction}
Let
$(\Chatset,{\honest}^{-},{\honest}^{+},n,\ecomp,\Chatcon,\varepsilon^{-},\varepsilon^{+})$
be the data of an \termref{term-admissible-input}{admissible input} (Definition~\ref{def:endpoint-rule}),
where $({\honest}^{-},{\honest}^{+})$ is its \termref{term-honest-envelope}{honest envelope}
(Definition~\ref{def:envelope}) --- an input $\honest$ in the
\termref{term-point-form}{point form} being read as the \termref{term-degenerate-envelope}{degenerate envelope}
$(\honest,\honest)$ --- and let $R$ be an \termref{term-admissible-endpoint-rule}{admissible endpoint rule}
(Definition~\ref{def:endpoint-rule}). With the \termref{term-guarded-endpoints}{guarded endpoints}
$k^{\pm}$ of $R$, \eqref{eq:guard}, define
\begin{equation}\label{eq:boxbounds}
  q^{-}(\hat c)\coloneqq
  \begin{cases}k^{-}(\hat c)/n,&\hat c\in\Chatcon,\\[2pt] 0,&\hat c\notin\Chatcon,\end{cases}
  \qquad
  q^{+}(\hat c)\coloneqq
  \begin{cases}k^{+}(\hat c)/n,&\hat c\in\Chatcon,\\[2pt] 1,&\hat c\notin\Chatcon.\end{cases}
\end{equation}
The \termdef{term-construction-map}{construction map} $\mathsf{Acc}_R$ --- written $\mathsf{Acc}$
where the rule is clear or is the default --- returns the
$\mathrm{polytope}$ geometry in \termref{term-canonical-box-form}{canonical box form} with these
$(\qlo,\qhi)$, i.e.\ the data \eqref{eq:boxrows}, and
\begin{equation}\label{eq:accset}
  \acSet\ \coloneqq\ \{\mathbf{q}\in\simplex{\Chatset}:\
  q^{-}(\hat c)\le q(\hat c)\le q^{+}(\hat c)\ \ \forall\hat c\in\Chatset\}.
\end{equation}
\end{definition}
\claimmeta{\metalabel{Uses}\mref{Def.}{def:taxonomy}, \mref{Def.}{def:envelope}, \mref{Def.}{def:endpoint-rule}, \mref{Def.}{def:interval}
\metasep \metalabel{Used by}\mref{Rem.}{rem:guard-activation}, \mref{Thm.}{thm:box-construction}, \mref{Rem.}{rem:selector}, \mref{Thm.}{thm:completeness-bound}, \mref{Lem.}{lem:widening}, \mref{Thm.}{thm:monotonicity-limits}, \mref{Prop.}{prop:bracketed-endpoints}
\metasep newly derived
\metasep implemented (effective: accepted)
\metasep \metalabel{Code}\coderef{https://github.com/fraunhofersg/certiqs-key/blob/v0.1.0/src/certiqs_key/statistics/construction.py}{certiqs_key.statistics.construction:construct}}

% claim statistics/binomial-box/thm:box-construction
\claimsummary{statistics/binomial-box/thm:box-construction}{For any admissible endpoint rule (count bounds meeting binomial tail budgets), the acceptance box has exact rational bounds and contains every law within the honest envelope's bounds. It is non-empty and compact and, by construction, an admissible acceptance set (\xref{meat-core/def:admissible-acceptance-set}).}
\claimbookmark{statistics/binomial-box/thm:box-construction}{Theorem}{thm:box-construction}{Properties of the constructed box}
\claimleadin{%
Whatever admissible rule produced the endpoints, the output is a valid
acceptance set that accepts every honest law of the envelope: the guard, not
the rule, places the envelope inside the box.
}
\begin{mainclaim}
\begin{theorem}[Properties of the constructed box]\label{thm:box-construction}
On every admissible input (Definition~\ref{def:endpoint-rule}) and for every
admissible endpoint rule $R$, under a \termref{term-conservative-evaluation}{conservative evaluation}
(Definition~\ref{def:conservative-evaluation}):
\begin{enumerate}[label=(\roman*),leftmargin=2.2em]
\item $\mathsf{Acc}_R$ is well defined, and all of its output data are exact
rationals: each $q^{\pm}(\hat c)$ is either an integer multiple of $1/n$ or one
of $0,1$, and the entries of $B$ lie in $\{0,\pm1\}$;
\item $0\le\qlo\le\qhi\le\mathbf{1}$ componentwise, so the pair
$(\qlo,\qhi)$ meets the conditions of the canonical box form of
Definition~\ref{def:taxonomy};
\item $\mathsf{H}\subseteq\acSet$, where $\mathsf{H}$ is the
\termref{term-envelope-set}{envelope set} \eqref{eq:envelope-set} of the input envelope: every
honest law the envelope admits is accepted, and for the degenerate envelope
this reads $\honest\in\acSet$. The envelope set is non-empty; consequently
$\acSet\neq\varnothing$ and, by Theorem~\ref{thm:geometry-validation}\,(iii),
$\sum_{\hat c}q^{-}(\hat c)\le1\le\sum_{\hat c}q^{+}(\hat c)$;
\item $\acSet$ is a non-empty, closed, convex, compact subset of
$\simplex{\Chatset}$; \xref{meat-core/def:admissible-acceptance-set} holds for the output of
$\mathsf{Acc}_R$ by construction, for every admissible input and every
admissible rule.
\end{enumerate}
\end{theorem}
\end{mainclaim}
\claimmeta{\metalabel{Uses}\mref{Def.}{def:taxonomy}, \mref{Thm.}{thm:geometry-validation}, \mref{Def.}{def:envelope}, \mref{Def.}{def:conservative-evaluation}, \mref{Def.}{def:endpoint-rule}, \mref{Lem.}{lem:guard}, \mref{Def.}{def:construction}
\metasep \metalabel{Assumes}\aref{A-ALPHABET}{ass:alphabet}
\metasep \metalabel{Used by}\mref{Rem.}{rem:envelope}, \mref{Thm.}{thm:honest-containment}, \mref{Thm.}{thm:monotonicity-limits}
\metasep newly derived
\metasep implemented (effective: accepted)
\metasep \metalabel{Code}\coderef{https://github.com/fraunhofersg/certiqs-key/blob/v0.1.0/src/certiqs_key/statistics/construction.py}{certiqs_key.statistics.construction:construct}
\metasep \metalabel{Oracles}\coderef{https://github.com/fraunhofersg/certiqs-key/blob/v0.1.0/tests/oracles/registry_statistics_acceptance.py}{registry_statistics_acceptance.py}}

\begin{proof}
\emph{(i).} Well-definedness: $R$ returns integers in $\{0,\dots,n\}$
(Definition~\ref{def:endpoint-rule}\,(a)), so the guarded endpoints
\eqref{eq:guard} exist and are unique, and
\eqref{eq:boxbounds}--\eqref{eq:accset} are explicit formulas in them. The
rationality statement is immediate: $k^{\pm}(\hat c)$ are integers in
$\{0,\dots,n\}$ and $n\ge1$.

\emph{(ii).} For $\hat c\in\Chatcon$, \eqref{eq:ordering} gives $0\le
k^{-}(\hat c)\le k^{+}(\hat c)\le n$, hence $0\le q^{-}(\hat c)\le q^{+}(\hat
c)\le1$. For $\hat c\notin\Chatcon$, $0\le0\le1\le1$.

\emph{(iii).} Let $\mathbf{p}\in\mathsf{H}$, let $\hat c\in\Chatcon$ and
write $p^{\mp}=p^{\mp}_{\hat c}$. By \eqref{eq:ordering} and
\eqref{eq:envelope-set},
$q^{-}(\hat c)=k^{-}(\hat c)/n\le\lfloor np^{-}\rfloor/n\le np^{-}/n=p^{-}
\le p(\hat c)$ and
$q^{+}(\hat c)=k^{+}(\hat c)/n\ge\lceil np^{+}\rceil/n\ge np^{+}/n=p^{+}
\ge p(\hat c)$, so $q^{-}(\hat c)\le p(\hat c)\le q^{+}(\hat c)$. For
$\hat c\notin\Chatcon$ the bounds are $0\le p(\hat c)\le1$, which hold because
$\mathbf{p}\in\simplex{\Chatset}$. Since also
$\mathbf{p}\in\simplex{\Chatset}$, all defining conditions of
\eqref{eq:accset} are met and $\mathbf{p}\in\acSet$; hence
$\mathsf{H}\subseteq\acSet$. For the degenerate envelope
$\mathsf{H}=\{\honest\}$ (Definition~\ref{def:envelope}).

Non-emptiness of $\mathsf{H}$: by Definition~\ref{def:envelope} the pair
$({\honest}^{-},{\honest}^{+})$ is a canonical box form with denotation
$\mathsf{H}$ whose bounds satisfy \eqref{eq:boxfeas}, so
Theorem~\ref{thm:geometry-validation}\,(iii), by which \eqref{eq:boxfeas} and
non-emptiness are equivalent, gives $\mathsf{H}\neq\varnothing$.
Hence $\acSet\supseteq\mathsf{H}$ is non-empty. The box sums follow from the
same equivalence applied to the non-empty $\acSet$, a canonical box by (ii).

\emph{(iv).} Theorem~\ref{thm:geometry-validation}\,(i) applied to the geometry of (ii),
whose \termref{term-validation-predicate}{validation predicate} holds by (iii).
\end{proof}

% claim statistics/binomial-box/rem:selector
\claimsummary{statistics/binomial-box/rem:selector}{Constraining fewer outcomes narrows each interval but leaves more coordinates bounded only by the simplex, so no choice of the constrained outcomes dominates. Completeness holds for every choice, and the unconstrained outcomes keep the bounds 0 and 1, so the box always has two rows per outcome.}
\claimbookmark{statistics/binomial-box/rem:selector}{Remark}{rem:selector}{The role of the constrained outcome set}
\claimleadin{%
Which outcomes are constrained is part of the input; the default, all of
them, is Requirement~\ref{req:request-defaults}.
}
\begin{remark}[The role of the constrained outcome set]\label{rem:selector}
Constraining fewer outcomes concentrates the budget on those that are
constrained, which narrows each interval; it also weakens the set, because the
unconstrained coordinates are then bounded only by the simplex. Both effects
are real and pull in opposite directions, so no choice of $\Chatcon$
dominates. Completeness holds for whatever $\Chatcon$ is supplied
(Theorem~\ref{thm:completeness-bound}), and the unconstrained coordinates receive the vacuous
bounds $[0,1]$ of \eqref{eq:boxbounds} rather than being omitted from $B$ ---
so the row count $m_{\mathrm{row}}=2|\Chatset|$ does not depend on $\Chatcon$
and a consumer never has to infer a missing row.
\end{remark}
\claimmeta{\metalabel{Uses}\mref{Def.}{def:construction}
\metasep \metalabel{Used by}\mref{Req.}{req:request-defaults}
\metasep newly derived
\metasep accepted}

\section{Completeness}\label{group-statistics/binomial-box/lem:envelope-dominance}
\sectionintro{A union bound combines the allocated tails into a bound on honest rejection, with no independence between outcome counts. A count whose per-round conditional success probability stays between two bounds has tails no heavier than the binomial laws at those bounds (\hyperlink{claim-statistics/binomial-box/lem:envelope-dominance}{Lemma~\ref*{lem:envelope-dominance}}), so independent rounds are not needed either. The main result, \hyperlink{claim-statistics/binomial-box/thm:completeness-bound}{Theorem~\ref*{thm:completeness-bound}}, bounds the honest run's probability of failing the test by \ensuremath{\varepsilon}*, the sum of those tails, and \ensuremath{\varepsilon}* by the budget; the remarks say how \ensuremath{\varepsilon}* is evaluated and reported, why soundness does not depend on it, and what an envelope buys and costs.}

% claim statistics/binomial-box/lem:envelope-dominance
\claimsummary{statistics/binomial-box/lem:envelope-dominance}{A count of successes whose conditional success probability in each round, given the past, lies in [a, b] has an upper tail no heavier than the binomial law at b and a lower tail no heavier than that at a. The box can therefore be built from binomial tails without independent rounds.}
\claimbookmark{statistics/binomial-box/lem:envelope-dominance}{Lemma}{lem:envelope-dominance}{Binomial dominance under bounded conditional success probabilities}
\claimleadin{%
This is the classical stochastic dominance of a sum of adapted Bernoulli
variables by a binomial law, derived here in full by a backward induction on
the conditional tail. Under Assumption~\ref{ass:hon-env} the indicators of
one outcome over the rounds of an honest run are such a sequence, with $a$
and $b$ the envelope's two values at that outcome (Theorem~\ref{thm:completeness-bound}).
}
\begin{lemma}[Binomial dominance under bounded conditional success probabilities]\label{lem:envelope-dominance}
Let $\mathcal{F}_0\subseteq\mathcal{F}_1\subseteq\dots\subseteq\mathcal{F}_n$
be a filtration on a probability space, $n\ge1$, and let
$\chi_1,\dots,\chi_n$ be $\{0,1\}$-valued random variables with $\chi_j$
$\mathcal{F}_j$-measurable. Let $0\le a\le b\le1$ and suppose that, for every
$j\in\{1,\dots,n\}$, almost surely
\begin{equation}\label{eq:dominance-hyp}
  a\ \le\ P_j\ \coloneqq\ \Pr\big[\chi_j=1\ \big|\ \mathcal{F}_{j-1}\big]\ \le\ b .
\end{equation}
Put $S\coloneqq\sum_{j=1}^{n}\chi_j$, and for $m\in\{0,\dots,n\}$,
$r\in[0,1]$ and an integer $k$ write
$G^{r}_{m}(k)\coloneqq\Pr[\mathrm{B}^{r}_{m}\ge k]$, where
$\mathrm{B}^{r}_{m}$ is binomial with parameters $(m,r)$ ($\mathrm{B}^{r}_{0}=0$).
Then for every integer $k$
\begin{equation}\label{eq:dominance}
  \Pr[S\ge k]\ \le\ \Pr\big[\mathrm{B}^{b}_{n}\ge k\big],
  \qquad
  \Pr[S\le k]\ \le\ \Pr\big[\mathrm{B}^{a}_{n}\le k\big],
\end{equation}
and equivalently
\begin{equation}\label{eq:dominance-strict}
  \Pr[S>k]\ \le\ \Pr\big[\mathrm{B}^{b}_{n}>k\big],
  \qquad
  \Pr[S<k]\ \le\ \Pr\big[\mathrm{B}^{a}_{n}<k\big].
\end{equation}
\end{lemma}
\claimmeta{\metalabel{Used by}\mref{Thm.}{thm:completeness-bound}
\metasep newly derived
\metasep accepted}

\begin{proof}
\emph{Two facts about the binomial tail.} (1) For $m\ge0$ and every integer
$k$, $G^{r}_{m}(k-1)-G^{r}_{m}(k)=\Pr[\mathrm{B}^{r}_{m}=k-1]\ge0$. (2) The
\emph{Pascal recursion of the binomial tail}: for $m\ge0$ and every integer
$k$,
\begin{equation}\label{eq:pascal-tail}
  G^{r}_{m+1}(k)\ =\ r\,G^{r}_{m}(k-1)+(1-r)\,G^{r}_{m}(k).
\end{equation}
Indeed $\mathrm{B}^{r}_{m+1}$ has the law of $X+Y$ with $Y$ binomial
$(m,r)$ and $X$ Bernoulli$(r)$ independent of $Y$, and conditioning on $X$
gives $\Pr[X+Y\ge k]=r\Pr[Y\ge k-1]+(1-r)\Pr[Y\ge k]$.

\emph{The upper tail, in conditional form.} For $j\in\{0,\dots,n\}$ put
$S_{>j}\coloneqq\sum_{i=j+1}^{n}\chi_i$ (so $S_{>0}=S$ and $S_{>n}=0$). We show
by backward induction on $j$ that for every integer $k$, almost surely,
\begin{equation}\label{eq:dominance-cond}
  \Pr\big[S_{>j}\ge k\ \big|\ \mathcal{F}_j\big]\ \le\ G^{b}_{n-j}(k).
\end{equation}
\emph{Base $j=n$.} $S_{>n}=0$, so the left side is $\mathbf{1}\{k\le0\}
=G^{b}_{0}(k)$.
\emph{Step from $j$ to $j-1$}, for $j\in\{1,\dots,n\}$, with $m\coloneqq n-j$.
Since $S_{>j-1}=\chi_j+S_{>j}$ and $\chi_j\in\{0,1\}$,
\begin{equation*}
  \mathbf{1}\{S_{>j-1}\ge k\}
  =\chi_j\,\mathbf{1}\{S_{>j}\ge k-1\}+(1-\chi_j)\,\mathbf{1}\{S_{>j}\ge k\}.
\end{equation*}
Conditioning on $\mathcal{F}_j$, with $\chi_j$ $\mathcal{F}_j$-measurable and
bounded, and using the induction hypothesis at $k-1$ and at $k$ together with
$\chi_j\ge0$ and $1-\chi_j\ge0$, gives almost surely
\begin{equation*}
  \Pr\big[S_{>j-1}\ge k\,\big|\,\mathcal{F}_j\big]
  \le\chi_j\,G^{b}_{m}(k-1)+(1-\chi_j)\,G^{b}_{m}(k).
\end{equation*}
Conditioning further on $\mathcal{F}_{j-1}\subseteq\mathcal{F}_j$ (the tower
property, and monotonicity of conditional expectation) gives almost surely
\begin{equation*}
  \Pr\big[S_{>j-1}\ge k\,\big|\,\mathcal{F}_{j-1}\big]
  \le G^{b}_{m}(k)+P_j\,\big(G^{b}_{m}(k-1)-G^{b}_{m}(k)\big)
  \le G^{b}_{m}(k)+b\,\big(G^{b}_{m}(k-1)-G^{b}_{m}(k)\big)
  =G^{b}_{m+1}(k),
\end{equation*}
the second inequality by $P_j\le b$ from \eqref{eq:dominance-hyp} and fact
(1), the equality by \eqref{eq:pascal-tail}. As $m+1=n-(j-1)$, this is
\eqref{eq:dominance-cond} at $j-1$. At $j=0$, taking expectations gives
$\Pr[S\ge k]\le G^{b}_{n}(k)$, the first inequality of \eqref{eq:dominance}.
Only the upper bound $b$ was used.

\emph{The lower tail.} Put $\chi'_j\coloneqq1-\chi_j$. These are
$\{0,1\}$-valued and adapted, and
$\Pr[\chi'_j=1\,|\,\mathcal{F}_{j-1}]=1-P_j\le1-a$ almost surely by
\eqref{eq:dominance-hyp}. The upper-tail statement applied to them with the
bound $1-a$ gives, for every integer $k$,
$\Pr[S\le k]=\Pr[\,n-S\ge n-k\,]\le\Pr[\mathrm{B}^{1-a}_{n}\ge n-k]$. Since
$n-\mathrm{B}^{a}_{n}$ is binomial with parameters $(n,1-a)$ (it counts the
failures), the right side is $\Pr[n-\mathrm{B}^{a}_{n}\ge n-k]
=\Pr[\mathrm{B}^{a}_{n}\le k]$, the second inequality of
\eqref{eq:dominance}.

\emph{The strict forms.} $S$ and the binomial variables are integer-valued, so
$\{S>k\}=\{S\ge k+1\}$ and $\{S<k\}=\{S\le k-1\}$, and likewise for
$\mathrm{B}^{b}_{n}$ and $\mathrm{B}^{a}_{n}$; \eqref{eq:dominance} at $k+1$
and at $k-1$ is \eqref{eq:dominance-strict}, and conversely.
\end{proof}

% claim statistics/binomial-box/thm:completeness-bound
\claimsummary{statistics/binomial-box/thm:completeness-bound}{If each honest round's conditional outcome law stays inside the envelope, an honest run fails the acceptance test with probability at most \ensuremath{\varepsilon}*, the sum of the binomial tails at the envelope's ends beyond the box's count bounds, and \ensuremath{\varepsilon}* is at most the completeness budget. No independence is needed.}
\claimbookmark{statistics/binomial-box/thm:completeness-bound}{Theorem}{thm:completeness-bound}{The completeness bound}
\claimleadin{%
The theorem covers a single honest law as a special case: Assumption A-HON,
the acceptance statistics being independent and identically distributed with
law $\honest$, is Assumption~\ref{ass:hon-env} for the
\termref{term-degenerate-envelope}{degenerate envelope} $(\honest,\honest)$
(Lemma~\ref{lem:hon-degenerate}). On an evaluation that is not conservative,
an upward-rounded $k^{-}$ or a downward-rounded $k^{+}$ can break the bound by
the reported completeness figure $\varepsilon^{\star}$: a figure is reported
that the set does not meet (Lemma~\ref{lem:widening}\,(v)).
}
\begin{mainclaim}
\begin{theorem}[The completeness bound]\label{thm:completeness-bound}
Assume a \termref{term-conservative-evaluation}{conservative evaluation}
(Definition~\ref{def:conservative-evaluation}), and assume A-HON-ENV
(Assumption~\ref{ass:hon-env}) for the
\termref{term-honest-envelope}{honest envelope} $({\honest}^{-},{\honest}^{+})$ of the input: on an
honest run there is a filtration
$\mathcal{F}_0\subseteq\dots\subseteq\mathcal{F}_n$ with $\Chat_j$
$\mathcal{F}_j$-measurable and
${\honest}^{-}(\hat c)\le\Pr[\Chat_j=\hat c\,|\,\mathcal{F}_{j-1}]\le
{\honest}^{+}(\hat c)$ almost surely, for every $j$ and $\hat c$. Let
$\acSet$ be the output of $\mathsf{Acc}_R$
(Definition~\ref{def:construction}) for an \termref{term-admissible-endpoint-rule}{admissible endpoint rule}
$R$ (Definition~\ref{def:endpoint-rule}). Then
\begin{equation}\label{eq:completeness}
  \Pr\big[\operatorname{freq}_{\Chat_1^n}\notin\acSet\big]
  \ \le\ \varepsilon^{\star}
  \ \coloneqq\ \sum_{\hat c\in\Chatcon}
  \Big(\Pr\big[\Ncount{\hat c}^{-}<k^{-}(\hat c)\big]
      +\Pr\big[\Ncount{\hat c}^{+}>k^{+}(\hat c)\big]\Big)
  \ \le\ \ecomp ,
\end{equation}
where $\Ncount{\hat c}^{-}$ and $\Ncount{\hat c}^{+}$ are binomial with
parameters $(n,{\honest}^{-}(\hat c))$ and $(n,{\honest}^{+}(\hat c))$
respectively (Definition~\ref{def:endpoint-rule}). Equivalently, the honest
run passes the acceptance test, the event $\OmAT$ of
\xref{meat-core/def:postprocessing-model}, with probability
$\Pr[\OmAT]=\Pr[\operatorname{freq}_{\Chat_1^n}\in\acSet]
\ge1-\varepsilon^{\star}\ge1-\ecomp$. For the degenerate envelope both
$\Ncount{\hat c}^{\mp}$ are binomial $(n,\honest(\hat c))$.
\end{theorem}
\end{mainclaim}
\claimmeta{\metalabel{Uses}\mref{Def.}{def:freq}, \mref{Lem.}{lem:freq-in-simplex}, \mref{Def.}{def:envelope}, \mref{Def.}{def:conservative-evaluation}, \mref{Def.}{def:endpoint-rule}, \mref{Lem.}{lem:guard}, \mref{Def.}{def:construction}, \mref{Lem.}{lem:envelope-dominance}
\metasep \metalabel{Assumes}\aref{A-HON-ENV}{ass:hon-env}, \aref{A-ALPHABET}{ass:alphabet}
\metasep \metalabel{Used by}\mref{Rem.}{rem:epsstar}, \mref{Rem.}{rem:sound-indep}, \mref{Rem.}{rem:envelope}, \mref{Thm.}{thm:monotonicity-limits}, \mref{Rem.}{rem:completeness-on-rules}, \mref{Req.}{req:completeness-report}, \mref{Req.}{req:honest-model-record}, \mref{Rem.}{rem:singleton-honest}
\metasep newly derived
\metasep implemented (effective: accepted)
\metasep \metalabel{Code}\coderef{https://github.com/fraunhofersg/certiqs-key/blob/v0.1.0/src/certiqs_key/statistics/construction.py}{certiqs_key.statistics.construction:report_completeness}
\metasep \metalabel{Oracles}\coderef{https://github.com/fraunhofersg/certiqs-key/blob/v0.1.0/tests/oracles/registry_statistics_acceptance.py}{registry_statistics_acceptance.py}}

\begin{proof}
\emph{Idea.} Each constrained outcome's count has tails no heavier than the
binomial laws at the envelope's two values (Lemma~\ref{lem:envelope-dominance}),
the test fails exactly when some count leaves its \termref{term-guarded-endpoints}{guarded endpoints}, and the
union bound adds the probabilities of those events with no independence;
the guard keeps each tail within its allocation.

\emph{Per-outcome tails.} Fix $\hat c\in\Chatcon$ and let $\chi_j$ be the
random variable that equals $1$ if $\Chat_j=\hat c$ and $0$ otherwise. Each
$\chi_j$ is $\mathcal{F}_j$-measurable, because $\Chat_j$ is, and
$\Pr[\chi_j=1\,|\,\mathcal{F}_{j-1}]=\Pr[\Chat_j=\hat c\,|\,\mathcal{F}_{j-1}]$
lies in $[{\honest}^{-}(\hat c),{\honest}^{+}(\hat c)]$ almost surely by
A-HON-ENV. By Definition~\ref{def:freq},
$\Ncount{\hat c}(\Chat_1^n)=\sum_{j=1}^{n}\chi_j$. Lemma~\ref{lem:envelope-dominance}
with $a={\honest}^{-}(\hat c)$ and $b={\honest}^{+}(\hat c)$, in its strict
form \eqref{eq:dominance-strict}, gives for every integer $k$
\begin{equation*}
  \Pr\big[\Ncount{\hat c}(\Chat_1^n)<k\big]\le\Pr\big[\Ncount{\hat c}^{-}<k\big],
  \qquad
  \Pr\big[\Ncount{\hat c}(\Chat_1^n)>k\big]\le\Pr\big[\Ncount{\hat c}^{+}>k\big].
\end{equation*}
No independence \emph{across} different $\hat c$ is claimed, and none is
used: the counts are in fact negatively related through
$\sum_{\hat c}\Ncount{\hat c}=n$. Nor is independence \emph{across rounds}
used: the per-round conditional laws may change with $j$ and with the past.

\emph{Event inclusion.} By Lemma~\ref{lem:freq-in-simplex}\,(i),
$\operatorname{freq}_{\Chat_1^n}\in\simplex{\Chatset}$ always, so by
\eqref{eq:accset} the event $\{\operatorname{freq}_{\Chat_1^n}\notin\acSet\}$
occurs exactly when some coordinate leaves its box bounds. For
$\hat c\notin\Chatcon$ the bounds are $[0,1]$ by \eqref{eq:boxbounds} and can
never be violated. For $\hat c\in\Chatcon$, and using $n>0$,
$\operatorname{freq}_{\Chat_1^n}(\hat c)<q^{-}(\hat c)$ is equivalent to
$\Ncount{\hat c}(\Chat_1^n)<k^{-}(\hat c)$, and
$\operatorname{freq}_{\Chat_1^n}(\hat c)>q^{+}(\hat c)$ to
$\Ncount{\hat c}(\Chat_1^n)>k^{+}(\hat c)$. Therefore
\begin{equation*}
  \big\{\operatorname{freq}_{\Chat_1^n}\notin\acSet\big\}
  \ =\ \bigcup_{\hat c\in\Chatcon}
  \Big(\big\{\Ncount{\hat c}(\Chat_1^n)<k^{-}(\hat c)\big\}
  \cup\big\{\Ncount{\hat c}(\Chat_1^n)>k^{+}(\hat c)\big\}\Big).
\end{equation*}

\emph{Aggregation.} A probability measure is finitely subadditive: for events
$A_1,\dots,A_M$, $\Pr[\bigcup_{i}A_i]\le\sum_{i}\Pr[A_i]$, with no
independence hypothesis. (Equivalently, in complement form,
$\Pr[\bigcap_i A_i^{c}]\ge1-\sum_i\Pr[A_i]$.) Applying it to the $2|\Chatcon|$
events above, and bounding each by the per-outcome tails at
$k=k^{-}(\hat c)$ and $k=k^{+}(\hat c)$, gives the first inequality of
\eqref{eq:completeness}.

\emph{Budget.} By Lemma~\ref{lem:guard}\,(ii) each term of
$\varepsilon^{\star}$ is at most the corresponding $\varepsilon^{\mp}(\hat c)$,
and the allocation constraint of an \termref{term-admissible-input}{admissible input}
(Definition~\ref{def:endpoint-rule}) sums these to at most $\ecomp$. The
reformulation in terms of $\Pr[\OmAT]$ is the complement, using
$\OmAT=\{\operatorname{freq}_{\hat c_1^n}\in\acSet\}$, the acceptance event
of \xref{meat-core/def:postprocessing-model} on the \termref{term-frequency-vector}{frequency vector} of
Definition~\ref{def:freq}. For the degenerate envelope,
$\Ncount{\hat c}^{-}$ and $\Ncount{\hat c}^{+}$ have the one law binomial
$(n,\honest(\hat c))$, and it is the law of $\Ncount{\hat c}(\Chat_1^n)$
itself: with $\mathrm{B}$ binomial $(n,\honest(\hat c))$, the per-outcome
tails, which are Lemma~\ref{lem:envelope-dominance} at
$a=b=\honest(\hat c)$, give for every integer $k$
$\Pr[\Ncount{\hat c}(\Chat_1^n)<k]\le\Pr[\mathrm{B}<k]$, that is
$\Pr[\Ncount{\hat c}(\Chat_1^n)\ge k]\ge\Pr[\mathrm{B}\ge k]$, and
$\Pr[\Ncount{\hat c}(\Chat_1^n)>k-1]\le\Pr[\mathrm{B}>k-1]$, that is
$\Pr[\Ncount{\hat c}(\Chat_1^n)\ge k]\le\Pr[\mathrm{B}\ge k]$; the two
upper tails agree at every $k$. So the terms of $\varepsilon^{\star}$ are
the tails of $\Ncount{\hat c}(\Chat_1^n)$ itself.
\end{proof}

% claim statistics/binomial-box/rem:epsstar
\claimsummary{statistics/binomial-box/rem:epsstar}{The figure \ensuremath{\varepsilon}* is exactly computable at a rational envelope and generally below the budget, as integrality, the guard and degenerate coordinates leave budget unspent. An upward-rounded evaluation clamped at the budget still bounds the honest failure probability; a downward-rounded one need not.}
\claimbookmark{statistics/binomial-box/rem:epsstar}{Remark}{rem:epsstar}{Exact evaluation of ε\textasciicircum{}* and its conservative report}
\claimleadin{%
Which figure is reported on which rule and evaluation is
Requirement~\ref{req:completeness-report}; the remark gives the mathematics
it rests on, and Remark~\ref{rem:completeness-on-rules} the figures a
closed-form rule can afford.
}
\begin{remark}[Exact evaluation of $\varepsilon^{\star}$ and its conservative report]
\label{rem:epsstar}
Every term of $\varepsilon^{\star}$ is one of the closed forms
\eqref{eq:tails-used} of Corollary~\ref{cor:tail-forms}, at the envelope's
lower value for a lower tail and at its upper value for an upper tail, so
$\varepsilon^{\star}$ is exactly computable at a rational envelope (for a
single law, at rational $\honest$; Lemma~\ref{lem:exact-strict-tails}) and is
generally strictly smaller than
$\ecomp$ --- the integrality of $k^{\pm}$ and the guard of
Remark~\ref{rem:guard} both leave budget unspent, and the degenerate-mass
convention of Definition~\ref{def:endpoint-rule}\,(c) leaves a side's budget
entirely unspent where that side's envelope value is $0$ or $1$
(Proposition~\ref{prop:rule-exact}\,(v) for the exact rule, read side-wise).
Where the tails cannot be evaluated exactly they are rounded \emph{upward}
(a \termref{term-conservative-evaluation}{conservative evaluation},
Definition~\ref{def:conservative-evaluation}): by \eqref{eq:completeness} an upward-rounded
evaluation of $\varepsilon^{\star}$ is still at least the honest failure
probability, whereas a downward-rounded one need not be, and would claim a
completeness the set does not have.

Upward rounding can carry the figure above $\ecomp$ --- in the degenerate
case where every tail meets its allocation exactly and the allocation sums to
exactly $\ecomp$, any strict over-estimate exceeds the budget and leaves the
interval $[0,\ecomp]$. Clamping closes this without weakening anything: for
every upward-rounded evaluation $\tilde\varepsilon\ge\varepsilon^{\star}$,
the figure $\min\{\ecomp,\tilde\varepsilon\}$ is again a valid bound, since
both $\ecomp$ and $\tilde\varepsilon$ are at least $\varepsilon^{\star}$,
which by \eqref{eq:completeness} is at least the honest failure probability;
and it lies in $[0,\ecomp]$ by construction. Exact evaluation needs no clamp,
because $\varepsilon^{\star}\le\ecomp$ is part of \eqref{eq:completeness}.
\end{remark}
\claimmeta{\metalabel{Uses}\mref{Cor.}{cor:tail-forms}, \mref{Lem.}{lem:exact-strict-tails}, \mref{Def.}{def:conservative-evaluation}, \mref{Def.}{def:endpoint-rule}, \mref{Prop.}{prop:rule-exact}, \mref{Rem.}{rem:guard}, \mref{Thm.}{thm:completeness-bound}
\metasep \metalabel{Used by}\mref{Rem.}{rem:completeness-on-rules}, \mref{Req.}{req:completeness-report}
\metasep newly derived
\metasep accepted}

% claim statistics/binomial-box/rem:sound-indep
\claimsummary{statistics/binomial-box/rem:sound-indep}{The honest-round assumptions enter only the completeness statements: if honest behaviour leaves the declared envelope, \ensuremath{\varepsilon}* is void and aborts may be more frequent. Keys from accepted runs are unaffected, since the acceptance set enters the security statement only as a fixed set.}
\claimbookmark{statistics/binomial-box/rem:sound-indep}{Remark}{rem:sound-indep}{Soundness does not depend on the honest-round assumptions}
\begin{remark}[Soundness does not depend on the honest-round assumptions]\label{rem:sound-indep}
Assumption~\ref{ass:hon-env}, and Assumption~\ref{ass:hon} which is its
degenerate case (Lemma~\ref{lem:hon-degenerate}), are statements about the
honest channel and are used only in the completeness statements, such as
Theorems~\ref{thm:completeness-bound} and~\ref{thm:monotonicity-limits}, which read A-HON-ENV, and
Theorem~\ref{thm:singleton-completeness}, which reads A-HON. Under the output-state model of the
test statistics (\xref{reductions/ass:tsm}) the round statistics are
state-dependent in general, so no such concentration is available against an
adversary --- and none is needed, because completeness is a usability
property. If the true honest behaviour leaves the declared envelope --- for a
single declared law, if the true honest channel differs from the declared
$\honest$ --- the figure $\varepsilon^{\star}$ is void and the protocol may
abort more often than advertised; the key produced on the runs that
\emph{do} accept is unaffected, because $\acSet$ enters the security
statement only as the fixed set $\hlow$ is an infimum over. Soundness is
unaffected in either case.
\end{remark}
\claimmeta{\metalabel{Uses}\mref{Thm.}{thm:completeness-bound}
\metasep \metalabel{Assumes}\aref{A-HON-ENV}{ass:hon-env}, \aref{A-HON}{ass:hon}, \extaref{../reductions/contract.pdf}{reductions/ass:tsm}{A-TSM}
\metasep \metalabel{Used by}\mref{Rem.}{rem:completeness-on-rules}
\metasep newly derived
\metasep accepted}

% claim statistics/binomial-box/rem:envelope
\claimsummary{statistics/binomial-box/rem:envelope}{An envelope certifies completeness for an honest channel that drifts within and between runs, and a finite family of honest laws enters through its entrywise minimum and maximum, which is already tight. The price is box width, which costs rate and never soundness.}
\claimbookmark{statistics/binomial-box/rem:envelope}{Remark}{rem:envelope}{Benefit and cost of an honest envelope}
\begin{remark}[Benefit and cost of an honest envelope]\label{rem:envelope}
\emph{What it covers.} Under Assumption~\ref{ass:hon-env} the conditional law
of each honest round, given the past, may be any law inside the envelope, and
may change from round to round. The construction therefore certifies
completeness for an honest channel that drifts within a run, that differs
from run to run (each run's conditional laws lying in the envelope), and
whose drift depends on past outcomes. A \emph{set} of candidate honest
statistics
$\{\mathbf{p}_1,\dots,\mathbf{p}_M\}\subseteq\simplex{\Chatset}\cap\mathbb{Q}^{\Chatset}$
--- several calibrations, the members of a simulated channel family, a
central value with a band around it --- enters through its entrywise minimum
and maximum, ${\honest}^{-}(\hat c)\coloneqq\min_i p_i(\hat c)$ and
${\honest}^{+}(\hat c)\coloneqq\max_i p_i(\hat c)$, which satisfy
Definition~\ref{def:envelope} because each $\mathbf{p}_i$ does (for the sums,
$\sum_{\hat c}{\honest}^{-}(\hat c)\le\sum_{\hat c}p_1(\hat c)=1\le
\sum_{\hat c}{\honest}^{+}(\hat c)$). The box then certifies every law of the
\termref{term-envelope-set}{envelope set} $\mathsf{H}$ (Theorem~\ref{thm:box-construction}\,(iii)), which
contains the convex hull of the family and in general strictly more; the
convex hull itself, as the set of certified laws, would need a joint
concentration argument over several coordinates at once, which this
submodule's per-outcome construction does not make. With the degenerate
envelope nothing changes (Lemma~\ref{lem:hon-degenerate}).

\emph{What it costs.} The box is built from the lower endpoint at
$p^{-}_{\hat c}$ and the upper endpoint at $p^{+}_{\hat c}$, so it widens by
the envelope's width on each coordinate: by Lemma~\ref{lem:guard}\,(i) the
interval $[k^{-}(\hat c),k^{+}(\hat c)]$ contains
$[\lfloor np^{-}_{\hat c}\rfloor,\lceil np^{+}_{\hat c}\rceil]$, whatever the
budget. A wider box lowers the consumer's $\hlow$, since an infimum over a
superset is smaller; the price is paid in rate, never in soundness
(Remark~\ref{rem:sound-indep}).

\emph{No tightening step is needed.} An envelope formed as the entrywise
minimum and maximum of a finite family is already tight in the sense that
matters here: for every $\hat c$, ${\honest}^{-}(\hat c)=p_i(\hat c)$ for some
member $i$, and then
${\honest}^{-}(\hat c)=1-\sum_{\hat c'\neq\hat c}p_i(\hat c')\ge
1-\sum_{\hat c'\neq\hat c}{\honest}^{+}(\hat c')$; symmetrically
${\honest}^{+}(\hat c)\le1-\sum_{\hat c'\neq\hat c}{\honest}^{-}(\hat c')$.
So no coordinate bound can be moved inwards using the simplex constraint
alone; indeed each bound is attained by a member of the family, which lies in
$\mathsf{H}$, so $\min_{\mathbf{p}\in\mathsf{H}}p(\hat c)={\honest}^{-}(\hat c)$
and $\max_{\mathbf{p}\in\mathsf{H}}p(\hat c)={\honest}^{+}(\hat c)$, and the
family's envelope is its own tightening. A declared envelope that is looser
than it could be --- one whose bounds could be moved inwards without losing a
law the channel may have --- is still valid; it is only conservative, costing
box width and hence rate, and the reported $\varepsilon^{\star}$ remains a
bound on the honest failure probability by Theorem~\ref{thm:completeness-bound}. If the true
honest behaviour leaves the declared envelope, $\varepsilon^{\star}$ is void
and soundness is unaffected (Remark~\ref{rem:sound-indep}).
\end{remark}
\claimmeta{\metalabel{Uses}\mref{Def.}{def:envelope}, \mref{Lem.}{lem:guard}, \mref{Thm.}{thm:box-construction}, \mref{Thm.}{thm:completeness-bound}
\metasep newly derived
\metasep accepted}

\section{Honest containment and infeasibility}\label{group-statistics/binomial-box/thm:honest-containment}
\sectionintro{The box contains every law of the \termref{term-envelope-set}{envelope set} for every budget and allocation, so a law of the envelope that the consuming declaration reaches gives a compatible pair and the consuming engine's infeasibility condition cannot arise (\hyperlink{claim-statistics/binomial-box/thm:honest-containment}{Theorem~\ref*{thm:honest-containment}}). The same pair makes the consuming engine's minimal divergence zero, so a positive lower bound on it shows that no law of the envelope is reached (\hyperlink{claim-statistics/binomial-box/rem:hlow-consequence}{Remark~\ref*{rem:hlow-consequence}}). The construction never fails on an \termref{term-admissible-input}{admissible input}, so every refusal comes from the input: a violated condition of an admissible input, or data not of the input's form.}

% claim statistics/binomial-box/thm:honest-containment
\claimsummary{statistics/binomial-box/thm:honest-containment}{Every law of the envelope set lies in the constructed box for every budget and allocation, so a reachable one gives a compatible pair and rules out the consuming engine's infeasibility condition. The construction never fails on an admissible input, so every refusal comes from the input.}
\claimbookmark{statistics/binomial-box/thm:honest-containment}{Theorem}{thm:honest-containment}{Honest containment and infeasibility}
\claimleadin{%
Theorem~\ref{thm:box-construction}\,(iii) places every law of the \termref{term-envelope-set}{envelope set} in
the box; this theorem draws two consequences, one for the consuming engine's
infeasibility condition and one for refusals. The one clause of an
\termref{term-admissible-input}{admissible input} that this module cannot decide, that the alphabet, its
enumeration and the round count are those of the consuming declaration, is
checked where both artefacts are available, in bundle validation
(\href{../../../spec/architecture.md\#36-bundle-validation-is-input-validation-not-a-module}{\nolinkurl{spec/architecture.md}\texttt{\#36-bundle-validation-is-input-validation-not-a-module}};
Requirement~\ref{req:matched-alphabet});
how the engine refuses on the other clauses is
Requirement~\ref{req:input-refusal}.
}
\begin{theorem}[Honest containment and infeasibility]\label{thm:honest-containment}
Let $\acSet$ be the output of $\mathsf{Acc}_R$ on an \termref{term-admissible-input}{admissible input}
(Definition~\ref{def:endpoint-rule}), for any
\termref{term-admissible-endpoint-rule}{admissible endpoint rule} $R$, under a \termref{term-conservative-evaluation}{conservative evaluation}
(Definition~\ref{def:conservative-evaluation}).
\begin{enumerate}[label=(\roman*),leftmargin=2.2em]
\item $\mathsf{H}\subseteq\acSet$, where $\mathsf{H}$ is the
\termref{term-envelope-set}{envelope set} \eqref{eq:envelope-set} of the input envelope, for every
admissible $\ecomp\in[0,1)$, every $\Chatcon$ and every admissible
allocation; for the \termref{term-degenerate-envelope}{degenerate envelope}, $\honest\in\acSet$.
\item Assume A-REACH (Assumption~\ref{ass:reach}): there exists
$\omega^{\mathrm{hon}}\in\Feas$, the feasible set of the consuming
declaration $\Decl$ (\xref{protocol-compiler/def:feasible}), with
$\mathbf{p}^{\mathrm{hon}}\coloneqq\Phimap[\omega^{\mathrm{hon}}]\in\mathsf{H}$
(for the degenerate envelope, $\Phimap[\omega^{\mathrm{hon}}]=\honest$). Then
the pair $(\mathbf{p}^{\mathrm{hon}},\omega^{\mathrm{hon}})$ satisfies
$\mathbf{p}^{\mathrm{hon}}\in\acSet$, $\omega^{\mathrm{hon}}\in\Feas$ and
$\mathbf{p}^{\mathrm{hon}}\ll\Phimap[\omega^{\mathrm{hon}}]$. Consequently the
condition
\begin{equation}\label{eq:infeasible}
  \text{\emph{no} }\mathbf{q}\in\acSet\text{ is absolutely continuous with
  respect to }\Phimap[\omega]\text{ for any }\omega\in\Feas
\end{equation}
is \emph{false} for this $\acSet$.
\item \emph{(refusals come from the input)} The \termref{term-validation-predicate}{validation predicate}
of Theorem~\ref{thm:geometry-validation}\,(iii) is decidable in closed form, and on every
admissible input it is satisfied by the output of $\mathsf{Acc}_R$
unconditionally, for every admissible rule. Hence the construction never
fails on an admissible input, and a refusal can arise only from an input that
is not admissible. On data of the form of an admissible input
(Definition~\ref{def:endpoint-rule}), the violations of its conditions that
are decidable from this module's own inputs are exactly:
\begin{enumerate}[label=(\arabic*),leftmargin=2.2em]
\item $|\Chatset|=0$;
\item $|\Chatset|$ not finite;
\item $n<1$;
\item an honest input in the \termref{term-point-form}{point form} with
$\honest\notin\simplex{\Chatset}$, including a coordinate outside
$\Chatset$;
\item an honest input in the envelope form violating \eqref{eq:envelope}:
some entry of ${\honest}^{-}$ or ${\honest}^{+}$ outside $[0,1]$, some
${\honest}^{-}(\hat c)>{\honest}^{+}(\hat c)$,
$\sum_{\hat c}{\honest}^{-}(\hat c)>1$, $\sum_{\hat c}{\honest}^{+}(\hat
c)<1$, or a coordinate outside $\Chatset$;
\item $\ecomp\notin[0,1)$;
\item some $\varepsilon^{\pm}(\hat c)<0$;
\item $\sum_{\hat c\in\Chatcon}(\varepsilon^{-}(\hat c)+\varepsilon^{+}(\hat
c))>\ecomp$;
\item $\Chatcon\not\subseteq\Chatset$;
\item some coordinate of $\honest$ or of ${\honest}^{\mp}$, or $\ecomp$, or
some $\varepsilon^{\pm}(\hat c)$ is not an exact rational, violating
Definition~\ref{def:conservative-evaluation};
\item the requested endpoint rule, or the requested evaluation, is none of
those of clause (viii) of Definition~\ref{def:endpoint-rule}.
\end{enumerate}
Exactly one clause of an admissible input is \emph{not} decidable from this
module's inputs: clause (ix) of Definition~\ref{def:endpoint-rule}, the
condition of Definition~\ref{def:matched-alphabet} that $\Chatset$, its
enumeration and $n$ are those of the same declaration $\Decl$ the
\exttermref{../foundations/contract.pdf}{term-consumer}{consumer} solves. That clause
is a relation between artefacts, not a property of an input.
\end{enumerate}
\end{theorem}
\claimmeta{\metalabel{Uses}\mref{Def.}{def:matched-alphabet}, \mref{Thm.}{thm:geometry-validation}, \mref{Def.}{def:envelope}, \mref{Def.}{def:conservative-evaluation}, \mref{Def.}{def:endpoint-rule}, \mref{Thm.}{thm:box-construction}
\metasep \metalabel{Assumes}\aref{A-REACH}{ass:reach}, \aref{A-ALPHABET}{ass:alphabet}
\metasep \metalabel{Used by}\mref{Rem.}{rem:hlow-consequence}, \mref{Req.}{req:input-refusal}
\metasep newly derived
\metasep implemented (effective: accepted)
\metasep \metalabel{Code}\coderef{https://github.com/fraunhofersg/certiqs-key/blob/v0.1.0/src/certiqs_key/statistics/construction.py}{certiqs_key.statistics.construction:validate_request}
\metasep \metalabel{Oracles}\coderef{https://github.com/fraunhofersg/certiqs-key/blob/v0.1.0/tests/oracles/registry_statistics_acceptance.py}{registry_statistics_acceptance.py}}

\begin{proof}
\emph{(i).} Theorem~\ref{thm:box-construction}\,(iii).

\emph{(ii).} By (i), $\mathbf{p}^{\mathrm{hon}}\in\mathsf{H}\subseteq
\acSet$. Absolute continuity
$\mathbf{p}^{\mathrm{hon}}\ll\Phimap[\omega^{\mathrm{hon}}]$ reads:
$p^{\mathrm{hon}}(\hat c)>0$ implies
$\Phimap[\omega^{\mathrm{hon}}](\hat c)>0$. The two vectors are equal by
definition of $\mathbf{p}^{\mathrm{hon}}$, so the implication is trivially
true. This exhibits a witness
$(\mathbf{q},\omega)=(\mathbf{p}^{\mathrm{hon}},\omega^{\mathrm{hon}})$ to the
negation of the universally quantified statement \eqref{eq:infeasible}; for
the degenerate envelope the witness is $(\honest,\omega^{\mathrm{hon}})$.

\emph{(iii).} Decidability in closed form is \eqref{eq:boxfeas};
unconditional satisfaction on every admissible input is
Theorem~\ref{thm:box-construction}\,(iii). The enumerated list is the negation, clause by
clause, of the conditions the clauses of an admissible input
(Definition~\ref{def:endpoint-rule}) place on data of its form, with the
single stated exception of the cross-artefact condition of
Definition~\ref{def:matched-alphabet}: an input
of that form for which none of the eleven violations (1)--(11) holds, and
whose alphabet and round count are matched to the consuming declaration, is
admissible, and then
(i)--(ii) apply.
\end{proof}

% claim statistics/binomial-box/rem:hlow-consequence
\claimsummary{statistics/binomial-box/rem:hlow-consequence}{Under \hyperlink{claim-statistics/ass:reach}{A-REACH} a box built from the honest envelope contains the statistics of a feasible state, which rules out an infinite value and a positive minimal divergence in the consuming engine. The assumption is never verified here; a positive bound on that divergence refutes it.}
\claimbookmark{statistics/binomial-box/rem:hlow-consequence}{Remark}{rem:hlow-consequence}{Scope of the statements on infeasibility and minimal divergence}
\begin{remark}[Scope of the statements on infeasibility and minimal divergence]
\label{rem:hlow-consequence}
The consuming engine's contract identifies \eqref{eq:infeasible} as the source of
$\hlow=+\infty$ (\xref{meat-core/rem:hlow-range}) and requires that value to
propagate rather than be clamped (\xref{meat-core/req:infinite-value}).
Theorem~\ref{thm:honest-containment}\,(ii) rules that
condition out for a set built by $\mathsf{Acc}$ from an
\termref{term-honest-envelope}{honest envelope} one of whose laws is reachable --- for a single law,
from a reachable honest distribution --- which is the precise sense in which
this module \emph{surfaces} infeasibility rather than passing it through: the
construction is ruled out as a source of \eqref{eq:infeasible}, so when the
consuming engine meets that condition on a set built by $\mathsf{Acc}$, the
contrapositive of Theorem~\ref{thm:honest-containment}\,(ii) identifies its
cause as a failure of A-REACH, an incompatibility of the declared honest
envelope with the declaration, instead of leaving an extended-real value with
no diagnosis.

The same pair has a finite counterpart. Under A-REACH,
Theorem~\ref{thm:honest-containment}\,(ii) gives
$\mathbf{p}^{\mathrm{hon}}\in\acSet$ with
$\mathbf{p}^{\mathrm{hon}}=\Phimap[\omega^{\mathrm{hon}}]$ and
$\omega^{\mathrm{hon}}\in\Feas$, so the minimal divergence between the
accepted distributions and the statistics of the feasible set, the quantity
$\sfD_{\min}$ of \xref{meat-core/prop:minimal-divergence}, is zero for a set
built by $\mathsf{Acc}$, and the consuming engine's value has no positive
penalty floor. Contrapositively, a positive lower bound on $\sfD_{\min}$,
which the consuming engine forms on the runs it checks
(\xref{meat-core/req:unreachable-acceptance-set}), shows that A-REACH fails
for the declared envelope and the declaration. The converse does not hold:
the box contains the envelope set and may meet the statistics of the
feasible set where no law of the envelope set does, so a bound that is not
positive says nothing about A-REACH.

The statement has two limits. First, it is conditional on A-REACH
(Assumption~\ref{ass:reach}), which this module never verifies, and which the
consuming engine can refute but never establish: whether some
law of the declared honest envelope is producible by the declared protocol
is a joint property of whoever declares the honest model and of the
declaration, and asserting it here would assert something this module cannot
see. Second, it excludes the condition \eqref{eq:infeasible}; whether $\hlow$
is then finite is a statement about the consuming engine's objective and
belongs to that contract, which is why it is quoted here and used nowhere in
a proof. At $\hlow=+\infty$ the consuming engine emits no key-length claim
and reports the mismatch between the declaration and the acceptance set
(\xref{meat-core/req:infinite-value}), and nothing here changes that.
\end{remark}
\claimmeta{\metalabel{Uses}\mref{Thm.}{thm:honest-containment}
\metasep \metalabel{Assumes}\aref{A-REACH}{ass:reach}
\metasep newly derived
\metasep accepted}

\section{Monotonicity and limiting cases}\label{group-statistics/binomial-box/def:monotone-rule}
\sectionintro{Lowering the box's lower bounds and raising its upper ones gives a superset that is failed no more often and has no larger infimum, while the reverse move can break the completeness figure (\hyperlink{claim-statistics/binomial-box/lem:widening}{Lemma~\ref*{lem:widening}}). Under a \termref{term-monotone-allocation-rule}{monotone allocation rule} and an endpoint rule \termref{term-monotone-in-the-allocation}{monotone in the allocation}, a larger budget gives a smaller or equal set and hence a larger or equal infimum for any consumer (\hyperlink{claim-statistics/binomial-box/thm:monotonicity-limits}{Theorem~\ref*{thm:monotonicity-limits}}). A zero budget gives the whole simplex away from degenerate coordinates, a one-letter alphabet a constant set, and for the exact rule at a fixed positive allocation the box closes in on the envelope's box at rate n\textasciicircum{}(\ensuremath{-}1/2) without being nested in the round count.}

% claim statistics/binomial-box/def:monotone-rule
\claimsummary{statistics/binomial-box/def:monotone-rule}{A monotone allocation rule allocates each completeness budget within that budget, every allocated tail budget non-decreasing in it; the uniform rule is one. It is the hypothesis on the allocation under which \hyperlink{claim-statistics/binomial-box/thm:monotonicity-limits}{Theorem~\ref*{thm:monotonicity-limits}} (i) compares the acceptance sets at two budgets.}
\claimbookmark{statistics/binomial-box/def:monotone-rule}{Definition}{def:monotone-rule}{Monotone allocation rule}
\claimleadin{%
Comparing the acceptance sets at two completeness budgets needs the
allocation at the larger budget to give every tail at least as much as the
allocation at the smaller one. Where a request names no allocation, the
uniform rule is the default (Requirement~\ref{req:request-defaults}).
}
\begin{definition}[Monotone allocation rule]\label{def:monotone-rule}
A \termdef{term-monotone-allocation-rule}{monotone allocation rule}
assigns to each budget $e\in[0,1)$ an allocation
$\varepsilon^{\pm}_{(e)}$ on $\Chatcon$ such that
$\sum_{\hat c}(\varepsilon^{-}_{(e)}(\hat c)+\varepsilon^{+}_{(e)}(\hat c))\le
e$ and each $e\mapsto\varepsilon^{\pm}_{(e)}(\hat c)$ is non-decreasing. The
\termdef{term-uniform-rule}{uniform rule}
$\varepsilon^{\pm}_{(e)}(\hat c)=e/(2|\Chatcon|)$ is monotone.
\end{definition}
\claimmeta{\metalabel{Used by}\mref{Thm.}{thm:monotonicity-limits}, \mref{Req.}{req:request-defaults}
\metasep newly derived
\metasep implemented
\metasep \metalabel{Code}\coderef{https://github.com/fraunhofersg/certiqs-key/blob/v0.1.0/src/certiqs_key/statistics/construction.py}{certiqs_key.statistics.construction:uniform_allocation}}

% claim statistics/binomial-box/lem:widening
\claimsummary{statistics/binomial-box/lem:widening}{Lowering the box's lower bounds and raising its upper ones gives a superset that is failed no more often and has no larger infimum, and integer endpoints so moved still meet their tail budgets. Moving them the other way can break the tail budgets and the reported completeness figure \ensuremath{\varepsilon}*.}
\claimbookmark{statistics/binomial-box/lem:widening}{Lemma}{lem:widening}{Widening the box keeps its guarantees}
\claimleadin{%
A \termref{term-conservative-evaluation}{conservative evaluation} (Definition~\ref{def:conservative-evaluation})
can only return a smaller lower and a larger upper endpoint than an exact one,
and the consuming engine's outward rounding of the emitted bounds moves them
the same way. The lemma says that every such move keeps what the construction
guarantees and costs only rate, and that the reverse move does not;
Proposition~\ref{prop:rule-exact}\,(iii) and Lemma~\ref{lem:guard}\,(ii) are
its counterparts for a rule's own interval.
}
\begin{lemma}[Widening the box keeps its guarantees]\label{lem:widening}
Let $\acSet$ be the output of $\mathsf{Acc}_R$
(Definition~\ref{def:construction}) on an \termref{term-admissible-input}{admissible input}
(Definition~\ref{def:endpoint-rule}) for an \termref{term-admissible-endpoint-rule}{admissible endpoint rule}
$R$, with box bounds $\qlo,\qhi$ \eqref{eq:boxbounds} and
\termref{term-guarded-endpoints}{guarded endpoints} $k^{\mp}$ \eqref{eq:guard}. Let
$\tilde{\mathbf{q}}^{-},\tilde{\mathbf{q}}^{+}\in\mathbb{R}^{\Chatset}$
satisfy $\tilde q^{-}(\hat c)\le q^{-}(\hat c)$ and
$\tilde q^{+}(\hat c)\ge q^{+}(\hat c)$ for every $\hat c\in\Chatset$, and put
\begin{equation}\label{eq:widened-set}
  \acSet'\ \coloneqq\ \{\mathbf{q}\in\simplex{\Chatset}:\
  \tilde q^{-}(\hat c)\le q(\hat c)\le\tilde q^{+}(\hat c)\ \
  \forall\hat c\in\Chatset\}.
\end{equation}
\begin{enumerate}[label=(\roman*),leftmargin=2.2em]
\item \emph{(containment)} $\acSet\subseteq\acSet'$.
\item \emph{(failure probability)} For all $\Chatset$-valued random
variables $\Chat_1,\dots,\Chat_n$ on one probability space,
\begin{equation*}
  \Pr\big[\operatorname{freq}_{\Chat_1^n}\notin\acSet'\big]
  \ \le\ \Pr\big[\operatorname{freq}_{\Chat_1^n}\notin\acSet\big].
\end{equation*}
\item \emph{(infima)} For every function
$\Psi:\simplex{\Chatset}\to[-\infty,+\infty]$,
\begin{equation}\label{eq:widened-inf}
  \inf_{\mathbf{q}\in\acSet'}\Psi(\mathbf{q})
  \ \le\ \inf_{\mathbf{q}\in\acSet}\Psi(\mathbf{q}).
\end{equation}
\item \emph{(integer endpoints)} Let
$\tilde k^{-},\tilde k^{+}:\Chatcon\to\mathbb{Z}$ satisfy
$0\le\tilde k^{-}(\hat c)\le k^{-}(\hat c)$ and
$k^{+}(\hat c)\le\tilde k^{+}(\hat c)\le n$ for every $\hat c\in\Chatcon$,
and let $\tilde{\mathbf{q}}^{\mp}$ be given by \eqref{eq:boxbounds} with
$\tilde k^{\mp}$ in place of $k^{\mp}$. Then $\tilde{\mathbf{q}}^{\mp}$
satisfy the hypothesis above, so (i)--(iii) hold; the pair
$(\tilde{\mathbf{q}}^{-},\tilde{\mathbf{q}}^{+})$ meets the conditions of the
\termref{term-canonical-box-form}{canonical box form} of Definition~\ref{def:taxonomy}, and its
denotation \eqref{eq:boxset} is $\acSet'$; and for every $\hat c\in\Chatcon$,
with $\Ncount{\hat c}^{\mp}$ as in Definition~\ref{def:endpoint-rule},
\begin{equation}\label{eq:widened-tails}
  \Pr\big[\Ncount{\hat c}^{-}<\tilde k^{-}(\hat c)\big]
  \le\Pr\big[\Ncount{\hat c}^{-}<k^{-}(\hat c)\big]
  \le\varepsilon^{-}(\hat c),
  \qquad
  \Pr\big[\Ncount{\hat c}^{+}>\tilde k^{+}(\hat c)\big]
  \le\Pr\big[\Ncount{\hat c}^{+}>k^{+}(\hat c)\big]
  \le\varepsilon^{+}(\hat c).
\end{equation}
Summed over $\Chatcon$, the tails at $\tilde k^{\mp}$ therefore total at most
the tails at $k^{\mp}$, and at most $\ecomp$.
\item \emph{(not the reverse)} Replacing $k^{-}(\hat c)$ by a larger integer
or $k^{+}(\hat c)$ by a smaller one keeps, in general, neither the budgets of
\eqref{eq:widened-tails} nor the tail sum at $k^{\mp}$ as a bound on the
probability that the \termref{term-frequency-vector}{frequency vector} leaves the set. Take
$\Chatset=\{\hat c_0,\hat c_1\}$, the \termref{term-degenerate-envelope}{degenerate envelope} at
$\honest=(\tfrac12,\tfrac12)$ (Definition~\ref{def:envelope}), $n=3$,
$\Chatcon=\{\hat c_0\}$,
$\varepsilon^{-}(\hat c_0)=\varepsilon^{+}(\hat c_0)=\tfrac25$,
$\ecomp=\tfrac45$ and the \termref{term-exact-binomial-quantile-rule}{exact binomial quantile rule} of
Definition~\ref{def:interval}. Then $k^{-}(\hat c_0)=1$, $k^{+}(\hat c_0)=2$,
and the tails at these endpoints sum to $\tfrac14$. Replacing
$k^{-}(\hat c_0)$ by $2$, or $k^{+}(\hat c_0)$ by $1$, gives a tail of
$\tfrac12>\tfrac25$ on the replaced side; and for $\Chat_1,\Chat_2,\Chat_3$
independent with law $\honest$, the frequency vector lies outside the set
\eqref{eq:accset} built from the replaced endpoints with probability
$\tfrac58>\tfrac14$, while it lies outside $\acSet$ with probability
$\tfrac14$.
\end{enumerate}
\end{lemma}
\claimmeta{\metalabel{Uses}\mref{Def.}{def:freq}, \mref{Def.}{def:taxonomy}, \mref{Def.}{def:envelope}, \mref{Def.}{def:endpoint-rule}, \mref{Lem.}{lem:guard}, \mref{Def.}{def:interval}, \mref{Def.}{def:construction}
\metasep \metalabel{Used by}\mref{Thm.}{thm:monotonicity-limits}
\metasep newly derived
\metasep accepted}

\begin{proof}
\emph{(i).} Let $\mathbf{q}\in\acSet$ and $\hat c\in\Chatset$. By
\eqref{eq:accset} and the hypothesis,
$\tilde q^{-}(\hat c)\le q^{-}(\hat c)\le q(\hat c)\le q^{+}(\hat c)\le
\tilde q^{+}(\hat c)$. Since also $\mathbf{q}\in\simplex{\Chatset}$, every
defining condition of \eqref{eq:widened-set} holds, and
$\mathbf{q}\in\acSet'$.

\emph{(ii).} The frequency vector $\operatorname{freq}_{\Chat_1^n}$ of
Definition~\ref{def:freq} is a function of the finitely valued string
$\Chat_1^n$, so $\{\operatorname{freq}_{\Chat_1^n}\in S\}$ is the finite
union of the events $\{\Chat_1^n=\hat c_1^n\}$ over the strings with
$\operatorname{freq}_{\hat c_1^n}\in S$, an event for every set $S$. By (i),
$\{\operatorname{freq}_{\Chat_1^n}\in\acSet\}\subseteq
\{\operatorname{freq}_{\Chat_1^n}\in\acSet'\}$; taking complements,
$\{\operatorname{freq}_{\Chat_1^n}\notin\acSet'\}\subseteq
\{\operatorname{freq}_{\Chat_1^n}\notin\acSet\}$, and a probability measure
is monotone under inclusion.

\emph{(iii).} For every $\mathbf{q}\in\acSet$, (i) gives
$\mathbf{q}\in\acSet'$, so the infimum over $\acSet'$ is at most
$\Psi(\mathbf{q})$. The infimum over $\acSet'$ is therefore a lower bound of
$\{\Psi(\mathbf{q}):\mathbf{q}\in\acSet\}$, hence at most its greatest lower
bound, which is \eqref{eq:widened-inf}. This holds in the extended reals for
every $\Psi$, with no continuity or measurability required, and also when
$\acSet$ is empty, its infimum then being $+\infty$.

\emph{(iv).} \emph{The hypothesis.} For $\hat c\in\Chatcon$,
\eqref{eq:boxbounds} gives
$\tilde q^{-}(\hat c)=\tilde k^{-}(\hat c)/n\le k^{-}(\hat c)/n=q^{-}(\hat c)$
and $\tilde q^{+}(\hat c)=\tilde k^{+}(\hat c)/n\ge k^{+}(\hat c)/n
=q^{+}(\hat c)$, using $n\ge1$; for $\hat c\notin\Chatcon$ both pairs of
bounds are $0$ and $1$. So $\tilde{\mathbf{q}}^{\mp}$ satisfy the hypothesis,
and (i)--(iii) apply.

\emph{The canonical box form.} Every $\tilde q^{\mp}(\hat c)$ is an integer
divided by $n$, or $0$ or $1$, hence rational. For $\hat c\in\Chatcon$,
\eqref{eq:ordering} of Lemma~\ref{lem:guard}\,(i) gives
$k^{-}(\hat c)\le k^{+}(\hat c)$, so
$0\le\tilde k^{-}(\hat c)\le k^{-}(\hat c)\le k^{+}(\hat c)\le
\tilde k^{+}(\hat c)$ and $0\le\tilde q^{-}(\hat c)\le\tilde q^{+}(\hat c)$;
for $\hat c\notin\Chatcon$, $0\le0\le1$. These are the conditions of the
canonical box form, and its denotation \eqref{eq:boxset} with these bounds is
\eqref{eq:widened-set}.

\emph{The tails.} By Lemma~\ref{lem:guard}\,(ii) the map
$k\mapsto\Pr[\Ncount{\hat c}^{-}<k]$ is non-decreasing and the map
$k\mapsto\Pr[\Ncount{\hat c}^{+}>k]$ is non-increasing on the integers, so
$\tilde k^{-}(\hat c)\le k^{-}(\hat c)$ and
$\tilde k^{+}(\hat c)\ge k^{+}(\hat c)$ give the first inequality of each
half of \eqref{eq:widened-tails}; the second is \eqref{eq:budget-met} of the
same lemma. Summing the first inequalities over $\hat c\in\Chatcon$ compares
the two tail sums, and summing the second ones bounds the sum at $k^{\mp}$ by
$\sum_{\hat c\in\Chatcon}(\varepsilon^{-}(\hat c)+\varepsilon^{+}(\hat c))$,
which the allocation of an admissible input keeps at most $\ecomp$
(Definition~\ref{def:endpoint-rule}\,(vi)).

\emph{(v).} With the degenerate envelope at $\tfrac12$, the two counts
$\Ncount{\hat c_0}^{-}$ and $\Ncount{\hat c_0}^{+}$ have one law; write
$\Ncount{}$ for a count binomial with parameters $(3,\tfrac12)$, whose law is
$(\tfrac18,\tfrac38,\tfrac38,\tfrac18)$ on $\{0,1,2,3\}$. The exact rule
takes \eqref{eq:kappaminus}: $\Pr[\Ncount{}<0]=0$ and
$\Pr[\Ncount{}<1]=\tfrac18$ are at most $\tfrac25$, while
$\Pr[\Ncount{}<k]\ge\Pr[\Ncount{}<2]=\tfrac12>\tfrac25$ for every $k\ge2$, so
$\kappa^{-}(\hat c_0)=1$. It takes \eqref{eq:kappaplus}:
$\Pr[\Ncount{}>2]=\tfrac18\le\tfrac25$ and $\Pr[\Ncount{}>3]=0$, while
$\Pr[\Ncount{}>k]\ge\Pr[\Ncount{}>1]=\tfrac12>\tfrac25$ for every $k\le1$, so
$\kappa^{+}(\hat c_0)=2$. The guard \eqref{eq:guard} is inert, since
$\lfloor\tfrac32\rfloor=1$ and $\lceil\tfrac32\rceil=2$, so
$k^{-}(\hat c_0)=1$ and $k^{+}(\hat c_0)=2$, and the tails at them are
$\Pr[\Ncount{}<1]+\Pr[\Ncount{}>2]=\tfrac18+\tfrac18=\tfrac14$. The replaced
tails are $\Pr[\Ncount{}<2]=\tfrac12$ and $\Pr[\Ncount{}>1]=\tfrac12$, each
above $\tfrac25$.

For the failure probabilities, let $\Chat_1,\Chat_2,\Chat_3$ be independent
with law $\honest$. The count of $\hat c_0$ in $\Chat_1^3$ is then binomial
with parameters $(3,\tfrac12)$, so it has the law of $\Ncount{}$. On a
two-letter alphabet a vector of $\simplex{\Chatset}$ is fixed by its
coordinate at $\hat c_0$, and the bounds at $\hat c_1$ are $0$ and $1$ by
\eqref{eq:boxbounds}, so a frequency vector lies in a set \eqref{eq:accset}
built from endpoints $(a,b)$ at $\hat c_0$ exactly when the count of
$\hat c_0$ lies in $\{a,\dots,b\}$. For $\acSet$, with $(a,b)=(1,2)$, the
failure probability is $\Pr[\Ncount{}\in\{0,3\}]=\tfrac14$. With $k^{-}$
replaced by $2$, $(a,b)=(2,2)$ and the failure probability is
$1-\Pr[\Ncount{}=2]=\tfrac58$; with $k^{+}$ replaced by $1$, $(a,b)=(1,1)$
and it is $1-\Pr[\Ncount{}=1]=\tfrac58$. Both exceed $\tfrac14$.
\end{proof}

% claim statistics/binomial-box/lem:cheb
\claimsummary{statistics/binomial-box/lem:cheb}{A binomial (n, p) count deviates from np by at least t \textgreater{} 0 with probability at most np(1 \ensuremath{-} p)/t\ensuremath{^{2}}. This gives the finite-sample width bound of the exact rule in \hyperlink{claim-statistics/binomial-box/thm:monotonicity-limits}{Theorem~\ref*{thm:monotonicity-limits}} (iv).}
\claimbookmark{statistics/binomial-box/lem:cheb}{Lemma}{lem:cheb}{Chebyshev's inequality for a binomial count}
\begin{technicalclaim}
\begin{lemma}[Chebyshev's inequality for a binomial count]\label{lem:cheb}
Let $\Ncount{}$ be binomial with parameters $(n,p)$ and let $t>0$. Then
$\Pr\big[|\Ncount{}-np|\ge t\big]\le np(1-p)/t^{2}$.
\end{lemma}
\end{technicalclaim}
\claimmeta{\metalabel{Used by}\mref{Thm.}{thm:monotonicity-limits}
\metasep newly derived
\metasep accepted}

\begin{proof}
Write $\Ncount{}=\sum_{j=1}^{n}\chi_j$ with $\chi_j$ independent
Bernoulli$(p)$; then $\mathbb{E}[\Ncount{}]=np$ and, by independence,
$\operatorname{Var}(\Ncount{})=\sum_{j}\operatorname{Var}(\chi_j)=np(1-p)$.
Consider the non-negative random variable $(\Ncount{}-np)^{2}$. Pointwise it is
at least $t^{2}$ wherever $|\Ncount{}-np|\ge t$ and at least $0$ elsewhere, so
it dominates the random variable taking the value $t^{2}$ wherever
$|\Ncount{}-np|\ge t$ and $0$ elsewhere. Taking expectations,
$t^{2}\,\Pr\big[|\Ncount{}-np|\ge t\big]
\le\mathbb{E}\big[(\Ncount{}-np)^{2}\big]
=\operatorname{Var}(\Ncount{})=np(1-p)$, which is the claim.
\end{proof}

% claim statistics/binomial-box/thm:monotonicity-limits
\claimsummary{statistics/binomial-box/thm:monotonicity-limits}{The acceptance set shrinks as the budget grows under monotone allocation and endpoint rules, so no infimum over it, h\_low included, decreases. Zero-budget and one-letter limits are explicit; at a fixed positive allocation the exact rule's set tends to the envelope's box as n\textasciicircum{}(\ensuremath{-}1/2), unnested in n.}
\claimbookmark{statistics/binomial-box/thm:monotonicity-limits}{Theorem}{thm:monotonicity-limits}{Monotonicity in the budget and limiting cases}
\claimleadin{%
Parts (i) and (ii) need the endpoint rule itself to be
\termref{term-monotone-in-the-allocation}{monotone in the allocation}, and which rules are is each rule's own
statement: the \termref{term-exact-binomial-quantile-rule}{exact rule} with its
sweep evaluation (Proposition~\ref{prop:rule-exact}\,(ii), and (vii) where
a tail is decided against an upper bound) and the Bernstein
rule (Proposition~\ref{prop:rule-bernstein}\,(ii)) are, while the enclosure
evaluation of the exact rule (Proposition~\ref{prop:bracketed-endpoints})
does not claim it, so (i) and (ii) are not asserted for it. The monotonicity
of $\hlow$ in (ii) is a trade-off, not an improvement: the same change of
budget makes the honest run abort no less often. A one-letter acceptance
alphabet, the case of (vi), imposes no constraint on the \exttermref{../foundations/contract.pdf}{term-consumer}{consumer} and
carries no information; it is a degenerate case to be recognised, not a
regime to be used.
}
\begin{theorem}[Monotonicity in the budget and limiting cases]\label{thm:monotonicity-limits}
Fix $\Chatset$, an \termref{term-honest-envelope}{honest envelope} $({\honest}^{-},{\honest}^{+})$
with \termref{term-envelope-set}{envelope set} $\mathsf{H}$ (Definition~\ref{def:envelope}; a single
law $\honest$ is the \termref{term-degenerate-envelope}{degenerate envelope}, $\mathsf{H}=\{\honest\}$),
$n$, $\Chatcon$ and an \termref{term-admissible-endpoint-rule}{admissible endpoint rule} $R$
(Definition~\ref{def:endpoint-rule}) under a \termref{term-conservative-evaluation}{conservative evaluation}
(Definition~\ref{def:conservative-evaluation}), write $p^{\mp}_{\hat c}$ for the
envelope's values as there, and write $\acSet(e)$ for the output of
$\mathsf{Acc}_R$ at budget $e$ under a fixed monotone allocation rule
(Definition~\ref{def:monotone-rule}).
\begin{enumerate}[label=(\roman*),leftmargin=2.2em]
\item \emph{(budget monotonicity: the set shrinks as the budget grows)} If
$R$ is monotone in the allocation and $0\le e_1\le e_2<1$, then
$\acSet(e_2)\subseteq\acSet(e_1)$.
\item \emph{(consequence for any consumer infimum)} Under the hypothesis of
(i), for every function $\Psi:\simplex{\Chatset}\to[-\infty,+\infty]$ and
$e_1\le e_2$,
\begin{equation}\label{eq:inf-mono}
  \inf_{\mathbf{q}\in\acSet(e_1)}\Psi(\mathbf{q})
  \ \le\ \inf_{\mathbf{q}\in\acSet(e_2)}\Psi(\mathbf{q}) .
\end{equation}
In particular, taking
$\Psi(\mathbf{q})=\inf_{\omega\in\Feas}O_{\alpha}(\mathbf{q},\omega)$, the
quantity $\hlow$, which \xref{meat-core/def:feasible} defines as
$\inf_{\mathbf{q}\in\acSet}\inf_{\omega\in\Feas}O_{\alpha}$, is
\emph{non-decreasing} in $\ecomp$: a weaker completeness guarantee buys a
larger or equal entropy bound.
\item \emph{(limit $\ecomp=0$, every rule)} At $e=0$ every
$\varepsilon^{\pm}(\hat c)=0$ and
\begin{equation*}
\begin{aligned}
  \acSet(0)=\big\{\mathbf{q}\in\simplex{\Chatset}:\
  &q(\hat c)=0\ \text{for }\hat c\in\Chatcon\text{ with }p^{+}_{\hat c}=0,\\
  &q(\hat c)=1\ \text{for }\hat c\in\Chatcon\text{ with }p^{-}_{\hat
  c}=1\big\},
\end{aligned}
\end{equation*}
with $\varepsilon^{\star}=0$. If $p^{+}_{\hat c}>0$ and $p^{-}_{\hat c}<1$ for
every $\hat c\in\Chatcon$ (for the degenerate envelope: if
$0<\honest(\hat c)<1$ for every $\hat c\in\Chatcon$), then
$\acSet(0)=\simplex{\Chatset}$: a zero completeness budget buys the whole
simplex and no constraint at all.
\item \emph{(limit $n\to\infty$, the exact rule)} Let $R$ be the exact rule
of Definition~\ref{def:interval}, let $\hat c\in\Chatcon$ and write
$p^{\mp}=p^{\mp}_{\hat c}$. Then
\begin{equation}\label{eq:width}
  0\ \le\ q^{+}(\hat c)-p^{+}\ \le\
  \sqrt{\frac{p^{+}(1-p^{+})}{n\,\varepsilon^{+}(\hat c)}}+\frac1n,
  \qquad
  0\ \le\ p^{-}-q^{-}(\hat c)\ \le\
  \sqrt{\frac{p^{-}(1-p^{-})}{n\,\varepsilon^{-}(\hat c)}}+\frac1n ,
\end{equation}
the first whenever $0<p^{+}<1$ and $\varepsilon^{+}(\hat c)>0$, the second
whenever $0<p^{-}<1$ and $\varepsilon^{-}(\hat c)>0$.
Consequently, if $\Chatcon=\Chatset$, $0<p^{\mp}_{\hat c}<1$ and
$\varepsilon^{\pm}(\hat c)>0$ for all $\hat c$, and the allocation is held
fixed as $n$ grows (for instance the uniform rule at a fixed $e>0$), then
every $\mathbf{q}\in\acSet$ satisfies
$p^{-}_{\hat c}-O(n^{-1/2})\le q(\hat c)\le p^{+}_{\hat c}+O(n^{-1/2})$
uniformly in $\hat c$: $\acSet$ lies coordinatewise within $O(n^{-1/2})$ of
the box $\prod_{\hat c}[p^{-}_{\hat c},p^{+}_{\hat c}]$, and it contains
$\mathsf{H}$ by Theorem~\ref{thm:box-construction}\,(iii). For the degenerate envelope,
$\sup_{\mathbf{q}\in\acSet}\|\mathbf{q}-\honest\|_{\infty}=O(n^{-1/2})$ and
$\acSet$ converges to $\{\honest\}$ in Hausdorff distance.
\item \emph{(no monotonicity in $n$, with a counterexample; the exact rule)}
For the exact rule, $\acSet$ is \emph{not} monotone in $n$, already for a
degenerate envelope. Take $\Chatset=\{\hat c_0,\hat c_1\}$, the degenerate
envelope at $\honest=(\tfrac12,\tfrac12)$, $\Chatcon=\{\hat c_0\}$ and
$\varepsilon^{-}(\hat c_0)=\varepsilon^{+}(\hat c_0)=\tfrac25$ (so
$\ecomp=\tfrac45$). Then
\begin{equation*}
\begin{aligned}
  n&=2:\quad [k^{-},k^{+}]=[1,1],\quad
   \acSet=\{(\tfrac12,\tfrac12)\};\\
  n&=3:\quad [k^{-},k^{+}]=[1,2],\quad
   \acSet=\{\mathbf{q}\in\simplex{\Chatset}:\tfrac13\le q(\hat
   c_0)\le\tfrac23\},
\end{aligned}
\end{equation*}
so $\acSet(n{=}2)\subsetneq\acSet(n{=}3)$: the set is strictly larger at the
larger round count.
\item \emph{(degenerate alphabet $|\Chatset|=1$, every rule)} If
$\Chatset=\{\hat c_0\}$ then, for every $n\ge1$, every $\ecomp\in[0,1)$,
every $\Chatcon\subseteq\Chatset$ and every admissible allocation,
\begin{equation*}
  {\honest}^{+}(\hat c_0)=1,
  \qquad
  \mathsf{H}=\acSet=\simplex{\Chatset}=\{\delta_{\hat c_0}\},
\end{equation*}
and $\varepsilon^{\star}=0$ whenever ${\honest}^{-}(\hat c_0)=1$, in
particular for the degenerate envelope, which is then
$\honest=\delta_{\hat c_0}$. The emitted set is therefore constant in all of
the construction's parameters on a one-letter alphabet, being at once the
whole simplex and a single point; for the degenerate envelope so are its box
data, and completeness is perfect. For a non-degenerate envelope
(${\honest}^{-}(\hat c_0)<1$, which \eqref{eq:envelope} permits here) the
lower box bound $q^{-}(\hat c_0)$ and the reported $\varepsilon^{\star}$ may be
positive and may vary with the parameters; the reported figure is still a
valid bound (by Theorem~\ref{thm:completeness-bound}), merely a loose one, the true honest
failure probability being $0$.
\end{enumerate}
\end{theorem}
\claimmeta{\metalabel{Uses}\mref{Def.}{def:envelope}, \mref{Def.}{def:conservative-evaluation}, \mref{Def.}{def:endpoint-rule}, \mref{Lem.}{lem:guard}, \mref{Def.}{def:interval}, \mref{Def.}{def:construction}, \mref{Thm.}{thm:box-construction}, \mref{Thm.}{thm:completeness-bound}, \mref{Def.}{def:monotone-rule}, \mref{Lem.}{lem:widening}, \mref{Lem.}{lem:cheb}
\metasep \metalabel{Assumes}\aref{A-ALPHABET}{ass:alphabet}
\metasep \metalabel{Used by}\mref{Rem.}{rem:n-nonmono}, \mref{Rem.}{rem:singleton-honest}
\metasep newly derived
\metasep accepted
\metasep \metalabel{Oracles}\coderef{https://github.com/fraunhofersg/certiqs-key/blob/v0.1.0/tests/oracles/registry_statistics_acceptance.py}{registry_statistics_acceptance.py}}

\begin{proof}
\emph{(i).} Fix $\hat c\in\Chatcon$, write $p^{\mp}=p^{\mp}_{\hat c}$ and
$\varepsilon^{\pm}_i\coloneqq\varepsilon^{\pm}_{(e_i)}(\hat c)$, so
$\varepsilon^{\pm}_1\le\varepsilon^{\pm}_2$ by
Definition~\ref{def:monotone-rule}. Since $R$ is monotone in the allocation
(Definition~\ref{def:endpoint-rule}),
$R^{-}(n,p^{-},\varepsilon^{-}_1)\le R^{-}(n,p^{-},\varepsilon^{-}_2)$ and
$R^{+}(n,p^{+},\varepsilon^{+}_1)\ge R^{+}(n,p^{+},\varepsilon^{+}_2)$. Both
guards in \eqref{eq:guard} are monotone in their first argument, their second
arguments $\lfloor np^{-}\rfloor$ and $\lceil np^{+}\rceil$ not depending on
the budget, whence $k^{-}_1(\hat c)\le k^{-}_2(\hat c)$ and
$k^{+}_1(\hat c)\ge k^{+}_2(\hat c)$,
that is, $[k^{-}_2,k^{+}_2]\subseteq[k^{-}_1,k^{+}_1]$. The endpoints
$k^{\mp}_1$ at the smaller budget are therefore integer endpoints of the kind
Lemma~\ref{lem:widening}\,(iv) admits for the output $\acSet(e_2)$:
$0\le k^{-}_1(\hat c)\le k^{-}_2(\hat c)$ and
$k^{+}_2(\hat c)\le k^{+}_1(\hat c)\le n$, the outer bounds by
\eqref{eq:ordering} at $e_1$ and the inner ones as just shown. The set
\eqref{eq:widened-set} they give is $\acSet(e_1)$, since both sets are built
by \eqref{eq:boxbounds} on the same $\Chatcon$, with the bounds $[0,1]$ off
$\Chatcon$ at both budgets. Lemma~\ref{lem:widening}\,(i) gives
$\acSet(e_2)\subseteq\acSet(e_1)$.

\emph{(ii).} \eqref{eq:inf-mono} is \eqref{eq:widened-inf} of
Lemma~\ref{lem:widening}\,(iii) applied to the inclusion of (i), with
$\acSet(e_2)$ in the role of $\acSet$ and $\acSet(e_1)$ in that of $\acSet'$,
and it holds
for arbitrary $\Psi$ with values in the extended reals, no continuity or
measurability required. The specialization substitutes the definition of
$\hlow$ as an infimum over $\acSet$ and changes nothing else.

\emph{(iii).} The allocation constraint at $e=0$ forces every
$\varepsilon^{\pm}(\hat c)=0$, these being non-negative with sum at most $0$.
Fix $\hat c\in\Chatcon$ and write $p^{\mp}=p^{\mp}_{\hat c}$; by
\eqref{eq:guard} each side is computed from its own envelope value alone.

\emph{Lower side}, $R^{-}=R^{-}(n,p^{-},0)$. If $p^{-}\in(0,1)$:
Definition~\ref{def:endpoint-rule}\,(b) gives $R^{-}=0$, so
$k^{-}(\hat c)=\min\{0,\lfloor np^{-}\rfloor\}=0$. If $p^{-}=0$:
Definition~\ref{def:endpoint-rule}\,(c) gives $R^{-}=0$ and
$k^{-}(\hat c)=\min\{0,0\}=0$. If $p^{-}=1$: (c) gives $R^{-}=n$ and
$k^{-}(\hat c)=\min\{n,n\}=n$. So the lower bound is $q^{-}(\hat c)=0$ unless
$p^{-}=1$, where it is $q^{-}(\hat c)=1$.

\emph{Upper side}, $R^{+}=R^{+}(n,p^{+},0)$. If $p^{+}\in(0,1)$: (b) gives
$R^{+}=n$, so $k^{+}(\hat c)=\max\{n,\lceil np^{+}\rceil\}=n$. If $p^{+}=0$:
(c) gives $R^{+}=0$ and $k^{+}(\hat c)=\max\{0,0\}=0$. If $p^{+}=1$: (c) gives
$R^{+}=n$ and $k^{+}(\hat c)=\max\{n,n\}=n$. So the upper bound is
$q^{+}(\hat c)=1$ unless $p^{+}=0$, where it is $q^{+}(\hat c)=0$.

On $\simplex{\Chatset}$ the bound $q(\hat c)\ge1$ is $q(\hat c)=1$ and the
bound $q(\hat c)\le0$ is $q(\hat c)=0$, and the bounds $[0,1]$ are vacuous;
assembling the cases in \eqref{eq:accset} gives the displayed set, and it is
the whole simplex when no $\hat c\in\Chatcon$ has $p^{+}_{\hat c}=0$ or
$p^{-}_{\hat c}=1$. Finally $\varepsilon^{\star}=0$. A lower term of
\eqref{eq:completeness} is $\Pr[\Ncount{\hat c}^{-}<0]=0$ when
$k^{-}(\hat c)=0$, and when $p^{-}=1$ it is $\Pr[\Ncount{\hat c}^{-}<n]=0$
because $\Ncount{\hat c}^{-}$, binomial $(n,1)$, equals $n$ almost surely. An
upper term is $\Pr[\Ncount{\hat c}^{+}>n]=0$ when $k^{+}(\hat c)=n$, and when
$p^{+}=0$ it is $\Pr[\Ncount{\hat c}^{+}>0]=0$ because
$\Ncount{\hat c}^{+}$, binomial $(n,0)$, equals $0$ almost surely.

\emph{(iv).} The left inequalities of \eqref{eq:width} are \eqref{eq:ordering}
(Lemma~\ref{lem:guard}\,(i)):
$q^{+}(\hat c)=k^{+}(\hat c)/n\ge\lceil np^{+}\rceil/n\ge p^{+}$ and
$q^{-}(\hat c)=k^{-}(\hat c)/n\le\lfloor np^{-}\rfloor/n\le p^{-}$.

For the upper one, $\kappa^{+}(\hat c)$ is the quantile \eqref{eq:kappaplus}
at $p^{+}$ and $\Ncount{\hat c}^{+}$ is binomial $(n,p^{+})$ (the
\termref{term-side-wise-reading}{side-wise reading} of Definition~\ref{def:endpoint-rule}). Set
$t\coloneqq\sqrt{np^{+}(1-p^{+})/\varepsilon^{+}(\hat c)}>0$ and
$k\coloneqq\lceil np^{+}+t\rceil$. Any integer $j>k$ satisfies $j>np^{+}+t$,
so
\begin{equation*}
  \Pr[\Ncount{\hat c}^{+}>k]\ \le\ \Pr[\Ncount{\hat c}^{+}-np^{+}>t]
  \ \le\ \Pr\big[|\Ncount{\hat c}^{+}-np^{+}|\ge t\big]
  \ \le\ \frac{np^{+}(1-p^{+})}{t^{2}}\ =\ \varepsilon^{+}(\hat c)
\end{equation*}
by Lemma~\ref{lem:cheb}. If $k\le n$ this exhibits an admissible index in
\eqref{eq:kappaplus}, so $\kappa^{+}(\hat c)\le k$ by minimality; if $k>n$
then $\kappa^{+}(\hat c)\le n<k$ anyway. In both cases
$\kappa^{+}(\hat c)\le\lceil np^{+}+t\rceil$. Since also
$\lceil np^{+}\rceil\le\lceil np^{+}+t\rceil$, the guard gives
$k^{+}(\hat c)\le\lceil np^{+}+t\rceil\le np^{+}+t+1$, hence
$q^{+}(\hat c)-p^{+}\le t/n+1/n$, which is the stated bound because
$t/n=\sqrt{p^{+}(1-p^{+})/(n\varepsilon^{+}(\hat c))}$.

Symmetrically, $\kappa^{-}(\hat c)$ is the quantile \eqref{eq:kappaminus} at
$p^{-}$ and $\Ncount{\hat c}^{-}$ is binomial $(n,p^{-})$. Set
$t'\coloneqq\sqrt{np^{-}(1-p^{-})/\varepsilon^{-}(\hat c)}>0$ and
$k'\coloneqq\max\{0,\lfloor np^{-}-t'\rfloor\}$. If $\lfloor np^{-}-t'\rfloor
\le0$ then $k'=0$ and $\Pr[\Ncount{\hat c}^{-}<0]=0\le\varepsilon^{-}(\hat
c)$. Otherwise $k'=\lfloor np^{-}-t'\rfloor\le np^{-}-t'$, so any integer
$j\le k'-1$ satisfies $j<np^{-}-t'$ and
$\Pr[\Ncount{\hat c}^{-}<k']\le\Pr[\Ncount{\hat c}^{-}-np^{-}<-t']\le
\Pr[|\Ncount{\hat c}^{-}-np^{-}|\ge t']\le\varepsilon^{-}(\hat c)$ by
Lemma~\ref{lem:cheb}. Either way $k'$ is admissible in
\eqref{eq:kappaminus}, so $\kappa^{-}(\hat c)\ge k'$. In the unclipped case
$k'=\lfloor np^{-}-t'\rfloor>np^{-}-t'-1$; in the clipped case
$\lfloor np^{-}-t'\rfloor\le0$ forces $np^{-}-t'<1$, hence $np^{-}-t'-1<0=k'$.
In both cases $\kappa^{-}(\hat c)\ge k'>np^{-}-t'-1$. As
$\lfloor np^{-}\rfloor>np^{-}-1\ge np^{-}-t'-1$, the guard gives
$k^{-}(\hat c)\ge np^{-}-t'-1$ and hence $p^{-}-q^{-}(\hat c)\le t'/n+1/n$.

For the consequence: $\acSet\subseteq\prod_{\hat c}[q^{-}(\hat
c),q^{+}(\hat c)]$ by \eqref{eq:accset}, so every $\mathbf{q}\in\acSet$ has
$p^{-}_{\hat c}-q(\hat c)\le p^{-}_{\hat c}-q^{-}(\hat c)$ and
$q(\hat c)-p^{+}_{\hat c}\le q^{+}(\hat c)-p^{+}_{\hat c}$, which
\eqref{eq:width} bounds by $O(n^{-1/2})$ uniformly
over the finite alphabet, both bounds of \eqref{eq:width} applying at every
$\hat c$ because $0<p^{\mp}_{\hat c}<1$ and $\varepsilon^{\pm}(\hat c)>0$
there, with fixed values; $\mathsf{H}\subseteq\acSet$ is
Theorem~\ref{thm:box-construction}\,(iii). For the degenerate envelope
$p^{-}_{\hat c}=p^{+}_{\hat c}=\honest(\hat c)$, so the same two bounds give
$\|\mathbf{q}-\honest\|_{\infty}\le\max_{\hat c}\max\{q^{+}(\hat
c)-\honest(\hat c),\,\honest(\hat c)-q^{-}(\hat c)\}=O(n^{-1/2})$ for every
$\mathbf{q}\in\acSet$, and $\honest\in\acSet$ by Theorem~\ref{thm:box-construction}\,(iii),
so the other half of the Hausdorff distance is $0$.

\emph{(v).} Direct computation with the degenerate envelope at $p=\tfrac12$,
so that $\Ncount{}\coloneqq\Ncount{\hat c_0}^{-}=\Ncount{\hat c_0}^{+}$ is
binomial $(n,\tfrac12)$, and the exact rule. At
$n=2$ the law of $\Ncount{}$ is $(\tfrac14,\tfrac12,\tfrac14)$ on
$\{0,1,2\}$: $\Pr[\Ncount{}<0]=0\le\tfrac25$ and
$\Pr[\Ncount{}<1]=\tfrac14\le\tfrac25$ but
$\Pr[\Ncount{}<2]=\tfrac34>\tfrac25$, so $\kappa^{-}=1$;
$\Pr[\Ncount{}>0]=\tfrac34>\tfrac25$ and $\Pr[\Ncount{}>1]=\tfrac14\le\tfrac25$,
so $\kappa^{+}=1$. The guard is inert ($\lfloor1\rfloor=\lceil1\rceil=1$), so
$[k^{-},k^{+}]=[1,1]$ and $q(\hat c_0)=\tfrac12$, which on a two-letter
alphabet forces $\mathbf{q}=(\tfrac12,\tfrac12)$. At $n=3$ the law is
$(\tfrac18,\tfrac38,\tfrac38,\tfrac18)$ on $\{0,1,2,3\}$:
$\Pr[\Ncount{}<1]=\tfrac18\le\tfrac25$ but
$\Pr[\Ncount{}<2]=\tfrac12>\tfrac25$, so $\kappa^{-}=1$;
$\Pr[\Ncount{}>1]=\tfrac12>\tfrac25$ and
$\Pr[\Ncount{}>2]=\tfrac18\le\tfrac25$, so $\kappa^{+}=2$. The guard is again
inert ($\lfloor\tfrac32\rfloor=1$, $\lceil\tfrac32\rceil=2$), so
$[k^{-},k^{+}]=[1,2]$ and $q(\hat c_0)\in[\tfrac13,\tfrac23]$. The two sets are
as displayed and the inclusion is strict.

\emph{(vi).} On $\Chatset=\{\hat c_0\}$ the only vector with non-negative
coordinates summing to $1$ is $\delta_{\hat c_0}$, so
$\simplex{\Chatset}=\{\delta_{\hat c_0}\}$. Condition \eqref{eq:envelope}
gives ${\honest}^{+}(\hat c_0)\le1\le{\honest}^{+}(\hat c_0)$, that is,
${\honest}^{+}(\hat c_0)=1$, and ${\honest}^{-}(\hat c_0)\le1$, so
$\delta_{\hat c_0}\in\mathsf{H}\subseteq\simplex{\Chatset}$ and
$\mathsf{H}=\{\delta_{\hat c_0}\}$; for the degenerate envelope also
${\honest}^{-}(\hat c_0)=1$, that is, $\honest=\delta_{\hat c_0}$. Since
$\acSet\subseteq\simplex{\Chatset}$ by \eqref{eq:accset} and
$\acSet\supseteq\mathsf{H}\neq\varnothing$ by Theorem~\ref{thm:box-construction}\,(iii),
$\acSet=\simplex{\Chatset}=\{\delta_{\hat c_0}\}$ whatever the parameters. For
$\varepsilon^{\star}$: if $\Chatcon=\varnothing$, \eqref{eq:boxbounds} gives
the vacuous bounds $[0,1]$ and $\varepsilon^{\star}$ is an empty sum, hence
$0$. If $\Chatcon=\{\hat c_0\}$, the upper side has
$p^{+}_{\hat c_0}=1$, so Definition~\ref{def:endpoint-rule}\,(c) gives
$R^{+}=n$, the guard gives $k^{+}(\hat c_0)=n$ and the upper term
$\Pr[\Ncount{\hat c_0}^{+}>n]$ is $0$. If moreover
${\honest}^{-}(\hat c_0)=1$, the same convention gives $R^{-}=n$,
$k^{-}(\hat c_0)=\min\{n,n\}=n$, the box bounds
$q^{-}(\hat c_0)=q^{+}(\hat c_0)=1$ independent of every parameter, and the
lower term $\Pr[\Ncount{\hat c_0}^{-}<n]$ is $0$ because
$\Ncount{\hat c_0}^{-}$, binomial $(n,1)$, equals $n$ almost surely; so
$\varepsilon^{\star}=0$. The two cases exhaust $\Chatcon\subseteq\Chatset$,
and neither used $n$, $\ecomp$, the allocation or anything of $R$ beyond
Definition~\ref{def:endpoint-rule}\,(c). When ${\honest}^{-}(\hat c_0)<1$ the
lower term is the tail of a binomial $(n,{\honest}^{-}(\hat c_0))$ below
$k^{-}(\hat c_0)$, which Lemma~\ref{lem:guard}\,(ii) bounds by
$\varepsilon^{-}(\hat c_0)$ and which need not vanish; the bound of
Theorem~\ref{thm:completeness-bound} still holds, the true failure probability being
$\Pr[\operatorname{freq}_{\Chat_1^n}\notin\{\delta_{\hat c_0}\}]=0$.
\end{proof}

% claim statistics/binomial-box/rem:n-nonmono
\claimsummary{statistics/binomial-box/rem:n-nonmono}{The exact rule's endpoints are integer quantiles on a lattice that refines with n, so its interval can widen when n grows by one, as \hyperlink{claim-statistics/binomial-box/thm:monotonicity-limits}{Theorem~\ref*{thm:monotonicity-limits}} (v) shows, while the n\textasciicircum{}(\ensuremath{-}1/2) width of (iv) stands. A consumer may rely on monotonicity in the budget, not in the round count.}
\claimbookmark{statistics/binomial-box/rem:n-nonmono}{Remark}{rem:n-nonmono}{Non-monotonicity of the exact rule in the round count}
\begin{remark}[Non-monotonicity of the exact rule in the round count]
\label{rem:n-nonmono}
The endpoints $\kappa^{\pm}$ are integer quantiles of a distribution whose
support lattice refines with $n$, so the attainable tail levels move in jumps
and the interval $[k^{-}/n,k^{+}/n]$ can widen when $n$ increases by one, as
Theorem~\ref{thm:monotonicity-limits}\,(v) exhibits. The asymptotic statement
\eqref{eq:width} is unaffected, the width still falling as $n^{-1/2}$, and
no claim of this module needs monotonicity in $n$. A consumer that assumed
``more rounds, tighter set'' would be assuming something false. The
monotonicity that \emph{is} available, and that a consumer may rely on, is
monotonicity in $\ecomp$, Theorem~\ref{thm:monotonicity-limits}\,(i), with the direction
stated there.
\end{remark}
\claimmeta{\metalabel{Uses}\mref{Thm.}{thm:monotonicity-limits}
\metasep newly derived
\metasep accepted}

\section{Elementary enclosures}\label{group-statistics/binomial-box/lem:log-enclosure}
\sectionintro{Rational bounds on logarithms from integer square roots, and two-sided Stirling bounds on factorials with the constant enclosed by Wallis's product, are the only transcendental ingredients of the closed-form rule and of the local evaluation. Every bound has an explicit direction and an explicit width.}

% claim statistics/binomial-box/lem:log-enclosure
\claimsummary{statistics/binomial-box/lem:log-enclosure}{For x \textgreater{} 0 and m \ensuremath{\ge} 1, m(1 \ensuremath{-} x\textasciicircum{}(\ensuremath{-}1/m)) \ensuremath{\le} ln x \ensuremath{\le} m(x\textasciicircum{}(1/m) \ensuremath{-} 1), with an explicit slack; for rational x and m = 2\textasciicircum{}j, j integer square roots give rational bounds of any prescribed width. The Bernstein rule and the enclosure evaluation take every logarithm this way.}
\claimbookmark{statistics/binomial-box/lem:log-enclosure}{Lemma}{lem:log-enclosure}{Rational bounds on the logarithm from integer square roots}
\claimleadin{%
No logarithm is taken from a library: every logarithm the Bernstein rule and
the local evaluation of the exact rule need is bounded from both sides by
rationals. The bounds are the tangent inequality $\ln t\le t-1$ at
$t=x^{\pm1/m}$, the identity by which Briggs tabulated logarithms through
repeated square roots before series methods, and their slack is the
second-order remainder of the exponential series; both are derived in full
here.
}
\begin{lemma}[Rational bounds on the logarithm from integer square roots]\label{lem:log-enclosure}
Let $x>0$ be real and $m\ge1$ an integer. Then
\begin{equation}\label{eq:log-root-bounds}
  m\big(1-x^{-1/m}\big)\ \le\ \ln x\ \le\ m\big(x^{1/m}-1\big),
\end{equation}
with equality on either side if and only if $x=1$, and both slacks are bounded
by
\begin{equation}\label{eq:log-root-slack}
  \max\Big\{m\big(x^{1/m}-1\big)-\ln x,\ \ \ln x-m\big(1-x^{-1/m}\big)\Big\}
  \ \le\ \frac{(\ln x)^{2}}{2m}\,e^{|\ln x|/m}.
\end{equation}
Consequently, if $l,u$ are rationals with $0<l\le x^{1/m}\le u$, then
\begin{equation}\label{eq:log-rational-bounds}
  m\Big(1-\frac1l\Big)\ \le\ \ln x\ \le\ m\,(u-1),
\end{equation}
and for rational $x=a/b$ ($a,b$ positive integers) and $m=2^{j}$, such
$l,u$ of any prescribed positive width are computed by $j$ integer square
roots: with $\operatorname{isqrt}(N)\coloneqq\lfloor\sqrt N\rfloor$
for an integer $N\ge0$, one has for every rational $y=c/d>0$
\begin{equation}\label{eq:isqrt-bounds}
  \frac{\operatorname{isqrt}(c\,d\,4^{s})}{d\,2^{s}}
  \ \le\ \sqrt{y}\ <\
  \frac{\operatorname{isqrt}(c\,d\,4^{s})+1}{d\,2^{s}},
\end{equation}
the two bounds differing by $1/(d\,2^{s})\le2^{-s}$, and iterating
\eqref{eq:isqrt-bounds} $j$ times, the lower bound always taken from the
previous lower bound and the upper from the previous upper, gives
$l\le x^{1/2^{j}}\le u$ whose width $u-l$ is an explicit rational the
caller reads off and drives below any prescribed bound by the precisions
$s$ of the steps.
\end{lemma}
\claimmeta{\metalabel{Used by}\mref{Lem.}{lem:stirling-enclosure}, \mref{Prop.}{prop:rule-bernstein}, \mref{Lem.}{lem:pmf-enclosure}, \mref{Prop.}{prop:bracketed-endpoints}, \mref{Rem.}{rem:local-evaluation}, \mref{Req.}{req:enclosure-substrate}
\metasep newly derived
\metasep implemented
\metasep \metalabel{Code}\coderef{https://github.com/fraunhofersg/certiqs-key/blob/v0.1.0/src/certiqs_key/statistics/logs.py}{certiqs_key.statistics.logs:log_bounds}}

\begin{proof}
\emph{The tangent inequality.} For $t>0$ put $f(t)\coloneqq t-1-\ln t$. Then
$f'(t)=1-1/t$, which is negative on $(0,1)$ and positive on $(1,\infty)$, so
$f$ attains its minimum at $t=1$, where $f(1)=0$. Hence
$\ln t\le t-1$ for every $t>0$, with equality if and only if $t=1$.

\emph{\eqref{eq:log-root-bounds}.} Apply the tangent inequality at
$t=x^{1/m}$: $\ln x^{1/m}=(\ln x)/m\le x^{1/m}-1$, i.e.\ $\ln x\le
m(x^{1/m}-1)$. Apply it at $t=x^{-1/m}$: $-(\ln x)/m\le x^{-1/m}-1$, i.e.\
$\ln x\ge m(1-x^{-1/m})$. Equality holds on either side exactly when the
argument $t$ equals $1$, that is, when $x=1$.

\emph{\eqref{eq:log-root-slack}.} Put $y\coloneqq(\ln x)/m$, so
$x^{1/m}=e^{y}$ and $x^{-1/m}=e^{-y}$. The two slacks are
$m(e^{y}-1-y)$ and $m(e^{-y}-1+y)$. From the exponential series,
$e^{\pm y}-1\mp y=\sum_{k\ge2}(\pm y)^{k}/k!$, and since $2/k!\le1/(k-2)!$
for $k\ge2$,
\begin{equation*}
  \big|e^{\pm y}-1\mp y\big|\ \le\ \sum_{k\ge2}\frac{|y|^{k}}{k!}
  \ =\ \frac{y^{2}}{2}\sum_{k\ge2}\frac{2\,|y|^{k-2}}{k!}
  \ \le\ \frac{y^{2}}{2}\sum_{i\ge0}\frac{|y|^{i}}{i!}
  \ =\ \frac{y^{2}}{2}\,e^{|y|}.
\end{equation*}
Multiplying by $m$ and substituting $y=(\ln x)/m$ gives
\eqref{eq:log-root-slack}.

\emph{\eqref{eq:log-rational-bounds}.} If $0<l\le x^{1/m}$ then
$x^{-1/m}\le1/l$, so $m(1-x^{-1/m})\ge m(1-1/l)$; if $x^{1/m}\le u$ then
$m(x^{1/m}-1)\le m(u-1)$. Combine with \eqref{eq:log-root-bounds}.

\emph{\eqref{eq:isqrt-bounds}.} For $y=c/d>0$,
$\sqrt{y}=\sqrt{c\,d\,4^{s}}\big/(d\,2^{s})$, and for the integer
$N\coloneqq c\,d\,4^{s}$ the definition of $\operatorname{isqrt}$ gives
$\operatorname{isqrt}(N)\le\sqrt N<\operatorname{isqrt}(N)+1$; divide by
$d\,2^{s}$. The difference of the two bounds is $1/(d\,2^{s})\le2^{-s}$
because $d\ge1$.

\emph{The iteration.} The square root is non-decreasing on $[0,\infty)$, so
from $l_i\le x^{1/2^{i}}\le u_i$ follows $\sqrt{l_i}\le x^{1/2^{i+1}}\le
\sqrt{u_i}$; applying the lower bound of \eqref{eq:isqrt-bounds} to
$y=l_i$ and the upper bound to $y=u_i$, at a precision $s_i$ of the caller's
choice, gives rationals $l_{i+1}\le x^{1/2^{i+1}}\le u_{i+1}$ with
$l_{i+1}>0$ whenever $l_i>0$ and $s_i$ is large enough that the lower bound
is positive. Starting from $l_0=u_0=x$, after $j$ steps
$l\coloneqq l_j\le x^{1/2^{j}}\le u_j\eqqcolon u$. For the width
$w_i\coloneqq u_i-l_i$: each step's exact roots satisfy
$\sqrt{u_i}-\sqrt{l_i}=w_i/(\sqrt{u_i}+\sqrt{l_i})$, and each of the two
roundings moves its bound by at most $1/(d\,2^{s_i})\le2^{-s_i}$, with $d$
the denominator of its own argument, because the two bounds of
\eqref{eq:isqrt-bounds} enclose the root and differ by $1/(d\,2^{s_i})$: the
lower bound satisfies $l_{i+1}>\sqrt{l_i}-2^{-s_i}$ and the upper bound
$u_{i+1}\le\sqrt{u_i}+2^{-s_i}$. The two add at most $2\cdot2^{-s_i}$ in
total, so
\begin{equation*}
  w_{i+1}\ \le\ \frac{w_i}{\sqrt{u_i}+\sqrt{l_i}}+2\cdot2^{-s_i}
  \ \le\ \frac{w_i}{2\sqrt{l_i}}+2\cdot2^{-s_i},
\end{equation*}
with $w_0=0$. To bound the factors $1/(2\sqrt{l_i})$ independently of the
precisions, write $x_i\coloneqq x^{1/2^{i}}$ and require of each $s_i$ that
$2^{-s_i}\le(1/\sqrt2-1/2)\,x_{i+1}$, a condition on $x$ alone that every
larger $s_i$ also meets. Then $l_i\ge x_i/2$ for every $i\le j$, by
induction: $l_0=x$, and $l_i\ge x_i/2$ gives
$l_{i+1}>\sqrt{l_i}-2^{-s_i}\ge x_{i+1}/\sqrt2-(1/\sqrt2-1/2)\,x_{i+1}
=x_{i+1}/2$, using $\sqrt{x_i}=x_{i+1}$. In particular every $l_i$ is
positive, and $1/(2\sqrt{l_i})\le1/(\sqrt2\,x_{i+1})\le K$ for $i<j$, where
$K\coloneqq\max\{1,\max_{i<j}1/(\sqrt2\,x_{i+1})\}$ depends on $x$ and $j$
alone. By induction
$w_j\le\sum_{i<j}K^{\,j-1-i}\,2\cdot2^{-s_i}$: finitely many terms, each
driven below any prescribed bound by its own $s_i$ without changing $K$. The
achieved width is the explicit rational $u_j-l_j$ a caller reads off.
\end{proof}

% claim statistics/binomial-box/lem:stirling-enclosure
\claimsummary{statistics/binomial-box/lem:stirling-enclosure}{Robbins' bounds enclose ln m! between Stirling's approximation plus 1/(12m + 1) and plus 1/(12m), and Wallis's partial products enclose 2\ensuremath{\pi} by rationals of any width. With rational bounds on ln m and ln(2\ensuremath{\pi}) this encloses ln m! by rationals of explicit width.}
\claimbookmark{statistics/binomial-box/lem:stirling-enclosure}{Lemma}{lem:stirling-enclosure}{Robbins' Stirling bounds with Wallis's product}
\claimleadin{%
A binomial mass point is a ratio of factorials, and the local evaluation of
the exact rule needs $\ln m!$ from both sides without a library logarithm.
Part~(i) is Stirling's formula with the two-sided error bounds of Robbins,
and part~(ii) is Wallis's product for $\pi/2$, which makes the constant
$\tfrac12\ln(2\pi)$ computable in exact rational arithmetic; both are derived
in full here, the constant being identified through the central binomial
coefficient.
}
\begin{technicalclaim}
\begin{lemma}[Robbins' Stirling bounds with Wallis's product]\label{lem:stirling-enclosure}
Put $C\coloneqq\tfrac12\ln(2\pi)$.
\begin{enumerate}[label=(\roman*),leftmargin=2.2em]
\item \emph{(Robbins' bounds)} For every integer $m\ge1$,
\begin{equation}\label{eq:robbins}
  \Big(m+\frac12\Big)\ln m-m+C+\frac{1}{12m+1}
  \ <\ \ln m!\ <\
  \Big(m+\frac12\Big)\ln m-m+C+\frac{1}{12m},
\end{equation}
equivalently
$\sqrt{2\pi m}\,(m/e)^{m}\,e^{1/(12m+1)}<m!<\sqrt{2\pi m}\,(m/e)^{m}\,e^{1/(12m)}$.
\item \emph{(Wallis's product)} For every integer $M\ge1$, the rational
\begin{equation}\label{eq:wallis-partial}
  P_M\ \coloneqq\ \prod_{k=1}^{M}\frac{(2k)^{2}}{(2k-1)(2k+1)}
\end{equation}
satisfies
\begin{equation}\label{eq:wallis-bounds}
  P_M\ <\ \frac{\pi}{2}\ <\ P_M\,\frac{2M+1}{2M},
  \qquad\text{hence}\qquad
  4P_M\ <\ 2\pi\ <\ 4P_M\Big(1+\frac{1}{2M}\Big),
\end{equation}
so that $2C=\ln(2\pi)$ is enclosed, through
Lemma~\ref{lem:log-enclosure} applied to the two rationals in
\eqref{eq:wallis-bounds}, by rationals of any prescribed width.
\item \emph{(rational enclosure of $\ln m!$)} Let $m\ge1$ and let
$\lambda^{-}\le\ln m\le\lambda^{+}$ and $c^{-}\le\ln(2\pi)\le c^{+}$ be
rationals. Then
\begin{equation}\label{eq:lnfact-enclosure}
  \Big(m+\frac12\Big)\lambda^{-}-m+\frac{c^{-}}{2}+\frac{1}{12m+1}
  \ <\ \ln m!\ <\
  \Big(m+\frac12\Big)\lambda^{+}-m+\frac{c^{+}}{2}+\frac{1}{12m},
\end{equation}
an enclosure by rationals whose width is at most
$(m+\tfrac12)(\lambda^{+}-\lambda^{-})+\tfrac12(c^{+}-c^{-})
+\tfrac{1}{12m(12m+1)}$. For $m=0$, $\ln 0!=0$ exactly.
\end{enumerate}
\end{lemma}
\end{technicalclaim}
\claimmeta{\metalabel{Uses}\mref{Lem.}{lem:log-enclosure}
\metasep \metalabel{Used by}\mref{Lem.}{lem:pmf-enclosure}, \mref{Rem.}{rem:local-evaluation}
\metasep newly derived
\metasep implemented
\metasep \metalabel{Code}\coderef{https://github.com/fraunhofersg/certiqs-key/blob/v0.1.0/src/certiqs_key/statistics/logs.py}{certiqs_key.statistics.logs:log_factorial_bounds}}

\begin{proof}
\emph{Idea.} The sequence $d_m=\ln m!-(m+\tfrac12)\ln m+m$ has differences
given by a series in $1/(2m+1)$. Bounding that series from both sides makes
$d_m-\frac{1}{12m}$ increasing and $d_m-\frac{1}{12m+1}$ decreasing, with a
common limit, and Wallis's product, proved from the sine-power integrals,
identifies the limit as $\tfrac12\ln(2\pi)$ through the central binomial
coefficient.

\emph{Step 1 (a series).} For real $|x|<1$,
\begin{equation}\label{eq:atanh-series}
  \ln\frac{1+x}{1-x}\ =\ 2\sum_{k\ge0}\frac{x^{2k+1}}{2k+1}.
\end{equation}
Indeed both sides vanish at $x=0$, and their derivatives agree:
$\frac{\mathrm{d}}{\mathrm{d}x}\ln\frac{1+x}{1-x}=\frac{1}{1+x}+\frac{1}{1-x}
=\frac{2}{1-x^{2}}=2\sum_{k\ge0}x^{2k}$, the geometric series converging
uniformly on $|x|\le\rho$ for every $\rho<1$, so it may be integrated
termwise from $0$ to $x$.

\emph{Step 2 (the differences of $d_m$).} For $m\ge1$ put
$d_m\coloneqq\ln m!-(m+\tfrac12)\ln m+m$. Then
\begin{equation*}
  d_m-d_{m+1}
  =\ln m!-\ln(m+1)!-\Big(m+\frac12\Big)\ln m+\Big(m+\frac32\Big)\ln(m+1)-1
  =\Big(m+\frac12\Big)\ln\frac{m+1}{m}-1,
\end{equation*}
using $\ln(m+1)!-\ln m!=\ln(m+1)$. Put $x\coloneqq1/(2m+1)\in(0,\tfrac13]$;
then $(1+x)/(1-x)=(2m+2)/(2m)=(m+1)/m$ and $m+\tfrac12=1/(2x)$, so by
\eqref{eq:atanh-series}
\begin{equation}\label{eq:dm-series}
  d_m-d_{m+1}=\frac{1}{2x}\cdot2\sum_{k\ge0}\frac{x^{2k+1}}{2k+1}-1
  =\sum_{k\ge0}\frac{x^{2k}}{2k+1}-1=\sum_{k\ge1}\frac{x^{2k}}{2k+1}.
\end{equation}

\emph{Step 3 (two bounds on the difference).} Every term of
\eqref{eq:dm-series} is positive, so keeping the first,
\begin{equation*}
  d_m-d_{m+1}\ >\ \frac{x^{2}}{3}=\frac{1}{3(2m+1)^{2}}
  \ >\ \frac{1}{12m+1}-\frac{1}{12m+13}=\frac{12}{(12m+1)(12m+13)},
\end{equation*}
where the last inequality is, after cross-multiplication,
$(12m+1)(12m+13)>36(2m+1)^{2}$, i.e.\
$144m^{2}+168m+13>144m^{2}+144m+36$, i.e.\ $24m>23$, true for $m\ge1$.
Bounding each term by $1/3$ times the geometric series,
\begin{equation*}
  d_m-d_{m+1}\ <\ \frac{x^{2}}{3}\sum_{k\ge0}x^{2k}=\frac{x^{2}}{3(1-x^{2})}
  =\frac{1}{3\big((2m+1)^{2}-1\big)}=\frac{1}{12m(m+1)}
  =\frac{1}{12m}-\frac{1}{12(m+1)}.
\end{equation*}

\emph{Step 4 (two monotone sequences with a common limit).} Put
$a_m\coloneqq d_m-\frac{1}{12m}$ and $b_m\coloneqq d_m-\frac{1}{12m+1}$. By
the upper bound of Step 3,
$a_{m+1}-a_m=-(d_m-d_{m+1})+\frac{1}{12m}-\frac{1}{12(m+1)}>0$, so $(a_m)$
is strictly increasing; by the lower bound,
$b_{m+1}-b_m=-(d_m-d_{m+1})+\frac{1}{12m+1}-\frac{1}{12m+13}<0$, so $(b_m)$
is strictly decreasing. Since $a_m<b_m$ for every $m$, $(a_m)$ is bounded
above by $b_1$ and $(b_m)$ below by $a_1$, so both converge, and
$b_m-a_m=\frac{1}{12m}-\frac{1}{12m+1}\to0$ forces a common limit $C_0$.
Strict monotonicity gives $a_m<C_0<b_m$ for every $m\ge1$, that is,
\begin{equation}\label{eq:robbins-C0}
  C_0+\frac{1}{12m+1}\ <\ d_m\ <\ C_0+\frac{1}{12m}\qquad(m\ge1).
\end{equation}

\emph{Step 5 (Wallis's product).} For an integer $n\ge0$ put
$I_n\coloneqq\int_0^{\pi/2}\sin^{n}x\,\mathrm{d}x$, so $I_0=\pi/2$ and
$I_1=1$. For $n\ge2$, integration by parts with $u=\sin^{n-1}x$ and
$\mathrm{d}v=\sin x\,\mathrm{d}x$ gives
$I_n=\big[-\sin^{n-1}x\cos x\big]_0^{\pi/2}
+(n-1)\int_0^{\pi/2}\sin^{n-2}x\cos^{2}x\,\mathrm{d}x
=(n-1)(I_{n-2}-I_n)$, whence $I_n=\frac{n-1}{n}I_{n-2}$. Iterating,
\begin{equation*}
  I_{2M}=\frac{\pi}{2}\prod_{k=1}^{M}\frac{2k-1}{2k},
  \qquad
  I_{2M+1}=\prod_{k=1}^{M}\frac{2k}{2k+1}\qquad(M\ge0).
\end{equation*}
On $(0,\pi/2)$ one has $0<\sin x<1$, so $\sin^{n+1}x<\sin^{n}x$ there and
$0<I_{n+1}<I_n$ for every $n$. Hence, for $M\ge1$,
$I_{2M+1}<I_{2M}<I_{2M-1}$. Now
$I_{2M}/I_{2M+1}=\frac{\pi}{2}\prod_{k=1}^{M}\frac{(2k-1)(2k+1)}{(2k)^{2}}
=\frac{\pi/2}{P_M}$, and $I_{2M+1}<I_{2M}$ gives $P_M<\pi/2$. Further
$I_{2M-1}=\prod_{k=1}^{M-1}\frac{2k}{2k+1}=I_{2M+1}\cdot\frac{2M+1}{2M}$, so
$I_{2M}<I_{2M-1}$ gives $\frac{\pi/2}{P_M}<\frac{2M+1}{2M}$, i.e.\
$\pi/2<P_M\frac{2M+1}{2M}$. Multiplying by $4$ gives the second display of
\eqref{eq:wallis-bounds}; this is (ii), and Lemma~\ref{lem:log-enclosure}
applied to the two rationals $4P_M$ and $4P_M(1+1/(2M))$ encloses
$\ln(2\pi)=2C$ by rationals, of width at most the sum of the root-bound
widths and $\ln(1+1/(2M))\le1/(2M)$, hence of any prescribed width for $M$
and the root precision large enough.

\emph{Step 6 (the constant).} We identify $C_0=C$. From
$\ln m!=d_m+(m+\tfrac12)\ln m-m$,
\begin{equation*}
  \ln\binom{2M}{M}=\ln(2M)!-2\ln M!
  =d_{2M}-2d_M+\Big(2M+\frac12\Big)\ln(2M)-2M-(2M+1)\ln M+2M
  =d_{2M}-2d_M+2M\ln2+\frac12\ln2-\frac12\ln M,
\end{equation*}
so $\binom{2M}{M}=e^{\,d_{2M}-2d_M}\,4^{M}\sqrt{2/M}$ and
\begin{equation}\label{eq:central-binomial}
  \Big(\frac{4^{M}}{\binom{2M}{M}}\Big)^{2}=e^{\,2(2d_M-d_{2M})}\,\frac{M}{2}.
\end{equation}
On the other hand $P_M$ is expressible through the same coefficient:
$\prod_{k=1}^{M}(2k)^{2}=4^{M}(M!)^{2}$,
$\prod_{k=1}^{M}(2k-1)=\frac{(2M)!}{2^{M}M!}$ and
$\prod_{k=1}^{M}(2k+1)=\frac{(2M+1)!}{2^{M}M!}$, so
\begin{equation*}
  P_M=\frac{4^{M}(M!)^{2}\cdot4^{M}(M!)^{2}}{(2M)!\,(2M+1)!}
  =\frac{16^{M}(M!)^{4}}{\big((2M)!\big)^{2}(2M+1)}
  =\frac{1}{2M+1}\Big(\frac{4^{M}}{\binom{2M}{M}}\Big)^{2}
  =e^{\,2(2d_M-d_{2M})}\,\frac{M}{2(2M+1)} .
\end{equation*}
As $M\to\infty$, $2d_M-d_{2M}\to2C_0-C_0=C_0$ by Step 4 and
$M/(2(2M+1))\to1/4$, while $P_M\to\pi/2$ by \eqref{eq:wallis-bounds}.
Hence $\pi/2=e^{2C_0}/4$, i.e.\ $e^{2C_0}=2\pi$ and $C_0=\tfrac12\ln(2\pi)=C$.
With this, \eqref{eq:robbins-C0} is \eqref{eq:robbins}, and exponentiating
gives the product form. This proves (i).

\emph{Step 7 (part (iii)).} In \eqref{eq:robbins} replace $\ln m$ by its bounds
(the coefficient $m+\tfrac12$ is positive) and $C=\tfrac12\ln(2\pi)$ by
$\tfrac12c^{\mp}$; the width is the sum of the two widths so introduced plus
$\frac{1}{12m}-\frac{1}{12m+1}=\frac{1}{12m(12m+1)}$. At $m=0$, $0!=1$.
\end{proof}

\section{The Bernstein rule}\label{group-statistics/binomial-box/lem:bernstein-tail}
\sectionintro{Bernstein's tail inequality gives count endpoints in closed form, with no tail evaluated. With upward-rounded rational bounds the rule is admissible, \termref{term-monotone-in-the-allocation}{monotone in the allocation} when every rounding is, and its unrounded half-width decreases strictly with the round count, though nesting of its integer-rounded boxes in the round count is not claimed (\hyperlink{claim-statistics/binomial-box/prop:rule-bernstein}{Proposition~\ref*{prop:rule-bernstein}}). A remark says which completeness figure a rule that evaluates no tail can report.}

% claim statistics/binomial-box/lem:bernstein-tail
\claimsummary{statistics/binomial-box/lem:bernstein-tail}{Either tail of a binomial (n, p) count at distance t \textgreater{} 0 from np has probability at most exp(\ensuremath{-}t\ensuremath{^{2}}/(2(v + t/3))) with v = np(1 \ensuremath{-} p), and every t at least a closed-form root t(L) brings both below e\textasciicircum{}(\ensuremath{-}L). The Bernstein endpoint rule evaluates that root.}
\claimbookmark{statistics/binomial-box/lem:bernstein-tail}{Lemma}{lem:bernstein-tail}{Bernstein's inequality for a binomial count}
\claimleadin{%
The lemma is Bernstein's inequality for a sum of independent, bounded,
centred random variables, specialised to a binomial count and derived in full
here; its closed-form deviation $t(L)$, the positive root of
$t^{2}-(2L/3)\,t-2vL=0$, is what the Bernstein endpoint rule evaluates.
}
\begin{technicalclaim}
\begin{lemma}[Bernstein's inequality for a binomial count]\label{lem:bernstein-tail}
Let $\Ncount{}$ be binomial with parameters $(n,p)$, $n\ge1$, $p\in[0,1]$,
and put $v\coloneqq np(1-p)$. For every real $t>0$,
\begin{equation}\label{eq:bernstein}
  \Pr\big[\Ncount{}-np\ge t\big]\ \le\ \exp\Big(-\frac{t^{2}}{2\,(v+t/3)}\Big),
  \qquad
  \Pr\big[\Ncount{}-np\le -t\big]\ \le\ \exp\Big(-\frac{t^{2}}{2\,(v+t/3)}\Big).
\end{equation}
Consequently, for every real $L>0$ and every
\begin{equation}\label{eq:bernstein-t}
  t\ \ge\ t(L)\ \coloneqq\ \frac{L}{3}+\sqrt{\Big(\frac{L}{3}\Big)^{2}+2vL},
\end{equation}
both tails in \eqref{eq:bernstein} are at most $e^{-L}$.
\end{lemma}
\end{technicalclaim}
\claimmeta{\metalabel{Uses}\mref{Def.}{def:freq}
\metasep \metalabel{Used by}\mref{Prop.}{prop:rule-bernstein}
\metasep newly derived
\metasep implemented (effective: accepted)
\metasep \metalabel{Code}\coderef{https://github.com/fraunhofersg/certiqs-key/blob/v0.1.0/src/certiqs_key/statistics/bernstein.py}{certiqs_key.statistics.bernstein:bernstein_deviation}}

\begin{proof}
\emph{Idea.} The exponential Markov inequality, with the
moment-generating-function bound $\mathbb{E}[e^{\lambda X}]\le
\exp\big(\sigma^{2}(e^{\lambda}-1-\lambda)\big)$ for a centred $X\le1$ and the
estimate $e^{\lambda}-1-\lambda\le\lambda^{2}/(2(1-\lambda/3))$, evaluated at
$\lambda=t/(v+t/3)$; the deviation $t(L)$ then solves a quadratic.

If $p\in\{0,1\}$ then $\Ncount{}=np$ with probability one, both
probabilities in \eqref{eq:bernstein} are $0$ for $t>0$, and there is nothing
to prove. Let $0<p<1$, so $v>0$.

\emph{Step 1 (a moment-generating-function bound).} Let $X$ be a real random
variable with $\mathbb{E}[X]=0$, $\operatorname{Var}(X)=\sigma^{2}<\infty$ and
$X\le1$ almost surely, and let $\lambda>0$. Put
$g(\lambda)\coloneqq e^{\lambda}-1-\lambda$. We claim
\begin{equation}\label{eq:mgf}
  \mathbb{E}\big[e^{\lambda X}\big]\ \le\ 1+\sigma^{2}g(\lambda)\ \le\
  \exp\big(\sigma^{2}g(\lambda)\big).
\end{equation}
The second inequality is $1+z\le e^{z}$. For the first, consider
$\psi(u)\coloneqq(e^{u}-1-u)/u^{2}$ for $u\neq0$, with $\psi(0)\coloneqq1/2$
(its limit). We show that $\psi$ is non-decreasing on $\mathbb{R}$. For
$u\neq0$, $\psi'(u)=N(u)/u^{3}$ with
$N(u)\coloneqq(e^{u}-1)u-2(e^{u}-1-u)=e^{u}(u-2)+u+2$. Now $N(0)=0$,
$N'(u)=e^{u}(u-1)+1$ with $N'(0)=0$, and $N''(u)=u\,e^{u}$, which has the sign
of $u$; hence $N'$ decreases on $(-\infty,0]$ and increases on $[0,\infty)$,
so $N'\ge N'(0)=0$ everywhere, so $N$ is non-decreasing, and with $N(0)=0$
the sign of $N(u)$ is the sign of $u$. Therefore $\psi'(u)=N(u)/u^{3}\ge0$ for
every $u\neq0$, and $\psi$ is non-decreasing (continuity at $0$ follows from
the series $\psi(u)=\sum_{k\ge2}u^{k-2}/k!$). For a real $x\le1$ and
$\lambda>0$ we have $\lambda x\le\lambda$, hence
$\psi(\lambda x)\le\psi(\lambda)$, i.e.\ for $x\neq0$
\begin{equation*}
  \frac{e^{\lambda x}-1-\lambda x}{\lambda^{2}x^{2}}\ \le\
  \frac{e^{\lambda}-1-\lambda}{\lambda^{2}},
  \qquad\text{so}\qquad
  e^{\lambda x}\ \le\ 1+\lambda x+x^{2}g(\lambda),
\end{equation*}
and the last inequality holds at $x=0$ too. Apply it to $X$ pointwise and
take expectations: $\mathbb{E}[e^{\lambda X}]\le1+\lambda\,\mathbb{E}[X]
+g(\lambda)\,\mathbb{E}[X^{2}]=1+\sigma^{2}g(\lambda)$, which is
\eqref{eq:mgf}.

\emph{Step 2 (exponential Markov).} Write $\Ncount{}=\sum_{j=1}^{n}\chi_j$ with
$\chi_j$ independent Bernoulli$(p)$ (Definition~\ref{def:freq}), and
$X_j\coloneqq\chi_j-p$. Then $\mathbb{E}[X_j]=0$,
$\operatorname{Var}(X_j)=p(1-p)$, and $X_j\le1-p\le1$. For $\lambda>0$,
Markov's inequality applied to the non-negative variable
$e^{\lambda\sum_jX_j}$ and independence give
\begin{equation*}
  \Pr\Big[\sum_{j}X_j\ge t\Big]
  \ \le\ e^{-\lambda t}\,\mathbb{E}\Big[e^{\lambda\sum_jX_j}\Big]
  \ =\ e^{-\lambda t}\prod_{j=1}^{n}\mathbb{E}\big[e^{\lambda X_j}\big]
  \ \le\ \exp\big(-\lambda t+v\,g(\lambda)\big),
\end{equation*}
by \eqref{eq:mgf} with $\sigma^{2}=p(1-p)$ and $np(1-p)=v$.

\emph{Step 3 (an elementary bound on $g$).} For integers $k\ge2$,
$k!\ge2\cdot3^{k-2}$ (induction: $2!=2$, and
$(k+1)!=(k+1)\,k!\ge3\cdot2\cdot3^{k-2}$). Hence for $0<\lambda<3$,
\begin{equation*}
  g(\lambda)=\sum_{k\ge2}\frac{\lambda^{k}}{k!}
  \ \le\ \frac{\lambda^{2}}{2}\sum_{k\ge2}\Big(\frac{\lambda}{3}\Big)^{k-2}
  \ =\ \frac{\lambda^{2}}{2\,(1-\lambda/3)}.
\end{equation*}

\emph{Step 4 (the choice of $\lambda$).} Put
$\lambda\coloneqq t/(v+t/3)$, which lies in $(0,3)$ because $v>0$. Then
$1-\lambda/3=v/(v+t/3)$, so Step 3 gives
$g(\lambda)\le\lambda^{2}(v+t/3)/(2v)=t^{2}/\big(2v(v+t/3)\big)$, and the
exponent of Step 2 is at most
\begin{equation*}
  -\lambda t+v\,g(\lambda)
  \ \le\ -\frac{t^{2}}{v+t/3}+\frac{t^{2}}{2(v+t/3)}
  \ =\ -\frac{t^{2}}{2(v+t/3)},
\end{equation*}
which is the first inequality of \eqref{eq:bernstein}. For the second, apply
Steps 1--4 to $-X_j=p-\chi_j$, which has the same mean and variance and
satisfies $-X_j\le p\le1$: $\Pr[\sum_j(-X_j)\ge t]=\Pr[\Ncount{}-np\le-t]$.

\emph{Step 5 (the closed-form deviation).} For $t>0$, the inequality
$t^{2}/(2(v+t/3))\ge L$ is equivalent to $t^{2}-(2L/3)\,t-2vL\ge0$. The
quadratic on the left has the two real roots $L/3\pm\sqrt{(L/3)^{2}+2vL}$, of
which $t(L)$ is the positive one and the other is $\le0$; so the inequality
holds for every $t\ge t(L)$, and then both bounds in \eqref{eq:bernstein}
are at most $e^{-L}$.
\end{proof}

% claim statistics/binomial-box/prop:rule-bernstein
\claimsummary{statistics/binomial-box/prop:rule-bernstein}{Bernstein's inequality with upward-rounded rational bounds on the logarithm and the root gives closed-form endpoints, an admissible rule, monotone in the allocation when every rounding is. Its unrounded half-width falls strictly in n as n\textasciicircum{}(\ensuremath{-}1/2), and it is never narrower than the exact rule.}
\claimbookmark{statistics/binomial-box/prop:rule-bernstein}{Proposition}{prop:rule-bernstein}{The Bernstein endpoint rule}
\claimleadin{%
The emitted interval of the
\termref{term-exact-binomial-quantile-rule}{exact rule} is not monotone in $n$
(Theorem~\ref{thm:monotonicity-limits}\,(v)). The Bernstein rule is the alternative whose
unrounded half-width decreases strictly with $n$; that its integer-rounded
boxes are nested in $n$ is not claimed. Its endpoints come in closed form
from Bernstein's inequality, with no tail evaluated, and every
logarithm and root is bounded by rationals rounded in the direction that can
only widen the interval, the direction a \termref{term-conservative-evaluation}{conservative evaluation}
(Definition~\ref{def:conservative-evaluation}) admits.
A request selects it (Requirement~\ref{req:rule-in-force}), and $j$ and $s$
are recorded with the rule (Requirement~\ref{req:honest-model-record}).
}
\begin{proposition}[The Bernstein endpoint rule]\label{prop:rule-bernstein}
Assume a \termref{term-conservative-evaluation}{conservative evaluation}
(Definition~\ref{def:conservative-evaluation}). Fix an integer $j\ge1$ and a
precision $s\ge0$ (parameters of the evaluation). The \emph{Bernstein rule}
$R_{\mathrm{B}}$ assigns to
$(n,p,\varepsilon^{-},\varepsilon^{+})$, with $v\coloneqq np(1-p)$:
\begin{itemize}[leftmargin=1.9em]
\item if $p=0$: $R^{-}_{\mathrm{B}}=R^{+}_{\mathrm{B}}=0$; if $p=1$:
$R^{-}_{\mathrm{B}}=R^{+}_{\mathrm{B}}=n$;
\item otherwise, for each sign separately: if $\varepsilon^{-}=0$ then
$R^{-}_{\mathrm{B}}=0$, and if $\varepsilon^{+}=0$ then
$R^{+}_{\mathrm{B}}=n$; if $\varepsilon^{\mp}>0$, let $L^{\mp}$ be the
rational upper bound $m(u-1)$ of \eqref{eq:log-rational-bounds} on
$\ln(1/\varepsilon^{\mp})$ with $m=2^{j}$ and $u$ the upper root enclosure
of \eqref{eq:isqrt-bounds} at precision $s$, let
\begin{equation}\label{eq:bernstein-endpoint-t}
  t^{\mp}\ \coloneqq\ \text{a rational }\ \ge\
  \frac{L^{\mp}}{3}+\sqrt{\Big(\frac{L^{\mp}}{3}\Big)^{2}+2vL^{\mp}}
\end{equation}
obtained by rounding the square root upward with \eqref{eq:isqrt-bounds},
and put
\begin{equation}\label{eq:bernstein-endpoints}
  R^{-}_{\mathrm{B}}\ \coloneqq\ \max\big\{0,\ \lfloor np-t^{-}\rfloor\big\},
  \qquad
  R^{+}_{\mathrm{B}}\ \coloneqq\ \min\big\{n,\ \lceil np+t^{+}\rceil\big\}.
\end{equation}
\end{itemize}
Then:
\begin{enumerate}[label=(\roman*),leftmargin=2.2em]
\item $R_{\mathrm{B}}$ is an \termref{term-admissible-endpoint-rule}{admissible endpoint rule} in the sense of
Definition~\ref{def:endpoint-rule}\,(a)--(c), computed from $n$, $p$ and the
allocation by finitely many operations on integers and rationals (integer
square roots included) and no others; Lemma~\ref{lem:guard} therefore
applies to it;
\item $R_{\mathrm{B}}$ is \termref{term-monotone-in-the-allocation}{monotone in the allocation}, provided every
rounding in $L^{\mp}$ and $t^{\mp}$ is a non-decreasing function of its
real argument. The integer-square-root rounding of \eqref{eq:isqrt-bounds}
is \emph{not} such a function of a rational argument as written (its excess over the
root depends on the denominator), but it becomes one when the argument is
first rounded up to the fixed dyadic grid $2^{-s}\mathbb{Z}$, i.e.\ replaced
by $\lceil y\,2^{s}\rceil/2^{s}$ before the root is taken; the evaluation
takes $u$ from \eqref{eq:isqrt-bounds} applied to that grid point (with
$d=2^{s}$ fixed) at every step of Lemma~\ref{lem:log-enclosure}'s
iteration, rounds the square root in \eqref{eq:bernstein-endpoint-t} the
same way, and takes $L^{\mp}$ and $t^{\mp}$ themselves as the least
multiples of $2^{-s}$ at or above the values so obtained. Every such step
is a composition of non-decreasing maps, and every rounding is upward, so
(i) is unaffected;
\item \emph{(half-width)} for $0<p<1$ and $\varepsilon^{\mp}>0$, the
\termref{term-guarded-endpoints}{guarded endpoints} of $R_{\mathrm{B}}$ satisfy
\begin{equation}\label{eq:bernstein-width}
  0\ \le\ q^{+}(\hat c)-p\ \le\ \frac{t^{+}}{n}+\frac1n,
  \qquad
  0\ \le\ p-q^{-}(\hat c)\ \le\ \frac{t^{-}}{n}+\frac1n,
\end{equation}
where, for the unrounded value of \eqref{eq:bernstein-endpoint-t},
\begin{equation}\label{eq:bernstein-halfwidth}
  \frac{t^{\mp}}{n}
  \ =\ \frac{L^{\mp}}{3n}+\sqrt{\Big(\frac{L^{\mp}}{3n}\Big)^{2}
        +\frac{2p(1-p)L^{\mp}}{n}}
\end{equation}
is strictly decreasing in $n$ and is $O(n^{-1/2})$ at fixed $p$ and
$\varepsilon^{\mp}$. The rule never emits an interval narrower than the raw
quantile interval of the exact rule of Definition~\ref{def:interval}
(Proposition~\ref{prop:rule-exact}\,(iii)).
\end{enumerate}
\end{proposition}
\claimmeta{\metalabel{Uses}\mref{Def.}{def:conservative-evaluation}, \mref{Def.}{def:endpoint-rule}, \mref{Lem.}{lem:guard}, \mref{Def.}{def:interval}, \mref{Lem.}{lem:log-enclosure}, \mref{Lem.}{lem:bernstein-tail}
\metasep \metalabel{Used by}\mref{Rem.}{rem:completeness-on-rules}
\metasep newly derived
\metasep implemented (effective: accepted)
\metasep \metalabel{Code}\coderef{https://github.com/fraunhofersg/certiqs-key/blob/v0.1.0/src/certiqs_key/statistics/rules.py}{certiqs_key.statistics.rules:BernsteinRule}}

\begin{proof}
\emph{(i).} Definition~\ref{def:endpoint-rule}\,(c) and (b) hold by the
first two clauses of the rule. For (a) at $p\in\{0,1\}$: the count is $0$
respectively $n$ almost surely, so with $R^{\mp}_{\mathrm{B}}=0$
respectively $n$ both tails of \eqref{eq:rule-admissible} vanish and the
values lie in $\{0,\dots,n\}$. For (a) at $0<p<1$: the values
\eqref{eq:bernstein-endpoints} lie in $\{0,\dots,n\}$ by the clipping. If
$\varepsilon^{+}=0$ then $R^{+}_{\mathrm{B}}=n$ and
$\Pr[\Ncount{}>n]=0\le\varepsilon^{+}$. If $\varepsilon^{+}>0$: by
Lemma~\ref{lem:log-enclosure}, $L^{+}\ge\ln(1/\varepsilon^{+})$, so
$e^{-L^{+}}\le\varepsilon^{+}$; by Lemma~\ref{lem:bernstein-tail} with
$L=L^{+}$, every $t\ge t(L^{+})$ has $\Pr[\Ncount{}-np\ge t]\le e^{-L^{+}}$,
and $t^{+}\ge t(L^{+})$ by \eqref{eq:bernstein-endpoint-t}. If
$R^{+}_{\mathrm{B}}=n$ the tail is $0$. Otherwise
$R^{+}_{\mathrm{B}}=\lceil np+t^{+}\rceil\ge np+t^{+}$, so the event
$\{\Ncount{}>R^{+}_{\mathrm{B}}\}$ is contained in
$\{\Ncount{}-np>t^{+}\}\subseteq\{\Ncount{}-np\ge t^{+}\}$ and
$\Pr[\Ncount{}>R^{+}_{\mathrm{B}}]\le e^{-L^{+}}\le\varepsilon^{+}$. The
lower side is the mirror image: if $R^{-}_{\mathrm{B}}=0$ the tail
$\Pr[\Ncount{}<0]$ is $0$; otherwise
$R^{-}_{\mathrm{B}}=\lfloor np-t^{-}\rfloor\le np-t^{-}$, so
$\{\Ncount{}<R^{-}_{\mathrm{B}}\}\subseteq\{\Ncount{}-np<-t^{-}\}
\subseteq\{\Ncount{}-np\le-t^{-}\}$, whose probability is at most
$e^{-L^{-}}\le\varepsilon^{-}$ by the second inequality of
\eqref{eq:bernstein}. The computation uses the rational arithmetic of the
inputs, the integer square roots of \eqref{eq:isqrt-bounds} for $u$ and for
the root in \eqref{eq:bernstein-endpoint-t} (applied to the rational
$(L^{\mp}/3)^{2}+2vL^{\mp}$), and floor, ceiling, maximum and minimum.

\emph{(ii).} Fix all arguments but $\varepsilon^{+}$ and let
$\varepsilon^{+}_1\le\varepsilon^{+}_2$, both positive. Then
$1/\varepsilon^{+}_1\ge1/\varepsilon^{+}_2$. With the argument first
rounded up to the grid, $y\mapsto\lceil y\,2^{s}\rceil$ is non-decreasing,
and for a \emph{fixed} denominator $d=2^{s}$ the bound
$u=(\operatorname{isqrt}(c\,d\,4^{s})+1)/(d\,2^{s})$ of
\eqref{eq:isqrt-bounds} is non-decreasing in the integer numerator $c$
because $\operatorname{isqrt}$ is non-decreasing on the integers; so each
step of the root iteration is non-decreasing in its input, and $u$ is
non-decreasing in $1/\varepsilon^{+}$. The final upward rounding to the
grid is non-decreasing too, so $L^{+}_1\ge L^{+}_2$; the map
$L\mapsto L/3+\sqrt{(L/3)^{2}+2vL}$ is increasing in $L\ge0$ and the upward
rounding of the root is monotone, so $t^{+}_1\ge t^{+}_2$; hence
$R^{+}_{\mathrm{B},1}\le R^{+}_{\mathrm{B},2}$ by
\eqref{eq:bernstein-endpoints}, which is the required non-increase of
$\varepsilon^{+}\mapsto R^{+}_{\mathrm{B}}$. At $\varepsilon^{+}_1=0$ the
value $n$ is the largest possible, so the monotonicity extends to the zero
budget. The lower side is the mirror image ($R^{-}_{\mathrm{B}}$
non-decreasing in $\varepsilon^{-}$, with the value $0$ at
$\varepsilon^{-}=0$ the smallest possible).

\emph{(iii).} The left inequalities of \eqref{eq:bernstein-width} are
\eqref{eq:ordering}. For the right ones: by \eqref{eq:guard} and
\eqref{eq:bernstein-endpoints}, $k^{+}(\hat c)=\max\{R^{+}_{\mathrm{B}},
\lceil np\rceil\}\le\lceil np+t^{+}\rceil\le np+t^{+}+1$ (using
$\lceil np\rceil\le\lceil np+t^{+}\rceil$ and $n\le\lceil np+t^{+}\rceil$
when the clip at $n$ is active), so $q^{+}(\hat c)-p\le(t^{+}+1)/n$;
symmetrically $k^{-}(\hat c)=\min\{R^{-}_{\mathrm{B}},\lfloor np\rfloor\}
\ge\lfloor np-t^{-}\rfloor\ge np-t^{-}-1$ when the clip at $0$ is inactive,
and $k^{-}(\hat c)\ge0>np-t^{-}-1$ when it is active (then
$\lfloor np-t^{-}\rfloor\le0$, so $np-t^{-}<1$); hence
$p-q^{-}(\hat c)\le(t^{-}+1)/n$. Dividing \eqref{eq:bernstein-endpoint-t} by
$n$ gives \eqref{eq:bernstein-halfwidth}; each of its two terms is strictly
decreasing in $n$ (the first as $1/n$, the second as the square root of a sum
of a term in $1/n^{2}$ and a term in $1/n$, both decreasing, at fixed
$p\in(0,1)$ and $L^{\mp}>0$), and it is $O(n^{-1/2})$ because the first term
is $O(n^{-1})$ and the radicand is $O(n^{-1})$. The comparison with the exact
rule: the exact rule's $\kappa^{+}$ is the minimum of the set in
\eqref{eq:kappaplus} (Definition~\ref{def:interval}), and
$R^{+}_{\mathrm{B}}\in\{0,\dots,n\}$ lies in that set by (i), so
$R^{+}_{\mathrm{B}}\ge\kappa^{+}$; likewise $\kappa^{-}$ is the maximum of the
set in \eqref{eq:kappaminus}, in which $R^{-}_{\mathrm{B}}$ lies by (i), so
$R^{-}_{\mathrm{B}}\le\kappa^{-}$; the guard preserves both inequalities.
\end{proof}

% claim statistics/binomial-box/rem:completeness-on-rules
\claimsummary{statistics/binomial-box/rem:completeness-on-rules}{The tightest completeness figure this submodule certifies on any admissible rule is \ensuremath{\varepsilon}*, but a rule that evaluates no tail can report the allocation sum or a clamped upward-rounded tail evaluation; both bound the honest failure probability. Neither choice changes the emitted set or soundness.}
\claimbookmark{statistics/binomial-box/rem:completeness-on-rules}{Remark}{rem:completeness-on-rules}{Completeness figures on a closed-form rule}
\begin{remark}[Completeness figures on a closed-form rule]\label{rem:completeness-on-rules}
Theorem~\ref{thm:completeness-bound} holds for every \termref{term-admissible-endpoint-rule}{admissible endpoint rule}, and
its $\varepsilon^{\star}$, the sum of the binomial tails at the
\termref{term-guarded-endpoints}{guarded endpoints}, is the tightest completeness figure this submodule
certifies for the emitted set, whatever rule produced the endpoints: every
figure this remark lists is at least $\varepsilon^{\star}$, while the honest failure
probability itself, which $\varepsilon^{\star}$ bounds by a union bound, may
be smaller. What differs between rules is
the cost of \emph{evaluating} it. On the exact rule of
Definition~\ref{def:interval} the tails are decided in the course of finding
the quantiles, and $\varepsilon^{\star}$ is reported exactly wherever they are
decided exactly; where they are decided against the upper bounds of
Proposition~\ref{prop:rule-exact}\,(vii) it is rounded upward and clamped
(Remark~\ref{rem:epsstar}). On a closed-form rule such as the Bernstein rule
of Proposition~\ref{prop:rule-bernstein}, the endpoints are found without
evaluating any tail, and two reports are admissible, both upward-rounded in
the sense of a \termref{term-conservative-evaluation}{conservative evaluation}
(Definition~\ref{def:conservative-evaluation}) and Remark~\ref{rem:epsstar}:
\begin{enumerate}[label=(\alph*),leftmargin=2.2em]
\item \emph{the allocation sum}
$\sum_{\hat c\in\Chatcon}\big(\varepsilon^{-}(\hat c)+\varepsilon^{+}(\hat
c)\big)$. By Lemma~\ref{lem:guard}\,(ii) each tail at a guarded endpoint is
at most its allocation, so the sum is at least $\varepsilon^{\star}$ and
hence, by \eqref{eq:completeness}, at least the honest failure probability;
and it lies in $[0,\ecomp]$ because the allocation of an
\termref{term-admissible-input}{admissible input} is non-negative and sums to at most $\ecomp$
(Definition~\ref{def:endpoint-rule}), so no clamp is needed. It costs nothing
beyond the inputs and is what a rule that evaluates no tail can always
report;
\item \emph{an upward-rounded evaluation of the tails at the emitted
endpoints}, $\min\{\ecomp,\tilde\varepsilon\}$ for any
$\tilde\varepsilon\ge\varepsilon^{\star}$ obtained by bounding each tail from
above (exactly, by the closed forms of the exact rule's evaluation, or by an
enclosure), which is the clamped figure of Remark~\ref{rem:epsstar}. At
$\tilde\varepsilon=\varepsilon^{\star}$ it is tighter than (a) by the budget
the guarded endpoints leave unspent,
$\sum_{\hat c\in\Chatcon}(\varepsilon^{-}(\hat c)+\varepsilon^{+}(\hat c))-\varepsilon^{\star}$,
which on a closed-form rule includes the slack of the rule's own tail bound
besides what integrality and the guard leave; an
upward-rounded $\tilde\varepsilon>\varepsilon^{\star}$ need not be at most
(a), since an upper bound on a tail at a guarded endpoint may exceed that
tail's allocation.
\end{enumerate}
Neither choice affects the emitted set, and neither affects soundness
(Remark~\ref{rem:sound-indep}); the difference is the advertised abort
probability of the honest run, which (a) overstates by the unspent budget.
Requirement~\ref{req:completeness-report} fixes which of the two is reported
on each route, and Requirement~\ref{req:honest-model-record} what is recorded
beside the figure.
\end{remark}
\claimmeta{\metalabel{Uses}\mref{Def.}{def:conservative-evaluation}, \mref{Def.}{def:endpoint-rule}, \mref{Lem.}{lem:guard}, \mref{Def.}{def:interval}, \mref{Thm.}{thm:completeness-bound}, \mref{Rem.}{rem:epsstar}, \mref{Rem.}{rem:sound-indep}, \mref{Prop.}{prop:rule-bernstein}
\metasep \metalabel{Used by}\mref{Rem.}{rem:local-evaluation}, \mref{Req.}{req:completeness-report}, \mref{Req.}{req:honest-model-record}
\metasep newly derived
\metasep accepted}

\section{The exact rule's local evaluation}\label{group-statistics/binomial-box/lem:pmf-enclosure}
\sectionintro{A binomial tail is enclosed in log space from a window of mass points and a geometric majorant of the rest, each mass point enclosed through Stirling's bounds. Endpoints bracketed by these enclosures form an \termref{term-admissible-endpoint-rule}{admissible endpoint rule} that relaxes the exact quantiles: equal to them where the bracket is decided, and otherwise, where the tail decays at the quantile, conservative by a number of counts bounded through the enclosure's slack and that decay (\hyperlink{claim-statistics/binomial-box/prop:bracketed-endpoints}{Proposition~\ref*{prop:bracketed-endpoints}}).}

% claim statistics/binomial-box/lem:pmf-enclosure
\claimsummary{statistics/binomial-box/lem:pmf-enclosure}{Rational bounds on ln m!, ln p and ln(1 \ensuremath{-} p) enclose the logarithm of a binomial mass point, with width the sum of the component widths. On the Stirling branch no intermediate is of the size of the mass point, of p\textasciicircum{}j or of a binomial coefficient.}
\claimbookmark{statistics/binomial-box/lem:pmf-enclosure}{Lemma}{lem:pmf-enclosure}{Rational enclosure of the logarithm of a binomial mass point}
\claimleadin{%
At large $n$ a binomial mass point such as $p^{n}$ is too small, and its
exact rational too long, for an implementation to hold, while its logarithm
is of moderate size. The logarithm of the law of Theorem~\ref{thm:binomial-tails} is a sum of
three log-factorials and two logarithms, and each is replaced by its
rational bound in the direction its sign requires.
}
\begin{technicalclaim}
\begin{lemma}[Rational enclosure of the logarithm of a binomial mass point]\label{lem:pmf-enclosure}
Let $n\ge1$, let $p\in(0,1)$ be rational, let $\Ncount{}$ be binomial with
parameters $(n,p)$ and write $\mathrm{pmf}(j)\coloneqq\Pr[\Ncount{}=j]
=\binom{n}{j}p^{j}(1-p)^{n-j}>0$ for $0\le j\le n$ (Theorem~\ref{thm:binomial-tails}),
and $F(m)\coloneqq\ln m!$ for an integer $m\ge0$. Then
\begin{equation}\label{eq:lnpmf}
  \ln\mathrm{pmf}(j)\ =\ F(n)-F(j)-F(n-j)+j\ln p+(n-j)\ln(1-p).
\end{equation}
Let $F^{-}(m)\le F(m)\le F^{+}(m)$ be rational bounds for
$m\in\{n,j,n-j\}$ --- those of Lemma~\ref{lem:stirling-enclosure}\,(iii)
for $m\ge1$, or $F^{\mp}(m)$ the root bounds of
Lemma~\ref{lem:log-enclosure} on the integer $m!$ where that integer is held
exactly, and $F^{\mp}(0)=0$ --- and let
$\lambda_p^{-}\le\ln p\le\lambda_p^{+}$ and
$\lambda_q^{-}\le\ln(1-p)\le\lambda_q^{+}$ be rational bounds
(Lemma~\ref{lem:log-enclosure}). Put
\begin{equation}\label{eq:pmf-enclosure}
\begin{aligned}
  \ell^{-}(j)&\coloneqq F^{-}(n)-F^{+}(j)-F^{+}(n-j)
              +j\,\lambda_p^{-}+(n-j)\,\lambda_q^{-},\\
  \ell^{+}(j)&\coloneqq F^{+}(n)-F^{-}(j)-F^{-}(n-j)
              +j\,\lambda_p^{+}+(n-j)\,\lambda_q^{+}.
\end{aligned}
\end{equation}
Then $\ell^{-}(j)\le\ln\mathrm{pmf}(j)\le\ell^{+}(j)$, both bounds are
rationals, and the width is the sum of the component widths,
\begin{equation}\label{eq:pmf-width}
  \ell^{+}(j)-\ell^{-}(j)
  =\sum_{m\in\{n,\,j,\,n-j\}}\big(F^{+}(m)-F^{-}(m)\big)
  +j\,(\lambda_p^{+}-\lambda_p^{-})+(n-j)\,(\lambda_q^{+}-\lambda_q^{-}).
\end{equation}
In particular, if every logarithm entering the bounds is enclosed to width
$\delta$ and $\ln(2\pi)$ to width $\delta_C$, the width is at most
$(3n+\tfrac32)\,\delta+\tfrac32\,\delta_C+\tfrac{3}{156}$. A log precision
of $2^{-s}$ with $s$ about $\log_2 n$ bits beyond the wanted precision brings
the terms in $\delta$ and $\delta_C$ within it; the remaining term, at most
$\tfrac{1}{12m(12m+1)}$ for each $m\in\{n,j,n-j\}$ taken on the Stirling
branch, is fixed by $m$, and no choice of $s$ reduces it. On the Stirling
branch every quantity in \eqref{eq:pmf-enclosure} is a rational whose size
is polynomial in the bit sizes of $n$, $p$ and the precision, and no intermediate has the magnitude of $\mathrm{pmf}(j)$,
$p^{j}$ or $\binom{n}{j}$ themselves; the exact-factorial branch is for
small $m$ only and holds the integer $m!$ by construction.
\end{lemma}
\end{technicalclaim}
\claimmeta{\metalabel{Uses}\mref{Thm.}{thm:binomial-tails}, \mref{Lem.}{lem:log-enclosure}, \mref{Lem.}{lem:stirling-enclosure}
\metasep \metalabel{Used by}\mref{Prop.}{prop:bracketed-endpoints}, \mref{Rem.}{rem:local-evaluation}, \mref{Req.}{req:enclosure-substrate}
\metasep newly derived
\metasep implemented
\metasep \metalabel{Code}\coderef{https://github.com/fraunhofersg/certiqs-key/blob/v0.1.0/src/certiqs_key/statistics/enclosure.py}{certiqs_key.statistics.enclosure:log_pmf_bounds}}

\begin{proof}
\eqref{eq:lnpmf} is the logarithm of the law, using
$\ln\binom{n}{j}=\ln n!-\ln j!-\ln(n-j)!$ and $p,1-p>0$. For the bounds,
replace each term of \eqref{eq:lnpmf} by the bound in the direction its
sign requires: $F(n)$, $j\ln p$ and $(n-j)\ln(1-p)$ enter with positive
coefficients ($j\ge0$, $n-j\ge0$) and are bounded below by their lower and
above by their upper bounds; $F(j)$ and $F(n-j)$ enter with a minus sign, so
a lower bound on $\ln\mathrm{pmf}(j)$ uses their upper bounds and an upper
bound uses their lower bounds. Subtracting the two displays of
\eqref{eq:pmf-enclosure} gives \eqref{eq:pmf-width}. For the estimate: on
the Stirling branch, by Lemma~\ref{lem:stirling-enclosure}\,(iii), the width
of $F^{\mp}(m)$ is at most
$(m+\tfrac12)\delta+\tfrac12\delta_C+\tfrac{1}{12m(12m+1)}$, and
$\frac{1}{12m(12m+1)}\le\frac{1}{156}$ for $m\ge1$; on the exact-factorial
branch it is the width of the root bounds on $\ln m!$, a logarithm entering
the bounds, so at most $\delta\le(m+\tfrac12)\delta$ for $m\ge1$; and for
$m=0$ it is $0$. So each of the three widths is at most
$(m+\tfrac12)\delta+\tfrac12\delta_C+\tfrac{1}{156}$; summing over
$m\in\{n,j,n-j\}$, whose values add up to $2n$, and adding
$j\delta+(n-j)\delta=n\delta$ gives at most
$(2n+\tfrac32)\delta+n\delta+\tfrac32\delta_C+\tfrac{3}{156}$. The size
statement holds because each term of \eqref{eq:pmf-enclosure} is a sum of
products of the given rationals with the integers $n$, $j$, $n-j$ and
$m+\tfrac12$; the rational $\frac{1}{12m+1}$ and its like are of the bit
size of $m$.
\end{proof}

% claim statistics/binomial-box/lem:tail-majorant
\claimsummary{statistics/binomial-box/lem:tail-majorant}{A binomial upper tail lies between a window of exact mass points and that window plus a geometric majorant of the rest, because consecutive mass ratios decrease. The same ratios give log-space bounds in units of one mass point, a per-count decay of the tail and, by reflection, the lower tail.}
\claimbookmark{statistics/binomial-box/lem:tail-majorant}{Lemma}{lem:tail-majorant}{The binomial tail from a window and a geometric majorant}
\claimleadin{%
The lemma is the classical geometric bound on a binomial tail through the
decreasing ratio of consecutive mass points, Feller's tail estimate for the
binomial law, derived in full here from the ratio identity, a telescoping
product and a geometric series.
}
\begin{technicalclaim}
\begin{lemma}[The binomial tail from a window and a geometric majorant]\label{lem:tail-majorant}
Let $n\ge1$, let $p\in(0,1)$ be rational, let $\Ncount{}$ be binomial with
parameters $(n,p)$, and write $\mathrm{pmf}(j)\coloneqq\Pr[\Ncount{}=j]$ for
$0\le j\le n$ (Theorem~\ref{thm:binomial-tails}) and $\mathrm{pmf}(j)\coloneqq0$ for
$j>n$.
\begin{enumerate}[label=(\roman*),leftmargin=2.2em]
\item \emph{(ratios)} For $0\le j\le n-1$,
\begin{equation}\label{eq:pmf-ratio}
  r_j\ \coloneqq\ \frac{\mathrm{pmf}(j+1)}{\mathrm{pmf}(j)}
  \ =\ \frac{(n-j)\,p}{(j+1)\,(1-p)},
\end{equation}
a positive rational, strictly decreasing in $j$; put $r_n\coloneqq0$.
\item \emph{(window plus majorant)} Let $0\le k\le n-1$ and let $T\ge0$ be
an integer with $k+T+1\le n$ and $r_{k+T+1}<1$ (which holds in particular
when $k+T+1=n$). Then, with $S_T\coloneqq\sum_{i=0}^{T-1}\mathrm{pmf}(k+1+i)$,
\begin{equation}\label{eq:tail-window}
  S_T\ \le\ \Pr[\Ncount{}>k]\ \le\
  S_T+\frac{\mathrm{pmf}(k+T+1)}{1-r_{k+T+1}} .
\end{equation}
If instead $k+T+1>n$, then $\Pr[\Ncount{}>k]=S_T$ exactly.
\item \emph{(in units of one mass point)} In the setting of (ii) with
$T\ge1$, put
\begin{equation}\label{eq:window-rationals}
  A_T\coloneqq\sum_{i=0}^{T-1}\ \prod_{l=k+1}^{k+i}r_l,
  \qquad
  B_T\coloneqq\frac{\prod_{l=k+1}^{k+T}r_l}{1-r_{k+T+1}}
\end{equation}
(the empty product is $1$), positive rationals computable from $n$, $p$, $k$
and $T$ alone. Then
\begin{equation}\label{eq:tail-log}
  \ln\mathrm{pmf}(k+1)+\ln A_T
  \ \le\ \ln\Pr[\Ncount{}>k]\ \le\
  \ln\mathrm{pmf}(k+1)+\ln(A_T+B_T).
\end{equation}
\item \emph{(per-count decay)} For $1\le k\le n$ with $r_k<1$,
\begin{equation}\label{eq:tail-decay}
  \Pr[\Ncount{}>k]\ \le\ r_k\,\Pr[\Ncount{}>k-1].
\end{equation}
\item \emph{(the lower tail)} $\Ncount{}'\coloneqq n-\Ncount{}$ is binomial
with parameters $(n,1-p)$, and $\Pr[\Ncount{}<k]=\Pr[\Ncount{}'>n-k]$ for
every integer $k$; so (i)--(iv) applied to $\Ncount{}'$ enclose the lower
tail of $\Ncount{}$.
\end{enumerate}
\end{lemma}
\end{technicalclaim}
\claimmeta{\metalabel{Uses}\mref{Thm.}{thm:binomial-tails}
\metasep \metalabel{Used by}\mref{Prop.}{prop:bracketed-endpoints}, \mref{Rem.}{rem:local-evaluation}, \mref{Req.}{req:enclosure-substrate}
\metasep newly derived
\metasep implemented
\metasep \metalabel{Code}\coderef{https://github.com/fraunhofersg/certiqs-key/blob/v0.1.0/src/certiqs_key/statistics/enclosure.py}{certiqs_key.statistics.enclosure:upper_tail_log_bounds}}

\begin{proof}
\emph{(i).} From the law,
$\frac{\mathrm{pmf}(j+1)}{\mathrm{pmf}(j)}
=\frac{\binom{n}{j+1}}{\binom{n}{j}}\cdot\frac{p}{1-p}
=\frac{n-j}{j+1}\cdot\frac{p}{1-p}$, using
$\binom{n}{j+1}=\binom{n}{j}\frac{n-j}{j+1}$. The factor $\frac{n-j}{j+1}$
is strictly decreasing in $j$ (numerator decreasing, denominator
increasing, both positive for $0\le j\le n-1$), and $\frac{p}{1-p}>0$.

\emph{(ii).} $\Pr[\Ncount{}>k]=\sum_{j=k+1}^{n}\mathrm{pmf}(j)
=S_T+\sum_{j=k+T+1}^{n}\mathrm{pmf}(j)$, and every term is non-negative, so
$S_T\le\Pr[\Ncount{}>k]$; if $k+T+1>n$ the second sum is empty and equality
holds. For the upper bound let $k+T+1\le n$ and $j\ge k+T+1$. Telescoping,
$\mathrm{pmf}(j)=\mathrm{pmf}(k+T+1)\prod_{l=k+T+1}^{j-1}r_l$, and each
factor satisfies $r_l\le r_{k+T+1}$ by (i) (for $l\le n-1$; and if the
product reaches $l=n$ the mass point is $0$ anyway), so
$\mathrm{pmf}(j)\le\mathrm{pmf}(k+T+1)\,r_{k+T+1}^{\,j-k-T-1}$. Summing over
$j\ge k+T+1$ and bounding the finite geometric sum by the full series with
ratio $r_{k+T+1}<1$,
$\sum_{j\ge k+T+1}\mathrm{pmf}(j)\le\mathrm{pmf}(k+T+1)/(1-r_{k+T+1})$.

\emph{(iii).} Telescoping again,
$\mathrm{pmf}(k+1+i)=\mathrm{pmf}(k+1)\prod_{l=k+1}^{k+i}r_l$ for
$0\le i\le T$, so $S_T=\mathrm{pmf}(k+1)\,A_T$ and
$\mathrm{pmf}(k+T+1)/(1-r_{k+T+1})=\mathrm{pmf}(k+1)\,B_T$. Substituting
into \eqref{eq:tail-window} and taking logarithms (all quantities positive:
$A_T\ge1$ since its first term is $1$, and $B_T>0$) gives
\eqref{eq:tail-log}. The rationals $A_T$, $B_T$ are finite sums and
products of the rationals \eqref{eq:pmf-ratio}.

\emph{(iv).} For every $j\ge k$ with $j\le n-1$, (i) gives
$\mathrm{pmf}(j+1)=r_j\,\mathrm{pmf}(j)\le r_k\,\mathrm{pmf}(j)$, and the
same inequality holds trivially for $j\ge n$ (both sides $0$). Summing over
$j\ge k$: $\Pr[\Ncount{}>k]=\sum_{j\ge k}\mathrm{pmf}(j+1)\le
r_k\sum_{j\ge k}\mathrm{pmf}(j)=r_k\Pr[\Ncount{}\ge k]
=r_k\Pr[\Ncount{}>k-1]$.

\emph{(v).} $\Pr[n-\Ncount{}=i]=\Pr[\Ncount{}=n-i]
=\binom{n}{n-i}p^{n-i}(1-p)^{i}=\binom{n}{i}(1-p)^{i}p^{n-i}$, the law of a
binomial with parameters $(n,1-p)$; and $\Ncount{}<k$ if and only if
$n-\Ncount{}>n-k$.
\end{proof}

% claim statistics/binomial-box/prop:bracketed-endpoints
\claimsummary{statistics/binomial-box/prop:bracketed-endpoints}{Endpoints bracketed by certified tail enclosures form an admissible rule whose box contains the exact rule's. Each equals the exact quantile where its bracket is decided; otherwise, where the tail decays at the quantile, the excess in counts is bounded through the enclosure's slack and that decay.}
\claimbookmark{statistics/binomial-box/prop:bracketed-endpoints}{Proposition}{prop:bracketed-endpoints}{The enclosure evaluation of the exact rule}
\claimleadin{%
The \termref{term-exact-binomial-quantile-rule}{exact rule} compares binomial tails
with their budgets. Where exact sums are too costly, the comparisons can be
decided against enclosures of the tails instead, which is what a
\termref{term-conservative-evaluation}{conservative evaluation} (Definition~\ref{def:conservative-evaluation})
admits: a comparison decided against an upper bound on the tail. A request selects this evaluation
(Requirement~\ref{req:rule-in-force}), and its endpoints are found by the one
search procedure Requirement~\ref{req:search-procedure} fixes;
Remark~\ref{rem:local-evaluation} discusses the search, the enclosures and
their cost. Through the closing sentences of (iii), the width bound
\eqref{eq:width} of Theorem~\ref{thm:monotonicity-limits}\,(iv), proved for
the exact rule, carries over to this evaluation wherever the bracket is
decided and, where it is not, up to the excess (iii) bounds, divided by $n$;
beyond that it is not asserted for this evaluation.
}
\begin{proposition}[The enclosure evaluation of the exact rule]\label{prop:bracketed-endpoints}
Assume a \termref{term-conservative-evaluation}{conservative evaluation}
(Definition~\ref{def:conservative-evaluation}). Let $n\ge1$, let
$p\in(0,1)$ and $\varepsilon^{+}\in(0,1)$ be rational, let
$\Ncount{}$ be binomial with parameters $(n,p)$, and let $\kappa^{+}$ be the
exact upper quantile of Definition~\ref{def:interval} at
$(n,p,\varepsilon^{+})$. Let
\begin{equation}\label{eq:tail-enclosure}
  \tau^{-}(k)\ \le\ \Pr[\Ncount{}>k]\ \le\ \tau^{+}(k),\qquad k=0,\dots,n,
\end{equation}
be an enclosure with $\tau^{-}(n)=\tau^{+}(n)=0$, each bound either a
rational or given through a rational bound on its logarithm, such as the
one Lemma~\ref{lem:tail-majorant}\,(ii)--(iii) yields with the mass point
enclosed by Lemma~\ref{lem:pmf-enclosure} and the logarithms by
Lemma~\ref{lem:log-enclosure}. Call $k$ \emph{certified admissible},
$\mathsf{C}^{+}(k)$, when the available upper bound proves
$\Pr[\Ncount{}>k]\le\varepsilon^{+}$: $\tau^{+}(k)\le\varepsilon^{+}$, or in
log space $U(k)\le\lambda^{-}$ for rationals $U(k)\ge\ln\Pr[\Ncount{}>k]$
and $\lambda^{-}\le\ln\varepsilon^{+}$; and \emph{certified inadmissible},
$\mathsf{D}^{+}(k)$, when the lower bound proves
$\Pr[\Ncount{}>k]>\varepsilon^{+}$: $\tau^{-}(k)>\varepsilon^{+}$, or in log
space $\Lambda(k)>\lambda^{+}$ for rationals
$\Lambda(k)\le\ln\Pr[\Ncount{}>k]$ and $\lambda^{+}\ge\ln\varepsilon^{+}$.
An index that is neither is \emph{undecided}. The predicate
$\mathsf{C}^{+}(n)$ always holds.

A \emph{bracketed upper endpoint} is an integer $\tilde k^{+}\in\{0,\dots,n\}$
with $\mathsf{C}^{+}(\tilde k^{+})$ and, if $\tilde k^{+}\ge1$,
$\lnot\mathsf{C}^{+}(\tilde k^{+}-1)$. The \emph{enclosure evaluation of the
exact rule} is the endpoint rule
$R_{\mathrm{encl}}$ whose $R^{+}_{\mathrm{encl}}$ is the bracketed upper
endpoint returned by a fixed deterministic search procedure that returns such
an endpoint for every input, whose $R^{-}_{\mathrm{encl}}$
is obtained in the same way for the lower tail through the mirror
$\Ncount{}'=n-\Ncount{}$ of Lemma~\ref{lem:tail-majorant}\,(v), and which
returns the values of Definition~\ref{def:endpoint-rule}\,(b) and (c) at a
zero budget and at $p\in\{0,1\}$. Then:
\begin{enumerate}[label=(\roman*),leftmargin=2.2em]
\item \emph{(admissible relaxation)} Every bracketed upper endpoint satisfies
$\Pr[\Ncount{}>\tilde k^{+}]\le\varepsilon^{+}$ and
$\kappa^{+}\le\tilde k^{+}$; mirror-wise every bracketed lower endpoint
$\tilde k^{-}$ satisfies $\Pr[\Ncount{}<\tilde k^{-}]\le\varepsilon^{-}$ and
$\tilde k^{-}\le\kappa^{-}$. Hence $R_{\mathrm{encl}}$ is an
\termref{term-admissible-endpoint-rule}{admissible endpoint rule}, Lemma~\ref{lem:guard} applies to it, and the
box
$\mathsf{Acc}_{R_{\mathrm{encl}}}$ emits (Definition~\ref{def:construction})
contains the box the exact rule emits: every $q^{-}$ is at most and every
$q^{+}$ at least the exact rule's.
\item \emph{(the bracket decided)} If $\tilde k^{+}\ge1$ and
$\mathsf{D}^{+}(\tilde k^{+}-1)$ holds, then $\tilde k^{+}=\kappa^{+}$; if
$\tilde k^{+}=0$ then $\tilde k^{+}=\kappa^{+}$. The mirror statement holds
for the lower side.
\item \emph{(conservatism in counts)} Suppose $\tilde k^{+}>\kappa^{+}$,
$\kappa^{+}\le n-2$ and $r_{\kappa^{+}+1}<1$, with $r_j$ the ratios
\eqref{eq:pmf-ratio}. Let $w\ge0$
be the slack of the failed certificate at $\tilde k^{+}-1$:
$w\coloneqq\ln\tau^{+}(\tilde k^{+}-1)-\ln\Pr[\Ncount{}>\tilde k^{+}-1]$, or
in log space
$w\coloneqq U(\tilde k^{+}-1)-\ln\Pr[\Ncount{}>\tilde k^{+}-1]
+(\ln\varepsilon^{+}-\lambda^{-})$. Then
\begin{equation}\label{eq:bracket-conservatism}
  \tilde k^{+}-\kappa^{+}\ <\ 1+\frac{w}{-\ln r_{\kappa^{+}+1}} .
\end{equation}
The emitted endpoint is therefore the exact quantile up to a number of
counts bounded by the enclosure's slack divided by the per-count decay of
the tail at the quantile. The mirror statement holds for the lower side,
through $\Ncount{}'$ (Lemma~\ref{lem:tail-majorant}\,(v)): with
$\kappa^{-}=n-\kappa'^{+}$ and $\tilde k^{-}=n-\tilde k'^{+}$, where the
primed quantities are those of $\Ncount{}'$ at level $\varepsilon^{-}$,
$\kappa^{-}-\tilde k^{-}<1+w'/(-\ln r'_{\kappa'^{+}+1})$ under the mirrored
hypotheses. Consequently, wherever the bracket is decided the emitted
endpoint is the exact rule's, and so are its guarded endpoint and its box
bound; wherever it is undecided the guarded endpoint exceeds the exact rule's
by at most $\tilde k^{+}-\kappa^{+}$ counts and the box bound $q^{+}$ the
exact rule's by at most $(\tilde k^{+}-\kappa^{+})/n$, which
\eqref{eq:bracket-conservatism} bounds under the hypotheses above;
mirror-wise for $k^{-}$ and $q^{-}$. Where the hypothesis $r_{\kappa^{+}+1}<1$ fails (the
exact quantile at or below the mean, reachable at a near-unit tail budget,
Proposition~\ref{prop:rule-exact}\,(vi)) the endpoint is still admissible by
(i) and only this bound is unavailable; at $\kappa^{+}=n-1$ the excess
$\tilde k^{+}-\kappa^{+}$ is at most one count.
\end{enumerate}
\end{proposition}
\claimmeta{\metalabel{Uses}\mref{Def.}{def:conservative-evaluation}, \mref{Def.}{def:endpoint-rule}, \mref{Lem.}{lem:guard}, \mref{Def.}{def:interval}, \mref{Prop.}{prop:rule-exact}, \mref{Def.}{def:construction}, \mref{Lem.}{lem:log-enclosure}, \mref{Lem.}{lem:pmf-enclosure}, \mref{Lem.}{lem:tail-majorant}
\metasep \metalabel{Used by}\mref{Rem.}{rem:local-evaluation}, \mref{Req.}{req:search-procedure}, \mref{Req.}{req:endpoint-exact-record}
\metasep newly derived
\metasep implemented (effective: accepted)
\metasep \metalabel{Code}\coderef{https://github.com/fraunhofersg/certiqs-key/blob/v0.1.0/src/certiqs_key/statistics/enclosure.py}{certiqs_key.statistics.enclosure:bracketed_upper_endpoint}}

\begin{proof}
\emph{(i).} $\mathsf{C}^{+}(\tilde k^{+})$ gives
$\Pr[\Ncount{}>\tilde k^{+}]\le\tau^{+}(\tilde k^{+})\le\varepsilon^{+}$, or
in log space $\Pr[\Ncount{}>\tilde k^{+}]\le e^{U(\tilde k^{+})}
\le e^{\lambda^{-}}\le\varepsilon^{+}$. So $\tilde k^{+}$ belongs to the set
in \eqref{eq:kappaplus}, and $\kappa^{+}\le\tilde k^{+}$ by the minimality
that defines $\kappa^{+}$ (Proposition~\ref{prop:rule-exact}\,(iii)). The
lower side is the same argument for $\Ncount{}'$: a certified index
$\tilde k'$ for the upper tail of $\Ncount{}'$ at level $\varepsilon^{-}$
gives $\Pr[\Ncount{}<n-\tilde k']=\Pr[\Ncount{}'>\tilde k']\le\varepsilon^{-}$
by Lemma~\ref{lem:tail-majorant}\,(v), so $\tilde k^{-}\coloneqq n-\tilde k'$
lies in the set of \eqref{eq:kappaminus} and $\tilde k^{-}\le\kappa^{-}$ by
maximality. Both values lie in $\{0,\dots,n\}$; with the conventions at a
zero budget and at $p\in\{0,1\}$ imposed by definition, $R_{\mathrm{encl}}$
satisfies Definition~\ref{def:endpoint-rule}\,(a)--(c). For the containment:
by \eqref{eq:guard} the \termref{term-guarded-endpoints}{guarded endpoints} are monotone in the rule's
endpoints, so $\tilde k^{-}\le\kappa^{-}$ and $\tilde k^{+}\ge\kappa^{+}$
give $k^{-}_{\mathrm{encl}}(\hat c)\le k^{-}_{\mathrm{ex}}(\hat c)$ and
$k^{+}_{\mathrm{encl}}(\hat c)\ge k^{+}_{\mathrm{ex}}(\hat c)$ at every
constrained outcome, hence the same order for the bounds \eqref{eq:boxbounds}
and the inclusion of the sets \eqref{eq:accset}.

\emph{(ii).} If $\tilde k^{+}=0$ then $\kappa^{+}\le0$ by (i), so
$\kappa^{+}=0$. Otherwise $\mathsf{D}^{+}(\tilde k^{+}-1)$ gives
$\Pr[\Ncount{}>\tilde k^{+}-1]\ge\tau^{-}(\tilde k^{+}-1)>\varepsilon^{+}$
(in log space
$\Pr[\Ncount{}>\tilde k^{+}-1]\ge e^{\Lambda(\tilde k^{+}-1)}
>e^{\lambda^{+}}\ge\varepsilon^{+}$), so $\tilde k^{+}-1$ is not in the set
of \eqref{eq:kappaplus}. That set is an up-set of $\{0,\dots,n\}$, because
$k\mapsto\Pr[\Ncount{}>k]$ is non-increasing (Lemma~\ref{lem:guard}\,(ii));
hence its minimum $\kappa^{+}$ exceeds $\tilde k^{+}-1$, i.e.\
$\kappa^{+}\ge\tilde k^{+}$, and with (i) equality. The lower side is the
mirror.

\emph{(iii).} Since $\tilde k^{+}>\kappa^{+}\ge0$, the index $\tilde k^{+}-1$
exists and is not certified admissible, so
$\tau^{+}(\tilde k^{+}-1)>\varepsilon^{+}$ (in log space
$U(\tilde k^{+}-1)>\lambda^{-}$), and by the definition of $w$,
\begin{equation*}
  \ln\Pr[\Ncount{}>\tilde k^{+}-1]\ =\ \ln\tau^{+}(\tilde k^{+}-1)-w
  \ >\ \ln\varepsilon^{+}-w
\end{equation*}
(in log space $\ln\Pr[\Ncount{}>\tilde k^{+}-1]=U(\tilde k^{+}-1)
-\big(U(\tilde k^{+}-1)-\ln\Pr[\Ncount{}>\tilde k^{+}-1]\big)
>\lambda^{-}-\big(w-(\ln\varepsilon^{+}-\lambda^{-})\big)
=\ln\varepsilon^{+}-w$). On the other hand $\tilde k^{+}-1\ge\kappa^{+}$,
and for every $l$ with $\kappa^{+}+1\le l\le\tilde k^{+}-1$ one has
$r_l\le r_{\kappa^{+}+1}<1$ by Lemma~\ref{lem:tail-majorant}\,(i), so
\eqref{eq:tail-decay} gives $\Pr[\Ncount{}>l]\le r_l\Pr[\Ncount{}>l-1]
\le r_{\kappa^{+}+1}\Pr[\Ncount{}>l-1]$; iterating from $l=\kappa^{+}+1$ to
$l=\tilde k^{+}-1$,
\begin{equation*}
  \Pr[\Ncount{}>\tilde k^{+}-1]\ \le\
  r_{\kappa^{+}+1}^{\,\tilde k^{+}-1-\kappa^{+}}\,\Pr[\Ncount{}>\kappa^{+}]
  \ \le\ r_{\kappa^{+}+1}^{\,\tilde k^{+}-1-\kappa^{+}}\,\varepsilon^{+},
\end{equation*}
the last step by \eqref{eq:exact-admissible}. Taking logarithms and
combining with the previous display,
$\ln\varepsilon^{+}-w<\ln\varepsilon^{+}+(\tilde k^{+}-1-\kappa^{+})
\ln r_{\kappa^{+}+1}$, i.e.\
$(\tilde k^{+}-1-\kappa^{+})(-\ln r_{\kappa^{+}+1})<w$, which is
\eqref{eq:bracket-conservatism} since $-\ln r_{\kappa^{+}+1}$ is positive and
finite: $r_{\kappa^{+}+1}<1$ by hypothesis, and $r_{\kappa^{+}+1}>0$ by
Lemma~\ref{lem:tail-majorant}\,(i) because $\kappa^{+}+1\le n-1$. The mirror
statement is the same argument for the upper tail of $\Ncount{}'$ at level
$\varepsilon^{-}$, translated by $k\mapsto n-k$: the exact lower quantile of
$\Ncount{}$ is $n$ minus the exact upper quantile of $\Ncount{}'$, because
$\Pr[\Ncount{}<k]=\Pr[\Ncount{}'>n-k]$ turns the maximum in
\eqref{eq:kappaminus} into the minimum in \eqref{eq:kappaplus} for
$\Ncount{}'$.

\emph{The closing sentences of (iii).} By \eqref{eq:guard} the guarded upper
endpoint is $\max\{\tilde k^{+},\lceil np\rceil\}$ for this evaluation and
$\max\{\kappa^{+},\lceil np\rceil\}$ for the exact rule, and by
\eqref{eq:boxbounds} the box bound is the guarded endpoint divided by $n$.
Where the bracket is decided, (ii) gives $\tilde k^{+}=\kappa^{+}$, so the two
guarded endpoints and the two bounds coincide. Otherwise, the map
$x\mapsto\max\{x,\lceil np\rceil\}$ is non-decreasing and changes by at most
$x'-x$ between $x\le x'$, so with $\kappa^{+}\le\tilde k^{+}$ from (i) the
guarded endpoint of this evaluation exceeds the exact rule's by at least $0$
and at most $\tilde k^{+}-\kappa^{+}$, and dividing by $n$ gives the bound on
$q^{+}$; \eqref{eq:bracket-conservatism} bounds $\tilde k^{+}-\kappa^{+}$
under the hypotheses of (iii). At $\kappa^{+}=n-1$, $\tilde k^{+}\le n$ gives
$\tilde k^{+}-\kappa^{+}\le1$. The lower side is the mirror, with
$x\mapsto\min\{x,\lfloor np\rfloor\}$ and $\tilde k^{-}\le\kappa^{-}$.
\end{proof}

% claim statistics/binomial-box/rem:local-evaluation
\claimsummary{statistics/binomial-box/rem:local-evaluation}{\hyperlink{claim-statistics/binomial-box/prop:bracketed-endpoints}{Proposition~\ref*{prop:bracketed-endpoints}} holds for every search, window and enclosure, so these are choices outside the proposition. A bisection finds a bracketed endpoint in about log\ensuremath{_{2}} n evaluations, and the window trades evaluation time for the last counts of the endpoint.}
\claimbookmark{statistics/binomial-box/rem:local-evaluation}{Remark}{rem:local-evaluation}{Evaluating the exact rule locally}
\begin{remark}[Evaluating the exact rule locally]\label{rem:local-evaluation}
Proposition~\ref{prop:bracketed-endpoints} holds for every bracketed
endpoint, whatever procedure found it, and for every enclosure of the form
\eqref{eq:tail-enclosure}. Four choices therefore lie outside the
proposition: the search, the window, the arithmetic substrate and the report
of the bracket verdict. Requirements~\ref{req:search-procedure},
\ref{req:enclosure-substrate} and~\ref{req:endpoint-exact-record} fix the
engine's search, substrate and report.

\emph{The search.} Certified admissibility $\mathsf{C}^{+}$ is a predicate
on $\{0,\dots,n\}$ that holds at $n$; a bisection that keeps
$\mathsf{C}^{+}$ true at its upper end and false at its lower end (starting
from $n$ and from the fictitious index $-1$), treating an undecided index as
not certified, returns a bracketed upper endpoint after about $\log_2 n$
evaluations of the enclosure. A downward scan from any index known to be
certified is the alternative; the endpoint of the Bernstein rule
(Proposition~\ref{prop:rule-bernstein}) is admissible and usually certified,
and is a natural start. Fixing the procedure, with its parameters, makes the
evaluation a rule in the sense of Definition~\ref{def:endpoint-rule};
Requirement~\ref{req:search-procedure} fixes it. Where the map
$k\mapsto\tau^{+}(k)$ is non-increasing (which it is up to the rounding of
the enclosure, since the tail itself decreases by the factor $r_k$ per count,
Lemma~\ref{lem:tail-majorant}\,(iv), and the enclosure's width varies
slowly), every bracketed endpoint is the least certified index, and the
evaluation is then \termref{term-monotone-in-the-allocation}{monotone in the allocation} as well (the set of certified
indices grows with $\varepsilon^{+}$ when the logarithm of
$\varepsilon^{+}$ is enclosed monotonically, as the fixed-grid root bounds of
Proposition~\ref{prop:rule-bernstein}\,(ii) are); the claims do not rest on
this.

\emph{The window.} In \eqref{eq:tail-log} the bracket's width is
$\ln(1+B_T/A_T)$, and $B_T/A_T$ is about $r^{T}/(1-r)$ with $r$ the ratio
near $k$; at a deviation of $t$ counts above the mean, $1-r_k$ is of the
order of $t/(np(1-p))$, so a window of $T$ of the order of
$(\ln\frac{1}{1-r}+c)/(1-r)$ terms leaves a relative slack of about
$e^{-c}$ in the tail, that is, about $e^{-c}$ nats in $w$. Together with the
per-count decay $-\ln r_{\kappa^{+}+1}$, which is of the same order
$t/(np(1-p))$, \eqref{eq:bracket-conservatism} says what the last counts
cost in evaluation time: the window trades time for the last counts, and
the exact quantile is reached exactly when the bracket is decided.

\emph{The substrate.} The Stirling remainders and the root bounds are
rationals of the size of a logarithm, while the products $\prod r_l$ inside
the rationals $A_T$ and $B_T$ of \eqref{eq:window-rationals}, held as exact
rationals, would have denominators that grow with the window. Every bound is
used in the direction a \termref{term-conservative-evaluation}{conservative evaluation}
(Definition~\ref{def:conservative-evaluation}) admits: a tail comparison is
decided against a bound that rounds the tail upward.
Requirement~\ref{req:enclosure-substrate} states the arithmetic of the
evaluation.

\emph{The verdict.} Whether the bracket was decided is
Proposition~\ref{prop:bracketed-endpoints}\,(ii)'s condition. The bound
\eqref{eq:bracket-conservatism} is not computable as it stands: $w$ involves
the unknown $\ln\Pr[\Ncount{}>\tilde k^{+}-1]$ and $r_{\kappa^{+}+1}$ the
unknown $\kappa^{+}$. With a certified-inadmissible index $k_{\mathsf{D}}$
below the bracket in hand and $0<r_{k_{\mathsf{D}}+2}<1$, and the enclosure
in log space, a computable bound on the excess $\tilde k^{+}-\kappa^{+}$ is
\eqref{eq:bracket-conservatism} with $-\ln r_{k_{\mathsf{D}}+2}$ in place of
$-\ln r_{\kappa^{+}+1}$ and
$W\coloneqq U(\tilde k^{+}-1)-\Lambda(\tilde k^{+}-1)+(\lambda^{+}-\lambda^{-})$
in place of $w$. Indeed $W\ge w$, because
$\Lambda(\tilde k^{+}-1)\le\ln\Pr[\Ncount{}>\tilde k^{+}-1]$ and
$\lambda^{+}\ge\ln\varepsilon^{+}$; and since $k_{\mathsf{D}}$ is certified
inadmissible, $\kappa^{+}\ge k_{\mathsf{D}}+1$, so the decreasing ratios give
$r_{\kappa^{+}+1}\le r_{k_{\mathsf{D}}+2}<1$ and
$-\ln r_{\kappa^{+}+1}\ge-\ln r_{k_{\mathsf{D}}+2}>0$. Where
$\tilde k^{+}>\kappa^{+}$ and $\kappa^{+}\le n-2$, the hypotheses of
Proposition~\ref{prop:bracketed-endpoints}\,(iii) then hold and the
substituted bound is at least its right-hand side; where
$\tilde k^{+}>\kappa^{+}=n-1$, the excess is $1$, below the substituted bound
because $W\ge w>0$, the certificate at $\tilde k^{+}-1\ge\kappa^{+}$ having
failed although that tail is at most $\varepsilon^{+}$; and where
$\tilde k^{+}=\kappa^{+}$ the excess is $0$.
Requirement~\ref{req:endpoint-exact-record} says what is recorded.
Completeness on this evaluation is the figure of
Remark~\ref{rem:completeness-on-rules}\,(b), from the upper enclosures
$\tau^{+}$ at the emitted guarded endpoints, clamped by $\ecomp$
(Requirement~\ref{req:completeness-report}); on the sweep evaluation no
bracket is formed: the exact rule is evaluated by the closed forms of
Definition~\ref{def:interval}, its tail comparisons decided exactly or
against the upper bounds of Proposition~\ref{prop:rule-exact}\,(vii), and
none of this arises.
\end{remark}
\claimmeta{\metalabel{Uses}\mref{Def.}{def:conservative-evaluation}, \mref{Def.}{def:endpoint-rule}, \mref{Def.}{def:interval}, \mref{Lem.}{lem:log-enclosure}, \mref{Lem.}{lem:stirling-enclosure}, \mref{Rem.}{rem:completeness-on-rules}, \mref{Lem.}{lem:pmf-enclosure}, \mref{Lem.}{lem:tail-majorant}, \mref{Prop.}{prop:bracketed-endpoints}
\metasep \metalabel{Used by}\mref{Req.}{req:search-procedure}, \mref{Req.}{req:enclosure-substrate}
\metasep newly derived
\metasep accepted}

\section{Requirements: the request, its refusal and the reports}\label{group-statistics/binomial-box/req:rule-in-force}
\sectionintro{The engine builds the box with the endpoint rule and tail evaluation the request selects, and takes its defaults where the request is silent. Before constructing anything it refuses an input that violates a clause its own inputs decide, reports the completeness figure its route names and never a downward-rounded one, and records beside it the input, the rule and the evaluation it was computed under.}

% claim statistics/binomial-box/req:rule-in-force
\claimsummary{statistics/binomial-box/req:rule-in-force}{The engine builds the box with the endpoint rule and tail evaluation the request selects: the exact rule evaluated by its closed forms, which is the default, the enclosure evaluation of the exact rule, or the Bernstein rule. It runs no other rule.}
\claimbookmark{statistics/binomial-box/req:rule-in-force}{Requirement}{req:rule-in-force}{The rule the request selects}
\claimleadin{%
Each rule of this submodule is an \termref{term-admissible-endpoint-rule}{admissible endpoint rule}, so every
result stated for all such rules holds for the box the engine builds,
whichever the request chooses; a result that needs more of the rule, such as
monotonicity in the allocation, holds on the rules whose propositions prove
it.
}
\begin{requirement}[The rule the request selects]\label{req:rule-in-force}
The engine builds the box of Definition~\ref{def:construction} with the rule
that the request's \iface{endpoint_rule} and \iface{tail_evaluation} select,
and with no other: \iface{endpoint_rule: exact} with
\iface{tail_evaluation: sweep}, the \termref{term-exact-binomial-quantile-rule}{exact binomial quantile rule} of
Definition~\ref{def:interval} with its tail comparisons decided on the
closed forms, exactly or against upper bounds on the tails as
Proposition~\ref{prop:rule-exact}\,(vii) admits, which is the default; \iface{endpoint_rule: exact} with
\iface{tail_evaluation: enclosure}, the enclosure evaluation
$R_{\mathrm{encl}}$ of the exact rule
(Proposition~\ref{prop:bracketed-endpoints}), an admissible relaxation of it;
and \iface{endpoint_rule: bernstein}, the Bernstein rule of
Proposition~\ref{prop:rule-bernstein}, on which \iface{tail_evaluation}
selects only how the completeness figure is computed
(Requirement~\ref{req:completeness-report}).
\end{requirement}
\claimmeta{\metalabel{Motivated by}\mref{Def.}{def:endpoint-rule}, \mref{Def.}{def:interval}, \mref{Prop.}{prop:rule-exact}, \mref{Def.}{def:construction}, \mref{Prop.}{prop:rule-bernstein}, \mref{Prop.}{prop:bracketed-endpoints}, \mref{Req.}{req:completeness-report}
\metasep implemented (effective: accepted)
\metasep \metalabel{Code}\coderef{https://github.com/fraunhofersg/certiqs-key/blob/v0.1.0/src/certiqs_key/statistics/rules.py}{certiqs_key.statistics.rules:resolve_rule}}

% claim statistics/binomial-box/req:request-defaults
\claimsummary{statistics/binomial-box/req:request-defaults}{Where a request names no constrained outcomes the engine constrains every outcome, and where it names no allocation it takes the uniform rule, which gives each tail of each constrained outcome the budget divided by twice the number of constrained outcomes.}
\claimbookmark{statistics/binomial-box/req:request-defaults}{Requirement}{req:request-defaults}{Defaults for the constrained outcomes and the allocation}
\claimleadin{%
No choice of the constrained outcomes dominates (Remark~\ref{rem:selector}),
and the \termref{term-uniform-rule}{uniform rule} is a \termref{term-monotone-allocation-rule}{monotone allocation rule}
(Definition~\ref{def:monotone-rule}); the engine fixes both only as defaults.
}
\begin{requirement}[Defaults for the constrained outcomes and the allocation]\label{req:request-defaults}
Where a request omits the constrained outcomes, the engine takes
$\Chatcon=\Chatset$ (\iface{constrained_outcomes}); where it omits the
allocation, the engine takes the uniform rule of
Definition~\ref{def:monotone-rule} at the request's budget,
$\varepsilon^{\pm}(\hat c)=\ecomp/(2|\Chatcon|)$ for every
$\hat c\in\Chatcon$ (\iface{budget_allocation}).
\end{requirement}
\claimmeta{\metalabel{Motivated by}\mref{Rem.}{rem:selector}, \mref{Def.}{def:monotone-rule}
\metasep implemented (effective: accepted)
\metasep \metalabel{Code}\coderef{https://github.com/fraunhofersg/certiqs-key/blob/v0.1.0/src/certiqs_key/statistics/construction.py}{certiqs_key.statistics.construction:resolve}}

% claim statistics/binomial-box/req:input-refusal
\claimsummary{statistics/binomial-box/req:input-refusal}{Before constructing anything the engine decides every condition of an admissible input its inputs determine; on a violation it reports the first by its code and emits no acceptance set, repairing nothing. Data not of the input's form, with the request defaults in place, are refused too.}
\claimbookmark{statistics/binomial-box/req:input-refusal}{Requirement}{req:input-refusal}{Refusal of an input that violates a decidable clause}
\claimleadin{%
By Theorem~\ref{thm:honest-containment}\,(iii) the construction never fails on an
\termref{term-admissible-input}{admissible input}, so every refusal comes from the input; this is
\xref{foundations/req:consumer-decides-exactly} for the module's own inputs.
The clause relating the alphabet and the round count to the consuming
declaration is not among those decided here: it is checked in bundle
validation
(\href{../../../spec/architecture.md\#36-bundle-validation-is-input-validation-not-a-module}{\nolinkurl{spec/architecture.md}\texttt{\#36-bundle-validation-is-input-validation-not-a-module}}).
}
\begin{requirement}[Refusal of an input that violates a decidable clause]\label{req:input-refusal}
Before it constructs anything, the engine decides exactly every condition of
an admissible input (Definition~\ref{def:endpoint-rule}) that its own inputs
determine on data of the form that definition fixes. On a violation it emits
\iface{infeasibility_report} naming the violated condition by its code, the
first violation in the order the interface fixes, and it emits no
\iface{acceptance_set}; it repairs no input and puts no default in place of a
violating value. The codes of the violations (1)--(11) of
Theorem~\ref{thm:honest-containment}\,(iii) are, in that order,
\iface{alphabet_empty}, \iface{alphabet_not_finite},
\iface{rounds_out_of_range}, \iface{honest_not_a_distribution},
\iface{honest_envelope_not_admissible}, \iface{budget_out_of_range},
\iface{allocation_negative}, \iface{allocation_exceeds_budget},
\iface{selector_not_a_subset}, \iface{inputs_not_rational} and
\iface{endpoint_rule_unknown}. Data that, with the defaults of
Requirements~\ref{req:request-defaults} and~\ref{req:rule-in-force} in place,
do not have that form are refused as well: the engine emits no
\iface{acceptance_set} for them and substitutes nothing for them.
\end{requirement}
\claimmeta{\metalabel{Motivated by}\mref{Def.}{def:endpoint-rule}, \mref{Thm.}{thm:honest-containment}, \mref{Req.}{req:rule-in-force}, \mref{Req.}{req:request-defaults}, \extref{../foundations/contract.pdf}{foundations/req:consumer-decides-exactly}
\metasep implemented (effective: accepted)
\metasep \metalabel{Code}\coderef{https://github.com/fraunhofersg/certiqs-key/blob/v0.1.0/src/certiqs_key/statistics/construction.py}{certiqs_key.statistics.construction:validate_request}}

% claim statistics/binomial-box/req:completeness-report
\claimsummary{statistics/binomial-box/req:completeness-report}{The engine reports \ensuremath{\varepsilon}* itself on the exact rule evaluated by its closed forms, and otherwise the figure its route's row names: the allocation sum or an upward-rounded evaluation of the tails clamped at the budget. It never reports a downward-rounded figure.}
\claimbookmark{statistics/binomial-box/req:completeness-report}{Requirement}{req:completeness-report}{The reported completeness figure}
\claimleadin{%
A completeness figure bounds how often an honest run aborts, so it may be
loose but never low: it is rounded in the direction in which
\xref{foundations/req:rounding-direction} rounds an emitted upper bound.
}
\begin{requirement}[The reported completeness figure]\label{req:completeness-report}
The engine reports under \iface{completeness_certified} the following figure.
On \iface{endpoint_rule: exact} with \iface{tail_evaluation: sweep} it
reports $\varepsilon^{\star}$ of Theorem~\ref{thm:completeness-bound} itself, and where it
cannot evaluate a tail exactly it rounds the tail upward and reports
$\min\{\ecomp,\tilde\varepsilon\}$ for the resulting
$\tilde\varepsilon\ge\varepsilon^{\star}$ (Remark~\ref{rem:epsstar}). On
every other route it reports the admissible figure of
Remark~\ref{rem:completeness-on-rules} that the route's row of
\iface{completeness_report} names: the allocation sum, or
$\min\{\ecomp,\tilde\varepsilon\}$ for an upward-rounded evaluation
$\tilde\varepsilon\ge\varepsilon^{\star}$ of the tails at the emitted
endpoints. It never reports a downward-rounded figure.
\end{requirement}
\claimmeta{\metalabel{Motivated by}\mref{Thm.}{thm:completeness-bound}, \mref{Rem.}{rem:epsstar}, \mref{Rem.}{rem:completeness-on-rules}, \extref{../foundations/contract.pdf}{foundations/req:rounding-direction}
\metasep implemented (effective: accepted)
\metasep \metalabel{Code}\coderef{https://github.com/fraunhofersg/certiqs-key/blob/v0.1.0/src/certiqs_key/statistics/construction.py}{certiqs_key.statistics.construction:report_completeness}}

% claim statistics/binomial-box/req:honest-model-record
\claimsummary{statistics/binomial-box/req:honest-model-record}{With every completeness figure it reports, the engine records what the figure was computed from: the honest model, n, the constrained outcomes, the allocation, the endpoint rule, the tail evaluation and the parameters in force. A completeness figure is meaningful only relative to that record.}
\claimbookmark{statistics/binomial-box/req:honest-model-record}{Requirement}{req:honest-model-record}{The honest-model record beside every completeness figure}
\claimleadin{%
The bound of Theorem~\ref{thm:completeness-bound} holds for every
\termref{term-admissible-endpoint-rule}{admissible endpoint rule},
but the figure reported for it depends on the input, on the rule and on how
its tails were evaluated (Remark~\ref{rem:completeness-on-rules}).
}
\begin{requirement}[The honest-model record beside every completeness figure]\label{req:honest-model-record}
With every completeness figure it reports, the engine records, in the record
\iface{honest_model}, the input and the evaluation under which the figure was
computed: the honest model (the law in the \termref{term-point-form}{point form}, or the
\termref{term-honest-envelope}{honest envelope}), the round count $n$, the constrained outcomes
$\Chatcon$ and the allocation; the endpoint rule and the tail evaluation; and
the parameters in force, for the Bernstein rule its $j$ and $s$
(Proposition~\ref{prop:rule-bernstein}) and, wherever the tails are
evaluated by enclosures, the parameters of those enclosures
(Proposition~\ref{prop:bracketed-endpoints}). A completeness figure is
meaningful only relative to that record.
\end{requirement}
\claimmeta{\metalabel{Motivated by}\mref{Thm.}{thm:completeness-bound}, \mref{Prop.}{prop:rule-bernstein}, \mref{Rem.}{rem:completeness-on-rules}, \mref{Prop.}{prop:bracketed-endpoints}
\metasep implemented (effective: accepted)
\metasep \metalabel{Code}\coderef{https://github.com/fraunhofersg/certiqs-key/blob/v0.1.0/src/certiqs_key/statistics/construction.py}{certiqs_key.statistics.construction:HonestModel}}

\section{Requirements: the enclosure evaluation}\label{group-statistics/binomial-box/req:search-procedure}
\sectionintro{On the enclosure evaluation of the exact rule the engine finds every bracketed endpoint by one fixed, total and deterministic search, and computes in log space with directed rounding, never forming a number of the size of a mass point. It records, per constrained outcome and side, whether the bracket was decided.}

% claim statistics/binomial-box/req:search-procedure
\claimsummary{statistics/binomial-box/req:search-procedure}{On the enclosure evaluation of the exact rule the engine finds every bracketed endpoint by one fixed, total and deterministic search with fixed parameters, the procedure the interface records. Fixing it is what makes the evaluation an endpoint rule.}
\claimbookmark{statistics/binomial-box/req:search-procedure}{Requirement}{req:search-procedure}{One fixed search for the enclosure evaluation}
\claimleadin{%
Proposition~\ref{prop:bracketed-endpoints} holds for every bracketed
endpoint, whatever procedure found it; fixing the procedure, with its
parameters, is what makes the evaluation a rule in the sense of
Definition~\ref{def:endpoint-rule} (Remark~\ref{rem:local-evaluation}).
}
\begin{requirement}[One fixed search for the enclosure evaluation]\label{req:search-procedure}
On the enclosure evaluation of the exact rule, the engine finds every
bracketed endpoint of Proposition~\ref{prop:bracketed-endpoints}, upper and
lower, by one fixed search procedure, total and deterministic, with fixed
parameters: the procedure \iface{search_procedure} records.
\end{requirement}
\claimmeta{\metalabel{Motivated by}\mref{Def.}{def:endpoint-rule}, \mref{Prop.}{prop:bracketed-endpoints}, \mref{Rem.}{rem:local-evaluation}
\metasep implemented (effective: accepted)
\metasep \metalabel{Code}\coderef{https://github.com/fraunhofersg/certiqs-key/blob/v0.1.0/src/certiqs_key/statistics/enclosure.py}{certiqs_key.statistics.enclosure:bracketed_upper_endpoint}}

% claim statistics/binomial-box/req:enclosure-substrate
\claimsummary{statistics/binomial-box/req:enclosure-substrate}{The enclosure evaluation computes in log space with directed rounding: the window's products are dyadic numbers rounded in the direction each bound requires, and the only numbers formed have the magnitude of logarithms. No mass point, power of p or binomial coefficient is ever formed.}
\claimbookmark{statistics/binomial-box/req:enclosure-substrate}{Requirement}{req:enclosure-substrate}{Log-space arithmetic with directed rounding}
\claimleadin{%
A tail comparison of the enclosure evaluation is decided against a bound that
rounds the tail upward, which is what a \termref{term-conservative-evaluation}{conservative evaluation}
(Definition~\ref{def:conservative-evaluation}) admits;
exact rationals for the window would have denominators that grow with the
window (Remark~\ref{rem:local-evaluation}).
}
\begin{requirement}[Log-space arithmetic with directed rounding]\label{req:enclosure-substrate}
On the enclosure evaluation the engine computes everything in log space with
directed rounding: the mass point through \eqref{eq:pmf-enclosure}, the window
through the rationals $A_T$ and $B_T$ of \eqref{eq:window-rationals}, whose
logarithms Lemma~\ref{lem:log-enclosure} encloses, and the budget through the
same lemma. It holds the products $\prod r_l$ inside $A_T$ and $B_T$ as
dyadic numbers with an explicit binary exponent, rounded in the direction
each bound requires (down for $A_T$ in the lower bound, up for $A_T+B_T$ in
the upper), not as exact rationals, and the Stirling remainders and the root
bounds as rationals of the size of a logarithm; the only numbers it ever
forms have the magnitude of logarithms.
\end{requirement}
\claimmeta{\metalabel{Motivated by}\mref{Def.}{def:conservative-evaluation}, \mref{Lem.}{lem:log-enclosure}, \mref{Lem.}{lem:pmf-enclosure}, \mref{Lem.}{lem:tail-majorant}, \mref{Rem.}{rem:local-evaluation}
\metasep implemented (effective: accepted)
\metasep \metalabel{Code}\coderef{https://github.com/fraunhofersg/certiqs-key/blob/v0.1.0/src/certiqs_key/statistics/dyadic.py}{certiqs_key.statistics.dyadic:round_dyadic}}

% claim statistics/binomial-box/req:endpoint-exact-record
\claimsummary{statistics/binomial-box/req:endpoint-exact-record}{On the enclosure evaluation the engine records, per constrained outcome and side, whether the bracket was decided, that is whether the emitted endpoint is proved equal to the exact quantile. It may add the bound on the conservatism in counts.}
\claimbookmark{statistics/binomial-box/req:endpoint-exact-record}{Requirement}{req:endpoint-exact-record}{The bracket verdict per outcome and side}
\claimleadin{%
The enclosure evaluation may emit an endpoint more conservative than the
exact quantile, and nothing in the emitted box shows it; the bracket verdict
of Proposition~\ref{prop:bracketed-endpoints}\,(ii) does.
}
\begin{requirement}[The bracket verdict per outcome and side]\label{req:endpoint-exact-record}
On the enclosure evaluation of the exact rule the engine records in the
diagnostic \iface{endpoint_exact}, per constrained outcome and per side,
whether the bracket was decided, that is whether the emitted endpoint is
proved equal to the exact quantile by
Proposition~\ref{prop:bracketed-endpoints}\,(ii). When it holds a
certified-inadmissible index $k_{\mathsf{D}}$ below the bracket with
$0<r_{k_{\mathsf{D}}+2}<1$, it may also report the bound
\eqref{eq:bracket-conservatism} in counts with $-\ln r_{k_{\mathsf{D}}+2}$ in
place of $-\ln r_{\kappa^{+}+1}$ and the computable upper bound
$U(\tilde k^{+}-1)-\Lambda(\tilde k^{+}-1)+(\lambda^{+}-\lambda^{-})$ on $w$
in place of $w$ (Remark~\ref{rem:local-evaluation}).
\end{requirement}
\claimmeta{\metalabel{Motivated by}\mref{Prop.}{prop:bracketed-endpoints}, \mref{Rem.}{rem:local-evaluation}
\metasep implemented (effective: accepted)
\metasep \metalabel{Code}\coderef{https://github.com/fraunhofersg/certiqs-key/blob/v0.1.0/src/certiqs_key/statistics/construction.py}{certiqs_key.statistics.construction:EndpointExact}}

\clearpage
\part{Singleton acceptance set}
\sectionintro{This submodule treats the acceptance test that accepts exactly one declared frequency vector. Its main result is that under independent honest rounds with one law the test's accept probability is the multinomial probability of that vector (\hyperlink{claim-statistics/singleton/thm:singleton-completeness}{Theorem~\ref*{thm:singleton-completeness}}). A vector off the count grid, or one that puts weight on an outcome the honest law never produces, is never accepted, and then no completeness budget below one exists; otherwise the least budget is one minus the accept probability. That budget is zero only when the honest law and the vector are the same point mass, and at any realistic alphabet and round count it is close to one, which is why the construction map never emits a singleton. An instance that declares one records its least budget, or that none exists, and the singleton reads the honest model in point form only.}

\section{Completeness of a singleton acceptance set}\label{group-statistics/singleton/thm:singleton-completeness}
\sectionintro{A test that accepts exactly one \termref{term-frequency-vector}{frequency vector} is far more restrictive than a box around the honest law. The main result, \hyperlink{claim-statistics/singleton/thm:singleton-completeness}{Theorem~\ref*{thm:singleton-completeness}}, computes its accept probability from the multinomial law, says when no completeness budget below one exists and when the least budget is zero; the remark explains why the \termref{term-construction-map}{construction map} never emits a singleton.}

% claim statistics/singleton/thm:singleton-completeness
\claimsummary{statistics/singleton/thm:singleton-completeness}{Under i.i.d. honest rounds a singleton acceptance set accepts with its point's multinomial probability. That is 0, with no completeness budget below one, if n times a coordinate is no integer or the point leaves the law's support; else the least budget is one minus it, zero only for one point mass.}
\claimbookmark{statistics/singleton/thm:singleton-completeness}{Theorem}{thm:singleton-completeness}{Completeness of a singleton acceptance set}
\claimleadin{%
A test that accepts exactly one \termref{term-frequency-vector}{frequency vector} is far more
restrictive than a box around the honest law. Under independent honest
rounds with one law its accept probability is a single multinomial mass, so
its completeness can be computed exactly; an \exttermref{../foundations/contract.pdf}{term-instance}{instance} that declares such
a set records the least budget of (ii), or that none exists
(Requirement~\ref{req:singleton-budget-record}).
}
\begin{mainclaim}
\begin{theorem}[Completeness of a singleton acceptance set]\label{thm:singleton-completeness}
Assume A-HON, let $\Chatset$ and $n$ be a
\termref{term-matched-alphabet-and-round-count}{matched alphabet and round count}
(Definition~\ref{def:matched-alphabet}), and let
$\acSet=\{\mathbf{q}^{\ast}\}$ be a singleton geometry with
$\mathbf{q}^{\ast}\in\simplex{\Chatset}$. Call a completeness
budget $\ecomp$ \emph{admissible for this geometry} when $\ecomp\in[0,1)$ and
$\Pr[\operatorname{freq}_{\Chat_1^n}\notin\acSet]\le\ecomp$.
\begin{enumerate}[label=(\roman*),leftmargin=2.2em]
\item \emph{(exact accept probability)} With the convention $0^{0}=1$,
\begin{equation}\label{eq:singleton-acc}
  \Pr\big[\operatorname{freq}_{\Chat_1^n}\in\acSet\big]
  =\begin{dcases}
     \frac{n!}{\prod_{\hat c\in\Chatset}\big(n\,q^{\ast}(\hat c)\big)!}
     \prod_{\hat c\in\Chatset}\honest(\hat c)^{\,n\,q^{\ast}(\hat c)},
     &\text{if }n\,q^{\ast}(\hat c)\in\mathbb{Z}_{\ge0}\ \forall\hat c,\\
     0,&\text{otherwise;}
   \end{dcases}
\end{equation}
\item \emph{(the least admissible budget, and when there is none)} the accept
probability \eqref{eq:singleton-acc} is $0$ if and only if either
$n\,q^{\ast}(\hat c)\notin\mathbb{Z}_{\ge0}$ for some $\hat c$, or
$\mathbf{q}^{\ast}\not\ll\honest$, that is, some $\hat c$ has
$q^{\ast}(\hat c)>0$ and $\honest(\hat c)=0$. In that case \emph{no}
$\ecomp\in[0,1)$ is admissible for this geometry, and no completeness
statement whatsoever can be made about it; a budget of $1$ lies outside
$[0,1)$. Otherwise the accept probability is strictly positive and the least
admissible completeness budget is
\begin{equation}\label{eq:singleton-budget}
  \ecomp=1-\Pr\big[\operatorname{freq}_{\Chat_1^n}\in\acSet\big]\ \in[0,1),
\end{equation}
which does lie in $[0,1)$;
\item \emph{(when the budget is zero)} the budget \eqref{eq:singleton-budget}
equals $0$ if and only if $\honest$ is a point mass $\delta_{\hat c_0}$ for
some $\hat c_0\in\Chatset$ and $\mathbf{q}^{\ast}=\delta_{\hat c_0}$; in
particular it equals $0$ whenever $|\Chatset|=1$, and it is strictly positive
whenever $\honest$ gives positive mass to two distinct outcomes.
\end{enumerate}
\end{theorem}
\end{mainclaim}
\claimmeta{\metalabel{Uses}\mref{Def.}{def:freq}, \mref{Def.}{def:matched-alphabet}
\metasep \metalabel{Assumes}\aref{A-HON}{ass:hon}, \aref{A-ALPHABET}{ass:alphabet}
\metasep \metalabel{Used by}\mref{Rem.}{rem:singleton-honest}, \mref{Req.}{req:singleton-budget-record}, \mref{Req.}{req:singleton-point-form}
\metasep newly derived
\metasep implemented (effective: accepted)
\metasep \metalabel{Code}\coderef{https://github.com/fraunhofersg/certiqs-key/blob/v0.1.0/src/certiqs_key/statistics/singleton.py}{certiqs_key.statistics.singleton:singleton_completeness}
\metasep \metalabel{Oracles}\coderef{https://github.com/fraunhofersg/certiqs-key/blob/v0.1.0/tests/oracles/registry_statistics_acceptance.py}{registry_statistics_acceptance.py}}

\begin{proof}
\emph{(i).} By Definition~\ref{def:freq}, $\operatorname{freq}_{\Chat_1^n}=
\mathbf{q}^{\ast}$ holds exactly when $\Ncount{\hat c}=n\,q^{\ast}(\hat c)$ for
every $\hat c$, which is impossible unless every $n\,q^{\ast}(\hat c)$ is a
non-negative integer; that gives the second branch. In the first branch write
$n_{\hat c}\coloneqq n\,q^{\ast}(\hat c)$, so $\sum_{\hat c}n_{\hat c}=n$. Under
A-HON each string $\hat c_1^n$ has probability
$\prod_{\hat c}\honest(\hat c)^{\,n_{\hat c}}$ whenever its count vector is
$(n_{\hat c})_{\hat c}$, and the number of such strings is the multinomial
coefficient $n!/\prod_{\hat c}n_{\hat c}!$; summing over them gives
\eqref{eq:singleton-acc}.

\emph{(ii).} If some $n\,q^{\ast}(\hat c)$ is not a non-negative integer the
second branch of \eqref{eq:singleton-acc} gives probability $0$. Otherwise the
first branch applies; its multinomial coefficient is a positive integer, and
under the convention $0^{0}=1$ each factor $\honest(\hat c)^{\,n_{\hat c}}$
with $n_{\hat c}=0$ equals $1$. The product therefore vanishes if and only if
some $\hat c$ has $n_{\hat c}\ge1$ and $\honest(\hat c)=0$; and $n_{\hat c}\ge1$
is equivalent to $q^{\ast}(\hat c)>0$, so this says precisely that
$\mathbf{q}^{\ast}\not\ll\honest$. That establishes the stated
characterization of a zero accept probability.

By the definition in the statement, a budget $\ecomp$ is admissible for this
geometry exactly when it satisfies both the completeness requirement
$\Pr[\operatorname{freq}_{\Chat_1^n}\notin\acSet]\le\ecomp$, that is,
$\ecomp\ge1-\Pr[\operatorname{freq}_{\Chat_1^n}\in\acSet]$, and the domain
$\ecomp\in[0,1)$ of the completeness budget. The admissible budgets
are therefore the interval
$[\,1-\Pr[\operatorname{freq}_{\Chat_1^n}\in\acSet],\,1)$, which is non-empty
if and only if $\Pr[\operatorname{freq}_{\Chat_1^n}\in\acSet]>0$, and which
then has least element \eqref{eq:singleton-budget}. When the accept
probability is $0$ the interval is empty: there is no admissible budget, and
in particular $\ecomp=1$ is not one.

\emph{(iii).} The budget \eqref{eq:singleton-budget} is $0$ exactly when the
accept probability is $1$. That probability equals $1$ if and only if
$\operatorname{freq}_{\Chat_1^n}=\mathbf{q}^{\ast}$ almost surely. Suppose
$\honest$ gives positive mass to two distinct outcomes $\hat c\neq\hat c'$.
The string $(\hat c,\hat c,\dots,\hat c)$ and the string
$(\hat c',\hat c,\dots,\hat c)$ both have positive probability under A-HON, and
their frequency vectors differ (they differ by $1/n$ in coordinate $\hat c$),
so $\operatorname{freq}_{\Chat_1^n}$ is not almost surely constant and the
probability is strictly less than $1$. Hence $\honest$ must be a point mass
$\delta_{\hat c_0}$; then $\operatorname{freq}_{\Chat_1^n}=\delta_{\hat c_0}$
almost surely and the probability is $1$ exactly when
$\mathbf{q}^{\ast}=\delta_{\hat c_0}$. If $|\Chatset|=1$ the only distribution
on $\Chatset$ is that point mass, and $\mathbf{q}^{\ast}$ equals it.
\end{proof}

% claim statistics/singleton/rem:singleton-honest
\claimsummary{statistics/singleton/rem:singleton-honest}{The consuming engine certifies a singleton's acceptance-side minimum, but its completeness figure 1 \ensuremath{-} Pr[freq = q*] is close to one at realistic alphabets and round counts, and no budget exists whenever n q* leaves the count grid. Its accept probability needs one honest law.}
\claimbookmark{statistics/singleton/rem:singleton-honest}{Remark}{rem:singleton-honest}{Certifiability and completeness of the singleton}
\begin{remark}[Certifiability and completeness of the singleton]
\label{rem:singleton-honest}
The consuming engine certifies the acceptance-side minimum of a singleton
geometry (\xref{meat-core/frank-wolfe/rem:acc-side}~(a)), so a singleton
acceptance set is certifiable. By Theorem~\ref{thm:singleton-completeness}
its completeness figure is $1-\Pr[\operatorname{freq}=\mathbf{q}^{\ast}]$,
which is close to $1$ for any realistic alphabet and round count; that is
why the \termref{term-construction-map}{construction map} never emits it. The degenerate case of
Theorem~\ref{thm:singleton-completeness}\,(ii), where the accept probability is exactly $0$
and \emph{no} admissible completeness budget exists, is not exotic: it is
reached whenever $n\,\mathbf{q}^{\ast}$ has a non-integer coordinate, which
for a fixed $\mathbf{q}^{\ast}$ happens at most round counts.
Requirement~\ref{req:singleton-budget-record} says what an instance that
declares a singleton set records.

A singleton is also the $n=2$ output of the counterexample in
Theorem~\ref{thm:monotonicity-limits}\,(v): the construction can produce a set with a single
point, and when it does, \eqref{eq:completeness} still bounds the honest
failure probability. The two statements agree, because there the budget
$\ecomp=\tfrac45$ was large enough to pay for it.

Unlike the box construction, the singleton reads its honest input in the
\termref{term-point-form}{point form} only: Theorem~\ref{thm:singleton-completeness} assumes A-HON, one law
$\honest$ for independent rounds, because its accept probability
\eqref{eq:singleton-acc} is the exact multinomial probability under that one
law and has no counterpart for an envelope of laws. A single law is the
\termref{term-degenerate-envelope}{degenerate envelope} (Definition~\ref{def:envelope});
Requirement~\ref{req:singleton-point-form} says how an envelope is read.
\end{remark}
\claimmeta{\metalabel{Uses}\mref{Def.}{def:envelope}, \mref{Thm.}{thm:completeness-bound}, \mref{Thm.}{thm:monotonicity-limits}, \mref{Thm.}{thm:singleton-completeness}
\metasep newly derived
\metasep accepted}

\section{Requirements}\label{group-statistics/singleton/req:singleton-budget-record}
\sectionintro{An instance that declares a singleton acceptance set records its least completeness budget, or that no admissible budget exists, and claims no completeness the set does not have. The singleton reads the honest model in \termref{term-point-form}{point form} only: a \termref{term-degenerate-envelope}{degenerate envelope} is read as its law, and a non-degenerate one is refused.}

% claim statistics/singleton/req:singleton-budget-record
\claimsummary{statistics/singleton/req:singleton-budget-record}{An instance that declares a singleton acceptance set records the least completeness budget of \hyperlink{claim-statistics/singleton/thm:singleton-completeness}{Theorem~\ref*{thm:singleton-completeness}}, or, where the honest run can never meet the set, that no admissible budget exists and never a budget of one. It claims no completeness the set does not have.}
\claimbookmark{statistics/singleton/req:singleton-budget-record}{Requirement}{req:singleton-budget-record}{The completeness budget of a declared singleton}
\claimleadin{%
The \termref{term-construction-map}{construction map} never emits a singleton
(Remark~\ref{rem:singleton-honest}), but an \exttermref{../foundations/contract.pdf}{term-instance}{instance} may declare one,
and Theorem~\ref{thm:singleton-completeness} gives its completeness exactly.
}
\begin{requirement}[The completeness budget of a declared singleton]\label{req:singleton-budget-record}
An instance that declares a singleton acceptance set $\{\mathbf{q}^{\ast}\}$
records the completeness budget $\ecomp$ of \eqref{eq:singleton-budget}, or,
in the degenerate case of Theorem~\ref{thm:singleton-completeness}\,(ii) where the accept
probability is exactly $0$, records that \emph{no} admissible completeness
budget exists rather than a budget of $1$, which is the honest reading of a
geometry the honest run can never satisfy. In either case it claims no
completeness the set does not have.
\end{requirement}
\claimmeta{\metalabel{Motivated by}\mref{Thm.}{thm:singleton-completeness}, \mref{Rem.}{rem:singleton-honest}
\metasep accepted}

% claim statistics/singleton/req:singleton-point-form
\claimsummary{statistics/singleton/req:singleton-point-form}{The singleton reads the honest model in the point form only: a degenerate envelope is read as its law, and a non-degenerate envelope is refused rather than reduced to one of its laws.}
\claimbookmark{statistics/singleton/req:singleton-point-form}{Requirement}{req:singleton-point-form}{Point-form honest input for the singleton}
\claimleadin{%
The accept probability of Theorem~\ref{thm:singleton-completeness} is a multinomial mass under
one law and has no counterpart for an envelope of laws
(Remark~\ref{rem:singleton-honest}).
}
\begin{requirement}[Point-form honest input for the singleton]\label{req:singleton-point-form}
The singleton takes its honest input in the \termref{term-point-form}{point form} only: the
engine reads a \termref{term-degenerate-envelope}{degenerate envelope} as its law, and refuses a
non-degenerate envelope with \iface{honest_envelope_not_degenerate} rather
than reducing it to one of its laws.
\end{requirement}
\claimmeta{\metalabel{Motivated by}\mref{Thm.}{thm:singleton-completeness}, \mref{Rem.}{rem:singleton-honest}
\metasep accepted
\metasep \metalabel{Code}\coderef{https://github.com/fraunhofersg/certiqs-key/blob/v0.1.0/src/certiqs_key/statistics/singleton.py}{certiqs_key.statistics.singleton:singleton_completeness}}
\clearpage
\bookmarksetupnext{level=part}
\section*{Interface}\addcontentsline{toc}{section}{Interface}
\sectionintro{The declared inputs and outputs of the module and its submodules: each key with its symbol, when it is present and the claims that fix it. The domains, types and refusal rules are in the interface files themselves.}
\subsection*{Inputs: Acceptance statistics (statistics)}
\noindent{\small From \href{interface.yaml}{\nolinkurl{interface.yaml}}.}\par\smallskip
{\footnotesize\begin{xltabular}{\linewidth}{@{}>{\raggedright\arraybackslash}p{4.4cm}>{\raggedright\arraybackslash}X>{\raggedright\arraybackslash}p{2.9cm}@{}}
\toprule
Key and symbol & Presence & Fixed by\\
\midrule
\endhead
\iface{acceptance_alphabet}\newline Chat & always & \mref{Def.}{def:matched-alphabet}, \mref{Req.}{req:matched-alphabet}, \mref{Def.}{def:taxonomy}, \mref{Thm.}{thm:monotonicity-limits}\\
\addlinespace
\iface{honest_distribution}\newline p\textasciicircum{}hon & always, in this form or as `honest\_envelope` (exactly one of the two) & \mref{Def.}{def:conservative-evaluation}, \mref{Ass.}{ass:hon}, \mref{Def.}{def:envelope}, \mref{Def.}{def:endpoint-rule}, \mref{Def.}{def:construction}\\
\addlinespace
\iface{honest_envelope}\newline (p\textasciicircum{}hon-, p\textasciicircum{}hon+) & instead of `honest\_distribution` (exactly one of the two) & \mref{Def.}{def:envelope}, \mref{Ass.}{ass:hon-env}, \mref{Def.}{def:conservative-evaluation}, \mref{Def.}{def:endpoint-rule}, \mref{Def.}{def:construction}, \mref{Thm.}{thm:honest-containment}, \mref{Req.}{req:input-refusal}\\
\addlinespace
\iface{rounds}\newline n & always & \mref{Def.}{def:matched-alphabet}, \mref{Ass.}{ass:alphabet}, \mref{Def.}{def:freq}, \mref{Thm.}{thm:box-construction}\\
\addlinespace
\iface{completeness_budget}\newline ecomp & always & \mref{Def.}{def:interval}, \mref{Thm.}{thm:monotonicity-limits}, \mref{Thm.}{thm:singleton-completeness}, \mref{Req.}{req:singleton-budget-record}\\
\addlinespace
\iface{constrained_outcomes}\newline Chat\_con & always; DEFAULT Chat\_con = Chat & \mref{Def.}{def:construction}, \mref{Rem.}{rem:selector}, \mref{Req.}{req:request-defaults}, \mref{Thm.}{thm:monotonicity-limits}\\
\addlinespace
\iface{budget_allocation}\newline (epsilon\textasciicircum{}-, epsilon\textasciicircum{}+) & always; DEFAULT the uniform rule epsilon\textasciicircum{}\{+-\}(chat) = ecomp / (2 \textbar{}Chat\_con\textbar{}) & \mref{Def.}{def:interval}, \mref{Def.}{def:monotone-rule}, \mref{Req.}{req:request-defaults}, \mref{Thm.}{thm:monotonicity-limits}\\
\addlinespace
\iface{endpoint_rule}\newline R & always; DEFAULT `exact` & \mref{Def.}{def:endpoint-rule}, \mref{Def.}{def:interval}, \mref{Prop.}{prop:rule-exact}, \mref{Prop.}{prop:rule-bernstein}, \mref{Req.}{req:rule-in-force}, \mref{Thm.}{thm:honest-containment}, \mref{Req.}{req:input-refusal}\\
\addlinespace
\iface{tail_evaluation} & always; DEFAULT `sweep` & \mref{Prop.}{prop:bracketed-endpoints}, \mref{Rem.}{rem:local-evaluation}, \mref{Prop.}{prop:rule-exact}, \mref{Cor.}{cor:tail-forms}, \mref{Lem.}{lem:exact-strict-tails}, \mref{Req.}{req:rule-in-force}\\
\bottomrule
\end{xltabular}}
\subsection*{Outputs: Acceptance statistics (statistics)}
\noindent{\small From \href{interface.yaml}{\nolinkurl{interface.yaml}}.}\par\smallskip
{\footnotesize\begin{xltabular}{\linewidth}{@{}>{\raggedright\arraybackslash}p{4.4cm}>{\raggedright\arraybackslash}X>{\raggedright\arraybackslash}p{2.9cm}@{}}
\toprule
Key and symbol & Presence & Fixed by\\
\midrule
\endhead
\iface{acceptance_set}\newline AcceptSet & on every run whose inputs lie in the input domain; absent -- not empty, not defaulted -- when `infeasibility\_report` is non-null & \mref{Def.}{def:taxonomy}, \mref{Req.}{req:geometry-tag}, \mref{Req.}{req:box-bounds-key}, \mref{Thm.}{thm:geometry-validation}, \mref{Def.}{def:construction}, \mref{Thm.}{thm:box-construction}, \mref{Rem.}{rem:two-tags}\\
\addlinespace
\iface{acceptance_support}\newline Chat\_acc & always & \mref{Thm.}{thm:box-oracle}, \mref{Req.}{req:acceptance-support}, \extref{../meat-core/contract.pdf}{meat-core/frank-wolfe/def:first-order-setting}\\
\addlinespace
\iface{completeness_certified}\newline epsilon\_star, or the upper bound on it the route reports & on every run that emits an acceptance\_set & \mref{Thm.}{thm:completeness-bound}, \mref{Rem.}{rem:epsstar}, \mref{Rem.}{rem:completeness-on-rules}, \mref{Req.}{req:completeness-report}\\
\addlinespace
\iface{honest_model} & on every run that emits an acceptance\_set & \mref{Req.}{req:honest-model-record}, \mref{Ass.}{ass:hon}, \mref{Ass.}{ass:hon-env}, \mref{Lem.}{lem:hon-degenerate}, \mref{Rem.}{rem:completeness-on-rules}, \mref{Rem.}{rem:local-evaluation}\\
\addlinespace
\iface{infeasibility_report} & always & \mref{Thm.}{thm:honest-containment}, \mref{Req.}{req:input-refusal}, \mref{Thm.}{thm:geometry-validation}, \mref{Thm.}{thm:box-construction}\\
\addlinespace
\iface{endpoint_exact} & present exactly under `endpoint\_rule = exact` with `tail\_evaluation = enclosure` (binomial-box/req:endpoint-exact-record). Absent under `sweep`, where no bracket is formed: the emitted endpoint is the exact quantile wherever the tails are evaluated exactly, and otherwise the endpoint of binomial-box/prop:rule-exact (vii); absent under `endpoint\_rule = bernstein`, where no bracket is formed. & \mref{Req.}{req:endpoint-exact-record}, \mref{Prop.}{prop:bracketed-endpoints}, \mref{Rem.}{rem:local-evaluation}\\
\addlinespace
\iface{assumption_ledger}\newline A-... & on every result that carries an acceptance\_set; absent with it when `infeasibility\_report` is non-null & \mref{Ass.}{ass:hon}, \mref{Ass.}{ass:hon-env}, \mref{Ass.}{ass:alphabet}, \mref{Ass.}{ass:reach}\\
\bottomrule
\end{xltabular}}
\clearpage
\bookmarksetupnext{level=part}
\section*{Claim index}\addcontentsline{toc}{section}{Claim index}
{\small\begin{xltabular}{\linewidth}{@{}>{\raggedright\arraybackslash}p{2.7cm}>{\raggedright\arraybackslash}Xr@{}}
\toprule
Claim & Title and identifier & Page\\
\midrule
\endhead
\mref{Definition}{def:freq} & Frequency vector and counts\newline{\footnotesize\color{metaGrey}\claimid{statistics/def:freq}} & \pageref{def:freq}\\
\addlinespace[2pt]
\mref{Lemma}{lem:freq-in-simplex} & Frequency vectors in the probability simplex\newline{\footnotesize\color{metaGrey}\claimid{statistics/lem:freq-in-simplex}} & \pageref{lem:freq-in-simplex}\\
\addlinespace[2pt]
\mref{Definition}{def:matched-alphabet} & Alphabet and round count matched to the consuming declaration\newline{\footnotesize\color{metaGrey}\claimid{statistics/def:matched-alphabet}} & \pageref{def:matched-alphabet}\\
\addlinespace[2pt]
\mref{Definition}{def:taxonomy} & Admissible geometries and the canonical box form\newline{\footnotesize\color{metaGrey}\claimid{statistics/def:taxonomy}} & \pageref{def:taxonomy}\\
\addlinespace[2pt]
\mref{Theorem}{thm:geometry-validation} & Validity of the admissible geometries\newline{\footnotesize\color{metaGrey}\claimid{statistics/thm:geometry-validation}} & \pageref{thm:geometry-validation}\\
\addlinespace[2pt]
\mref{Remark}{rem:two-tags} & The two tags and the consumer's certified geometries\newline{\footnotesize\color{metaGrey}\claimid{statistics/rem:two-tags}} & \pageref{rem:two-tags}\\
\addlinespace[2pt]
\mref{Theorem}{thm:binomial-tails} & Binomial tails in closed form\newline{\footnotesize\color{metaGrey}\claimid{statistics/thm:binomial-tails}} & \pageref{thm:binomial-tails}\\
\addlinespace[2pt]
\mref{Corollary}{cor:tail-forms} & The two strict binomial tails\newline{\footnotesize\color{metaGrey}\claimid{statistics/cor:tail-forms}} & \pageref{cor:tail-forms}\\
\addlinespace[2pt]
\mref{Lemma}{lem:exact-strict-tails} & Exact strict tails at a rational success probability\newline{\footnotesize\color{metaGrey}\claimid{statistics/lem:exact-strict-tails}} & \pageref{lem:exact-strict-tails}\\
\addlinespace[2pt]
\mref{Theorem}{thm:box-oracle} & Acceptance support and the exact box oracle\newline{\footnotesize\color{metaGrey}\claimid{statistics/thm:box-oracle}} & \pageref{thm:box-oracle}\\
\addlinespace[2pt]
\mref{Remark}{rem:box-oracle-consumer} & The closed-form box minimum for the consuming engine\newline{\footnotesize\color{metaGrey}\claimid{statistics/rem:box-oracle-consumer}} & \pageref{rem:box-oracle-consumer}\\
\addlinespace[2pt]
\mref{Definition}{def:envelope} & Honest envelope and its envelope set\newline{\footnotesize\color{metaGrey}\claimid{statistics/def:envelope}} & \pageref{def:envelope}\\
\addlinespace[2pt]
\mref{Lemma}{lem:hon-degenerate} & The envelope assumption at a degenerate envelope\newline{\footnotesize\color{metaGrey}\claimid{statistics/lem:hon-degenerate}} & \pageref{lem:hon-degenerate}\\
\addlinespace[2pt]
\mref{Definition}{def:conservative-evaluation} & Exact data and conservative evaluation\newline{\footnotesize\color{metaGrey}\claimid{statistics/def:conservative-evaluation}} & \pageref{def:conservative-evaluation}\\
\addlinespace[2pt]
\mref{Assumption}{ass:hon} & Honest i.i.d. rounds\newline{\footnotesize\color{metaGrey}\claimid{statistics/ass:hon}} & \pageref{ass:hon}\\
\addlinespace[2pt]
\mref{Assumption}{ass:hon-env} & Honest rounds inside a declared envelope\newline{\footnotesize\color{metaGrey}\claimid{statistics/ass:hon-env}} & \pageref{ass:hon-env}\\
\addlinespace[2pt]
\mref{Assumption}{ass:alphabet} & An acceptance value recorded in every round\newline{\footnotesize\color{metaGrey}\claimid{statistics/ass:alphabet}} & \pageref{ass:alphabet}\\
\addlinespace[2pt]
\mref{Assumption}{ass:reach} & An honest law produced by a feasible state\newline{\footnotesize\color{metaGrey}\claimid{statistics/ass:reach}} & \pageref{ass:reach}\\
\addlinespace[2pt]
\mref{Requirement}{req:geometry-tag} & Acceptance sets under exactly the two geometry tags\newline{\footnotesize\color{metaGrey}\claimid{statistics/req:geometry-tag}} & \pageref{req:geometry-tag}\\
\addlinespace[2pt]
\mref{Requirement}{req:box-bounds-key} & Recorded bounds of a canonical box\newline{\footnotesize\color{metaGrey}\claimid{statistics/req:box-bounds-key}} & \pageref{req:box-bounds-key}\\
\addlinespace[2pt]
\mref{Requirement}{req:acceptance-support} & The emitted acceptance support\newline{\footnotesize\color{metaGrey}\claimid{statistics/req:acceptance-support}} & \pageref{req:acceptance-support}\\
\addlinespace[2pt]
\mref{Requirement}{req:matched-alphabet} & Alphabet and round count taken from the consuming declaration\newline{\footnotesize\color{metaGrey}\claimid{statistics/req:matched-alphabet}} & \pageref{req:matched-alphabet}\\
\addlinespace[2pt]
\mref{Definition}{def:endpoint-rule} & Admissible inputs, endpoint rules and guarded endpoints\newline{\footnotesize\color{metaGrey}\claimid{statistics/binomial-box/def:endpoint-rule}} & \pageref{def:endpoint-rule}\\
\addlinespace[2pt]
\mref{Lemma}{lem:guard} & The honest-containment guard for every rule\newline{\footnotesize\color{metaGrey}\claimid{statistics/binomial-box/lem:guard}} & \pageref{lem:guard}\\
\addlinespace[2pt]
\mref{Definition}{def:interval} & The exact binomial quantile rule\newline{\footnotesize\color{metaGrey}\claimid{statistics/binomial-box/def:interval}} & \pageref{def:interval}\\
\addlinespace[2pt]
\mref{Proposition}{prop:rule-exact} & Properties of the exact binomial quantile rule\newline{\footnotesize\color{metaGrey}\claimid{statistics/binomial-box/prop:rule-exact}} & \pageref{prop:rule-exact}\\
\addlinespace[2pt]
\mref{Remark}{rem:guard} & Benefit and cost of the honest-containment guard\newline{\footnotesize\color{metaGrey}\claimid{statistics/binomial-box/rem:guard}} & \pageref{rem:guard}\\
\addlinespace[2pt]
\mref{Remark}{rem:guard-activation} & Regimes in which the guard acts\newline{\footnotesize\color{metaGrey}\claimid{statistics/binomial-box/rem:guard-activation}} & \pageref{rem:guard-activation}\\
\addlinespace[2pt]
\mref{Definition}{def:construction} & The construction map\newline{\footnotesize\color{metaGrey}\claimid{statistics/binomial-box/def:construction}} & \pageref{def:construction}\\
\addlinespace[2pt]
\mref{Theorem}{thm:box-construction} & Properties of the constructed box\newline{\footnotesize\color{metaGrey}\claimid{statistics/binomial-box/thm:box-construction}} & \pageref{thm:box-construction}\\
\addlinespace[2pt]
\mref{Remark}{rem:selector} & The role of the constrained outcome set\newline{\footnotesize\color{metaGrey}\claimid{statistics/binomial-box/rem:selector}} & \pageref{rem:selector}\\
\addlinespace[2pt]
\mref{Lemma}{lem:envelope-dominance} & Binomial dominance under bounded conditional success probabilities\newline{\footnotesize\color{metaGrey}\claimid{statistics/binomial-box/lem:envelope-dominance}} & \pageref{lem:envelope-dominance}\\
\addlinespace[2pt]
\mref{Theorem}{thm:completeness-bound} & The completeness bound\newline{\footnotesize\color{metaGrey}\claimid{statistics/binomial-box/thm:completeness-bound}} & \pageref{thm:completeness-bound}\\
\addlinespace[2pt]
\mref{Remark}{rem:epsstar} & Exact evaluation of $\varepsilon^{\star}$ and its conservative report\newline{\footnotesize\color{metaGrey}\claimid{statistics/binomial-box/rem:epsstar}} & \pageref{rem:epsstar}\\
\addlinespace[2pt]
\mref{Remark}{rem:sound-indep} & Soundness does not depend on the honest-round assumptions\newline{\footnotesize\color{metaGrey}\claimid{statistics/binomial-box/rem:sound-indep}} & \pageref{rem:sound-indep}\\
\addlinespace[2pt]
\mref{Remark}{rem:envelope} & Benefit and cost of an honest envelope\newline{\footnotesize\color{metaGrey}\claimid{statistics/binomial-box/rem:envelope}} & \pageref{rem:envelope}\\
\addlinespace[2pt]
\mref{Theorem}{thm:honest-containment} & Honest containment and infeasibility\newline{\footnotesize\color{metaGrey}\claimid{statistics/binomial-box/thm:honest-containment}} & \pageref{thm:honest-containment}\\
\addlinespace[2pt]
\mref{Remark}{rem:hlow-consequence} & Scope of the statements on infeasibility and minimal divergence\newline{\footnotesize\color{metaGrey}\claimid{statistics/binomial-box/rem:hlow-consequence}} & \pageref{rem:hlow-consequence}\\
\addlinespace[2pt]
\mref{Definition}{def:monotone-rule} & Monotone allocation rule\newline{\footnotesize\color{metaGrey}\claimid{statistics/binomial-box/def:monotone-rule}} & \pageref{def:monotone-rule}\\
\addlinespace[2pt]
\mref{Lemma}{lem:widening} & Widening the box keeps its guarantees\newline{\footnotesize\color{metaGrey}\claimid{statistics/binomial-box/lem:widening}} & \pageref{lem:widening}\\
\addlinespace[2pt]
\mref{Lemma}{lem:cheb} & Chebyshev's inequality for a binomial count\newline{\footnotesize\color{metaGrey}\claimid{statistics/binomial-box/lem:cheb}} & \pageref{lem:cheb}\\
\addlinespace[2pt]
\mref{Theorem}{thm:monotonicity-limits} & Monotonicity in the budget and limiting cases\newline{\footnotesize\color{metaGrey}\claimid{statistics/binomial-box/thm:monotonicity-limits}} & \pageref{thm:monotonicity-limits}\\
\addlinespace[2pt]
\mref{Remark}{rem:n-nonmono} & Non-monotonicity of the exact rule in the round count\newline{\footnotesize\color{metaGrey}\claimid{statistics/binomial-box/rem:n-nonmono}} & \pageref{rem:n-nonmono}\\
\addlinespace[2pt]
\mref{Lemma}{lem:log-enclosure} & Rational bounds on the logarithm from integer square roots\newline{\footnotesize\color{metaGrey}\claimid{statistics/binomial-box/lem:log-enclosure}} & \pageref{lem:log-enclosure}\\
\addlinespace[2pt]
\mref{Lemma}{lem:stirling-enclosure} & Robbins' Stirling bounds with Wallis's product\newline{\footnotesize\color{metaGrey}\claimid{statistics/binomial-box/lem:stirling-enclosure}} & \pageref{lem:stirling-enclosure}\\
\addlinespace[2pt]
\mref{Lemma}{lem:bernstein-tail} & Bernstein's inequality for a binomial count\newline{\footnotesize\color{metaGrey}\claimid{statistics/binomial-box/lem:bernstein-tail}} & \pageref{lem:bernstein-tail}\\
\addlinespace[2pt]
\mref{Proposition}{prop:rule-bernstein} & The Bernstein endpoint rule\newline{\footnotesize\color{metaGrey}\claimid{statistics/binomial-box/prop:rule-bernstein}} & \pageref{prop:rule-bernstein}\\
\addlinespace[2pt]
\mref{Remark}{rem:completeness-on-rules} & Completeness figures on a closed-form rule\newline{\footnotesize\color{metaGrey}\claimid{statistics/binomial-box/rem:completeness-on-rules}} & \pageref{rem:completeness-on-rules}\\
\addlinespace[2pt]
\mref{Lemma}{lem:pmf-enclosure} & Rational enclosure of the logarithm of a binomial mass point\newline{\footnotesize\color{metaGrey}\claimid{statistics/binomial-box/lem:pmf-enclosure}} & \pageref{lem:pmf-enclosure}\\
\addlinespace[2pt]
\mref{Lemma}{lem:tail-majorant} & The binomial tail from a window and a geometric majorant\newline{\footnotesize\color{metaGrey}\claimid{statistics/binomial-box/lem:tail-majorant}} & \pageref{lem:tail-majorant}\\
\addlinespace[2pt]
\mref{Proposition}{prop:bracketed-endpoints} & The enclosure evaluation of the exact rule\newline{\footnotesize\color{metaGrey}\claimid{statistics/binomial-box/prop:bracketed-endpoints}} & \pageref{prop:bracketed-endpoints}\\
\addlinespace[2pt]
\mref{Remark}{rem:local-evaluation} & Evaluating the exact rule locally\newline{\footnotesize\color{metaGrey}\claimid{statistics/binomial-box/rem:local-evaluation}} & \pageref{rem:local-evaluation}\\
\addlinespace[2pt]
\mref{Requirement}{req:rule-in-force} & The rule the request selects\newline{\footnotesize\color{metaGrey}\claimid{statistics/binomial-box/req:rule-in-force}} & \pageref{req:rule-in-force}\\
\addlinespace[2pt]
\mref{Requirement}{req:request-defaults} & Defaults for the constrained outcomes and the allocation\newline{\footnotesize\color{metaGrey}\claimid{statistics/binomial-box/req:request-defaults}} & \pageref{req:request-defaults}\\
\addlinespace[2pt]
\mref{Requirement}{req:input-refusal} & Refusal of an input that violates a decidable clause\newline{\footnotesize\color{metaGrey}\claimid{statistics/binomial-box/req:input-refusal}} & \pageref{req:input-refusal}\\
\addlinespace[2pt]
\mref{Requirement}{req:completeness-report} & The reported completeness figure\newline{\footnotesize\color{metaGrey}\claimid{statistics/binomial-box/req:completeness-report}} & \pageref{req:completeness-report}\\
\addlinespace[2pt]
\mref{Requirement}{req:honest-model-record} & The honest-model record beside every completeness figure\newline{\footnotesize\color{metaGrey}\claimid{statistics/binomial-box/req:honest-model-record}} & \pageref{req:honest-model-record}\\
\addlinespace[2pt]
\mref{Requirement}{req:search-procedure} & One fixed search for the enclosure evaluation\newline{\footnotesize\color{metaGrey}\claimid{statistics/binomial-box/req:search-procedure}} & \pageref{req:search-procedure}\\
\addlinespace[2pt]
\mref{Requirement}{req:enclosure-substrate} & Log-space arithmetic with directed rounding\newline{\footnotesize\color{metaGrey}\claimid{statistics/binomial-box/req:enclosure-substrate}} & \pageref{req:enclosure-substrate}\\
\addlinespace[2pt]
\mref{Requirement}{req:endpoint-exact-record} & The bracket verdict per outcome and side\newline{\footnotesize\color{metaGrey}\claimid{statistics/binomial-box/req:endpoint-exact-record}} & \pageref{req:endpoint-exact-record}\\
\addlinespace[2pt]
\mref{Theorem}{thm:singleton-completeness} & Completeness of a singleton acceptance set\newline{\footnotesize\color{metaGrey}\claimid{statistics/singleton/thm:singleton-completeness}} & \pageref{thm:singleton-completeness}\\
\addlinespace[2pt]
\mref{Remark}{rem:singleton-honest} & Certifiability and completeness of the singleton\newline{\footnotesize\color{metaGrey}\claimid{statistics/singleton/rem:singleton-honest}} & \pageref{rem:singleton-honest}\\
\addlinespace[2pt]
\mref{Requirement}{req:singleton-budget-record} & The completeness budget of a declared singleton\newline{\footnotesize\color{metaGrey}\claimid{statistics/singleton/req:singleton-budget-record}} & \pageref{req:singleton-budget-record}\\
\addlinespace[2pt]
\mref{Requirement}{req:singleton-point-form} & Point-form honest input for the singleton\newline{\footnotesize\color{metaGrey}\claimid{statistics/singleton/req:singleton-point-form}} & \pageref{req:singleton-point-form}\\
\bottomrule
\end{xltabular}}

\clearpage
\bookmarksetupnext{level=part}
\printindex

\clearpage\phantomsection\bookmarksetupnext{level=part}%
\addcontentsline{toc}{section}{References}
\bibliographystyle{plain}
\bibliography{../../../refs/references}

\end{document}
